\documentclass[11pt,a4paper]{article}
\usepackage[T1]{fontenc}
\usepackage[utf8]{inputenc}
\usepackage{lmodern}
\usepackage[margin=26mm,headheight=15pt]{geometry}
\usepackage{amsmath,amssymb,amsthm,mathtools,aliascnt,mathrsfs}
\usepackage{microtype,booktabs,longtable,array,enumitem}
\usepackage[dvipsnames]{xcolor}
\usepackage{fancyhdr,xurl}
\usepackage[colorlinks=true,linkcolor=MidnightBlue,citecolor=MidnightBlue,urlcolor=MidnightBlue]{hyperref}
\usepackage[nameinlink,noabbrev]{cleveref}
\hypersetup{pdftitle={Definable Surreal Numbers and Omnific Integers},pdfauthor={Merged research report for the Surreal project},pdfsubject={Ambient definability, HOD, reversible and bounded omnific coding, a transcendence dichotomy, the maximal initial core, real parameters, countable assembly, strong symmetries, and definable operations}}
\pagestyle{fancy}\fancyhf{}
\fancyhead[L]{\small Definable surreals and omnific integers}
\fancyhead[R]{\small\thepage}
\renewcommand{\headrulewidth}{0.3pt}
\setlength{\parskip}{3pt}
\setlength{\emergencystretch}{2em}
\setlist{itemsep=3pt,topsep=4pt,leftmargin=2em}
\numberwithin{equation}{section}
\allowdisplaybreaks[2]
\newtheorem{theorem}{Theorem}[section]
\newaliascnt{lemma}{theorem}
\newtheorem{lemma}[lemma]{Lemma}
\aliascntresetthe{lemma}
\crefname{lemma}{lemma}{lemmas}
\Crefname{lemma}{Lemma}{Lemmas}
\newaliascnt{proposition}{theorem}
\newtheorem{proposition}[proposition]{Proposition}
\aliascntresetthe{proposition}
\crefname{proposition}{proposition}{propositions}
\Crefname{proposition}{Proposition}{Propositions}
\newaliascnt{corollary}{theorem}
\newtheorem{corollary}[corollary]{Corollary}
\aliascntresetthe{corollary}
\crefname{corollary}{corollary}{corollaries}
\Crefname{corollary}{Corollary}{Corollaries}
\theoremstyle{definition}
\newaliascnt{definition}{theorem}
\newtheorem{definition}[definition]{Definition}
\aliascntresetthe{definition}
\crefname{definition}{definition}{definitions}
\Crefname{definition}{Definition}{Definitions}
\newaliascnt{example}{theorem}
\newtheorem{example}[example]{Example}
\aliascntresetthe{example}
\crefname{example}{example}{examples}
\Crefname{example}{Example}{Examples}
\newaliascnt{question}{theorem}
\newtheorem{question}[question]{Question}
\aliascntresetthe{question}
\crefname{question}{question}{questions}
\Crefname{question}{Question}{Questions}
\theoremstyle{remark}
\newaliascnt{remark}{theorem}
\newtheorem{remark}[remark]{Remark}
\aliascntresetthe{remark}
\crefname{remark}{remark}{remarks}
\Crefname{remark}{Remark}{Remarks}
\newcommand{\No}{\mathbf{No}}
\newcommand{\Oz}{\mathbf{Oz}}
\newcommand{\Ord}{\mathrm{Ord}}
\newcommand{\OD}{\mathrm{OD}}
\newcommand{\HOD}{\mathrm{HOD}}
\newcommand{\ZFC}{\mathrm{ZFC}}
\newcommand{\N}{\mathbb N}
\newcommand{\Z}{\mathbb Z}
\newcommand{\Q}{\mathbb Q}
\newcommand{\R}{\mathbb R}
\newcommand{\C}{\mathbb C}
\newcommand{\bd}{\operatorname{b}}
\newcommand{\ct}{\operatorname{ct}}
\newcommand{\st}{\operatorname{st}}
\newcommand{\supp}{\operatorname{supp}}
\newcommand{\NF}{\operatorname{NF}}
\newcommand{\dcl}{\operatorname{dcl}}
\newcommand{\acl}{\operatorname{acl}}
\newcommand{\Frac}{\operatorname{Frac}}
\newcommand{\TC}{\operatorname{TC}}
\newcommand{\dom}{\operatorname{dom}}
\newcommand{\ran}{\operatorname{ran}}
\newcommand{\sgn}{\operatorname{sgn}}
\newcommand{\Sat}{\operatorname{Sat}}
\newcommand{\RC}{\operatorname{rcl}}
\newcommand{\Fix}{\operatorname{Fix}}
\newcommand{\Def}{\mathrm{Def}}
\newcommand{\Code}{\mathsf C}
\newcommand{\Decode}{\mathsf D}
\newcommand{\Two}{\mathsf T}
\newcommand{\Mc}{\mathsf Q}
\newcommand{\esc}{\operatorname{esc}}
\newcommand{\SetDec}{\operatorname{Set}}
\newcommand{\Core}{\operatorname{Core}}
\newcommand{\Cell}{\mathcal U}
\newcommand{\KH}{K_{\HOD}}
\newcommand{\IH}{I_{\HOD}}
\newcommand{\ps}{\mathrel{\prec_s}}
\newcommand{\pseq}{\mathrel{\preceq_s}}
\newcommand{\res}{\mathbin{\upharpoonright}}
\newcommand{\cut}[2]{\left\{\,#1\;\middle|\;#2\,\right\}}
\newcommand{\Lexp}{\mathcal L_{\exp}}
\newcommand{\Lor}{\mathcal L_{\mathrm{or}}}
\newcommand{\repoSHA}{\texttt{9a385d3957bdfe3d9ea79f9a524751c90bd2c894}}
\newcommand{\repopath}[1]{\nolinkurl{#1}}
\newcommand{\src}[1]{\textup{\textsf{\small[#1]}}}
\newcommand{\mergetag}{\textup{\textsf{\small[merge]}}}
\newcommand{\lbl}[1]{\nolinkurl{#1}}
\newcommand{\fn}[1]{\nolinkurl{#1}}
\newcolumntype{P}[1]{>{\raggedright\arraybackslash}p{#1}}
\newcommand{\writetag}{\textup{\textsf{\small[write]}}}
\newenvironment{writenote}[1][]{\par\medskip\begingroup\small\noindent\textbf{[write]}\if\relax\detokenize{#1}\relax\else~\textbf{#1.}\fi~\ignorespaces}{\par\endgroup\medskip}
\newenvironment{dupstatement}[1]{\par\medskip\noindent\textbf{#1.}\ \begingroup\itshape\ignorespaces}{\endgroup\par\medskip}
% Parts II and III (batch 93)
\newcommand{\ROD}{\mathrm{ROD}}
\newcommand{\COD}{\mathrm{COD}}
\newcommand{\RD}{\mathrm{RD}}
\newcommand{\HR}{\HOD_{\R}}
\newcommand{\KR}{K_{\R}}
\newcommand{\IR}{I_{\R}}
\newcommand{\Kco}{K_{\mathrm{co}}}
\newcommand{\Ico}{I_{\mathrm{co}}}
\newcommand{\Hcup}{H_{\cup}}
\newcommand{\Asm}{\mathsf A_{\omega}}
\newcommand{\Zc}{\mathcal Z}
\newcommand{\Qrange}{\mathcal Q}
\newcommand{\Spec}{\operatorname{Spec}}
\newcommand{\Szero}{\mathfrak S_0}
\newcommand{\Somega}{\mathfrak S_\Omega}
\newcommand{\eps}{\varepsilon}
\newcommand{\smallrel}{\mathrel{\prec_{\!A}}}
\newcommand{\seq}[1]{\langle #1\rangle}
\newcommand{\cf}{\operatorname{cf}}
\newcommand{\cof}{\operatorname{cof}}
\begin{document}
\hypersetup{pageanchor=false}
\begin{titlepage}
\centering
\vspace*{0mm}
{\Large Merged research report for the Surreal project\par}
\vspace{7mm}
{\huge\bfseries Definable Surreal Numbers\\[3mm] and Omnific Integers\par}
\vspace{5mm}
{\Large Ambient definability, reversible and bounded omnific coding,\\[1mm]
hereditary ordinal definability, a transcendence dichotomy,\\[1mm]
and the maximal initial core\par}
\vspace{5mm}
{\large Part~I built from two manuscripts dated September 23, 2026;\\[1mm]
Parts~II and~III from three manuscripts dated October 4, 2026\par}
\vspace{4mm}
\begin{minipage}{\textwidth}
\begin{abstract}
A surreal number cannot be called definable without specifying a language,
a universe, and permitted parameters. We develop ambient set-theoretic
notions, keep them apart from definability in the ordered field and its
exponential expansion, and give an algebraic theory of the resulting
numbers. For a fixed parameter set, the ambiently definable surreals form a
real closed field closed under the Conway monomial map, the exponential and
the omnific floor; their omnific integers form an integer part whose
fraction field is exactly that field, with a denominator chosen from the
same parameters. Over a transitive model the definable field is an
elementary exponential subfield. Ordinal-definable surreals are exactly
the surreals of $\HOD$.

Two reversible omnific codes compress every surreal: a sign code with
leading term $\omega$, zero constant term and coefficients in
$\{1,2,3\}$, landing in the omnific interval $\omega<z<\omega+\omega^{1/2}$,
and the monomial code $x\mapsto\omega^{(1+x/\sqrt{1+x^2})/2}$, landing
between $0$ and $\omega$ and localizing to every infinitely wide omnific
interval. Consequently ordinal definability of all omnific integers in one
such interval is equivalent to $V=\HOD$, and the parameter-free version in
a model is equivalent to pointwise definability. If $V\ne\HOD$, the first
interval contains an explicit $\Ord$-indexed algebraically independent
family over the $\HOD$ surreals. If $\rho$ is the first ordinal of a
transitive model $M$ not definable from $A$, then $\rho$ is an epsilon
number and the definable surreals of birthday below $\rho$ form the largest
initial subclass of the definable field; it is again an elementary
exponential subfield and the fraction field of its own omnific integers.
Explicit normal forms show in both directions that definability of a whole
series differs from definability of its terms. We also treat
complexity-bounded escape ordinals, homogeneous forcing that adds no reals,
relativization, the surcomplex pairs, language dependence, and the finite
quotients of the definable omnific ring.

Part~II (added 4 October 2026) lets the real parameter vary: the
resulting field is captured by one real exactly when $\HOD_{\R}$
satisfies Choice; closure under countable Hahn sums is equivalent to
real--ordinal definability of countable ordinal sequences, holds in a
Cohen model, fails after Namba forcing over $L$, and so is independent of
ZFC; and the common fixed field of the strong automorphisms fixing all
reals and ordinals is the field of normal forms over the rational span of
the ordinals. Part~III (added the same day) interprets set theory
parameter-free in the field with the omnific predicate and the omega map,
proves absoluteness of normal forms between inner models without a
same-reals hypothesis, and characterizes definable sections of the real
decoder by correctness of $\omega_1$ in $\HOD$.

Part~I merges two independently written manuscripts and Part~II two
more; shared results are printed once with their sources named, and every
source's limitations are kept. It is an AI-assisted draft: it has not been
refereed, and nothing in it has been formalized in Lean.
\end{abstract}
\end{minipage}
\vfill
\begin{minipage}{\textwidth}
\footnotesize
Sources: manuscript 04 (the base text) and manuscript 07, both pinned to
repository commit \texttt{9a385d3}; their titles are in
\cref{dsn:sub:provenance}. The first non-$\HOD$ birthday is printed by
citation of the universes report. The provenance of every section is
recorded in \cref{dsn:app:provenance}; non-claims are collected in
\cref{dsn:app:nonclaims}. Finite checks accompany the report and test
finite identities only. Parts~II and~III (batch~93, manuscripts 02, 05 and
06) record their provenance in \cref{dsn:rp:sub:provenance,dsn:op:sub:provenance}.
\end{minipage}
\end{titlepage}
\hypersetup{pageanchor=true}
\setcounter{tocdepth}{2}
\begingroup\small\setlength{\parskip}{0pt}
\tableofcontents
\endgroup
{\footnotesize\noindent\textbf{[write]} Since 4 October 2026 the report has three Parts:
Part~\ref{dsn:part:rp} (from \cref{dsn:rp:sec:scope}; batch~93, manuscripts 02
and 05, labels \texttt{dsn:rp:} and \texttt{dsn:ca:}) treats real parameters,
countable assembly and strong symmetries, and Part~\ref{dsn:part:op} (from
\cref{dsn:op:sec:scope}; manuscript 06, labels \texttt{dsn:op:}) definable
operations and an interpretation of set theory. Part~\ref{dsn:part:one}'s
numbering and labels are unchanged.\par}
\clearpage

\part{Ambient definability, reversible and bounded omnific coding, hereditary ordinal definability, a transcendence dichotomy, and the maximal initial core}
\label{dsn:part:one}

\section{The question, the results, and how this report was assembled}
\label{dsn:sec:intro}

\subsection{The question}

The familiar phrase ``definable real number'' conceals several choices.
A positive root of a rational polynomial is definable in an ordered field;
a computable real can instead be defined by a statement about its
algorithm; and a set-theoretic definition can refer to ordinals, models,
and forcing extensions. These are different notions, not alternative
notations for one fixed countable field. The user-specified reference on
definable reals is a useful starting point \cite{WikiDef}; it is not used
as a proof source. The foundational analysis and the
pointwise-definability results used below come from
Hamkins--Linetsky--Reitz \cite{HLR}.

For surreals the distinction becomes essential, because the sign sequence
representing a number is a set while the totality of all surreals is a
proper class internally. The pure ordered field has no parameter-free
definable infinite element, whereas the first infinite ordinal $\omega$
has a very short parameter-free definition in set theory. An omnific
integer can simultaneously have elementary-looking normal-form
coefficients and contain the information of an arbitrary set of ordinals.
Our goal is to make these facts into a coherent theory rather than to
announce an unspecified collection of ``describable'' numbers.

The central object is not a syntactic list of expressions and not a
putative truth predicate inside the universe under discussion. Fix a model
$M$ of ZFC and an external collection $A$ of allowed parameters. An element
of $M$ is \emph{$A$-definable over $M$} when a standard first-order formula
of set theory, using a finite tuple from $A$, uniquely specifies it in $M$
(\cref{dsn:def:ambient}). We write $D^M_A$ for the definable closure,
$K^M_A$ for its surreals and $I^M_A$ for its omnific integers. When the
universe $V$ itself is meant, the same words are a metatheoretic scheme.

\subsection{Main results}\label{dsn:sub:main}

Let $\No$ be represented by ordinal-length sign sequences, and let $\Oz$
be the classical omnific-integer ring. Write
\begin{equation}\label{dsn:eq:maincell}
 \Cell=\{z\in\Oz:\ \omega<z<\omega+\omega^{1/2},\ \ct(z)=0\}.
\end{equation}
The exponentiation $x\mapsto\omega^x$ throughout this report is the
\emph{Conway monomial map}, also written $\Omega(x)$; it is not
$x\mapsto\exp(x\log\omega)$ (\cref{dsn:sub:normal}).

\paragraph{The definable field \src{04, 07}.}
For every parameter family the definable surreals $K_A$ form a real closed
field closed under $\Omega$, $\Omega^{-1}$ on monomials, $\exp$, $\log$
and the omnific floor; $I_A=K_A\cap\Oz$ is an integer part and
$K_A=\Frac(I_A)$ with a denominator chosen by one rule from the same
parameters (\cref{dsn:thm:hull}). Over a transitive model
$K^M_A\prec\No^M$ in the exponential language (\cref{dsn:thm:elementary}).
The ordinal-definable surreals are the surreal field of $\HOD$
(\cref{dsn:thm:HOD}).

\paragraph{Reversible compression \src{04, 07}.}
There is a parameter-free set-theoretically definable injection
$\Code:\No\to\Cell$ with a parameter-free definable inverse on its image,
preserving the definability degree for every parameter family
(\cref{dsn:thm:code}, source 04). The monomial code
$\Mc(x)=\omega^{\tau(x)}$, $\tau(x)=\frac12(1+x/\sqrt{1+x^2})$, is a
parameter-free definable order embedding of $\No$ into
$\Oz\cap(0,\omega)$ with the same property; it localizes to every
infinitely wide omnific interval, and there is a definable class function
$\esc$ on the ordinals with values in $\Oz\cap(0,\omega)$ and unbounded
birthdays (\cref{dsn:thm:bounded-code,dsn:thm:local-code,dsn:thm:bdescape},
source 07). Neither code is asserted to be a ring homomorphism or to be
definable in the pure ordered field.

\paragraph{One interval tests the universe \src{04, 07}.}
Every element of $\Cell$ is ordinal definable if and only if $V=\HOD$, and
the same holds for $\Oz\cap(0,\omega)$ and for the monomials $\omega^a$
with $0<a<1$. In the external semantics of a model $M\models\ZFC$, every
element of $\Cell^M$ is parameter-free definable in $M$ if and only if $M$
is pointwise definable (\cref{dsn:thm:globalOD,dsn:thm:pointwise}). Thus a
fixed leading asymptotic term, or a fixed numerical bound, does not bound
set-theoretic information.

\paragraph{A transcendence dichotomy \src{04}.}
Either $\No=\KH$, or $\Cell$ contains an $\Ord$-indexed algebraically
independent family over $\KH$ (\cref{dsn:thm:dichotomy}). The proof is
explicit: starting with any $t\in(0,1)\setminus\KH$, use
\begin{equation}\label{dsn:eq:introfam}
 z_\alpha=\omega+\omega^{\,t\omega^{-(\alpha+1)}}
 \qquad(\alpha\in\Ord).
\end{equation}
The obstruction to a polynomial relation is separation of normal-form
supports into distinct additive cosets of $\KH$.

\paragraph{The maximal initial core \src{07}.}
Let $M$ be a transitive model. If some ordinal of $M$ is not
$A$-definable and $\rho$ is the least such ordinal, taken externally, then
$\rho$ is an epsilon number, and
\[
 \Core^M_A=K^M_A\cap(\No_{<\rho})^M
\]
is the largest initial subclass of $K^M_A$. It is an elementary
exponential subfield of $\No^M$, and it is the fraction field of its own
omnific integers, with the birthday-controlled denominator
$\omega^{\bd(x)+1}$ (\cref{dsn:thm:core,dsn:thm:core-fractions}).

\paragraph{Two uniformity gaps \src{04, 07}.}
A fixed definable support and individually definable coefficients need not
determine a definable coefficient assignment
(\cref{dsn:thm:termwise,dsn:thm:bit-code}); conversely a uniformly
definable omnific normal form can have a nondefinable term
(\cref{dsn:prop:collector,dsn:thm:reverse-gap}). The correct invariant is
the entire canonical normal-form graph (\cref{dsn:prop:nf-invariant}).

\paragraph{Further results.}
The first birthday of a non-ordinal-definable surreal is the first
subset disagreement between $\HOD$ and $V$ (\cref{dsn:thm:firstbirthday},
the case $M=\HOD$ of a theorem of the universes report);
complexity-bounded definitions admit escape ordinals
(\cref{dsn:thm:sigmaescape}); no definable set contains its own definable
hull (\cref{dsn:thm:noncollection}); homogeneous forcing produces the
dichotomy without adding reals (\cref{dsn:thm:forcing}); and the
definable omnific ring has only the finite quotients $\Z/n\Z$
(\cref{dsn:thm:finite-quotients}).

\subsection{What is prior work, and what is proposed here}

The sign construction, normal forms, the monomial map, real closedness,
the exponential, and the omnific integer part are classical
\cite{Conway,Gonshor,vdDE,vdDEerr,LM}. Quantifier elimination, o-minimality,
ordinal definability, HOD, partial truth predicates, and Mostowski
collapse are standard tools \cite{Marker,vddTame,Jech,HLR}. Van den Dries
and Ehrlich supply the elementary exponential-field and epsilon-cutoff
results \cite{vdDE,vdDEerr}; Bournez and Guilmant restate the corrected
cutoff theorem \cite{BG}. We do not claim these tools, or their routine
closure consequences, as new.

Hamkins proved in 2012 that the existence of a definable
ordinal-indexed dense image in $\No$, and the density of the
ordinal-definable surreals, are equivalent to $V=\HOD$
\cite{HamkinsDensity}. This density theorem is prior work; the bounded
omnific tests below strengthen the form of the hypothesis, not the
underlying idea (\cref{dsn:thm:hamkins}). Source 04 did not cite it; the
merge adds the credit.

Chen, Hamkins, and Yang have announced that the surreal ordered field
expanded by birthday order is bi-interpretable with set theory
\cite{CHY}. The collection credits that announcement in its birthday-cutoff
report (\lbl{hset:thm:chy}). In particular, the general idea of encoding
well-founded membership graphs by surreal signs is \emph{not} a new
contribution of this report. No theorem below needs an unverified version
of that announced bi-interpretation.

The candidate-original contributions are: the particular reversible
omnific codes with fixed leading data or fixed numerical bounds, and their
one-interval tests for $V=\HOD$ and for pointwise definability (04, 07);
the explicit \emph{proper-class transcendence dichotomy in the same
interval} (04); the local coding, bounded birthday escape and bounded
noncollection (07); the two explicit normal-form uniformity gaps (04,
07); and the maximal initial elementary exponential core with its
birthday-controlled denominators (07). The birthday, complexity and
forcing results are refinements and applications with less ambitious
originality claims. Neither source's literature search located these
precise statements; neither search was exhaustive, and no theorem is
described as settling a named open problem in the published literature.
Mathematical correctness and bibliographic novelty are separate questions.

\subsection{How this report was assembled}\label{dsn:sub:provenance}

The report is built from two manuscripts, both dated 23 September 2026 and
both pinned to repository commit \repoSHA{} (35 commits before the
placement commit \texttt{a4dcb91}); they are items 04 and 07 of batch 27
and keep those numbers here.
\begin{itemize}
\item \textbf{04}: \emph{Definable Surreal Numbers and Omnific Integers:
Ambient definability, reversible coding, hereditary ordinal definability,
and a transcendence dichotomy}. It is the base text. It contributes the
general set-model framework, the sign code and the two-term code, the
one-interval tests (including the pointwise test for arbitrary set
models), support-subfield amplification and the $\HOD$ dichotomy,
localization, the termwise counterexample and the collector, the first
non-$\HOD$ birthday, the $\Sigma_n$ escape, forcing, relativization and
the surcomplex companion.
\item \textbf{07}: \emph{Definable Surreal Numbers and Omnific Integers:
Ambient truth, bounded coding, and the maximal initial core}. It works over
an externally fixed transitive model throughout. It contributes the
exponential elementarity of the definable field, the cut denominator rule,
the valuation theorem, the monomial and local codes, the birthday escape
and bounded noncollection, the bit code and the reverse gap, the maximal
initial core and its denominators, Hamkins's density theorem with the
topology corollary, the language-dependence results, and the arithmetic of
the definable omnific ring.
\end{itemize}
The source manuscripts are not shipped. Their code, recorded data and
audits are, under the prefixes \fn{04-hod-transcendence-} and
\fn{07-initial-core-}.

\paragraph{Why one report, and why 04 is the base.}
The two manuscripts have the same title and the same spine: the definable
field and its integer part with same-parameter denominators, the $\HOD$
identification, bounded tests for $V=\HOD$ and for pointwise
definability, the gap between whole normal forms and their terms,
noncollection, and the definable closure of the pure field. Each has a
large result the other lacks. On the shared theorems 04 has the weaker
hypotheses: its pointwise test holds for arbitrary set models, including
ill-founded ones, and its field theory is stated both for $V$ as a scheme
and for set models. 07's stronger conclusions (exponential elementarity,
the core) are proved under 07's own transitive-model hypothesis, and they
keep it here. The two manuscripts do not contradict each other or any
report in the collection.

\paragraph{Printed once.} Each of the following is printed once, crediting
both sources: unique-output closure (\cref{dsn:lem:closure}); the omnific
floor (\cref{dsn:prop:floor}); the definable closure of the pure field
(\cref{dsn:thm:purefield}); the definable field, integer part and fractions
(\cref{dsn:thm:hull}, with both denominator rules); closure under whole
definable cuts and sums (\cref{dsn:prop:uniform}); the $\HOD$
identification and the normal-form criterion
(\cref{dsn:thm:HOD,dsn:lem:HODsupport}); the bounded $V=\HOD$ test
(\cref{dsn:thm:globalOD}); the pointwise test (\cref{dsn:thm:pointwise});
noncollection (\cref{dsn:thm:noncollection}); the set-coding lemma
(\cref{dsn:lem:setcoding}); and the complexification
(\cref{dsn:prop:surcomplex}).

\paragraph{Kept as second routes.} 04's least-ordinal denominator and
07's cut denominator (\cref{dsn:thm:hull}); 04's sign code and 07's bit code
for the fixed-support phenomenon (\cref{dsn:thm:termwise,dsn:thm:bit-code});
04's collector and 07's reverse gap for the converse phenomenon
(\cref{dsn:prop:collector,dsn:thm:reverse-gap}); 04's two-term code and
07's monomial code (\cref{dsn:prop:twotermcode,dsn:thm:bounded-code}).

\paragraph{Cited, not reprinted.} 04's first-birthday theorem is the case
$M=\HOD$, $N=V$ of \lbl{univ:prop:beta} in the universes report, with the
same proof; it is printed by citation (\cref{dsn:thm:firstbirthday}).
07's arithmetic section restricts results of the omnific Diophantine and
quotient reports to definable fragments; the full-class antecedents are
named at each statement (\cref{dsn:sec:arithmetic}).

\paragraph{Added by the merge}, each tagged \mergetag{} with a complete
proof or an explicit citation:
\begin{itemize}
\item the fix of a gap in 04's proof that a proper $\KH$ is not
intrinsically definable (\cref{dsn:prop:notdefinable});
\item the remark on o-minimality of $\No^M_{\exp}$ used tacitly by 07
(\cref{dsn:rem:ominimal});
\item the extension of 07's bounded conditions of the pointwise test to
arbitrary set models (\cref{dsn:thm:pointwise}) and of 07's bounded
noncollection to arbitrary parameter families
(\cref{dsn:thm:noncollection});
\item the answer from the collection to both sources' language questions
(\cref{dsn:prop:pair}), and the re-scoped \cref{dsn:q:languages};
\item the relation to \lbl{odg:q:size} (\cref{dsn:rem:size}).
\end{itemize}
\Cref{dsn:app:stale} lists every stale statement corrected.

\begin{writenote}[Added 4 October 2026, batch~93]
This subsection describes Part~\ref{dsn:part:one}. Parts~\ref{dsn:part:rp}
and~\ref{dsn:part:op} were added from three manuscripts of batch~93 (02 and 05,
merged, and 06), each with its own provenance subsection
(\cref{dsn:rp:sub:provenance,dsn:op:sub:provenance}). Their notes in this Part
are the dated notes after \cref{dsn:rem:largecardinal} and in
\cref{dsn:app:provenance}; no statement, proof, number or label of this Part
changed.
\end{writenote}

\subsection{Conventions and renamed symbols}\label{dsn:sub:conventions}

The two sources used several letters for different objects. The
following convention is fixed before any use. Symbols of 04, the base, are
kept where possible and 07's are renamed; the few local renamings in 04's
text are listed as well.
\begin{itemize}
\item $D^M_A$ is the definable closure of $A$ in $(M,\in)$ (07's
$\Def_M(A)$); $K^M_A=D^M_A\cap\No^M$ are the definable surreals (07's
$D_M(A)$); $I^M_A=K^M_A\cap\Oz^M$ (07's $I_M(A)$); $O^M_A$ the definable
ordinals; $\rho^M_A$ the first omitted ordinal (07's $\rho_M(A)$).
\item $\Core^M_A$ is 07's maximal initial core (07's $K_M(A)$). The
letter $K$ never denotes the core.
\item $\Pi$ is the purely infinite ideal (07's $\mathcal J$); $\Pi_A=
K_A\cap\Pi$ (07's $\mathcal J_D$); $R_A=K_A\cap\R$ (07's $R_D$).
\item $\Code$ is 04's sign code, $\Decode$ its inverse; $\Mc$ and
$\Mc_{a,b}$ are 07's monomial and local codes (07's $C$, $C_{a,b}$).
\item $\Two$ is 04's two-term code, built on the rational map $h$;
$\esc$ is 07's birthday escape (07's $T$). The rational map of 04's
collector is renamed $\zeta$, and the width of 07's local code is $\ell$
(07's $h$).
\item $e_\beta=\omega^{-(\beta+1)}$ (04) and $\iota_\beta=
1/(\underline\beta+2)$ (07's $e_\beta$).
\item $B_n$ is 04's set of $\Sigma_n$-definable surreals; 07's set of
surreals of birthday at most $\alpha$ is written $\No_{\le\alpha}$ (07's
$B_\alpha$), its subsets of ordinals $X\subseteq\alpha$ (07's $B$), and
04's ordinal code of a set is $X_a$ (04's $B$).
\item 04's decoder of sets from signs is $\SetDec$ (04's $E$); $E_\alpha(X)$
is 07's bit code. In the set-coding lemma 04's index $\rho$ is written
$\varrho$, and in the forcing lemma 04's set $A\subseteq\theta$ is written
$Y$, since $\rho$ and $A$ are reserved.
\item $\lambda$ is a normal-form length and $\lambda_H$ the first
non-$\HOD$ birthday; an epsilon number is written $\epsilon$, and an
infinitesimal $\varepsilon$ (07 used $\lambda$ for epsilon numbers and for
other ordinals).
\item $\omega^x=\Omega(x)$ is the Conway map (07 wrote a bold $\omega$ for
the surreal $\omega$); $\underline\alpha$ is the surreal ordinal $\alpha$
when ordinal and surreal arithmetic might be confused.
\item $\prec$ is elementary substructure; $\ps$ is the proper-prefix
(simplicity) relation.
\end{itemize}
Other reports use some of these letters differently: the automorphisms
report writes $\mathcal U_{\Oz}$ for a group, whereas $\Cell$ here is an
interval of $\Oz$; the universes report writes $\beta(M,N)$ for the first
new birthday, which is $\lambda_H$ when $M=\HOD$ and $N=V$.

\subsection{Repository snapshot and dependency boundary}

Both manuscripts inspected the repository at commit \texttt{9a385d3} through a
GitHub connector: 04 read the catalogue and the introductions to the
birthday-cutoff and omnific Diophantine reports \cite{Repo,RepoHset,RepoOdg};
07 read the README, the catalogue and the detailed guides of the universes,
birthday and omnific Diophantine reports \cite{RepoUniv}. The catalogue
describes AI-assisted research drafts and separates written arguments from
Lean coverage. This report follows that distinction. It imports no
nonclassical Diophantine, quotient, automorphism, or bi-interpretation
claim as a proof dependency; the relations to other reports stated below
are credits and cross-references, checked at the placement commit.
\Cref{dsn:app:inspected} records the limited audit scope of each source.

\section{Foundations and canonical representations}
\label{dsn:sec:foundations}

\subsection{Universes, models, and class notation}

Our mathematical statements are made in ZFC, with proper-class notation
understood schematically. Each individual surreal is a \emph{set}.
A class function used below is specified by one set-theoretic formula;
all its values and all sums at which it is evaluated are set-sized.
An $\Ord$-indexed independent family means a definable class assignment
whose every finite subfamily is algebraically independent. It does not
mean that the entire family is a set or that its class sum exists.

When a literal satisfaction relation is needed, fix a set model
$M\models\ZFC$ in an ambient metatheory, and interpret $\No^M$, $\Oz^M$,
and the relevant definitions internally. This conditional framework does
not assert that ZFC proves the existence of such an $M$. The model-theoretic
statements explicitly identified as external apply to arbitrary set
models. For nonstandard models, their ``ordinals,'' ``finite sequences,''
and surreal codes are internal objects, not necessarily actual well-founded
sequences. Where a proof uses a least \emph{externally} chosen ordinal
(\cref{dsn:sec:core}) or external o-minimality (\cref{dsn:thm:elementary}),
$M$ is assumed transitive, as in source 07; then its ordinals and
finite sequences are genuine, and a surreal of $M$ is literally an
external sign sequence, although $M$ may have only some of the subsets of
a given ordinal.

\subsection{Signs, simplicity, and cuts}

A canonical surreal is a function
\[
 s:\alpha\longrightarrow\{-,+\},\qquad\alpha\in\Ord.
\]
Its birthday is $\bd(s)=\alpha$. Comparison is by the first disagreement,
with a missing sign ordered between $-$ and $+$. The ordinal $\alpha$
itself is represented by the all-plus sequence of length $\alpha$, written
$\underline\alpha$ when needed. The standard real embedding and the
arithmetic operations are parameter-free definable in set theory
\cite{Conway,Gonshor}. We write $y\ps x$ when $y$ is a proper prefix of
$x$, and $y\pseq x$ when equality is allowed.

A subclass $F\subseteq\No$ is \emph{initial} if it contains every prefix of
each of its elements. This is not the same as being an initial segment of
the numerical order. The notation $\No_{<\lambda}$ always refers to
birthdays, not numerical magnitude.

For sets $L,R$ of surreal numbers with $L<R$, there is a unique simplest
separator $\cut LR$. It is a prefix of every separator of the cut. Both the
restriction that $L,R$ be sets in the universe of construction and the
word ``simplest'' are essential. These standard facts are part of the
Conway--Gonshor simplicity theory; see also \cite{EhrlichKaplan}.

\subsection{Normal forms and two different exponentials}\label{dsn:sub:normal}

Every nonzero surreal has a unique Conway normal form
\begin{equation}\label{dsn:eq:normal}
 x=\sum_{\xi<\lambda}r_\xi\omega^{g_\xi},
 \qquad r_\xi\in\R\setminus\{0\},\quad
 g_\xi>g_\eta\ \ (\xi<\eta<\lambda).
\end{equation}
The length $\lambda$ is an ordinal, and the support
$\supp(x)=\{g_\xi:\xi<\lambda\}$ is reverse well-ordered and set-sized.
Conversely, any such coefficient/exponent sequence has a surreal sum.
Addition and multiplication are Hahn addition and convolution. We use
$\NF(x)$ for the \emph{whole indexed sequence}
$\xi\mapsto(g_\xi,r_\xi)$, including its length, not merely the
collection of its individual entries; for zero it is the empty function.

The uniqueness theorem supplies parameter-free set-theoretic operations
$x\mapsto\NF(x)$ and evaluation of valid normal forms. It supplies no
algorithm for arbitrary infinite codes and no assertion of topological
convergence of finite partial sums. The factorisation paper
\cite[Sections 1--2]{LM} is a convenient modern reference for the
normal-form convention and the distinction between series fields and
omnific subrings.

We write $\Omega(a)=\omega^a$ for the Conway omega map. It is an
increasing isomorphism from the additive group of surreals to the
multiplicative group of canonical monomials. Distinct exponents give
distinct Archimedean scales: if $a<b$, then $n\omega^a<\omega^b$ for every
positive ordinary integer $n$. The inverse $\Omega^{-1}$ is defined on
monomials. Gonshor's function $\exp$ is a different function, extending
ordinary real exponentiation with base $e$. Throughout, $\Lexp$ means the
ordered-field language with this $\exp$, \emph{not} with $\Omega$, and
$\Lor$ the ordered-field language. In general one must not replace
$\omega^a$ by $\exp(a\log\omega)$. All coding constructions below use the
Conway omega map.

\begin{lemma}[Ordinal domination \src{04}]\label{dsn:lem:ordinalbound}
Every set $S\subseteq\No$ is strictly bounded in absolute value by a
surreal ordinal. Moreover, the least ordinal $\delta$ satisfying
$|s|<\delta$ for every $s\in S$ is definable from $S$.
\end{lemma}
\begin{proof}
A sign sequence of length $\alpha$ lies between the all-minus and
all-plus sequences of length $\alpha$, hence $|s|\leq\bd(s)$.
Replacement gives the set of birthdays. An ordinal strictly exceeding
all of them gives the required bound. The qualifying ordinals form a
nonempty definable class, so it has a least member, uniquely specified
from $S$.
\end{proof}

\subsection{Omnific integers and the correct floor formula}

\begin{definition}\label{dsn:def:Oz}
An \emph{omnific integer} is a surreal whose normal form has no negative
exponents and whose coefficient at exponent $0$ is an ordinary integer.
Thus
\begin{equation}\label{dsn:eq:Oz}
 \Oz=\Pi\oplus\Z,\qquad
 \Pi=\{x\in\No:\supp(x)\subseteq\No_{>0}\}.
\end{equation}
The empty support is permitted, so $0\in\Pi$. The map
$\ct:\Oz\to\Z$ takes the coefficient at exponent $0$.
\end{definition}

This is the standard ring, not the ring generated by the ordinals.
For example, $\sqrt2\,\omega$, $\omega^{1/2}$ and $\omega+17$ are omnific,
but $\omega+\pi$, $\omega+1/2$ and $\omega^{-1}$ are not. A positive
omnific integer may be much smaller than $\omega$ while exceeding every
ordinary integer. The direct sum in \eqref{dsn:eq:Oz} is additive, not a
direct product of unital rings. Positive exponents stay positive under
addition, so $\Pi$ is an ideal and $\ct$ is a ring retraction.

\begin{proposition}[Integer part \src{04, 07}]\label{dsn:prop:floor}
The ring $\Oz$ is discretely ordered with least positive element $1$.
For every $x\in\No$ there is exactly one $z\in\Oz$ such that
$z\leq x<z+1$. This $z=\lfloor x\rfloor_{\Oz}$ is parameter-free
definable from $x$ in set theory, and $\lfloor x\rfloor_{\Oz}\pseq x$.
\end{proposition}
\begin{proof}
Write, by splitting the normal form,
$x=P+r+\varepsilon$, where $P\in\Pi$, $r\in\R$, and
$\varepsilon$ is infinitesimal. Put
\begin{equation}\label{dsn:eq:floor}
 \lfloor x\rfloor_{\Oz}=P+
 \begin{cases}
 r-1,&r\in\Z\ \text{and}\ \varepsilon<0,\\
 \lfloor r\rfloor,&\text{otherwise}.
 \end{cases}
\end{equation}
A nonzero positive element of $\Pi$ is larger than every real, since
its leading exponent is positive. Hence every positive element of
$\Oz$ is at least $1$. Formula~\eqref{dsn:eq:floor} directly verifies the
two inequalities. If two omnific integers had the required property,
their nonzero difference would have absolute value less than $1$,
contradicting discreteness. Both the normal-form formula and the unique
integer-part property are set-theoretic definitions.

For the prefix assertion (07), use the standard simplicity
characterization $z=\cut{z-1}{z+1}$ for $z\in\Oz$. The inequalities place
$x$ strictly between $z-1$ and $z+1$, and the simplest separator is a
prefix of every separator.
\end{proof}

\begin{remark}
The correction when $r$ is integral and $\varepsilon<0$ is essential:
$\lfloor\omega-\omega^{-1}\rfloor_{\Oz}=\omega-1$, not $\omega$.
The proposition is classical \cite{Conway,LM}; its proof is included
to make subsequent definability closure entirely explicit. The floor and
its correction are also treated in the omnific Diophantine report
\cite{RepoOdg}.
\end{remark}

\subsection{The imported exponential-field theorem}\label{dsn:sub:imported}

Source 07 uses the following classical input. The surreal exponential
field is an elementary extension of the real exponential field, and it is
o-minimal. If $\epsilon$ is an epsilon number, then $\No_{<\epsilon}$,
with the induced exponential, is an elementary substructure of $\No$ and
is closed under logarithms of positive elements. These are consequences of
\cite[Theorem 2.1 and Corollary 5.5]{vdDE}; the latter is stated even in
the richer restricted-analytic exponential language. The original paper
must be read with its erratum \cite{vdDEerr}, as the birthday-cutoff report
insists; Bournez and Guilmant prove that $\No_\epsilon$ is stable under
exponential and logarithm exactly when $\epsilon$ is an epsilon number
\cite{BG}, which is the closure input used in \cref{dsn:thm:core}, and
elementarity then follows from model completeness of the real exponential
field with logarithm, which is what \cite[Corollary 5.5]{vdDE} packages.
Source 07 cited the original paper without the erratum; the merge adds it.
Our conclusions about definable fragments use only the countable language
$\Lexp$.

\begin{remark}[o-minimality inside a transitive model \mergetag]
\label{dsn:rem:ominimal}
Source 07 uses tacitly that $\No^M_{\exp}$ is o-minimal as an external
structure when $M$ is transitive. This holds formula by formula. The
o-minimality of the real exponential field and the elementary equivalence
of the surreal and real exponential fields are theorems of ZFC, and by
uniform finiteness each formula $\varphi(x,\bar y)$ has a standard bound
$N_\varphi$ such that ZFC proves that every fibre $\varphi(\No_{\exp},\bar
b)$ is a union of at most $N_\varphi$ points and intervals. Applied inside
$M$, this statement is first order with a standard numeral, so it holds
externally of $\No^M_{\exp}$.
\end{remark}

\begin{lemma}[Classical birthday controls \src{07}]\label{dsn:lem:birthday-controls}
If $r\omega^a$ is a nonzero term in the normal form of $x$, then
\[
 \bd(a)\le\bd(x),\qquad \bd(r\omega^a)\le\bd(x).
\]
The normal-form length is at most $\bd(x)$, and for every surreal $a$,
\[
 \bd(a)\le\bd(\Omega(a))\le\omega^{\bd(a)},
\]
where the power on the right is ordinal exponentiation.
\end{lemma}
These are the estimates of \cite[Lemmas 4.1--4.2 and the proof of
Proposition 4.7]{vdDE}, imported as established results rather than
reproved from the sign-expansion analysis. Independently, the sign order
immediately gives $|a|\le\underline{\bd(a)}$, because the all-plus and
all-minus sequences are the extremal sequences of a fixed length. For the
merge, the upper bound $\bd(\Omega(a))\le\omega^{\bd(a)}$ was re-derived
from Gonshor's sign expansion of $\omega^a$; the exponent bound
$\bd(a)\le\bd(x)$ was not re-checked, and
\cref{dsn:thm:core-fractions} uses only its consequence $a\ge-\underline{\bd(x)}$.

\section{Which notion of definability?}
\label{dsn:sec:definitions}

\subsection{Ambient definability with parameters}

\begin{definition}[\src{04, 07}]\label{dsn:def:ambient}
Let $M\models\ZFC$ be a set model, and let $A$ be an external collection
of allowed elements of $M$. An element $a\in M$ is \emph{ambiently
$A$-definable in $M$} if some standard first-order formula
$\varphi(v,\bar u)$ and some finite tuple $\bar p$ from $A$ satisfy
\[
 M\models\forall v\,[\varphi(v,\bar p)\leftrightarrow v=a].
\]
Write $D_A^M=\dcl_{(M,\in)}(A)$, and define
\[
 K_A^M=D_A^M\cap\No^M,\qquad I_A^M=D_A^M\cap\Oz^M,\qquad
 O_A^M=D_A^M\cap\Ord^M.
\]
When $A=\varnothing$, say \emph{parameter-free definable}. For a finite
tuple $p$ write $K^M_p$. Superscripts are omitted when the universe is
fixed.
\end{definition}

For a set model, $D_A^M$ is an external set, definable using the external
satisfaction relation of $M$. In particular,
\begin{equation}\label{dsn:eq:count}
 |D_A^M|\leq\max(\aleph_0,|A|).
\end{equation}
This estimate counts standard formulas and finite parameter tuples. It is
external: it does not put into $M$ a list of all the elements definable in
$M$. Even though $\No^M$ and $K^M_A$ are sets externally, $\No^M$ is a
proper class internally, and $K^M_A$ need not be represented by a set of
$M$.

For the intended universe $V$, the same terminology is a metatheoretic
scheme: for each particular formula we can ask whether it uniquely
defines a given object. We do not add a first-order truth predicate for
$V$, and do not assume that the full collection $D_\varnothing^V$ is
an internally available set or class.

\begin{lemma}[Unique-output closure \src{04, 07}]\label{dsn:lem:closure}
Suppose a set-theoretic formula defines a unique output $F(\bar x)$
from each input in its domain. If every component of $\bar a$ lies
in $D_A$, then $F(\bar a)\in D_A$.
\end{lemma}
\begin{proof}
Choose definitions for the finitely many components of $\bar a$.
Existentially quantify those components, assert their defining formulas,
and assert the graph formula for $F$. This is a single formula using
only finitely many parameters from $A$, with a unique solution.
\end{proof}

This elementary lemma is the mechanism behind all positive closure results
in the report. It does not assert that $D^M_A\prec M$: general
set-theoretic existential statements need not have a definable witness
(07).

\subsection{Ordinal parameters}

\begin{definition}[\src{04, 07}]\label{dsn:def:OD}
A surreal or omnific integer is \emph{ordinal definable} if its canonical
set code is in $\OD$. Define
\[
 \KH=\No\cap\OD,\qquad \IH=\Oz\cap\OD.
\]
The standard internal definition of $\OD$ says that an object is uniquely
definable over some set structure $(V_\theta,\in)$ using finitely many
ordinal parameters below $\theta$. Reflection identifies this with the
usual notion of definability in $V$ from ordinal parameters. The extra
ordinal $\theta$ is allowed; no truth predicate for $V$ is required.
Inside a model $M$ the same definition gives $\KH^M=\No^M\cap\OD^M$.
\end{definition}

Both $\OD$ and $\HOD$ are definable classes of ZFC, unlike a uniform
truth-based relation between arbitrary formulas and their unique
solutions in $V$. $\HOD$ is the transitive class of sets whose transitive
closures, including the sets themselves, consist of ordinal-definable
sets; it is an inner model of ZFC with a canonical definable well-order
\cite{Jech,GoldbergPoveda}. Every ordinal is ordinal definable, using
itself as an allowed parameter. This alone shows why $\KH$ need not
be a countable collection.

\subsection{Intrinsic ordered-field definability}

For a specified language $\mathcal L$ on $\No^M$, \emph{intrinsic
definability} means $\dcl_{(\No^M,\mathcal L)}(P)$ for a set $P$ of surreal
parameters. The quantifiers then range over surreals, and only the symbols
of $\mathcal L$ may be used. The three notions --- ambient, ordinal and
intrinsic --- should never be conflated. The set-theoretic definition of
$\omega$ is parameter-free, but $\omega$ is not definable without
parameters in the pure ordered field, or even in its exponential reduct
(\cref{dsn:prop:exp-no-omega}). Conversely, allowing every surreal as a
parameter trivializes definability of individual surreals.

\begin{theorem}[Classical language contrast \src{04, 07}]\label{dsn:thm:purefield}
In the pure ordered-field language, the parameter-free definable
surreals are exactly the real algebraic numbers. Their intersection
with $\Oz$ is exactly $\Z$. More generally, in a real closed surreal
field $F$ (for example $F=\No^M$) and for any external parameter set
$P\subseteq F$,
\[
 \dcl_F(P)=\RC_F\bigl(\Q(P)\bigr),
\]
the relative real algebraic closure of $\Q(P)$ in $F$.
\end{theorem}
\begin{proof}
Quantifier elimination for real closed fields reduces a definable
singleton to a Boolean combination of polynomial sign conditions over the
field generated by finitely many parameters \cite{Marker,vddTame}. If a
point is transcendental over that field, it is not a zero of any of the
finitely many nonzero polynomials appearing in the formula, so their signs
are constant on a small interval around it and the formula cannot isolate
that point. Thus the point is algebraic over $\Q(P)$. Conversely, an
algebraic element is specified as the $k$th real root of its polynomial,
for a fixed finite $k$, with coefficients given by field terms in finitely
many parameters. Real algebraic elements of $\No$ lie in its standard real
subfield. A real is omnific exactly when it is an integer, by
\eqref{dsn:eq:Oz}.
\end{proof}

For the full proper-class field, the argument is a formula-by-formula
application of the same quantifier-elimination theorem. In particular,
$\omega$ is \emph{ambiently} parameter-free definable, but is not
parameter-free definable in the pure field. Thus ``the parameter-free
definable omnific integers'' can mean just $\Z$ when definability is taken
in the ordered field, but can include $\omega$, $\omega^{1/2}$ and
arbitrarily elaborate bounded codes when taken in set theory; both
assertions are correct in their respective languages. Nor does adding the
word ``omnific'' to our prose add a predicate to that language. The
structure $(\No;+,-,\cdot,<,0,1,\Oz)$ is a different structure; what the
collection now proves about its definable closure is in
\cref{dsn:prop:pair}.

\subsection{Provable descriptions and computability}

For a recursively axiomatized theory $T$, one can restrict to formulas
$\varphi$ for which $T\vdash\exists!x\,(\No(x)\wedge\varphi(x))$.
These are \emph{$T$-provably unique descriptions}, or provable names.
Their interpretation in a model of $T$ need not be invariant between
different models. For example the formula which selects $0$ if CH holds
and $1$ otherwise has a provably unique output in ZFC, but its value
depends on the model (07). Transitivity, syntactic shortness, and provable
uniqueness are separate issues from absoluteness; the statement ``$T$
proves uniqueness'' is not a replacement for choosing the semantic
universe.

A genuine effective presentation of a surreal sign sequence gives an
ambient definition when the presentation uniquely determines a
well-founded sequence. The converse fails already for real numbers:
the characteristic real of a noncomputable but arithmetically definable
set is definable. We therefore do not identify definability with any
one of the computable-surreal representations catalogued by the
collection, or with a numerical approximation algorithm, and we provide no
algorithm deciding whether an arbitrary formula uniquely defines a
surreal. When a construction is called uniform, its defining set-theoretic
formula is fixed; no internal truth predicate is thereby supplied.

\section{The field and integer part of ambiently definable numbers}
\label{dsn:sec:hulls}

Unless stated otherwise the statements of this section hold for $V$ as a
scheme and for any set model $M$; the superscript $M$ refers to such a
model, whose internal operations on surreals and normal forms are used.
Put $\Pi_A=K_A\cap\Pi$ and $R_A=K_A\cap\R$.

\begin{theorem}[Definable field and integer part \src{04, 07}]\label{dsn:thm:hull}
The ambiently $A$-definable surreals $K_A$ form a real closed field. It is
closed under $\Omega$, under $\Omega^{-1}$ on monomials, under $\exp$ and
positive logarithms, and under the omnific floor. The ring $I_A=K_A\cap\Oz$
is a discretely ordered integer part of $K_A$, and
\begin{equation}\label{dsn:eq:frac}
 K_A=\Frac(I_A).
\end{equation}
More precisely, every nonzero $x\in K_A$ can be written $x=p/q$ with
$p,q\in\Pi_A$ and $q>0$ a monomial, the pair $(p,q)$ being chosen by one
parameter-free definable rule from $x$; two such rules are
\begin{equation}\label{dsn:eq:denominators}
 q=\omega^{\delta}\ \ \text{(04)},\qquad
 q=\omega^{b},\ b=\cut{\{0\}\cup\{-a:a\in\supp x\}}{\varnothing}\ \ \text{(07)},
\end{equation}
where $\delta$ is the least positive ordinal with $\delta+g>0$ for all
$g\in\supp x$. The field $K_A$ is an elementary substructure of the ambient
surreal ordered field in the pure ordered-field language.
\end{theorem}
\begin{proof}
The constants $0,1$, addition, subtraction, multiplication, inverse away
from zero, $\Omega$, $\Omega^{-1}$, $\exp$ and $\log$ are uniquely definable
set-theoretic operations, so \cref{dsn:lem:closure} proves the field
axioms and the closure assertions by restriction. Every positive element
has a unique positive square root in $\No$, so that square root belongs to
$K_A$. For an odd-degree polynomial with coefficients in $K_A$, real
closedness of $\No$ gives at least one root; its finite set of roots has a
least member, uniquely definable from the coefficient tuple. That root lies
in $K_A$. The square-root and odd-degree-root criterion proves real
closedness. This is an application of the classical real-closed-field
theorem, not a new proof of real closedness of $\No$.

The floor from \cref{dsn:prop:floor} belongs to $D_A$ whenever its input
does. Hence every $x\in K_A$ has its omnific floor in $I_A$. Discreteness
and the subring property restrict from $\Oz$.

For the fraction-field assertion, let $0\ne x\in K_A$, and put
$S=\supp(x)$. The whole normal form and $S$ are definable from $x$.
By \cref{dsn:lem:ordinalbound} there exists a positive ordinal $\delta$
with $\delta+g>0$ for $g\in S$. Choose the \emph{least} such positive
ordinal (04). It is definable from $x$, so $\delta\in K_A$, and
$q=\omega^\delta$ and $p=qx$ belong to $K_A$. Alternatively (07), form the
set cut $b$ in \eqref{dsn:eq:denominators}; then $b>0$ and $b+a>0$ for
every $a\in S$, and $b$ is definable from $x$. In either case the normal
form of $p=qx$ is obtained by adding the exponent of $q$ to every exponent
of $x$, without changing the coefficients. Its exponents are strictly
positive, so $p,q\in\Pi_A$, and $x=p/q$. The zero case is $0=0/1$. The
reverse inclusion follows because $K_A$ is a field. Finally, every
inclusion of real closed fields is elementary by quantifier elimination
\cite{Marker}.
\end{proof}

\begin{remark}
The proof does not suppose that each exponent of an $A$-definable
normal form is $A$-definable. It only uses definability of the
\emph{entire support} to choose one bound. Confusing these
two assertions would invalidate the fraction-field argument.
The equality $\Frac(\Oz)=\No$ and the support-clearing method are not
new: they are \lbl{odg:thm:fractions} in the omnific Diophantine report
\cite{RepoOdg}. The point here is that the denominator can be chosen
canonically enough to preserve the \emph{same} allowed parameters; merely
invoking an unspecified denominator would not establish this conclusion.
\end{remark}

Source 07 states the closure under $\Omega$, $\exp$ and the floor, and the
cut denominator, for a transitive model; the proofs use only
\cref{dsn:lem:closure}, so they hold in the generality of 04's statement.
The same applies to \cref{dsn:prop:nf-invariant,dsn:thm:valuation} below.

\begin{lemma}[Definable selection in an o-minimal field \src{07}]\label{dsn:lem:selection}
Let $F$ be an o-minimal expansion of an ordered field. For each formula
$\varphi(y,\bar x)$ there is a parameter-free definable partial function
$s_\varphi(\bar x)$ whose domain is exactly
\[
 \{\bar a:F\models\exists y\,\varphi(y,\bar a)\}
\]
and whose value satisfies $\varphi$. A substructure closed under these
functions is elementary in $F$. Intersections of substructures closed under
all of them are elementary as well.
\end{lemma}
\begin{proof}
A nonempty definable subset of the line is a finite union of points and
intervals. Its leftmost component and its finite endpoints are definable
from its parameters. Select its unique point if it is a point. For an
interval choose its left endpoint if that endpoint belongs to it;
otherwise use the midpoint when both endpoints are finite, the finite
endpoint minus or plus $1$ in the one-sided unbounded cases, and $0$ for
the whole line. The alternatives and endpoints are uniformly definable for
the given formula. Thus this is a definable selection, not a choice made
separately with new parameters for each fibre.

Apply the Tarski--Vaught criterion to the selected witness. One-coordinate
selection suffices for all existential formulas; tuples can be handled
successively. Closure of an intersection follows because the same
single-valued selectors are used in every substructure.
\end{proof}

\begin{theorem}[Exponential elementarity \src{07}]\label{dsn:thm:elementary}
Let $M$ be a transitive set model of ZFC. In the ordered exponential
language,
\[
 (K^M_A,+,\cdot,<,\exp)\prec(\No^M,+,\cdot,<,\exp).
\]
In particular the same holds in the ordered-field reduct.
\end{theorem}
\begin{proof}
By \cref{dsn:sub:imported} and \cref{dsn:rem:ominimal} the right-hand
structure is o-minimal. Each fixed $\Lexp$-formula translates into a fixed
formula of set theory: restrict its quantifiers to $\No$, and replace its
function symbols by their definable graphs. Thus a selector of
\cref{dsn:lem:selection}, evaluated at a tuple from $K^M_A$, is an ambient
uniquely definable output. It lies in $K^M_A$ by \cref{dsn:lem:closure}.
The selection lemma gives elementarity. This translation is performed for
each standard formula separately. The proof never combines all
translations into a truth predicate inside $M$.
\end{proof}

\begin{remark}
A language with a symbol for every restricted real analytic function can
contain nondefinable real information in its symbols. The countability and
ambient-definability statements here concern $\Lexp$, not that arbitrary
uncountable expansion (07).
\end{remark}

\begin{corollary}[Open induction \src{04}]\label{dsn:cor:openinduction}
The nonnegative part of $I_A$ satisfies induction for every
quantifier-free formula in the ordered-ring language, with parameters
from $I_A$.
\end{corollary}
\begin{proof}
A quantifier-free formula with parameters in $I_A$ defines a
semialgebraic subset of the real closed field $K_A$, hence a finite
union of intervals and points. Every nonempty intersection of such a
set with $(I_A)_{\geq0}$ has a least member. Indeed, the integer-part
property supplies the first integer at or immediately above each
finite left endpoint, with a one-unit adjustment for a strict endpoint;
an interval unbounded to the left is treated by the lower bound $0$.
Take the least of the finitely many nonempty pieces.
If an induction instance failed, the set of counterexamples would
therefore have a least member. It is not $0$, and its predecessor is
nonnegative and satisfies the formula, contradicting the inductive
step. This is the standard integer-part route to open induction
\cite{Shepherdson,LM}.
\end{proof}

For the full ring $\Oz$ open induction, its failure beyond open formulas,
and the failure of Peano arithmetic are in the omnific Diophantine report
(\lbl{odg:rem:foursquares}, \lbl{odg:thm:induction}); the corollary is
the definable-fragment instance.

Related (batch 86, added 3 October 2026): open induction is sharp here. By
\cref{dsn:prop:ct}, $I_A$ is a discretely ordered ring with division with
remainder by every ordinary integer and with the ring retraction $\ct$ onto
$\Z$. The discrete-initial-subgroups report (its source~03; AI-assisted,
unrefereed, not formalized) shows that such a ring satisfies the single
bounded induction instance of its dyadic standard-cut formula
$\mathsf{Std}^{\mathrm{dy}}$ only if it is $\Z$ (\lbl{isg:sc:thm:boundary});
so the nonnegative part of $I_A$ satisfies open induction but not
$\mathsf{I}\Delta_0$, as for $\Oz$. (That report proves its theorems for
set-sized rings and reads class applications as in
\lbl{isg:sc:sec:classes}; the same reading applies to $I_A$.) The
\fn{polish-models-of-omnific-arithmetic} report proves open induction by the
route of the corollary above for two finite principal-part integer parts:
$\Z\oplus\bigoplus_{q>0}\R X^q$, in a Borel presentation
(\lbl{pma:bor:thm:arithmetic}), and its countable subring with real
algebraic coefficients, inside a Polish ordered field
(\lbl{pma:ser:prop:iopen}).

\begin{proposition}[Finite ambiguity does not enlarge the hull \src{04}]
\label{dsn:prop:algebraicity}
An element of $\No$ belonging to a finite ambiently $A$-definable set
of surreals is itself ambiently $A$-definable. In particular, ambient
model-theoretic algebraicity and definability coincide on the surreal
sort. Every surreal algebraic over $K_A$ already belongs to $K_A$.
\end{proposition}
\begin{proof}
Order the finite candidate set by the surreal order. Each member is
its $k$th member for a fixed finite $k$; this property is expressible
by a formula with the same parameters. For the final assertion, the
roots of a nonzero polynomial over $K_A$ form such a finite definable
set. Equivalently, a real closed field has no proper algebraic ordered
extension.
\end{proof}

This finite-ordering argument is a particularly transparent case of
model-theoretic algebraicity; broader set-theoretic questions about
algebraicity and implicit definability are studied in
Hamkins--Leahy \cite{HamkinsLeahy}. It concerns a genuinely finite
candidate set, not an internally nonstandard finite set viewed from
outside a nonstandard model.

\begin{proposition}[Whole-family closure \src{04, 07}]\label{dsn:prop:uniform}
If $L,R\in D_A$ are sets of surreals with $L<R$, their simplest separator
belongs to $K_A$. If an entire set-coded, strongly summable family of
surreals belongs to $D_A$ --- in particular a valid normal form whose
\emph{whole indexed sequence} is $A$-definable --- its canonical sum
belongs to $K_A$.
\end{proposition}
\begin{proof}
In both cases the output is a unique parameter-free definable function of
the input set codes; apply \cref{dsn:lem:closure}. For a cut, both option
sets must be jointly available as $A$-definable sets, not merely have
individually $A$-definable elements. Strong summability is the
normal-form condition that the union of supports is reverse well ordered
and each exponent receives only finitely many nonzero contributions.
\end{proof}

The hypothesis concerns the \emph{whole input code}. An arbitrary set of
individually definable entries need not satisfy it
(\cref{dsn:sec:uniformity}). This distinction is central, not a technical
qualification at the edge of the theory.

\begin{proposition}[The normal-form invariant \src{07}]\label{dsn:prop:nf-invariant}
For every $x\in\No^M$,
\[
 x\in K^M_A\quad\Longleftrightarrow\quad\NF(x)\in D^M_A.
\]
If $x\in K^M_A$ and $\xi<\dom\NF(x)$ is an $A$-definable ordinal, then
$g_\xi$ and $r_\xi$ are $A$-definable.
\end{proposition}
\begin{proof}
The normal-form map and its inverse are unique parameter-free definable
constructions. Evaluation of a set-coded function at an input is also
uniquely definable. Apply \cref{dsn:lem:closure} in both directions.
\end{proof}

Without the definability of the index, the final assertion need not hold
(\cref{dsn:prop:collector,dsn:thm:reverse-gap}). Without the definability
of the \emph{whole} graph, its converse need not hold either
(\cref{dsn:thm:termwise,dsn:thm:bit-code}).

\subsection{Value group and residue field}

Define $v(x)=-g_0$ for nonzero $x$ in \eqref{dsn:eq:normal}, and
$v(0)=\infty$. This is the normal-form valuation; it uses the negative of
the largest Conway exponent. Let
\[
 \mathcal O_A=\{x\in K_A:|x|<n\text{ for some }n\in\N_{>0}\},\quad
 \mathfrak m_A=\{x\in K_A:|x|<1/n\text{ for every }n\in\N_{>0}\}.
\]
Here the quantification over $\N$ is ambient set-theoretic, not a claim
that $\N$ is intrinsically definable in an o-minimal field.

\begin{theorem}[Value group and residue field \src{07}]\label{dsn:thm:valuation}
The value group of $K_A$ is canonically the ordered additive group
$(K_A,+)$: $v(K_A^\times)=K_A$. Standard part induces an isomorphism
$\mathcal O_A/\mathfrak m_A\simeq R_A$. Every finite $x\in K_A$ has a
unique decomposition $x=r+\varepsilon$ with $r\in R_A$ and
$\varepsilon\in\mathfrak m_A$.
\end{theorem}
\begin{proof}
The leading exponent is a unique definable function of a nonzero normal
form, so $v(x)\in K_A$. Conversely $v(\omega^{-a})=a$ for every $a\in K_A$,
and the monomial lies in $K_A$.

For finite $x$, its real constant term is its standard part; this remains
true when there is no constant term, in which case it is zero. This unique
real is definable from $x$, hence belongs to $R_A$. The ordinary
normal-form multiplication and addition laws show that standard part is a
surjective ring homomorphism with kernel exactly the infinitesimals.
Uniqueness of the decomposition follows because no nonzero real is
infinitesimal.
\end{proof}

This self-similarity of the value group should not be mistaken for a claim
that $K_A$ contains every Hahn series whose entries belong to $K_A$;
\cref{dsn:sec:uniformity} gives explicit obstructions.

\section{The ordinal-definable surreal field is the HOD field}
\label{dsn:sec:HOD}

The canonical sign representation makes hereditary definability much
more transparent than it is for arbitrary sets. In this section $V$ is the
universe of construction; every statement can equally be read inside a
model, with $V$, $\OD$ and $\HOD$ relativized.

\begin{theorem}[HOD identification \src{04, 07}]\label{dsn:thm:HOD}
With canonical sign-sequence representatives,
\begin{equation}\label{dsn:eq:HODid}
 \KH=\No\cap\OD=\No\cap\HOD=\No^{\HOD},
 \qquad
 \IH=\Oz\cap\OD=\Oz\cap\HOD=\Oz^{\HOD}.
\end{equation}
These are equalities of sign-sequence classes. The class $\KH$ is an
initial subclass of $\No$ and a real closed field that contains every
ordinal and is closed under $\Omega$, $\Omega^{-1}$ on monomials, $\exp$,
positive logarithms and the omnific floor. The ring $\IH$ is an integer
part with $\Frac(\IH)=\KH$.
\end{theorem}
\begin{proof}
Let $s:\alpha\to\{-,+\}$ be ordinal definable. Every element in its
transitive closure, other than $s$ itself, is built from ordinals and the
two fixed sign symbols by finitely many fixed set constructions, so it is
ordinal definable. Since $s$ itself is ordinal definable, $s\in\HOD$.
The converse is part of the definition of $\HOD$. $\HOD$ contains every
ordinal and is transitive, so being an ordinal-length sign sequence is
absolute, and its surreal sign sequences are exactly the ambient sign
sequences that belong to $\HOD$. This gives the carrier equalities for
$\No$ in \eqref{dsn:eq:HODid}.

The closure assertions follow by substituting ordinal definitions into the
unique-output constructions of \cref{dsn:thm:hull}; the ordinals used are
permitted parameters. The floor and the least-ordinal denominator give the
integer-part and fraction claims for $\IH=\Oz\cap\OD$. Initiality: a prefix
$x\res\beta$ is definable from $x$ and the ordinal $\beta$ (07).

It remains to compare with the internal arithmetic of $\HOD$. Comparison is
computed from the same signs; canonical options are the same initial
segments; and the recursive arithmetic on set-sized option closures is
absolute between a transitive inner model and the ambient universe, so the
ordered-field structure of $\No^{\HOD}$ is the restriction of that of $\No$
(this part is \lbl{univ:prop:absolute} in the universes report, for any
same-ordinal inner model). For the monomial map one may use its recursive
cut with integer and dyadic multipliers, so no new reals enter the
recursion (04); source 07 asserts the same compatibility for the normal
form, the exponential and the purely infinite/constant splitting. The
universes report proves the arithmetic and cut parts only, so for these
further operations the recursion argument just sketched is the source.
For the omnific carriers, an $\OD$ surreal has an $\OD$ normal-form graph
whose entries at ordinal indices are $\OD$, hence the whole graph is in
$\HOD$ (\cref{dsn:lem:HODsupport}) and its canonical interpretation agrees
in the two universes; this gives $\Oz\cap\HOD=\Oz^{\HOD}$ (07).
\end{proof}

When $V$ is replaced by a transitive model $M$, \cref{dsn:thm:elementary}
with all ordinals of $M$ as the external parameter set also shows that
$\KH^M$ is an elementary exponential subfield of $\No^M$ (07). This is
stronger than merely being closed under $\exp$.

\begin{lemma}[Support of an ordinal-definable number \src{04, 07}]
\label{dsn:lem:HODsupport}
If $x\in\KH$ has normal form \eqref{dsn:eq:normal}, then for every
$\xi<\lambda$,
\[
 r_\xi\in\R\cap\HOD,\qquad g_\xi\in\KH.
\]
In fact, $\NF(x)\in\HOD$, and
\begin{equation}\label{dsn:eq:NFcriterion}
 x\in\KH\quad\Longleftrightarrow\quad\NF(x)\in\HOD.
\end{equation}
The converse with only ``every coefficient and exponent is $\OD$'' is false
whenever there is a non-$\OD$ subset of an ordinal.
\end{lemma}
\begin{proof}
The entire normal form is uniquely definable from $x$, so is ordinal
definable. An individual entry is uniquely definable from $x$ and the
ordinal index $\xi$. Its real coefficient is ordinal definable, hence
hereditarily ordinal definable under the usual real coding. Its
surreal exponent belongs to $\HOD$ by \cref{dsn:thm:HOD}. The length,
indices, and the fixed tuple coding are also hereditarily ordinal
definable. Hence the whole sequence belongs to $\HOD$. Conversely,
$\NF(x)\in\HOD$ implies that it is ordinal definable, so its unique
surreal value is ordinal definable. The counterexamples are
\cref{dsn:thm:termwise} and $E_\alpha(X)$ of \cref{dsn:thm:bit-code} for a
non-$\OD$ $X\subseteq\alpha$.
\end{proof}

The last equivalence requires the whole sequence to lie in $\HOD$. The
strictly weaker condition ``every coefficient and exponent lies in
$\HOD$'' is insufficient, even with a $\HOD$ support and rational
coefficients. In particular, $\KH$ is not being identified with the
unrestricted ambient Hahn field containing \emph{all} externally assembled
series with coefficients and exponents in $\HOD$.

\begin{corollary}[No algebraic defect \src{04}]\label{dsn:cor:noalgebraic}
Every $x\in\No\setminus\KH$ is transcendental over $\KH$. The
inclusion $\KH\subseteq\No$ is elementary in the ordered-field
language, whether or not $V=\HOD$.
\end{corollary}
\begin{proof}
Apply \cref{dsn:prop:algebraicity} with ordinal parameters,
and then quantifier elimination.
\end{proof}

\begin{proposition}[The HOD field is not intrinsically definable when proper \src{04}; proof corrected \mergetag]
\label{dsn:prop:notdefinable}
If $\KH\ne\No$, no ordered-field formula, even with finitely many
surreal parameters, defines $\KH$ inside $\No$.
\end{proposition}
\begin{proof}
It suffices to show that an infinite definable subfield $F$ is the whole
field; $\KH$ is infinite. An infinite definable subset of a real closed
field is a finite union of points and intervals, so $F$ contains an open
interval $(a,b)$. The endpoints need not lie in $F$, so choose
$a'<b'$ inside $(a,b)$; they belong to $F$, and so does $c=b'-a'>0$. For
$0<y<c$ the point $a'+y$ lies in $(a',b')\subseteq F$, so $y\in F$; hence
$(0,c)\subseteq F$. For any $x>0$, the number $y=1/(x+2/c)$ lies in
$(0,c)$, hence in $F$, and then $x=1/y-2/c\in F$. Negatives and zero
follow, so $F$ is the entire field.
\end{proof}

Source 04 took $c$ from ``differences'' of the interval without placing
the endpoints in the subfield; the choice of $a',b'$ inside the interval
closes that gap. The same fact, that a proper real closed subfield has
empty interior, is used in \cref{dsn:prop:hod-topology}.

\section{Reversible omnific codes}
\label{dsn:sec:code}

\subsection{A sign code in a fixed omnific interval \src{04}}

For each ordinal $\beta$, viewed as a surreal ordinal, define
\begin{equation}\label{dsn:eq:eps}
 e_\beta=\omega^{-(\beta+1)}.
\end{equation}
The $e_\beta$ form a strictly decreasing class of positive infinitesimal
surreals. In particular $0<e_\beta\leq\omega^{-1}<1/2$.
Every $e_\beta$ is ordinal definable. For a sign sequence
$s:\alpha\to\{-,+\}$ put
\[
 c_s(\beta)=\begin{cases}1,&s(\beta)=-,\\2,&s(\beta)=+.\end{cases}
\]

\begin{definition}[Self-delimiting omnific code \src{04}]\label{dsn:def:code}
Define
\begin{equation}\label{dsn:eq:code}
 \Code(s)=\omega+
       \sum_{\beta<\alpha}c_s(\beta)\omega^{e_\beta}
       +3\omega^{e_\alpha},\qquad\alpha=\bd(s).
\end{equation}
The coefficient $3$ is an end marker, not a third sign value.
\end{definition}

\begin{theorem}[Reversible fixed-leading-term coding \src{04}]\label{dsn:thm:code}
Formula~\eqref{dsn:eq:code} defines a parameter-free set-theoretic injection
$\Code:\No\to\Cell$. Its image is set-theoretically definable, and
there is a parameter-free definable inverse $\Decode$ on that image.
For every model $M$ and every allowed parameter family $A$,
\begin{equation}\label{dsn:eq:degree}
 s\in D_A^M\quad\Longleftrightarrow\quad\Code^M(s)\in D_A^M.
\end{equation}
In particular, $s$ and $\Code(s)$ are mutually definable and have the
same ordinal-definability status.
\end{theorem}
\begin{proof}
The exponents in \eqref{dsn:eq:code} occur in strictly decreasing order:
first $1$, then $e_\beta$ for $\beta\leq\alpha$. They form a set,
because $\alpha$ is an ordinal. Each coefficient is nonzero and real,
so the expression is a valid Conway normal form. Every exponent is
positive. Thus the value is an omnific integer with constant term zero.
Its leading term is exactly $\omega$.

The tail $\Code(s)-\omega$ is positive. Its leading exponent is
$e_0=\omega^{-1}$, including the empty-sign case, in which the leading
tail coefficient is $3$. Since $e_0<1/2$, comparison of normal forms
gives $\Code(s)-\omega<\omega^{1/2}$. Hence $\Code(s)\in\Cell$.

To decode, take the unique normal form of $z$. Verify that its leading
term is $\omega$, that its remaining exponents are exactly
$e_\beta$ for an initial ordinal segment $\beta\leq\alpha$, that
the last coefficient is $3$, and that every preceding coefficient is
$1$ or $2$. The map $\beta\mapsto e_\beta$ is injective, so the end
marker uniquely determines $\alpha$. At each $\beta<\alpha$, recover
$-$ from coefficient $1$ and $+$ from coefficient $2$. This produces
exactly one sign sequence $s$ and satisfies $\Decode(\Code(s))=s$.

These are set-theoretic definitions: the assertion that there exists
an ordinal $\alpha$ and the specified indexed normal form is a
first-order assertion about sets. It neither quantifies over class
sequences nor needs a field-intrinsic way to recognize $\omega$.
The unique-output closure lemma applied to both maps proves
\eqref{dsn:eq:degree}. The same proof is a theorem internal to any model;
external definability uses standard formulas for the two maps.
\end{proof}

\begin{example}
The first three codes are
\[
 \Code(0)=\omega+3\omega^{\omega^{-1}},\qquad
 \Code(1)=\omega+2\omega^{\omega^{-1}}+3\omega^{\omega^{-2}},
 \qquad
 \Code(-1)=\omega+\omega^{\omega^{-1}}+3\omega^{\omega^{-2}}.
\]
The code for $0$ is not $0$. The coding theorem makes no assertion of
arithmetic compatibility.
\end{example}

\begin{corollary}[Congruence and leading-data blindness \src{04}]\label{dsn:cor:blind}
Every code has leading exponent $1$, leading coefficient $1$, constant
term $0$, and coefficients from $\{1,2,3\}$. Moreover, for every
ordinary positive integer $n$,
\[
 \Code(s)\in n\Oz.
\]
These invariants therefore do not bound ambient definability complexity.
\end{corollary}
\begin{proof}
Division by an ordinary $n>0$ scales the real coefficients and leaves
all exponents positive, so the quotient is still omnific. The remaining
assertions follow from the normal form. Since every surreal is
recoverable from its code, the listed common invariants cannot
separate definability degrees.
\end{proof}

\begin{proposition}[A complementary two-term code \src{04}]\label{dsn:prop:twotermcode}
There is also a parameter-free definable order isomorphism
\[
 \Two:\No\longrightarrow
 \Cell_2:=\{\omega+\omega^g:0<g<1/2\}\subseteq\Cell.
\]
Every value has exactly two nonzero normal-form terms, both with
coefficient $1$, and is mutually ambiently definable with its input.
\end{proposition}
\begin{proof}
The rationally defined map
\[
 h(x)=\frac14\left(1+\frac{x}{1+|x|}\right)
\]
is an increasing bijection from $\No$ onto $(0,1/2)$. For its
inverse, put $u=4g-1$ and use $x=u/(1-|u|)$. Therefore
$\Two(x)=\omega+\omega^{h(x)}$ has the stated normal form,
is strictly increasing, and has exactly the stated image. Extracting
the lower exponent from the unique normal form and applying the
inverse of $h$ recovers $x$ by a parameter-free set-theoretic
operation.
\end{proof}

The two codes have different purposes. The elementary two-term
code allows an arbitrary surreal exponent to carry all the input
information. The sign code $\Code$ instead fixes its exponents in
an ordinal-definable template, so that all information is stored in
the assignment of a finite coefficient alphabet. It is the latter
feature that proves the failure of termwise $\HOD$ closure
(\cref{dsn:thm:termwise}).

\begin{remark}[An interval is still a proper class \src{04}]
Although $\Cell$ is bounded by two explicitly given surreals, it is a
proper class in the ambient universe. Otherwise its subset
$\ran(\Code)$ would be a set and Replacement applied to $\Decode$
would make all of $\No$ a set. Ordered boundedness is not a size bound
for the full surreal field.
\end{remark}

\subsection{A monomial code below omega \src{07}}

\begin{lemma}[\src{07}]\label{dsn:lem:tau}
In every real closed ordered field $F$, the map
\[
 \tau(x)=\frac12\left(1+\frac{x}{\sqrt{1+x^2}}\right)
\]
is an order isomorphism $F\to(0,1)_F$. Its inverse is
\begin{equation}\label{dsn:eq:tauinv}
 \tau^{-1}(t)=\frac{2t-1}{\sqrt{1-(2t-1)^2}}.
\end{equation}
Both maps are parameter-free definable in the ordered-field language.
\end{lemma}
\begin{proof}
Since $1+x^2>x^2$, we have $|x|<\sqrt{1+x^2}$, so $0<\tau(x)<1$. If
$u\in(-1,1)$, put $x=u/\sqrt{1-u^2}$. Direct calculation gives
\[
 1+x^2=\frac1{1-u^2},\qquad
 \frac{x}{\sqrt{1+x^2}}=u.
\]
This proves the inverse formula. To check monotonicity without analytic
assumptions, the function $x/\sqrt{1+x^2}$ preserves sign. For $0\le x<y$,
comparison of its squares reduces to
\[
 x^2(1+y^2)<y^2(1+x^2),
\]
which is equivalent to $x^2<y^2$. The negative case follows by oddness, and
opposite signs are immediate.
\end{proof}

\begin{theorem}[Bounded omnific coding \src{07}]\label{dsn:thm:bounded-code}
The function
\begin{equation}\label{dsn:eq:mcode}
 \Mc(x)=\Omega(\tau(x))=\omega^{\tau(x)}
\end{equation}
is a parameter-free ambient-definable order embedding of $\No$ into
$\Oz\cap(0,\omega)$. Its image is exactly $\{\omega^a:0<a<1\}$, and its
inverse there is $\tau^{-1}\circ\Omega^{-1}$. For every positive ordinary
integer $n$,
\begin{equation}\label{dsn:eq:code-bounds}
 n<\Mc(x)<\omega/n.
\end{equation}
For any $M$ and $A$ as in \cref{dsn:def:ambient},
\begin{equation}\label{dsn:eq:code-def}
 x\in K^M_A\quad\Longleftrightarrow\quad \Mc(x)\in I^M_A.
\end{equation}
More strongly, $\dcl_{(M,\in)}(x)=\dcl_{(M,\in)}(\Mc(x))$.
\end{theorem}
\begin{proof}
The exponent $\tau(x)$ lies strictly between $0$ and $1$. A one-term normal
form with a positive exponent is purely infinite, hence omnific.
Comparison of the exponents $0$, $\tau(x)$, and $1$ gives
\eqref{dsn:eq:code-bounds}. Both $\tau$ and $\Omega$ are increasing, so
their composite is an order embedding. The displayed inverse follows from
uniqueness of the monomial exponent and \cref{dsn:lem:tau}. Each direction
is a unique ambient definition, so \cref{dsn:lem:closure} gives
\eqref{dsn:eq:code-def}. Substitution of these definitions also gives the
equality of definable closures, with or without additional parameters.
\end{proof}

This theorem does not assert that the interval is order-isomorphic to all
of $\No$: the image consists only of selected monomials. Nor does it give
an ordered-field-definable map; such a map into this discrete set would
contradict o-minimality. Its use of the ambient omega map is
indispensable. The code $\Mc$ and 04's two-term code $\Two$ rest on the
same idea, an algebraic compression of $\No$ into an interval followed by
the omega map; $\Mc$ has a single monomial term and lands below $\omega$,
$\Two$ has two terms and lands in $\Cell$.

\begin{theorem}[Local omnific coding \src{07}]\label{dsn:thm:local-code}
Let $a<b$ be omnific integers. If $\ell=b-a$ is positive infinite and
$\beta>0$ is its leading normal-form exponent, then
\begin{equation}\label{dsn:eq:local-code}
 \Mc_{a,b}(x)=a+\omega^{\beta\tau(x)}
\end{equation}
is an order embedding $\No\to(a,b)\cap\Oz$, ambient-definable over $a,b$,
with an inverse definable over the same parameters. Consequently $x$ and
$\Mc_{a,b}(x)$ are interdefinable over $a,b$. If $b-a$ is finite,
$(a,b)\cap\Oz$ is finite, so no injection from all surreals into it
exists.
\end{theorem}
\begin{proof}
The leading coefficient of $\ell$ is positive. Since
$0<\beta\tau(x)<\beta$, the monomial in \eqref{dsn:eq:local-code} is
positive infinite but of strictly smaller Archimedean order than $\ell$.
It is therefore strictly between $0$ and $\ell$. The sum is omnific, and
monotonicity follows from $\beta>0$. For a point $z$ in the image, recover
$\gamma=\Omega^{-1}(z-a)$, then $\tau(x)=\gamma/\beta$, then $x$ from
\eqref{dsn:eq:tauinv}. The leading exponent $\beta$ is definable from
$a,b$, so no additional parameter was introduced. If $\ell$ is finite it is
an ordinary positive integer by \eqref{dsn:eq:Oz}. Any omnific difference
strictly between $0$ and $\ell$ is then one of $1,\ldots,\ell-1$, so the
interval contains exactly those finitely many translates of $a$.
\end{proof}

The relevant condition is infinite \emph{width}, not infinite endpoints.
For example $(\omega,\omega+3)$ contains just two omnific integers, while
$(\omega,2\omega)$ can encode every surreal by the theorem. Compare
\cref{dsn:cor:localize}, which places an independent family at every
definable infinite scale.

\subsection{A uniformly definable bounded family with unbounded birthdays \src{07}}

An order embedding into a numerically bounded interval need not preserve
birthdays. Nevertheless its definable inverse lets us force birthdays to
escape every prescribed ordinal bound.

\begin{theorem}[Canonical birthday escape \src{07}]\label{dsn:thm:bdescape}
There is a parameter-free set-theoretically definable class function
$\esc:\Ord\to\Oz\cap(0,\omega)$ such that $\bd(\esc(\alpha))>\alpha$ for
every ordinal $\alpha$. It requires no global choice and no truth
predicate.
\end{theorem}
\begin{proof}
For a fixed ordinal $\alpha$, the class $\No_{\le\alpha}=\{z\in\No:
\bd(z)\le\alpha\}$ is a set: its members are sign sequences of bounded
ordinal length. By Separation and Replacement, the inverse image
\[
 Q_\alpha=\{\Mc^{-1}(z):z\in\No_{\le\alpha}\cap\ran(\Mc)\}
\]
is a set. Define
\begin{equation}\label{dsn:eq:bdescape}
 u_\alpha=\cut{Q_\alpha}{\varnothing},\qquad
 \esc(\alpha)=\Mc(u_\alpha).
\end{equation}
Every construction is uniquely and uniformly definable from $\alpha$.
Suppose $\bd(\esc(\alpha))\le\alpha$. Then $\esc(\alpha)\in
\No_{\le\alpha}\cap\ran(\Mc)$, so its unique inverse $u_\alpha$ belongs to
$Q_\alpha$. But the defining cut makes $u_\alpha$ strictly greater than
every member of $Q_\alpha$, including itself. This contradiction proves the
strict birthday bound.
\end{proof}

\begin{corollary}[\src{07}]\label{dsn:cor:escape-def}
For every $A$-definable ordinal $\alpha$ of $M$, there is an $A$-definable
omnific integer $z$ with $0<z<\omega$ and $\bd(z)>\alpha$. The
ordinal-definable omnific integers in $(0,\omega)$ have birthdays unbounded
in $\Ord$ in every universe of ZFC, independently of whether $V=\HOD$.
\end{corollary}
\begin{proof}
Apply \cref{dsn:thm:bdescape} internally and \cref{dsn:lem:closure}. With
ordinal parameters allowed, every $\esc(\alpha)$ is ordinal-definable.
\end{proof}

The last statement does not imply density of ordinal-definable surreals.
Birthday cofinality of a bounded family and density in the numerical order
are different properties; Hamkins's density theorem has the substantially
stronger hypothesis discussed in \cref{dsn:thm:hamkins}. The escape
$\esc$ is also different from 04's $\Sigma_n$ escape of
\cref{dsn:thm:sigmaescape}, which escapes a definability level rather than
a birthday.

\section{One interval detects the entire set-theoretic universe}
\label{dsn:sec:global}

We first spell out the standard set-coding step needed for the global
consequences. This step is not an original bi-interpretation claim.

\begin{lemma}[Every set is recoverable from some surreal \src{04, 07}]
\label{dsn:lem:setcoding}
There is a parameter-free partial set-theoretic function $\SetDec$ on
surreal sign sequences such that every set is $\SetDec(s)$ for some $s$.
No parameter-free selection of such an $s$ from the set being coded is
asserted. In particular, if every surreal of a model $M$ is parameter-free
definable over $M$, then every set of $M$ is.
\end{lemma}
\begin{proof}
For a set $a$, Choice gives a bijection $f:\kappa\to\TC(\{a\})$
for some ordinal $\kappa$. Transport membership to a well-founded
extensional relation $R$ on $\kappa$, and record the unique index
$\varrho$ with $f(\varrho)=a$.

Fix a parameter-free class pairing of ordinals $p:\Ord^2\to\Ord$. One
construction well-orders ordinal pairs first by their maximum and then
lexicographically, and uses the order type of each predecessor set. Every
predecessor collection is a set, so this defines a class bijection. (07
uses instead the injection $(\xi,\eta)\mapsto\kappa\cdot\xi+\eta$ with
ordinal arithmetic.) Form the set of ordinals
\[
 X_a=\{p(0,\kappa),p(1,\varrho)\}
       \cup\{p(2,p(\xi,\eta)):(\xi,\eta)\in R\}.
\]
Let $\lambda=\sup\{\gamma+1:\gamma\in X_a\}$ and take the sign
sequence of length $\lambda$ whose positive positions are exactly
$X_a$. From this sequence recover $X_a$, then $\kappa,\varrho,R$, and
finally the value at $\varrho$ of the unique Mostowski collapse. These
operations are parameter-free and unique whenever the code is valid. For
invalid codes leave $\SetDec$ undefined. The last sentence follows by
\cref{dsn:lem:closure} applied inside $M$. Mostowski collapse is standard
\cite{Jech}; the same pointed-graph mechanism is present in the announced
surreal/set-theory dictionary \cite{CHY}.
\end{proof}

The use of Choice in this proof is existential. It does \emph{not} produce
a canonical definable well-ordering of every transitive closure from its
underlying set. In particular, the lemma must not be strengthened to a
definable injection of arbitrary sets into omnific integers without further
assumptions (07).

\begin{theorem}[The ordinal-definability test \src{04, 07}]\label{dsn:thm:globalOD}
The following statements are equivalent in ZFC:
\begin{enumerate}[label=\textup{(\roman*)}]
\item $V=\HOD$;
\item every surreal is ordinal definable;
\item every omnific integer is ordinal definable;
\item every member of $\Cell$ is ordinal definable \src{04};
\item every value of the code $\Code$ is ordinal definable \src{04};
\item every omnific integer in $(0,\omega)$ is ordinal definable \src{07};
\item every monomial $\omega^a$ with $0<a<1$ is ordinal definable \src{07}.
\end{enumerate}
\end{theorem}
\begin{proof}
The implications (i)$\Rightarrow$(ii)$\Rightarrow$(iii)$\Rightarrow$(iv)
$\Rightarrow$(v) and (iii)$\Rightarrow$(vi)$\Rightarrow$(vii) are immediate
from the inclusions. If (v) holds, recover each $s$ from $\Code(s)$ by
\cref{dsn:thm:code}; if (vii) holds, recover each $x$ from $\Mc(x)$ by
\cref{dsn:thm:bounded-code}. Either way every surreal is ordinal definable.
For any set $a$, choose a surreal $s$ with $\SetDec(s)=a$ by
\cref{dsn:lem:setcoding}. Its canonical sign code lies in $\HOD$ by
\cref{dsn:thm:HOD}. More explicitly, the associated ordinal set $X_a$ lies
in $\HOD$, and $\HOD$ decodes the same well-founded extensional relation
and its Mostowski collapse. Thus $a\in\HOD$. Since $a$ was arbitrary,
$V=\HOD$.
\end{proof}

The equivalence of (i) and (ii) is in substance part of Hamkins's 2012
answer (\cref{dsn:thm:hamkins}); source 04 did not cite it. The bounded
conditions (iv)--(vii) are the forms proposed by the two sources.

\begin{corollary}[Two terms already detect HOD \src{04}]
The equivalent conditions in \cref{dsn:thm:globalOD} remain
equivalent if condition \textup{(iv)} is replaced by ``every member
of $\Cell_2$ is ordinal definable.''
\end{corollary}
\begin{proof}
Use the two mutually definable maps in \cref{dsn:prop:twotermcode} to
recover every surreal from such a two-term omnific integer, and then apply
\cref{dsn:thm:globalOD}.
\end{proof}

\begin{remark}
There is no Choice-based definable selection hidden in the converse.
We choose a code only to prove its existence; hypothesis (ii) says
\emph{every} such surreal code is ordinal definable. This is much
weaker as a proof requirement than constructing a canonical code from
an arbitrary set in a universe not satisfying $V=\HOD$.
\end{remark}

\begin{theorem}[Hamkins, 2012; included for comparison \src{07}]\label{dsn:thm:hamkins}
The ordinal-definable surreals are dense in $\No$ if and only if $V=\HOD$.
Hamkins further proves equivalence with a parameter-free definable map
$\Ord\to\No$ having dense image, and with definable surjectivity and
bijectivity in that setting \cite{HamkinsDensity}.
\end{theorem}
\begin{proof}[Proof of the density equivalence]
If $x$ is a sign sequence which is not $\OD$, let $x^-$ and $x^+$ be its
extensions by one minus sign and one plus sign. Every number strictly
between $x^-$ and $x^+$ extends $x$: an earlier disagreement, or
termination before the length of $x$, puts it outside this interval. If
such a number $y$ were $\OD$, then the prefix $y\res\bd(x)$ would be $\OD$
using the ordinal parameter $\bd(x)$. That prefix is $x$, a contradiction.
Thus density forces every surreal to be $\OD$, and
\cref{dsn:thm:globalOD} yields $V=\HOD$. The converse is immediate.
\end{proof}

The proof just given is the prefix argument from Hamkins's earlier result,
not a new density proof claimed here. The bounded tests use the explicit
omnific codes, whereas the density criterion is already known.

\begin{proposition}[Specialization of the old-field theorem \src{07}]\label{dsn:prop:hod-topology}
If $V\ne\HOD$, the subclass $\No^{\HOD}$ is closed and nowhere dense in the
fine order topology of $\No$.
\end{proposition}
\begin{proof}
For every $x\notin\No^{\HOD}$, the interval $(x^-,x^+)$ in the previous
proof contains no $\HOD$ surreal. Hence the complement is open, and
$\No^{\HOD}$ is closed.

For completeness, any proper real closed subfield $F\subsetneq\No$ has
empty interior. Suppose an open interval were contained in $F$, and choose
$c<d$ in that interval. For arbitrary $y\in\No$, $z=c+(d-c)\tau(y)$ lies
between $c$ and $d$, so $z\in F$. Then $(z-c)/(d-c)\in F$, and the algebraic
inverse \eqref{dsn:eq:tauinv}, using the positive square root in the real
closed field $F$, implies $y\in F$. This would give $F=\No$. Apply this to
the proper field $\No^{\HOD}$ and combine empty interior with closedness.
\end{proof}

The universes report proves the more general assertion for proper
same-ordinal inner models, including its surcomplex analogue
(\lbl{univ:thm:nowheredense}; source 07 cited it as ``Theorem 14.3'' of
that report). \Cref{dsn:prop:hod-topology} is the case $M=\HOD$ of that
theorem and is not counted as a new result. It also illustrates the
difference between fine surreal topology and the usual topology on the
embedded real line.

\begin{theorem}[The pointwise-definability test \src{04, 07}; items (v)--(vi) for arbitrary set models \mergetag]\label{dsn:thm:pointwise}
Let $M\models\ZFC$ be any set model. In external semantics, the
following are equivalent:
\begin{enumerate}[label=\textup{(\roman*)}]
\item $M$ is pointwise definable;
\item every member of $\No^M$ is parameter-free definable in $M$;
\item every member of $\Oz^M$ is parameter-free definable in $M$;
\item every member of $\Cell^M$ is parameter-free definable in $M$ \src{04};
\item every omnific integer $z$ of $M$ with $0<z<\omega$ is parameter-free
definable in $M$ \src{07};
\item every monomial $\omega^a$ of $M$ with $0<a<1$ is parameter-free
definable in $M$ \src{07}.
\end{enumerate}
\end{theorem}
\begin{proof}
The implications (i)$\Rightarrow$(ii)$\Rightarrow$(iii)$\Rightarrow$(iv)
and (iii)$\Rightarrow$(v)$\Rightarrow$(vi) are immediate. Each
$s\in\No^M$ is recovered by a fixed standard formula from
$\Code^M(s)\in\Cell^M$, and each $x\in\No^M$ from the monomial
$\Mc^M(x)$, so (iv) and (vi) each imply (ii). Internally to $M$, every set
has a code as in \cref{dsn:lem:setcoding}. For any fixed $a\in M$, choose
such a code $s\in M$. A standard formula defines $s$ by (ii), and composing
it with the standard decoder formula defines $a$. This does not require
$M$'s internal relation to be externally well-founded: the collapse
operation is the one that $M$ proves exists uniquely.
\end{proof}

Source 04 proved (i)--(iv) for arbitrary set models; source 07 proved
(i)--(iii), (v) and (vi) for transitive $M$. The argument for
(vi)$\Rightarrow$(ii) uses only internal formulas, so 04's proof covers
07's conditions for arbitrary set models; this extension is the merge's.

Pointwise definable models of ZFC exist whenever ZFC is consistent;
Hamkins--Linetsky--Reitz also prove strong extension theorems and treat
nonstandard as well as transitive models \cite{HLR}. Their existence
theorem is prior work, and no existence assumption is needed for the
equivalence. The new feature of \cref{dsn:thm:pointwise} is the
restriction of the test to a single fixed-leading-term or bounded omnific
interval. It also explains why an external counting argument cannot be
converted into an internal proof that some member of $\Oz^M$ is
undefinable: the formulas that define particular elements are not
presented to $M$ as a surjective internal list. In a pointwise definable
model every set can be definable while the model still satisfies its own
uncountability assertions.

\section{From one missing number to proper-class transcendence}
\label{dsn:sec:transcendence}

This section contains the strongest algebraic result, due to source 04.
We first isolate a support argument that does not involve definability at
all.

\begin{lemma}[Disjoint support cosets \src{04}]\label{dsn:lem:cosets}
Let $K\subseteq\No$ be a subfield, and let $G\subseteq(\No,+)$ be an
additive subgroup such that
\[
 \supp(a)\subseteq G\qquad(0\ne a\in K).
\]
Suppose $(g_i)_{i\in J}$ is a family such that every nonzero finitely
supported integer combination satisfies
\begin{equation}\label{dsn:eq:modind}
 \sum_i m_i g_i\notin G.
\end{equation}
Then the family $(\omega^{g_i})_{i\in J}$ is algebraically independent
over $K$. The assertion applies to a class family when the condition
is read on each finite subfamily.
\end{lemma}
\begin{proof}
Take a nonzero polynomial in finitely many distinct variables,
\[
 P(Y_1,\ldots,Y_n)=\sum_{\nu\in F}a_\nu Y_1^{\nu_1}\cdots
 Y_n^{\nu_n},
\]
where $F\subseteq\N^n$ is finite and the displayed coefficients are
nonzero. Since $\omega^u\omega^v=\omega^{u+v}$, the normal-form
support of the term at multi-index $\nu$ after substitution is
\begin{equation}\label{dsn:eq:supportshift}
 \supp(a_\nu)+\sum_{j=1}^n\nu_jg_j
 \ \subseteq\ G+\sum_{j=1}^n\nu_jg_j.
\end{equation}
For $\nu\ne\mu$, these cosets are disjoint. Otherwise subtracting an
exponent in their intersection would give
$\sum_j(\nu_j-\mu_j)g_j\in G$, contrary to \eqref{dsn:eq:modind}.
Every displayed term is nonzero, and its coefficients cannot be
cancelled by any other term. The finite sum is therefore nonzero.
This proves the required absence of polynomial relations. No infinite
polynomial, class sum, or interchange of summations is used.
\end{proof}

Related (batch 89, added 4 October 2026): the case $K=\R$, $G=\{0\}$ is
proved again, without this lemma, as Theorem~43.1 (\lbl{hset:kw:thm:double})
of \repopath{foundations-and-computation/birthday-cutoffs-and-hereditary-sets},
Part~VII (batch-89 manuscript~04; AI-assisted, unrefereed, not formalized):
for an injection $j:X\to\No$ the numbers $t_x=\omega^{\omega^{j(x)}}$ are
algebraically independent over $\R$ and are positive infinite omnific
integers, in ZF. That part uses it to show that embedding the free rings
$\Z[T_X]$ into $\Oz$, or the free fields $\R(T_X)$ into $\No$, for every set
$X$ is exactly the first Kinna--Wagner principle (its Theorem~43.5,
\lbl{hset:kw:thm:equivalences}); its Remark~40.4 records the overlap with
this lemma.

Related (batch 90, added 4 October 2026): a second case, with irrational
coefficients in the exponents, is Lemma~3.2 (\lbl{csf:lem:independence}) of
\repopath{foundations-and-computation/cantor-families-of-surreal-subfields}
(batch-90 manuscript~04; AI-assisted, unrefereed, not formalized): take $K=k$,
the field of real algebraic numbers, $G=\{0\}$ and $g_{q,b}=b\,\omega^q$ for
$q\in\Q$ and $b\in\{1\}\cup\{\sqrt p:p\text{ prime}\}$; nonzero integer
combinations of the $g_{q,b}$ are nonzero by the $\Q$-linear independence of
square roots of squarefree integers (its Lemma~3.1,
\lbl{csf:lem:radicals}). That report proves the lemma directly and uses the
monomials $\omega^{b\omega^q}$ as a transcendence basis of a countable real
closed field whose relative real closures $K_A$ form a closed Cantor family
of subfields on which embeddability is a complete analytic quasiorder (its
Theorem~7.4, \lbl{csf:thm:universality}); a dated note after its Lemma~3.2
credits this lemma.

\begin{theorem}[Support-subfield amplification \src{04}]\label{dsn:thm:amplify}
Let $K$ be a proper subfield of $\No$ satisfying
\begin{equation}\label{dsn:eq:Khyp}
 \omega\in K,\qquad e_\alpha\in K\ (\alpha\in\Ord),\qquad
 \supp(a)\subseteq K\ (0\ne a\in K),
\end{equation}
where $e_\alpha=\omega^{-(\alpha+1)}$.
For any $t\in(0,1)\setminus K$, the class family
\begin{equation}\label{dsn:eq:indfam}
 y_\alpha=\omega^{t e_\alpha},\qquad
 z_\alpha=\omega+y_\alpha\qquad(\alpha\in\Ord)
\end{equation}
is algebraically independent over $K$, both for the $y$ family and
for the $z$ family. Each $z_\alpha$ lies in $\Cell$ and has exactly
two nonzero normal-form terms, both with coefficient $1$.
Such a $t$ exists for every proper $K$ satisfying \eqref{dsn:eq:Khyp}.
\end{theorem}
\begin{proof}
First, the family $(e_\alpha)$ is linearly independent over $\Q$.
Indeed, a finite rational linear combination is already a normal
form with distinct exponents $-(\alpha+1)$; after deleting zero
coefficients it cannot vanish. Hence for distinct
$\alpha_1,\ldots,\alpha_n$ and integers $m_j$ not all zero,
\[
 c=\sum_{j=1}^n m_j e_{\alpha_j}\in K\setminus\{0\}.
\]
If $\sum_jm_jte_{\alpha_j}=tc$ belonged to $K$, division by $c$
would give $t\in K$, a contradiction. Thus the exponents
$te_\alpha$ satisfy \eqref{dsn:eq:modind} for the additive subgroup
$G=(K,+)$. \Cref{dsn:lem:cosets} proves algebraic independence of
the $y_\alpha$.

Translation by $\omega\in K$ is a polynomial automorphism on each
finite tuple of variables. Therefore a nonzero relation among the
$z_\alpha=\omega+y_\alpha$ would give a nonzero relation among the
$y_\alpha$, which is impossible.

Finally,
\[
 0<te_\alpha<e_\alpha\leq\omega^{-1}<1/2.
\]
Consequently $\omega^{te_\alpha}$ is a positive, purely infinite
omnific monomial strictly below $\omega^{1/2}$. The exponents $1$
and $te_\alpha$ are distinct and positive, so the asserted two-term
normal form is exact and $z_\alpha\in\Cell$.

For existence, take $x\in\No\setminus K$. Then $u=|x|>0$ is not
in $K$, since otherwise $x=u$ or $x=-u$ would be in $K$. Set
$t=u/(1+u)$. It lies strictly between $0$ and $1$. If $t\in K$,
then $u=t/(1-t)\in K$, again a contradiction.
\end{proof}

\begin{theorem}[The HOD-surreal transcendence dichotomy \src{04}]
\label{dsn:thm:dichotomy}
Exactly one of the following occurs:
\begin{enumerate}[label=\textup{(\roman*)}]
\item $V=\HOD$ and $\No=\KH$;
\item $V\ne\HOD$, and for any $t\in(0,1)\setminus\KH$, the family
\[
 \left(\omega+\omega^{\,t\omega^{-(\alpha+1)}}\right)_{\alpha\in\Ord}
\]
is algebraically independent over $\KH$ and is contained in $\Cell$.
\end{enumerate}
In case \textup{(ii)}, for every cardinal $\kappa$ there is an
algebraically independent subset of $\Cell$ of cardinality $\kappa$
over $\KH$. In particular, the transcendence defect has no set-sized
cardinal bound.
\end{theorem}
\begin{proof}
\Cref{dsn:thm:globalOD} identifies $V\ne\HOD$ with $\KH\ne\No$. The field
$\KH$ contains $\omega$ and every $e_\alpha$ because these are ordinal
definable. Its support property is \cref{dsn:lem:HODsupport}. Apply
\cref{dsn:thm:amplify}. Restricting the class family to $\alpha<\kappa$
gives a set family by Replacement. Algebraic independence implies
distinctness, hence that set has cardinality $\kappa$.
\end{proof}

\begin{remark}[What ``proper-class transcendence'' means \src{04}]
The theorem gives a concrete class injection with finite polynomial
independence, definable from one parameter $t$. It does not choose a
transcendence basis for $\No$, invoke Global Choice, or assign an
ordinary cardinal to a proper class. Each individual $z_\alpha$ is
transcendental over $\KH$, but this pointwise assertion alone would
be much weaker than the simultaneous algebraic independence proved
above. In a set model, the assertions have their internal cardinal
and ordinal meaning, not an assertion that its externally set-sized
carrier is literally a proper class in the metatheory.
\end{remark}

\begin{corollary}[Localization at other definable scales \src{04}]
\label{dsn:cor:localize}
Assume $V\ne\HOD$. For every $p\in\IH$, every positive
$\gamma\in\KH$, and every $t\in(0,1)\setminus\KH$, the family
\begin{equation}\label{dsn:eq:localfam}
 w_\alpha=p+\omega^{\gamma t e_\alpha}\qquad(\alpha\in\Ord)
\end{equation}
is algebraically independent over $\KH$ and lies in
$(p,p+\omega^\gamma)\cap\Oz$.
\end{corollary}
\begin{proof}
For every nonzero finite integer combination $c$ of the $e_\alpha$,
we have $c\in\KH\setminus\{0\}$. If $\gamma tc\in\KH$, division
by $\gamma c$ would put $t$ in $\KH$. The coset proof therefore
applies to the exponents $\gamma te_\alpha$. Translation by
$p\in\KH$ preserves independence. The exponents are positive and
less than $\gamma$, so the added monomials are omnific and lie
strictly between $0$ and $\omega^\gamma$.
\end{proof}

The widths in this corollary are infinite monomial widths. It would
be false to replace them indiscriminately by positive real widths:
an interval of width less than $1$ contains at most one omnific
integer. The result is localization at an arbitrary definable
\emph{infinite scale}, not density of omnific integers in $\No$.

\section{Uniform definability is not termwise definability}
\label{dsn:sec:uniformity}

There are two different pitfalls. Individually definable terms need
not form a definable sequence; conversely, a definable sequence need
not have definable entries at arbitrary ordinal indices. Each source
gives an explicit example of each; both pairs are kept, since the
mechanisms differ.

\begin{theorem}[Definable support and terms do not suffice \src{04}]
\label{dsn:thm:termwise}
Assume $V\ne\HOD$. There is a valid normal form $q$ satisfying all
of the following:
\begin{enumerate}[label=\textup{(\roman*)}]
\item $q\in\Cell\setminus\IH$;
\item its support, as a whole set, belongs to $\HOD$;
\item every coefficient is in $\{1,2,3\}$, and every individual
monomial term belongs to $\KH$;
\item the whole indexed coefficient/exponent sequence does not
belong to $\HOD$.
\end{enumerate}
Thus $\KH$ is not closed under arbitrary ambient Hahn-summable
families of its elements, even when the support belongs to $\HOD$.
\end{theorem}
\begin{proof}
Choose $s\notin\KH$ by \cref{dsn:thm:globalOD}, and let
$\alpha=\bd(s)$ and $q=\Code(s)$. The reversible code gives
$q\notin\IH$. Its support is the set
\[
 \{1\}\cup\{e_\beta:\beta\leq\alpha\},
\]
which is definable from the ordinal $\alpha$ and whose members all
belong to $\HOD$. Hence the support itself belongs to $\HOD$. Every
individual coefficient is one of three fixed rationals, and every
individual exponent is ordinal definable, so every individual term
is ordinal definable. If the whole normal form belonged to $\HOD$,
\cref{dsn:lem:HODsupport} would give $q\in\KH$, a contradiction.
The summable family here is simply the family of its distinct
monomial terms, so no additional summability issue occurs.
\end{proof}

\begin{remark}
In this example even the \emph{set of coefficient values} is a
finite $\HOD$ set. What fails is the definability of which coefficient
is assigned to which exponent. Saying ``the coefficients are
rational'' erases precisely the information carrying the original
sign sequence.
\end{remark}

\begin{theorem}[Uniform subset coding by omnific normal forms \src{07}]\label{dsn:thm:bit-code}
For an ordinal $\alpha$ and $X\subseteq\alpha$, put $\iota_\beta=
1/(\underline\beta+2)$ and
\begin{equation}\label{dsn:eq:bit-code}
 E_\alpha(X)=\omega^{3/4}
  +\sum_{\beta<\alpha}(1+\chi_X(\beta))\,\omega^{\iota_\beta}.
\end{equation}
Then this is a well-defined omnific integer satisfying
\begin{equation}\label{dsn:eq:bit-bounds}
 \omega^{3/4}\le E_\alpha(X)<2\omega^{3/4}<\omega,
\end{equation}
with strictness in the first inequality when $\alpha>0$. The pair
$(\alpha,X)$ and $E_\alpha(X)$ are mutually ambient-definable without
additional parameters. In particular,
\[
 E_\alpha(X)\in K^M_A\quad\Longleftrightarrow\quad \alpha,X\in D^M_A,
\]
and the analogous equivalence holds with ordinal definability in place of
$A$-definability.
\end{theorem}
\begin{proof}
The ordinal embedding is strictly increasing, so the exponents
$\iota_\beta$ are strictly decreasing and positive. Their largest possible
value is $1/2$, below the marker exponent $3/4$. Thus \eqref{dsn:eq:bit-code}
is already a decreasing normal form, all of whose exponents are positive.
It is omnific. Its leading term has coefficient $1$ at exponent $3/4$, so
its lower terms are strictly smaller in Archimedean order; this proves
\eqref{dsn:eq:bit-bounds}.

Starting from $(\alpha,X)$, Replacement and the normal-form construction
define the number uniquely. Starting from the number, take its canonical
normal form, remove the first, marker term, and take the order type of the
remaining decreasing support. This is $\alpha$, including the empty case.
The coefficient at the $\beta$th remaining place is $2$ precisely when
$\beta\in X$. Hence $X$ is recovered uniquely. The definability
equivalences follow from composition.
\end{proof}

The marker in \eqref{dsn:eq:bit-code} is necessary for this simple uniform
empty-case decoder and bound. No cancellation can obscure a bit: the
exponents are distinct and the coefficients are never zero. The bit code
and 04's sign code $\Code$ exploit the same phenomenon; $\Code$ is
self-delimiting through its coefficient alphabet $\{1,2,3\}$, while
$E_\alpha$ uses a marker term and lands below $\omega$.

\begin{corollary}[Individually definable entries are insufficient \src{07}]\label{dsn:cor:entries-insufficient}
Suppose $X\subseteq\omega$ belongs to $M$ but is not $A$-definable. Then
$E_\omega(X)$ is not $A$-definable, although its support is parameter-free
definable and every one of its real coefficients and exponents is
individually parameter-free definable. If $X$ is not ordinal-definable,
the entire number is not ordinal-definable either.
\end{corollary}
\begin{proof}
The support is the fixed set $\{3/4\}\cup\{1/(n+2):n\in\N\}$. Its entries
are rational, and each coefficient is either $1$ or $2$. Nondefinability of
the sum follows from the decoder in \cref{dsn:thm:bit-code}.
\end{proof}

There is no contradiction in saying that every coefficient is definable
while the coefficient assignment is not. The statement ``the coefficient
at $n$ is $1$'' or ``it is $2$'' may require a nonuniform choice of which
true statement to use. The missing information is exactly the subset $X$.
For arbitrary $X\subseteq\alpha$, all exponents in \eqref{dsn:eq:bit-code}
are ordinal-definable and the support is definable from $\alpha$; a
non-$\OD$ $X$ still gives a non-$\OD$ sum. Thus even allowing ordinal
parameters does not turn pointwise definability of data into uniform
definability.

\begin{proposition}[A definable integer can encode nondefinable entries \src{04}]
\label{dsn:prop:collector}
There is a parameter-free ambiently definable $q_*\in\Cell$ whose
normal-form coefficients canonically encode every member of
$\R\cap\HOD$. In any universe or fixed transitive model in which
some member of $\R\cap\HOD$ is not parameter-free definable, $q_*$
has a normal-form coefficient that is not parameter-free definable.
The last hypothesis is essential and is not asserted in every model.
\end{proposition}
\begin{proof}
The set $R_H=\R\cap\HOD$ is parameter-free definable. Restrict the
canonical set-like $\HOD$ well-order to it, and let
$(r_\xi)_{\xi<\lambda}$ be the unique increasing enumeration. The
entire enumeration and its ordinal order type $\lambda$ are
parameter-free definable. Define the injective real function
\[
 \zeta(r)=\frac32+\frac{r}{2(1+|r|)}\in(1,2).
\]
Its inverse on $(1,2)$ is parameter-free definable: if
$u=2\zeta(r)-3$, then $r=u/(1-|u|)$. Form
\begin{equation}\label{dsn:eq:collector}
 q_*=\omega+\sum_{\xi<\lambda}\zeta(r_\xi)\omega^{e_\xi}.
\end{equation}
This is a valid normal form and is parameter-free definable from the
whole enumeration. The set $R_H$ is nonempty, and the leading tail
exponent is $e_0<1/2$, so $q_*\in\Cell$. Every $r\in R_H$ occurs
in exactly one coefficient, after applying $\zeta$.
If that $r$ is not parameter-free definable, neither is $\zeta(r)$,
since the displayed inverse would otherwise define $r$.
\end{proof}

This does not contradict unique-output closure. Extracting the
coefficient at position $\xi$ uses $\xi$ as an additional parameter;
a parameter-free definition of $q_*$ need not define every such
ordinal position. Ordinal definability allows that parameter, which
is why all coefficients of an ordinal-definable surreal do belong
to $\HOD$.

\begin{theorem}[The reverse uniformity gap \src{07}]\label{dsn:thm:reverse-gap}
Suppose $\mu\in\Ord^M$ is $A$-definable and there is $\beta<\mu$ which is
not $A$-definable. Then
\begin{equation}\label{dsn:eq:uniform-series}
 U_\mu=\sum_{\gamma<\mu}\omega^{\iota_\gamma},\qquad
 \iota_\gamma=\frac1{\underline\gamma+2},
\end{equation}
belongs to $I^M_A\cap(0,\omega)$, but its exponent $\iota_\beta$ and its
corresponding monomial are not $A$-definable.
\end{theorem}
\begin{proof}
The entire normal-form family in \eqref{dsn:eq:uniform-series} is
uniformly defined from $\mu$. Hence so is its sum, which is a positive
omnific integer with leading exponent $1/2$. If $\iota_\beta$ were
$A$-definable, taking its reciprocal and subtracting $2$ would make the
surreal ordinal $\underline\beta$, and therefore the set-theoretic ordinal
$\beta$, $A$-definable. If the monomial were $A$-definable, applying
$\Omega^{-1}$ would give the same contradiction.
\end{proof}

This theorem is conditional on the stated ordinal configuration. One must
not assume that an ordinal which is uncountable \emph{inside} a countable
transitive model necessarily has an externally undefinable predecessor;
pointwise definable models are a reminder that it need not.
\Cref{dsn:sec:core} gives an exact criterion for when the reverse gap can
occur. The collector (04) obtains a nondefinable coefficient from real
information in $\HOD$; the reverse gap (07) obtains a nondefinable
exponent from ordinal information. The two are different mechanisms.

\paragraph{A compact lesson for proposed definitions \src{07}.}
Defining a ``definable surreal'' to mean ``a normal form with definable
coefficients and exponents'' produces an unstable notion unless a
representation and a uniformity condition are specified.
\Cref{dsn:thm:termwise,dsn:thm:bit-code} show that it can admit numbers
with an undefinable information sequence.
\Cref{dsn:prop:collector,dsn:thm:reverse-gap} show that it can reject
uniformly definable numbers. Using the entire canonical graph, as in
\cref{dsn:prop:nf-invariant}, avoids both errors and agrees exactly with
ambient definability.

\section{The maximal initial definable core}
\label{dsn:sec:core}

Throughout this section, which is source 07's, $M$ is an \emph{externally
fixed transitive} set model of ZFC and $A\subseteq M$; the proofs use a
least externally omitted ordinal, which need not exist for a nonstandard
model. We write $K=K^M_A$, $I=I^M_A$ and $O=O^M_A=\Ord^M\cap D^M_A$. The
next elementary observation links the simplicity tree to definability of
ordinals.

\begin{proposition}[Initiality criterion \src{07}]\label{dsn:prop:initiality}
The subclass $K^M_A$ is initial in $\No^M$ if and only if $O^M_A$ is
downward closed in $\Ord^M$.
\end{proposition}
\begin{proof}
If $K^M_A$ is initial and $\alpha\in O^M_A$, the all-plus surreal
$\underline\alpha$ is in $K^M_A$. Each $\underline\beta$ for $\beta<\alpha$
is a prefix, hence definable. Its birthday is $\beta$, so $\beta\in O^M_A$.
Conversely, let $x\in K^M_A$. Its birthday $\alpha$ is definable from $x$,
so $\alpha\in O^M_A$. If $\beta<\alpha$, downward closure makes $\beta$
$A$-definable. The prefix $x\res\beta$ is uniquely definable from
$x,\beta$ and therefore belongs to $K^M_A$.
\end{proof}

\subsection{The first omitted ordinal is an epsilon number}

Assume $O^M_A\ne\Ord^M$. Since $M$ is transitive, the external nonempty set
$\Ord^M\setminus O^M_A$ has a least element. Denote it by
\begin{equation}\label{dsn:eq:rho}
 \rho^M_A=\min\bigl(\Ord^M\setminus O^M_A\bigr).
\end{equation}
This is a metatheoretic definition. It does not say that the minimum is
defined from $A$ in $M$.

\begin{lemma}[Cantor-normal-form closure \src{07}]\label{dsn:lem:rho-epsilon}
The ordinal $\rho=\rho^M_A$ satisfies $\omega^\rho=\rho$, where the power
is ordinal exponentiation. In particular $\rho>\omega$.
\end{lemma}
\begin{proof}
The ordinals $0$ and all finite ordinals are parameter-free definable.
Suppose a nonzero $\rho$ were not a fixed point of ordinal exponentiation
with base $\omega$. Its finite Cantor normal form would be
\[
 \rho=\omega^{\gamma_0}n_0+\cdots+\omega^{\gamma_{k-1}}n_{k-1},
 \quad \gamma_0>\cdots>\gamma_{k-1},\quad n_j\in\N_{>0},
\]
with every $\gamma_j<\rho$. Indeed, the leading exponent is at most
$\rho$; equality would force $\rho=\omega^\rho$, and would leave no further
term or coefficient greater than one. All exponents in this finite
expression are $A$-definable by minimality of $\rho$. Ordinal
exponentiation, finite ordinal sums, and finite coefficients are uniquely
definable operations. The finite expression would therefore make $\rho$
$A$-definable, a contradiction. The least epsilon number is greater than
$\omega$ (and is itself parameter-free definable), so the final assertion
follows.
\end{proof}

The argument is ordinary ordinal definability, not a new theorem about
Cantor normal form. Its role is to identify exactly the cutoff at which the
classical epsilon-field theorem becomes available.

\subsection{The core theorem}

\begin{theorem}[Maximal initial elementary exponential core \src{07}]\label{dsn:thm:core}
Let $M$ be transitive, and suppose $\rho=\rho^M_A$ exists. Define
\begin{equation}\label{dsn:eq:core}
 \Core^M_A=K^M_A\cap(\No_{<\rho})^M.
\end{equation}
Then:
\begin{enumerate}[label=\textup{(\roman*)}]
\item $\Core^M_A$ is initial, and every initial subclass of $K^M_A$ is
contained in it.
\item Its surreal ordinals are exactly the ordinals below $\rho$.
\item It is an elementary exponential subfield of $\No^M$, hence real
closed and closed under positive logarithms.
\item $\Core^M_A\cap\Oz^M$ is an integer part of $\Core^M_A$.
\item $\Core^M_A\cap\R^M=K^M_A\cap\R^M$.
\end{enumerate}
If $O^M_A=\Ord^M$, the same initiality and elementary-field conclusions
hold with $\Core^M_A=K^M_A$.
\end{theorem}
\begin{proof}
Let $x\in\Core^M_A$, and put $\alpha=\bd(x)<\rho$. Every $\beta<\alpha$ is
$A$-definable by the definition of $\rho$. Hence $x\res\beta$ is
$A$-definable and has birthday below $\rho$. This proves initiality.

No $A$-definable surreal has birthday exactly $\rho$, since the birthday is
definable from the number. If an initial subclass $S\subseteq K^M_A$
contained $x$ with birthday greater than $\rho$, it would also contain
$x\res\rho$. The birthday of that prefix is $\rho$, giving the same
contradiction. Therefore every such $S$ is contained in \eqref{dsn:eq:core}.
This proves (i), and (ii) follows because every ordinal below $\rho$ is
definable and the all-plus sequence has that ordinal as birthday.

By \cref{dsn:lem:rho-epsilon}, $\rho$ is an epsilon number. Thus the
classical cutoff theorem (\cref{dsn:sub:imported}) gives
\[
 (\No_{<\rho})^M_{\exp}\prec\No^M_{\exp}.
\]
\Cref{dsn:thm:elementary} gives $(K^M_A)_{\exp}\prec\No^M_{\exp}$. Fix the
parameter-free definable selection functions of \cref{dsn:lem:selection}
for the common ambient exponential field. Each elementary substructure is
closed under them: for parameters in it, the uniquely defined selected
value belongs to it by elementarity. Their intersection is therefore closed
under the same selectors. By \cref{dsn:lem:selection}, the intersection in
\eqref{dsn:eq:core} is elementary. This proves (iii). This step uses common
definable selectors; it is not an unsupported assertion that arbitrary
intersections of elementary submodels are elementary.

For (iv), the omnific floor of $x\in\Core^M_A$ belongs to $K^M_A$ by
\cref{dsn:thm:hull} and is a prefix of $x$ by \cref{dsn:prop:floor}.
Initiality therefore places the floor in $\Core^M_A$. Finally every
embedded real has birthday at most $\omega<\rho$, giving (v).

If all ordinals of $M$ are $A$-definable, \cref{dsn:prop:initiality} makes
$K^M_A$ initial, and \cref{dsn:thm:elementary} gives the remaining field
assertion.
\end{proof}

\begin{corollary}[Parameter monotonicity \src{07}]\label{dsn:cor:core-monotone}
If $A\subseteq A'\subseteq M$, then $\Core^M_A\subseteq\Core^M_{A'}$,
interpreting the core as in \cref{dsn:thm:core} in the no-omission case.
When both first omitted ordinals exist, $\rho^M_A\le\rho^M_{A'}$.
\end{corollary}
\begin{proof}
The initial class $\Core^M_A$ is contained in $K^M_{A'}$, so the
maximality of $\Core^M_{A'}$ gives the inclusion. Every ordinal below
$\rho^M_A$ is also $A'$-definable, giving the ordinal inequality.
\end{proof}

\begin{theorem}[Birthday-controlled omnific denominators \src{07}]\label{dsn:thm:core-fractions}
Let $\epsilon\in\Ord^M$ be any epsilon number, and put
\[
 K_{A,<\epsilon}=K^M_A\cap(\No_{<\epsilon})^M,\qquad
 I_{A,<\epsilon}=I^M_A\cap(\No_{<\epsilon})^M.
\]
Then $\Frac(I_{A,<\epsilon})=K_{A,<\epsilon}$. For nonzero
$x\in K_{A,<\epsilon}$ with $\alpha=\bd(x)$, the explicit choice
\begin{equation}\label{dsn:eq:ordinal-denominator}
 q=\omega^{\underline{\alpha+1}},\qquad p=qx
\end{equation}
gives $p,q\in\Pi^M\cap K_{A,<\epsilon}$ and $x=p/q$. The rule is definable
from $x$ alone, without $\epsilon$ as an additional parameter. In
particular,
\[
 \Frac(\Core^M_A\cap\Oz^M)=\Core^M_A.
\]
\end{theorem}
\begin{proof}
For each exponent $a$ of $x$, \cref{dsn:lem:birthday-controls} gives
$\bd(a)\le\alpha$, hence $a\ge-\underline\alpha$. Therefore
\[
 a+\underline{\alpha+1}\ge
 -\underline\alpha+(\underline\alpha+1)=1>0.
\]
All exponents of $p$ are positive, so $p$ and $q$ are purely infinite. The
number $q$ is the surreal ordinal corresponding to the ordinal
$\omega^{\alpha+1}$. As $\epsilon$ is an epsilon number and
$\alpha<\epsilon$, we have $\alpha+1<\epsilon$ and
$\omega^{\alpha+1}<\epsilon$. Thus $\bd(q)<\epsilon$. The epsilon-cutoff
field is closed under multiplication, so $p=qx$ also has birthday below
$\epsilon$. Both $p$ and $q$ are definable from $x$, and hence belong to
$K^M_A$. The zero case is $0/1$. The core assertion follows by taking
$\epsilon=\rho^M_A$ when that ordinal exists, and by \cref{dsn:thm:hull}
in the no-omission case.
\end{proof}

\begin{corollary}[Normal-form heredity and valuation of the core \src{07}]\label{dsn:cor:core-normal}
The field $\Core=\Core^M_A$ is closed under $\Omega$ and its inverse on
monomials. Every exponent, coefficient, and monomial term in the normal
form of an element of $\Core$ belongs to $\Core$. Its normal-form value
group is $(\Core,+)$, and its residue field is $R_A=K^M_A\cap\R^M$.
\end{corollary}
\begin{proof}
Suppose first that $\rho=\rho^M_A$ exists. The bound
$\bd(\Omega(a))\le\omega^{\bd(a)}<\rho$ for $a\in\Core$, together with
ambient definability, proves closure under $\Omega$. The opposite birthday
inequality gives closure under its inverse on monomials.

If $x\in\Core$, its normal-form length is at most $\bd(x)<\rho$. Every
index in that length is therefore $A$-definable.
\Cref{dsn:prop:nf-invariant} makes each coefficient and exponent
definable. Their birthdays, and those of the corresponding terms, are
below $\rho$ by \cref{dsn:lem:birthday-controls}; coefficients are real and
hence have birthday at most $\omega$. Thus they lie in $\Core$. In the
no-omission case, every ordinal index is definable and the same reasoning
requires no birthday cutoff. Now repeat the proof of
\cref{dsn:thm:valuation}: leading exponents stay in $\Core$,
$\omega^{-a}\in\Core$ realizes every value $a\in\Core$, and standard part
stays in the common real subfield from \cref{dsn:thm:core}(v).
\end{proof}

\subsection{What the core theorem does and does not say}

The full ambient-definable field need not be initial: a definable high
ordinal can have an undefinable predecessor. Nevertheless it always has
the canonical initial elementary field just constructed. The reverse
normal-form gap of \cref{dsn:thm:reverse-gap} occurs exactly when
$O^M_A$ is not downward closed, equivalently when $K^M_A$ is not initial.

The ordinal $\rho$ is not an additional allowed parameter in the
definition of $K^M_A$. It indexes an externally described intersection.
Calling it ``the first undefinable ordinal'' in the metatheory does not
turn it into a definition inside $M$. This is precisely the distinction
that prevents a Berry-style contradiction.

The fraction-field conclusion required the additional birthday argument of
\cref{dsn:thm:core-fractions}; it was not automatic from initiality or real
closedness. In general an initial real closed subfield need not be the
fraction field of its omnific integer part: the real algebraic subfield has
omnific integer part $\Z$. Here it is the epsilon cutoff and the
definability of the ordinal denominator that remove this obstruction.

The ordinal $\rho^M_A$ and 04's threshold $\lambda_H$
(\cref{dsn:sec:firstbirthday}) are different invariants: $\rho$ is the
first ordinal not definable from $A$, and $\lambda_H$ is the first birthday
of a non-ordinal-definable surreal. With ordinal parameters every ordinal
is definable, so the core of the ordinal-definable field is the whole
field $\KH$, which is initial (\cref{dsn:thm:HOD}).

\begin{proposition}[Failure of finite-parameter continuity \src{07}]\label{dsn:prop:parameter-discontinuity}
Suppose $M$ contains an ordinal $\mu$ which is uncountable in the external
metatheory. Let $A=\{\beta:\beta\le\mu\}\subseteq M$. Then
\[
 \bigcup_{A_0\subseteq A\ \mathrm{finite}}\Core^M_{A_0}
 \ \subsetneq\ \Core^M_A.
\]
\end{proposition}
\begin{proof}
For finite $A_0$, only countably many ordinals are externally
$A_0$-definable. Some ordinal below $\mu$ is therefore omitted, so
$\rho^M_{A_0}<\mu$. The surreal ordinal $\underline\mu$ is not in
$\Core^M_{A_0}$, since its birthday exceeds that cutoff. On the other hand
every ordinal at most $\mu$ is itself an allowed parameter in $A$. Thus the
first omitted ordinal for $A$, if there is one, is greater than $\mu$, and
$\underline\mu\in\Core^M_A$. Parameter monotonicity gives the inclusion,
and this element proves it strict.
\end{proof}

This is a conditional statement about externally uncountable transitive
models. It does not apply merely because $M$ calls an ordinal uncountable.
It shows that, although each individual definition uses finitely many
parameters, requiring \emph{all} prefixes to remain definable is a
genuinely non-finitary condition.

\section{The first birthday of a non-ordinal-definable surreal}
\label{dsn:sec:firstbirthday}

When $V\ne\HOD$, there is a genuine definable first place where
ordinal-definable signs cease to exhaust the construction. It is
not the birthday of a canonically selected nondefinable number.

\begin{definition}[\src{04}]
Assuming $V\ne\HOD$, define
\begin{equation}\label{dsn:eq:firstlambda}
 \lambda_H=\min\{\alpha\in\Ord:
                 \mathcal P(\alpha)\nsubseteq\HOD\}.
\end{equation}
Existence follows from the set-of-ordinals coding in
\cref{dsn:lem:setcoding}. All terms in this definition are
first-order definable classes or set constructions.
\end{definition}

$\HOD$ is a parameter-free definable transitive inner class of $V$
containing every ordinal and satisfying ZFC, so the pair $M=\HOD$, $N=V$
satisfies the standing hypothesis (H) of the universes report, and
$\lambda_H=\beta(\HOD,V)$ in that report's notation.

\begin{theorem}[Exact birthday threshold \src{04}; the case $M=\HOD$ of \lbl{univ:prop:beta}]\label{dsn:thm:firstbirthday}
Assume $V\ne\HOD$. Then
\begin{equation}\label{dsn:eq:firstbd}
 \lambda_H=\min\{\bd(s):s\in\No\setminus\KH\}.
\end{equation}
The ordinal $\lambda_H$ is an infinite cardinal in both $\HOD$ and $V$.
It is itself a parameter-free definable ordinal, and therefore a
parameter-free definable omnific integer.
\end{theorem}
\begin{proof}
The equality \eqref{dsn:eq:firstbd} and the two cardinality assertions are
\lbl{univ:prop:beta} of the universes report
\cite{RepoUniv} (from that report's source 08), applied with $M=\HOD$ and
$N=V$ and combined with $\KH=\No^{\HOD}$ (\cref{dsn:thm:HOD}); that report
also proves $\delta(\HOD,V)\le\operatorname{cf}(\lambda_H)$ for its
omitted-cut invariant $\delta$. Its proof is the argument source 04 gives:
signs of length $\alpha$ and subsets of $\alpha$ are interdefinable from
$\alpha$; a $\HOD$ bijection with a smaller ordinal would pull a new subset
back; and a $V$-bijection with a smaller cardinal would give a $\HOD$ code
of a well-order of type $\lambda_H$. The birthday-cutoff report proves the
equality and the inner-model cardinality for transitive models $M\subsetneq N$
of ZFC with the same ordinals as \lbl{hset:thm:firstbirthday}, which covers
$\HOD^N\subseteq N$ for a transitive set model $N$. The remaining assertion is 04's:
\eqref{dsn:eq:firstlambda} is a parameter-free unique definition of
$\lambda_H$, since $\HOD$ is parameter-free definable, and ordinal surreals
are omnific.
\end{proof}

\begin{corollary}[No definable witness selection \src{04}]
The set of non-ordinal-definable signs of length $\lambda_H$ is a
nonempty parameter-free definable set, but none of its members is
ordinal definable. In particular, there is no ordinal-definable
selection of one of those signs, or of one of their omnific codes.
\end{corollary}
\begin{proof}
The set is defined using the definable predicate $\HOD$ and the
parameter-free ordinal $\lambda_H$. Its nonemptiness follows from
the definition of that ordinal. A selected ordinal-definable member
would contradict the defining condition. An ordinal-definable code
would yield such a member by the decoder.
\end{proof}

\subsection{A current set-theoretic restriction on the threshold}

The threshold is the first cardinal where $V$ and $\HOD$ disagree about
subsets. It therefore inherits restrictions already proved in set
theory; these are \emph{not} new surreal theorems in substance.
Benhamou, Cummings, Goldberg, Hayut, and Poveda prove that a singular
strong-limit cardinal of uncountable cofinality cannot be the first
such disagreement, and, from a measurable cardinal, obtain a model
whose first disagreement is $\aleph_\omega$
\cite{BCHGP2026}. Hence their first theorem gives:

\begin{corollary}[Translation of a cited preprint theorem \src{04}]
If $\lambda_H$ is singular and strong limit, then
$\operatorname{cf}(\lambda_H)=\omega$.
\end{corollary}
\begin{proof}
An infinite singular cardinal has infinite cofinality. Uncountable
cofinality is excluded by the cited compactness theorem, so its
cofinality must be $\omega$.
\end{proof}

The preprint was submitted on August 25, 2026. It is important not to
replace the cardinal conclusion of \cref{dsn:thm:firstbirthday}
by an unsupported regular-cardinal claim: the cited relative
consistency result explicitly permits a singular first threshold.
Source 04 credited only ``the older repository discussion of the first new
birthday across universes'' in the birthday-cutoff report; the exact
theorem is \lbl{univ:prop:beta}, and the same statement for transitive
same-ordinal models is \lbl{hset:thm:firstbirthday}.

\section{Complexity-bounded definability and escape ordinals}
\label{dsn:sec:complexity}

The unrestricted phrase ``all parameter-free definable numbers'' is
external. A fixed syntactic complexity bound, by contrast, admits an
internal treatment. This gives a precise version of the familiar
intuition that definitions of bounded complexity can be diagonalized
against without creating a paradox.

\subsection{The bounded hierarchy \src{04}}

Fix a standard natural number $n\geq1$. Use a standard syntactic
coding of formulas in the L\'evy hierarchy, with padding so that
$\Sigma_n$ formulas are also admitted at every higher level. There
is a partial truth formula $\Sat_{\Sigma_n}(e,x)$ for this fixed
fragment. The existence of each fixed-fragment truth definition is
standard; it does not provide one truth formula for all values of
$n$ simultaneously \cite{Jech,HLR}.

\begin{definition}[\src{04}]
Let $S_n$ be the collection of unique solutions of parameter-free
$\Sigma_n$ formulas, and put $B_n=S_n\cap\No$. More explicitly,
$x\in S_n$ means that some code $e$ of a formula with one free
variable satisfies
\[
 \Sat_{\Sigma_n}(e,x)\ \wedge\
 \forall y\bigl(\Sat_{\Sigma_n}(e,y)\longrightarrow y=x\bigr).
\]
Each formula is required to have a unique solution among all sets,
not just among the surreals.
\end{definition}

\begin{lemma}[\src{04}]
For every fixed standard $n$, both $S_n$ and $B_n$ are
parameter-free definable sets of ZFC.
\end{lemma}
\begin{proof}
For a formula code having a unique solution, send the code to that
solution; send every other code to a fixed default set. The
fixed-fragment truth formula defines this total function on a subset
of $\omega$, so Replacement gives its range as a set. Separation
removes the default value when necessary and restricts to actual
unique solutions. A further Separation restricts to surreal codes.
All definitions use the fixed standard integer $n$ only through its
finite numeral, so there are no free parameters.
\end{proof}

\begin{theorem}[Definability-complexity escape \src{04}]\label{dsn:thm:sigmaescape}
For every fixed standard $n\geq1$, let
\begin{equation}\label{dsn:eq:sigmaescape}
 \delta_n=\sup\{\bd(x)+1:x\in B_n\},
 \qquad d_n=\text{the all-plus surreal of length }\delta_n.
\end{equation}
Then $d_n$ is a parameter-free definable omnific integer, but
$d_n\notin B_n$. There is a standard $m>n$ with $d_n\in B_m$.
Consequently the hierarchy $(B_n)$ does not eventually stabilize.
\end{theorem}
\begin{proof}
Replacement gives the set of birthdays and the ordinal supremum in
\eqref{dsn:eq:sigmaescape}. The all-plus sequence is then uniquely defined,
so $d_n$ is parameter-free definable and is an ordinal surreal.
If $d_n\in B_n$, the definition of the supremum would imply
\[
 \bd(d_n)+1=\delta_n+1\leq\delta_n,
\]
which is impossible. On the other hand, the parameter-free formula
just given for $d_n$ is a particular finite formula. After conversion
to prenex form and adding dummy quantifiers, it belongs to some
$\Sigma_m$ fragment. Thus $d_n\in B_m$. By nesting and its failure
to lie in $B_n$, we must have $m>n$.
\end{proof}

\begin{remark}[No unproved complexity estimate \src{04}]
The theorem does not assert $m=n+1$, a closed expression for
$\delta_n$, or a computable method that outputs the full transfinite
sign sequence of $d_n$. The syntactic complexity of the normal-form,
birthday, and sign-coding formulas depends on coding conventions.
Only existence of a finite higher level is used. Nor is the sequence
of escape ordinals asserted to be cofinal in all ordinals.
\end{remark}

\subsection{Why the separate definitions cannot be collected uniformly}

\begin{theorem}[Noncollection \src{04, 07}; bounded clause for arbitrary $A$ \mergetag]
\label{dsn:thm:noncollection}
Let $M$ be a model or the universe, and $A$ an allowed parameter family.
\begin{enumerate}[label=\textup{(\roman*)}]
\item There is no ambiently $A$-definable set $S\subseteq\No$ such that
$K_A\subseteq S$. In particular, no ambiently $A$-definable map
$\omega\to\No$ has range containing all of $K_A$ \src{04, 07}.
\item No ambiently $A$-definable set contains every $A$-definable
omnific integer in $(0,\omega)$ \src{07}.
\end{enumerate}
\end{theorem}
\begin{proof}
(i) If such an $S$ were given, the simplest strict upper bound
$u=\cut S\varnothing$ is uniquely definable from $S$, so $u\in K_A$, but
$u>S$, so $u\notin S$ (07). Source 04 uses instead the all-plus surreal of
the least ordinal strictly above all birthdays in $S$, whose birthday
exceeds that of every member of $S$. For a definable map on $\omega$,
Replacement gives an $A$-definable range set, reducing to the first
assertion.

(ii) Restrict a putative set $S$ to $\ran(\Mc)$ and apply $\Mc^{-1}$
pointwise. Separation and Replacement produce an $A$-definable set of
surreals $S'=\{\Mc^{-1}(z):z\in S\cap\ran(\Mc)\}$. Every $x\in K_A$ has
$\Mc(x)\in I_A\cap(0,\omega)$ by \cref{dsn:thm:bounded-code}, so the assumed
property of $S$ would give $x\in S'$. This contradicts (i).
\end{proof}

Source 07 stated the theorem for a finite parameter tuple $p$ of a
transitive $M$; the proof uses only that the sets involved are defined from
finitely many allowed parameters, so it holds for any $A$. Thus external
countability does not even give an internally definable list of the
definable omnific integers in this bounded interval. The theorem concerns
\emph{sets} of $M$, not a universal assertion that $K_A$ cannot be a
definable class: in a pointwise definable model it is simply $\No^M$,
which is a definable class (07).

\begin{corollary}[Failure of a uniform bounded-truth collector \src{04}]
\label{dsn:cor:nouniformBn}
In an $\omega$-standard model of ZFC there is no single
parameter-free formula that, for each $n\geq1$, uniquely defines
the set $B_n$ above, uniformly in $n$.
\end{corollary}
\begin{proof}
A uniform formula would, by Replacement, give a parameter-free set
$B=\bigcup_{n\geq1}B_n$. Every parameter-free definable surreal is
uniquely defined by a finite formula and hence belongs to some
$B_n$. Thus $B$ would contain $K_\varnothing$, contrary to
\cref{dsn:thm:noncollection}. The $\omega$-standardness
hypothesis ensures that internal natural indices are precisely the
standard finite levels discussed here.
\end{proof}

This is the exact point at which the false ``least undefinable
surreal'' argument tries to use more truth than it has. Separate
formulas defining each fixed $B_n$ do not constitute one internal
formula defining the sequence of all these sets. The argument also
does not claim that the membership class of definable elements can
never accidentally be definable: in a pointwise definable model it
is simply the entire universe. What fails there is an internal
surjective set-sized list of those elements, not the external fact
that every one of them has some definition.

\section{Forcing can change the theory without adding reals}
\label{dsn:sec:forcing}

We now give a concrete source of models exhibiting the dichotomy.
The forcing facts are standard, but proofs of the relevant
homogeneity and closure steps are included. The result connects
ambient definability of real numbers to genuinely transfinite
surreal information without conflating the two.

\begin{lemma}[Homogeneity over a constructible ground \src{04}]
\label{dsn:lem:forcingHOD}
Let $W$ be a transitive model of ZFC satisfying $V=L$, and let
$G$ be generic for an ordinal-definable weakly homogeneous forcing
$\mathbb P\in W$. Then, in $W[G]$,
\[
 \HOD^{W[G]}=W.
\]
The assertion has the usual relative-model interpretation; it does
not assume that a generic filter over the actual universe is a set
in that universe.
\end{lemma}
\begin{proof}
First let $Y\subseteq\theta$ be ordinal definable in $W[G]$.
Membership in $Y$ is expressed there by one formula
$\psi(\beta,\bar\alpha)$ with a finite ordinal tuple. For each
$\beta<\theta$, weak homogeneity implies that the top condition
decides $\psi(\check\beta,\check{\bar\alpha})$.
Indeed, if one condition forced the sentence and another forced its
negation, an automorphism could make their images compatible, while
fixing all check names for ordinals. Compatible conditions cannot
force contradictory sentences. Thus
\[
 Y=\{\beta<\theta:1_{\mathbb P}\Vdash
                 \psi(\check\beta,\check{\bar\alpha})\}\in W
\]
by the definability of the forcing relation.

Every $\HOD$ set is coded within $\HOD$ by a set of ordinals: use its
canonical $\HOD$ well-order on a transitive closure and the pointed
membership-relation code. Such a code is itself ordinal definable,
so the previous paragraph places it in $W$. The absolute Mostowski
collapse in $W$ then recovers the original set. This proves
$\HOD^{W[G]}\subseteq W$.
Conversely, forcing does not change the ordinals or the constructible
hierarchy, so $L^{W[G]}=W$. Every constructible set is ordinal
definable in $W[G]$ using its position in the canonical constructible
well-order, and so are all members of its transitive closure.
Therefore $W\subseteq\HOD^{W[G]}$.
These are the standard homogeneity/$\HOD$ arguments \cite{Jech}.
\end{proof}

\begin{theorem}[Prescribed regular first birthday \src{04}]\label{dsn:thm:forcing}
Let $W\models\ZFC+V=L$ be transitive, and let $\kappa$ be an
infinite regular cardinal of $W$. Force with
\[
 \operatorname{Add}(\kappa,1)^W
 =\{p:p\text{ is a partial map }\kappa\to2,
                    \ |\dom(p)|^W<\kappa\},
\]
ordered by extension, and let $V=W[G]$. Then
\begin{enumerate}[label=\textup{(\roman*)}]
\item $\HOD^V=W$ and $\KH^V=\No^W$ as ordered fields;
\item the least birthday of a surreal outside $\KH^V$ is exactly
$\lambda_H^V=\kappa$;
\item $\Cell^V$ contains the explicit ordinal-indexed independent
family of \cref{dsn:thm:dichotomy} over the old field $\No^W$;
\item if $\kappa>\omega$, then $\R^V=\R^W$, and every real of $V$
is ordinal definable in $V$, despite \textup{(ii)} and \textup{(iii)}.
\end{enumerate}
\end{theorem}
\begin{proof}
For two conditions $p,q$, flip the bits on the set of common-domain
coordinates at which they disagree. This is a ground-model
automorphism sending $p$ to a condition compatible with $q$.
Thus the forcing is weakly homogeneous and ordinal definable, so
\cref{dsn:lem:forcingHOD} applies. The field identification follows
from \cref{dsn:thm:HOD} and the absoluteness of canonical surreal
arithmetic.

Regularity of $\kappa$ makes the union of a descending chain of
length less than $\kappa$ a condition. Given a name for a sequence
of ordinals of length $\mu<\kappa$, recursively strengthen a
condition to decide each coordinate, taking lower bounds at limit
stages and at the end. The resulting ground-model sequence is
forced to equal the name. Hence the forcing adds no subsets of any
$\mu<\kappa$. For $\kappa=\omega$, the corresponding assertion
about finite subsets is immediate. This closure also prevents a
collapse of the cardinal $\kappa$ by a shorter sequence.

The union $g=\bigcup G$ is a total subset of $\kappa$, and is new:
for every ground-model $a\subseteq\kappa$, the set of conditions
already disagreeing with $a$ at some coordinate is dense. Since
$\HOD^V=W$, the first non-$\HOD$ subset therefore occurs exactly at
$\kappa$. \Cref{dsn:thm:firstbirthday} proves (ii).
The new sign sequence with positive positions $g$ lies outside
$\KH^V$; normalize its absolute value as in
\cref{dsn:thm:amplify} to obtain the parameter $t$ for (iii).

When $\kappa>\omega$, the no-new-short-sequences conclusion in
particular gives no new subsets of $\omega$, hence no new reals.
All old reals belong to $W=\HOD^V$, which proves (iv).
\end{proof}

\begin{example}[All reals ordinal definable, but many omnific integers not \src{04}]
Take $\kappa=\omega_1^W$ in \cref{dsn:thm:forcing}. In the
extension, every real is ordinal definable; every surreal of birthday
less than $\omega_1^W$ is in the $\HOD$-surreal field; and a non-$\HOD$
surreal first appears at birthday $\omega_1^W$. Nevertheless a
single fixed interval of purely infinite omnific integers already
contains independent sets of every internal cardinality over the
old surreal field. This is a strict separation between the real
question and the surreal question, not a consequence of adding an
undefinable real.
\end{example}

\begin{example}[A Cohen real \src{07}]\label{dsn:ex:cohen}
Let a transitive ground universe satisfying $V=L$ be extended by a Cohen
real $c$. In the extension $\HOD=L$ (this is \cref{dsn:lem:forcingHOD} for
Cohen forcing, which is ordinal-definable and weakly homogeneous; the
standard argument concerns $\OD$ subsets of ordinals and does \emph{not}
assert that homogeneous forcing adds no new $\OD$ sets of arbitrary
kind). Consequently
\[
 z_c=\omega^{3/4}+\sum_{n<\omega}(1+\chi_c(n))\,\omega^{1/(n+2)}
\]
is a non-$\OD$ omnific integer in $(0,\omega)$. The displayed support is
parameter-free definable and every coefficient is either $1$ or $2$. This
realizes \cref{dsn:cor:entries-insufficient} concretely, relative to the
existence of the indicated ground model and generic extension. More
generally, if an extension adds a non-$\OD$ subset $X$ of an ordinal
$\alpha$, \cref{dsn:thm:bit-code} creates a non-$\OD$ omnific integer below
$\omega$, even when that forcing adds no real numbers. This is an immediate
conditional application of the code, not a claim that every forcing
extension has the same $\HOD$.
\end{example}

\begin{remark}[What is, and is not, absolute \src{04}]
The forcing theorem compares old sign sequences and their arithmetic.
It does not assert that every old parameter-free definition keeps
its truth or uniqueness in an extension. An ambient definition can
mention the continuum, a generic object, or the set-theoretic
universe itself. Its truth can change even when the particular sign
sequence and all arithmetic performed on old inputs stay unchanged.
Nor is the relation $V=\HOD$ preserved by arbitrary forcing.
\end{remark}

\section{Relativization and a surcomplex companion}
\label{dsn:sec:relativize}

\subsection{An allowed set of ordinal information \src{04}}

Let $p$ be a fixed set of ordinals. Write $\OD(p)$ for definitions
using $p$ together with finitely many ordinals, and $\HOD(p)$ for
hereditary membership in this class. Here the parameter is the
\emph{single set} $p$, not permission to use its arbitrary elements
as unrelated parameters. Since $p$ is a set of ordinals, it belongs
to $\HOD(p)$; this is an inner model of ZFC with its corresponding
canonical definable well-order \cite{Jech}.

\begin{proposition}[Relativized theory \src{04}]\label{dsn:prop:relative}
Put $K_p=\No\cap\OD(p)$ and $I_p=\Oz\cap\OD(p)$. Then
\[
 K_p=\No\cap\HOD(p)=\No^{\HOD(p)},\qquad
 K_p=\Frac(I_p).
\]
The field $K_p$ is real closed and $I_p$ is an integer part.
Every member of $\Cell$ is in $\OD(p)$ if and only if
$V=\HOD(p)$. If $V\ne\HOD(p)$, the family
\[
 \omega+\omega^{\,t\omega^{-(\alpha+1)}}\quad(\alpha\in\Ord)
\]
is algebraically independent over $K_p$ for every
$t\in(0,1)\setminus K_p$.
\end{proposition}
\begin{proof}
The hereditary sign-code argument only uses that ordinals and the
fixed sign tags are hereditarily definable; permitting $p$ does not
change it. Unique-output closure gives the field, floor, and fraction
claims. The reversible code uses no extra parameter, while the
membership-graph collapse works inside $\HOD(p)$ because that inner
model contains $p$. Finally every normal-form entry of an
$\OD(p)$ surreal is definable from $p$ and ordinal parameters, so
its support lies in $K_p$. The support-subfield amplification proof
applies without change.
\end{proof}

This relativization describes precisely how adding allowed information
changes the theory. It is different from allowing \emph{all} surreal
parameters, which makes every individual surreal trivially definable
by equality to itself.

\subsection{Canonical surcomplex pairs \src{04, 07}}

For completeness, define the surcomplex field by pairs
$\No(i)=\No\times\No$ with $i=(0,1)$ and the usual quadratic
extension operations. This fixed coordinate representation names a
real axis and an imaginary unit.

\begin{proposition}[Definable surcomplex closure \src{04, 07}]
\label{dsn:prop:surcomplex}
In the canonical pair representation, the ambiently $A$-definable
surcomplex numbers are exactly $K_A(i)$, and they form an algebraically
closed field. The ordinal-definable ones are exactly $\KH(i)$.
If $V\ne\HOD$, the real omnific family in
\cref{dsn:thm:dichotomy} remains algebraically independent over
$\KH(i)$ inside $\No(i)$ \src{04}.
\end{proposition}
\begin{proof}
Pairing and the two coordinate projections are parameter-free
set-theoretic operations. Thus a pair is $A$-definable precisely when
both coordinates are. The quadratic extension of a real closed field
by $i$ is algebraically closed. For the independence assertion,
write a purported complex-coefficient polynomial relation as
$P_0(\bar z)+iP_1(\bar z)=0$, with polynomials over $\KH$.
At real inputs both values are real, so both are zero; a nonzero
relation would contradict real algebraic independence.
\end{proof}

There is no conflict with complex conjugation or with the failure of
a pure field to choose an imaginary unit. The proposition uses
\emph{ambient canonical coordinates} with their distinguished orientation.
An abstract algebraically closed-field language without those coordinates
has a different definable closure, just as the abstract real field does;
coordinate-based definitions are not silently transferred to that reduct.

\section{Changing the language changes the answer}
\label{dsn:sec:languages}

The pure ordered field is \cref{dsn:thm:purefield}.

\subsection{Adding the ordinary exponential still does not define omega \src{07}}

\begin{proposition}[\src{07}]\label{dsn:prop:exp-no-omega}
In the exponential-field reduct,
\[
 \dcl_{\No^M_{\exp}}(\varnothing)\subseteq\R^M.
\]
In particular $\omega$ is not parameter-free definable in that reduct.
\end{proposition}
\begin{proof}
The embedded real exponential field is elementary in the surreal
exponential field by the classical theorem recalled in
\cref{dsn:sub:imported}. A parameter-free formula having a unique solution
in the larger structure must have a solution in the elementary
substructure, and uniqueness makes these the same solution. Thus the
solution is real.
\end{proof}

We make no attempt to identify the exact parameter-free definable closure
in the real exponential field. In particular, this argument supplies no
conclusion about whether any particular familiar transcendental real
beyond those directly named by the language is definable there.

\subsection{Simplicity already defines omega \src{07}}

\begin{proposition}[\src{07}]\label{dsn:prop:simplicity-omega}
In the structure $(\No^M;<,\ps)$, the ordinal subclass and the element
$\omega$ are parameter-free definable.
\end{proposition}
\begin{proof}
The condition
\[
 \operatorname{OrdSign}(x):\quad
 \forall y\,(y\ps x\longrightarrow y<x)
\]
holds exactly for all-plus sequences. It holds for an ordinal because each
proper prefix is smaller. If $x$ has a minus sign, the prefix ending just
before that sign is greater than $x$, so the condition fails. On this
definable ordinal subclass, the numerical order is the usual ordinal order.
A nonzero ordinal is a limit exactly when it has no greatest ordinal
predecessor; this is first-order expressible by
\[
 x>0\quad\wedge\quad
 \forall y\bigl(\operatorname{OrdSign}(y)\wedge y<x
 \longrightarrow
 \exists z\, (\operatorname{OrdSign}(z)\wedge y<z<x)\bigr).
\]
The least such ordinal is $\omega$.
\end{proof}

The formula $\operatorname{OrdSign}$ is not new: the birthday-cutoff report
uses the same formula, ``every proper prefix of $x$ lies below $x$'', in
the proof of \lbl{hset:prop:nondef}, which shows that neither birthday nor
prefix simplicity is definable in the ordered-field reduct or any o-minimal
expansion. Source 07 did not credit it. Much stronger interpretations
become possible after adding birthday structure: the birthday-cutoff report
develops birthday-enriched interpretations and bi-interpretations
\cite{RepoHset}, and Chen--Hamkins--Yang's related work has been publicly
announced by Hamkins \cite{CHY}. None of those interpretation results is
needed to prove the coding and core theorems above, and none is asserted
here as a new result.

\subsection{The field with the omnific predicate: what the collection proves \mergetag}

Both sources asked what becomes definable in
$\mathfrak S=(\No;+,-,\cdot,<,0,1,\Oz)$, the field with only the omnific
predicate added (04's Question 1, 07's question on intermediate
languages). At their pin the collection did not address this. It now
contains results that largely answer it; we record them with proofs from
the cited theorems.

\begin{proposition}[The pair $(\No,\Oz)$ \mergetag]\label{dsn:prop:pair}
\leavevmode
\begin{enumerate}[label=\textup{(\roman*)}]
\item $\dcl_{\mathfrak S}(\varnothing)\subseteq\R$. In particular no
infinite or infinitesimal surreal, and so neither $\omega$ nor any code
value of \cref{dsn:sec:code}, is parameter-free definable in $\mathfrak S$.
\item $\R$ and $\Z$ are parameter-free definable subsets of $\mathfrak S$.
Hence every real that is parameter-free definable in
$(\R;+,\cdot,<,\Z)$ belongs to $\dcl_{\mathfrak S}(\varnothing)$; since
that structure interprets second-order arithmetic, these include every real
whose binary expansion is parameter-free definable in second-order
arithmetic.
\item For every set $P\subseteq\No$, neither the class of Conway monomials
nor the graph of $x\mapsto\omega^x$ is definable in $\mathfrak S$ with
parameters from $P$.
\end{enumerate}
\end{proposition}
\begin{proof}
(i) Every field automorphism of $\No$ preserves the order, since the
positive elements are the nonzero squares; so a field automorphism
preserving $\Oz$ is an automorphism of $\mathfrak S$. Fix a formula
$\varphi(v)$ with a unique solution $x$ in $\mathfrak S$. By induction on
the subformulas of this fixed formula, every class automorphism $\sigma$ of
$\mathfrak S$ preserves $\varphi$, so $\sigma(x)=x$. By
\lbl{opa:thm:fixed} of the automorphisms report \cite{RepoOpa}, the common
fixed field of all field automorphisms of $\No$ preserving $\Oz$ is $\R$.
Hence $x\in\R$. (The dilation $\mathsf S_2$ of \lbl{odg:def:ex:dilation},
$\sum_g a_g\omega^g\mapsto\sum_g a_g\omega^{2g}$, already suffices: it is a
field automorphism preserving $\Oz$, and it fixes $x$ only if
$2\supp x=\supp x$; for a reverse well-ordered set this forces
$\supp x\subseteq\{0\}$, since a nonzero member $g$ would give the
increasing sequence $g,2g,4g,\ldots$ or $g,g/2,g/4,\ldots$ inside the
support.)

(ii) By \lbl{odg:def:thm:realrecovery} and \lbl{odg:def:cor:internal} of the
omnific Diophantine report \cite{RepoOdg}, the pure ring $\Oz$
parameter-free interprets $\No$ as its fraction field and defines in it the
subfield $\R$ by a formula whose quantifiers range over $\Oz$. In
$\mathfrak S$ the map $(a,b)\mapsto a/b$ on $\Oz\times(\Oz\setminus\{0\})$
is definable, so relativizing that formula to $\Oz$ defines $\R\subseteq\No$
without parameters. Then $\Z=\R\cap\Oz$. A parameter-free definition in
$(\R;+,\cdot,<,\Z)$ relativizes to one in $\mathfrak S$. That structure
defines $\N$ and codes subsets of $\N$ by binary expansions of reals, which
gives the last clause. (The pure ring $\Oz$ already defines $\Z$ and $\N$,
\lbl{odg:def:cor:arithmetic}.)

(iii) This is \lbl{opa:thm:parameters} of the automorphisms report.
\end{proof}

So among ambiently definable surreals only reals can become parameter-free
definable after adding the omnific predicate, and the codes of
\cref{dsn:sec:code}, which recognize ordinal-indexed monomials and
normal-form coefficients through set theory, have no counterpart there.
The exact identification of $\dcl_{\mathfrak S}(\varnothing)$ inside $\R$ is
not known; it is squeezed between the reals definable in
$(\R;+,\cdot,<,\Z)$ and $\R$. The expansion by the omega map is not
addressed by these results. \Cref{dsn:q:languages} is re-scoped
accordingly.

\section{Arithmetic of the definable omnific ring}
\label{dsn:sec:arithmetic}

These results of source 07 apply to $I=I_A$ and, with the same proofs, to
the $\OD$ omnific class. Their full-class antecedents are in the omnific
Diophantine and quotient reports \cite{RepoOdg,RepoOsq}, named at each
statement below; the arguments are included to verify that the required
division and support operations preserve ambient definability. They are
restricted duplicates, not new results.

\begin{proposition}[The constant-term retraction \src{07}]\label{dsn:prop:ct}
There is a split exact sequence of additive groups and a ring retraction
\[
 0\longrightarrow\Pi_A\longrightarrow I_A
   \xrightarrow{\ \ct\ }\Z\longrightarrow0,
 \qquad I_A=\Pi_A\oplus\Z.
\]
The ideal $\Pi_A$ is divisible as an additive group. For every positive
ordinary integer $n$,
\begin{equation}\label{dsn:eq:nI}
 nI_A=\Pi_A\oplus n\Z,\qquad I_A/nI_A\simeq\Z/n\Z.
\end{equation}
Furthermore $\Pi_A^2=\Pi_A$.
\end{proposition}
\begin{proof}
The purely infinite part and the constant coefficient are uniquely
definable from the normal form, so both remain in $K_A$. Products of terms
with positive exponents still have positive exponents; therefore constant
term is a ring homomorphism and $\Pi_A$ is an ideal. For $p\in\Pi_A$,
division by $n$ changes coefficients but not positive support, so
$p/n\in\Pi_A$. This proves divisibility and \eqref{dsn:eq:nI}.

For the final assertion take nonzero $p\in\Pi_A$ with exponent set
$S\subseteq\No_{>0}$. The simplest separator $\delta=\cut{0}{S}$ is
positive and strictly below every exponent of $p$. It is definable from
$p$. Factor $p=\omega^\delta\bigl(p\omega^{-\delta}\bigr)$. Both factors
have strictly positive normal-form exponents and belong to $K_A$, so both
are in $\Pi_A$. Thus $\Pi_A\subseteq\Pi_A^2$; the reverse containment
follows because it is an ideal. The zero case is automatic.
\end{proof}

The full-class statement is \lbl{odg:thm:finitequotients} ($n\Oz=\Pi+n\Z$,
$\Oz/n\Oz\cong\Z/n\Z$).

\begin{theorem}[All finite quotients \src{07}]\label{dsn:thm:finite-quotients}
Every unital ring homomorphism from $I_A$ to a finite ring factors through
$\ct:I_A\to\Z$. Every finite-index ideal of $I_A$ is $nI_A$ for a unique
positive integer $n$. Consequently the profinite completion of $I_A$ is
canonically $\widehat{I_A}\simeq\widehat\Z$. Also $\ct$ is the unique
unital ring homomorphism $I_A\to\Z$.
\end{theorem}
\begin{proof}
Let $f:I_A\to R$ and suppose the additive group of $R$ has finite exponent
$m$. For $p\in\Pi_A$ choose $q\in\Pi_A$ with $p=mq$. Then
$f(p)=mf(q)=0$. Thus $f$ factors through $I_A/\Pi_A\simeq\Z$.
If an ideal $J$ has finite index, apply this argument to $I_A\to I_A/J$ to
get $\Pi_A\subseteq J$. The image of $J$ in $\Z$ is $n\Z$ for a unique
$n\ge1$. Hence $J=\Pi_A\oplus n\Z=nI_A$. The compatible quotient
identifications in \eqref{dsn:eq:nI} identify the inverse system of all
finite quotients with that of $\Z$. For a homomorphism $f:I_A\to\Z$, the
image of each $p\in\Pi_A$ is divisible by every positive integer because
$p$ is. The only such integer is zero. Again $f$ factors through $\ct$,
and unitality forces its restriction to $\Z$ to be the identity.
\end{proof}

For the full ring these are \lbl{odg:thm:finitequotients},
\lbl{odg:cor:charideals} and \lbl{odg:eq:profinite} of the omnific
Diophantine report, and the much stronger universal theorem
\lbl{osq:thm:universal} of the quotient report, which factors every
homomorphism into a set-sized ring through $\ct$.
Related (batch 34): the surcomplex-automorphism report proves the
finite-quotient statement, and the transfer of solvability of integer
polynomial systems through $\ct$, uniformly for its integer-part rings $R_F$
over set-sized real closed fields $F$, with $\Oz=R_{\R}$ and $2^{\mathfrak c}$
pairwise nonisomorphic rings among them (\lbl{saut:gr:prop:finite},
\lbl{saut:gr:prop:dio}); it cites the theorem above for the rings $I_A$.

\begin{proposition}[A size-sensitive nontransfer \src{07}; already in the quotient report]\label{dsn:prop:size-boundary}
For a transitive set model $M$ and finite $A$, the externally countable
ring $I^M_A$ admits a homomorphism to a set-sized ring which does not
factor through constant term.
\end{proposition}
\begin{proof}
Take the identity map $I^M_A\to I^M_A$. Its target is a set externally, and
is countable by \eqref{dsn:eq:count}. It cannot factor through constant
term because $\omega\in I^M_A$, $\omega\ne0$, and $\ct(\omega)=0$.
\end{proof}

There is no contradiction with the full-class theorem. Its source is the
entire proper class of omnific integers in the same universe in which
target size is measured. Here the source is a fragment, and the external
set $I^M_A$ need not be an internal set of $M$. The finite-target theorem
survives because its proof uses divisibility rather than a proper-class
size obstruction. Source 07 presented this ``external size-boundary
warning'' as supplied by it; it is already in the quotient report, in the
paragraph ``The universe-relative reading'' of \lbl{osq:sub:foundations}
(from that report's manuscripts 04, 10 and 11), which observes that a
set-sized fragment viewed in a larger universe has its own identity map as
a set-sized target.

\begin{proposition}[Non-Noetherian behaviour \src{07}]\label{dsn:prop:nonnoetherian}
The ring $I_A$ is not Noetherian. In fact it has the strict ascending chain
\[
 (\omega)\subsetneq(\omega^{1/2})\subsetneq(\omega^{1/3})\subsetneq\cdots.
\]
\end{proposition}
\begin{proof}
Every displayed monomial is parameter-free ambient-definable. Since
$\omega^{1/n}=\omega^{1/(n+1)}\omega^{1/n-1/(n+1)}$ and the second factor is
omnific, the ideals are increasing. Equality would put
$\omega^{1/(n+1)}/\omega^{1/n}=\omega^{-1/(n(n+1))}$ in $I_A$. This is a
nonzero infinitesimal, hence is not omnific. Therefore each inclusion is
strict.
\end{proof}

The full-ring chain $(\omega)\subsetneq(\omega^{1/2})\subsetneq
(\omega^{1/4})\subsetneq\cdots$ is \lbl{odg:prop:notnormal}.

\begin{proposition}[Elementary Diophantine transfer \src{07}]\label{dsn:prop:equations}
A finite system of polynomial equations with ordinary integer coefficients
has a solution in $I_A$ if and only if it has a solution in $\Z$. This
equivalence need not persist when nonvanishing conditions are added.
\end{proposition}
\begin{proof}
Apply the ring retraction $\ct$ coordinatewise to any solution in $I_A$.
Conversely every integer solution already belongs to $I_A$. For the
failure with nonvanishing, the equation $X^2=2Y^2$ has the nonzero omnific
solution $X=\sqrt2\omega$, $Y=\omega$, both parameter-free
ambient-definable, whereas the only ordinary integer solution is $(0,0)$.
\end{proof}

The full-ring statements are \lbl{odg:thm:transfer} and
\lbl{odg:def:thm:ideal}, where $\exists y\,x^2=2y^2$ defines $\Pi$ in
$\Oz$. In particular the nonnegative part of $I_A$ is not a model of Peano
arithmetic: PA proves that the displayed equation has no nonzero solution.
This is distinct from open induction (\cref{dsn:cor:openinduction}), the
familiar weaker theory associated with integer parts of real closed fields
\cite{Shepherdson}; no stronger induction assertion is being inferred from
finite-quotient agreement with $\Z$.

\begin{remark}[Relation to \lbl{odg:q:size} \mergetag]\label{dsn:rem:size}
The omnific Diophantine report asks (\lbl{odg:q:size}) which set-sized
exponent groups or universe-relative constructions admit analogues of the
common-monomial and set-sized quotient theorems. The definable fragments
$K_A$ and $I_A$ supply one instance: same-parameter per-element monomial
clearing holds (\cref{dsn:thm:hull,dsn:thm:core-fractions}); finite
quotients remain $\Z/n\Z$ (\cref{dsn:thm:finite-quotients}); and the
set-sized quotient theorem fails through the identity map
(\cref{dsn:prop:size-boundary}, already observed in the quotient report).
This is partial information only. The question is about general set-sized
exponent groups and constructions, and it stays open.
\end{remark}

\section{Foundational boundaries, further questions, and formalization}
\label{dsn:sec:questions}

\subsection{From set models to an intended universe \src{07}}

The set-model convention is a clean semantic setting, not a theorem of ZFC
that such a model exists. All theorems quantified over such $M$ are
conditional statements. The set-theoretic coding constructions themselves,
including $\Code$, $\Two$, $\Mc$, $\Mc_{a,b}$, $E_\alpha(X)$ and $\esc$, are
ordinary ZFC-definable class constructions and do not require the existence
of a set model.

For an intended universe $V$, one may discuss parameter-free definability
externally, or work in a stronger metatheory supplying the relevant truth
apparatus. What one may not do is count all true definitions of $V$ using
an unmentioned truth predicate and then invoke Replacement \emph{inside
that same $V$}. The $\OD$/$\HOD$ part avoids this difficulty because $\OD$
and $\HOD$ have internal first-order definitions. The noncollection theorem
is formulated precisely so that no hidden internal list of definitions is
used.

For nonstandard models of ZFC, \cref{dsn:def:ambient} still uses standard
external formulas, and the closure-under-definitions principle continues to
hold; so do the pointwise test and the coding theorems, whose proofs are
internal. However, the ordinal and sign arguments about a \emph{least
externally omitted ordinal} need not transfer unchanged: the model's
ordinals may not be externally well founded. \Cref{dsn:thm:core} and its
corollaries are not claimed for such models. The pointwise-definability
literature treats nonstandard models as well as transitive cases
\cite{HLR}; these distinctions should be retained when importing its
results.

\subsection{Questions}

The questions below are further projects, not omitted hypotheses or
unfinished steps in the proofs above. They are not asserted to be
previously published open problems.

\begin{question}[Languages between fields and set theory \src{04, 07}; re-scoped]\label{dsn:q:languages}
Determine the parameter-free definable closure of
$\mathfrak S=(\No;+,-,\cdot,<,0,1,\Oz)$ inside $\R$, or of well-specified
set-sized models of its first-order theory. How much of ambient
definability is captured by an expansion naming the omega map, which bounded
coding maps become definable there, and what are its intrinsic definable
closures?
\end{question}
Source 04 asked the first part in the form ``which ambiently definable
surreals become intrinsically definable after adding only the omnific
predicate''; source 07 asked about $(\No;<,+,\cdot,\Oz)$ and the omega-map
expansion. By \cref{dsn:prop:pair} (from the automorphisms and omnific
Diophantine reports) only reals can, all reals definable in
$(\R;+,\cdot,<,\Z)$ do, and the monomials are not definable with set
parameters; the exact real closure and the omega-map expansion remain open.
The pure ordered field defines only real algebraic elements without
parameters (\cref{dsn:thm:purefield}); order plus simplicity already
defines $\omega$ (\cref{dsn:prop:simplicity-omega}); birthday-enriched
interpretation results supply a much stronger endpoint, but do not settle
every intermediate language by themselves.

\begin{question}[Birthday cost of reversible omnific coding \src{04}]\label{dsn:q:cost-birthday}
Obtain sharp upper and lower bounds for $\bd(\Code(s))$ in terms of
$\bd(s)$, and compare alternative reversible codes with the same
leading term and coefficient alphabet. Can a useful optimality
statement be proved for a specified class of such codes?
\end{question}
The construction controls support order type and the leading scale,
not the exact surreal birthday of its nested monomials. No unstated
birthday estimate is used in the independence theorem.
\Cref{dsn:lem:birthday-controls} bounds $\bd(\omega^a)$ by
$\omega^{\bd(a)}$, which is a starting point for upper bounds.

\begin{question}[Intermediate support fields \src{04}]\label{dsn:q:support}
Classify naturally occurring proper subfields satisfying
\eqref{dsn:eq:Khyp}, especially those arising from relative definability
or canonical inner models. Determine when stronger structural
invariants than the independent family can be read from the
quotient additive group $\No/K$.
\end{question}
The support-subfield theorem already gives one invariant uniformly:
a single element outside such a field produces an ordinal-indexed
independent family in one fixed omnific interval. Instances are $\KH$ and
$K_p=\No\cap\OD(p)$ (\cref{dsn:prop:relative}). For a proper same-ordinal
inner model $M$ as in the universes report, $\No^M$ would satisfy
\eqref{dsn:eq:Khyp} if the omega map and normal forms were absolute between
$M$ and $V$; \lbl{univ:prop:absolute} proves absoluteness of the arithmetic
and of cuts only, so this instance is not claimed from that report. For
inner models with the same reals it is now supplied by the large-cardinal
report (\cref{dsn:rem:largecardinal}).

\begin{question}[Possible first omitted ordinals \src{07}]\label{dsn:q:rho}
For which transitive $M$ and parameter sets $A$ can a specified epsilon
number occur as $\rho^M_A$? What additional closure constraints follow
from definable ordinal functions beyond finite Cantor normal form?
\end{question}
There are immediate further constraints: $O^M_A$ is closed under every
parameter-free definable ordinal-valued function on definable inputs. For
example, if an ordinal index is definable, its indexed epsilon number is
definable. Determining the exact resulting external cut is a set-theoretic
problem, not a matter of adding a few field operations.

\begin{question}[Effective syntactic costs \src{07}]\label{dsn:q:cost}
After fixing a proof system and a precise coding of formulas, how
efficiently can one translate definitions through $\Mc$, $E_\alpha$, and the
denominator rule? Which translations preserve a chosen level of the
set-theoretic formula hierarchy?
\end{question}
Interdefinability is proved here by fixed formula substitution. It does
not automatically preserve quantifier complexity, and no optimal
description-length theorem is asserted. This question links the present
semantic theory to the collection's separate work on computable surreals
without identifying the two notions.

\begin{remark}[The same-reals instance, from the large-cardinal report]
\label{dsn:rem:largecardinal}
The report
\fn{docs/foundations-and-computation/large-cardinal-embeddings-and-normal-forms}
proves, in its absoluteness lemma (\lbl{lce:lem:absolute}, from three
manuscripts of batch 30), that for transitive inner classes $M\subseteq V$ of
$\ZFC$ with the same ordinals and the same reals the inclusion
$\No^M\subseteq\No$ preserves the field operations, the omega map and
canonical normal forms. Its \lbl{lce:rem:dsnsupport} draws the consequence
for \cref{dsn:q:support}: if moreover $M\ne V$, then $K=\No^M$ satisfies
\eqref{dsn:eq:Khyp} and is proper (a transitive inner model of $\ZFC$
containing every set of ordinals is $V$), so
\cref{dsn:thm:amplify} applies to it. The targets of elementary embeddings
$j:V\to M$ with a critical point, such as the ultrapower by a normal
measure, are instances. Inner models that lack some real, such as $L$ when
$\R\not\subseteq L$, or $\HOD$ when it lacks a real, are not covered, and the
classification asked in \cref{dsn:q:support} stays open. This remark, and
the sentence after \cref{dsn:q:support} that points here, were added when
the large-cardinal report was written; nothing else in this report changed.
\end{remark}

\begin{writenote}[Status of the questions after batch~93, 4 October 2026]
Parts~\ref{dsn:part:rp} and~\ref{dsn:part:op} bear on four of these
questions; all stay open as re-scoped here.
\emph{\Cref{dsn:q:languages}.} Its omega-map half is answered at the level
of an interpretation: in $(\No;+,-,\cdot,<,0,1,\Oz,\Omega)$ every normal-form
coefficient and truncation is parameter-free definable
(\cref{dsn:op:thm:support}), and $(V,\in)$ has a parameter-free quotient
interpretation with representatives in the bounded monomial interval
(\cref{dsn:op:thm:V-interpretation}); source~02's
\cref{dsn:rp:prop:omega-code} is the case of the unary code. The definable
closures of that expansion, of $\mathfrak S$ and of the language with the
omega map alone, and a bi-interpretation, remain open
(\cref{dsn:op:q:dcl,dsn:op:q:biinterpretation} and the questions of
\cref{dsn:rp:sec:agenda}).
\emph{\Cref{dsn:q:support}.} The instance this remark leaves open is
settled: for every transitive inner model $M\ne V$ of $\ZFC$, with or
without all reals, $\No^M$ satisfies \eqref{dsn:eq:Khyp}, because canonical
operations and normal forms are absolute without a same-reals hypothesis
(\cref{dsn:op:thm:all-reals-free-absoluteness,dsn:op:cor:inner-support-fields}),
and \cref{dsn:thm:amplify} applies (\cref{dsn:op:cor:inner-independent});
so do $L$ and $\HOD$ when they lack reals. The real-parameter field $\KR$ is a
further instance when $V\ne\HOD_{\R}$ (\cref{dsn:rp:rem:support}). The
classification stays open (\cref{dsn:op:q:support}).
\emph{\Cref{dsn:q:cost}.} Answered in part: the canonical operations,
normal-form evaluation, $\Oz$ and the floor have $\Delta_1$ graphs over ZFC,
and $x$ and $\Mc(x)$ have the same $\Sigma_n$, $\Pi_n$ and $\Delta_n$
singleton status
(\cref{dsn:op:thm:canonical-delta-one,dsn:op:prop:definition-complexity-transfer,dsn:op:cor:code-levels}).
Source~06's proposed additive description-length bound still requires the
explicit syntax accounting of \cref{dsn:op:q:length}; optimality and the
other codes also remain open. \emph{Review clarification (4 October 2026):}
the published summary said ``with an additive description-length bound.''
That wording overstated the proof status: the credited source~06 sketch is
retained as \cref{dsn:op:q:length}, not refuted or promoted to a proved bound.
\emph{\Cref{dsn:q:cost-birthday}} is refined by
\cref{dsn:op:q:cost} and by the birthday-cost questions of
\cref{dsn:rp:sec:agenda}, and \cref{dsn:q:rho} is extended by
\cref{dsn:op:q:support}.
\end{writenote}

\subsection{A realistic Lean decomposition}

A formalization should separate the general Hahn-field mathematics
from satisfaction and internal models (04). The shortest algebraic route
is to formalize \cref{dsn:lem:cosets} for set-sized Hahn series and
finite polynomials first: coefficient supports in different cosets
cannot cancel. Next formalize the finite linear independence of
monomials, the normalized-parameter argument, and the finite form
of \cref{dsn:thm:amplify}. These steps require no model-theoretic
truth predicates or proper-class objects.

The set-theoretic track then needs a precise universe convention,
canonical sign codes, unique normal forms, and definable class maps.
$\HOD$ identification requires a model-of-ZFC interface, with standard
external formulas kept distinct from internal formula codes.
The statements about $B_n$ should be separate theorems for each
fixed fragment, not a falsely uniform truth evaluator. Source 07 proposes
three interfaces: a sign/normal-form interface with definable canonical
operations; an external definability-closure interface for set structures;
and a model-theoretic selection interface for o-minimal exponential fields.
The bounded code uses only the first two, while the core theorem
additionally uses the published epsilon-cutoff theorem and the intersection
argument. The repository documents sign, Hahn, size, and normal-form
interfaces, but neither source nor this report maps the new statements to
checked declarations, claims that the required interfaces already exist in
Lean, or claims to have built or extended the Lean code \cite{Repo}.

\section{Conclusion}

There is no single language-independent class of ``the definable
surreals.'' Once the ambient universe and permitted parameters are
fixed, however, the basic algebra is robust: definable surreals form
a real closed field, exponential and elementary over a transitive model,
and definable omnific integers are an integer part with exactly that
fraction field. Ordinal parameters reveal a particularly natural inner
field, the canonical surreal field of $\HOD$.

The strongest results are not counting observations. Explicit normal-form
maps compress every surreal definability degree into a fixed-leading-term
omnific interval, or into single monomials below $\omega$. If $\HOD$ fails
to contain the universe, one missing surreal generates an ordinal-indexed
algebraically independent family in that same interval. The failure of the
definable field to be initial is controlled exactly by the first omitted
ordinal, and the resulting maximal initial fragment is still an elementary
exponential field and the fraction field of its omnific integers. Uniform
coefficient assignments, rather than small coefficient alphabets, numerical
size or ordinary congruences, carry the missing information. These
conclusions coexist with, rather than evade, the traditional cautions about
truth, parameters, $\HOD$, and pointwise definable models.

\appendix
\crefalias{section}{appendix}
\crefalias{subsection}{subappendix}
\section{Provenance, proof scope, and reproducibility}
\label{dsn:app:provenance}

\begin{writenote}[4 October 2026]
This appendix concerns Part~\ref{dsn:part:one}. The provenance, files and
checks of Parts~\ref{dsn:part:rp} and~\ref{dsn:part:op} are in
\cref{dsn:rp:sub:provenance,dsn:op:sub:provenance} and in the README; their
notation is in \cref{dsn:rp:sub:notation,dsn:op:sub:notation}.
\end{writenote}

\subsection{Where each section comes from}

\begin{longtable}{@{}P{0.30\textwidth}P{0.64\textwidth}@{}}
\caption{Provenance of the sections.}\label{dsn:tab:sections}\\
\toprule
Section & Sources\\
\midrule
\endfirsthead
\toprule
Section & Sources\\
\midrule
\endhead
\ref{dsn:sec:intro} Introduction & 04's introduction and 07's scope
section, merged; provenance and conventions added.\\
\ref{dsn:sec:foundations} Foundations & 04 \S2 with 07 \S2 (simplicity,
initial classes, two exponentials, floor prefix, imported cutoff theorem
and birthday controls); \cref{dsn:rem:ominimal} merge.\\
\ref{dsn:sec:definitions} Definability & 04 \S3 with 07 \S3; field dcl
printed once.\\
\ref{dsn:sec:hulls} Field and integer part & 04 \S4 with 07 \S4 (selection,
$\Lexp$ elementarity, cut denominator, valuation, normal-form invariant).\\
\ref{dsn:sec:HOD} HOD & 04 \S5 with 07 \S8.1; proof of
\cref{dsn:prop:notdefinable} corrected.\\
\ref{dsn:sec:code} Codes & 04 \S6 and 07 \S\S5.1--5.3.\\
\ref{dsn:sec:global} Universe tests & 04 \S7 with 07 \S\S5.4, 8.2, 8.3.\\
\ref{dsn:sec:transcendence} Transcendence & 04 \S8.\\
\ref{dsn:sec:uniformity} Uniformity & 04 \S9 with 07 \S6.\\
\ref{dsn:sec:core} Core & 07 \S7.\\
\ref{dsn:sec:firstbirthday} First birthday & 04 \S10, printed by citation of
the universes report.\\
\ref{dsn:sec:complexity} Complexity & 04 \S11 with 07 \S5.5.\\
\ref{dsn:sec:forcing} Forcing & 04 \S12 with 07 \S8.4.\\
\ref{dsn:sec:relativize} Relativization & 04 \S13 with 07 \S9.4.\\
\ref{dsn:sec:languages} Languages & 07 \S\S9.2--9.3; \cref{dsn:prop:pair}
merge.\\
\ref{dsn:sec:arithmetic} Arithmetic & 07 \S10 with credits added.\\
\ref{dsn:sec:questions} Questions & 04 \S14 and 07 \S11.\\
Appendices & 04's appendices and 07's audit section and gallery.\\
\bottomrule
\end{longtable}

\subsection{Claim ledger}

\begin{longtable}{@{}P{0.27\textwidth}P{0.10\textwidth}P{0.23\textwidth}P{0.32\textwidth}@{}}
\caption{Status of the statement families.}\label{dsn:tab:ledger}\\
\toprule
Statement family & Source & Status & Dependency and qualification\\
\midrule
\endfirsthead
\toprule
Statement family & Source & Status & Dependency and qualification\\
\midrule
\endhead
Signs, normal forms, real closedness, omega map, exponential, floor &
04, 07 & Classical & Conway, Gonshor, van den Dries--Ehrlich (with
erratum), L'Innocente--Mantova; elementary proofs supplied where useful.\\
Exponential elementarity and epsilon cutoffs & 07 & Published input &
van den Dries--Ehrlich, erratum, Bournez--Guilmant. Not new.\\
Ambient field and integer-part closure & 04, 07 & Standard-method
consequence & Unique-output closure, ordered finite root selection, and a
uniformly chosen denominator. No priority claim for the closure
principle; $\Frac(\Oz)=\No$ is \lbl{odg:thm:fractions}.\\
$K^M_A\prec\No^M$ in $\Lexp$ & 07 & Application & Fixed-formula translation
and o-minimal definable selection; transitive $M$.\\
Global birthday/set-theory interpretation & 04, 07 & Prior announced work
& Chen--Hamkins--Yang; not used as a black-box dependency and not
claimed.\\
HOD/sign-code identification & 04, 07 & Foundational application &
Hereditariness of a sign code, inner-model recursion, canonical normal
forms. No claim to have introduced $\HOD$-valued surreal fields.\\
Fixed-interval sign code & 04 & Proposed contribution &
\Cref{dsn:thm:code}; search did not locate this precise code.\\
Bounded and local monomial codes & 07 & Proposed contribution &
\Cref{dsn:thm:bounded-code,dsn:thm:local-code}, with two-way parameter
preservation.\\
One-interval universe tests & 04, 07 & Proposed consequences &
\Cref{dsn:thm:globalOD,dsn:thm:pointwise}; standard pointed set coding
followed by the omnific codes. Density form is Hamkins 2012; pointwise
model existence is HLR.\\
Support-subfield amplification and HOD dichotomy & 04 & Principal proposed
contribution & \Cref{dsn:thm:amplify,dsn:thm:dichotomy}; finite
support-coset proof. No transcendence-basis claim.\\
Two uniformity gaps & 04, 07 & Proposed explicit formulations &
\Cref{dsn:thm:termwise,dsn:thm:bit-code,dsn:prop:collector,dsn:thm:reverse-gap};
the general distinction between pointwise and uniform definability is
standard.\\
Maximal initial elementary core & 07 & Principal proposed contribution &
\Cref{dsn:thm:core,dsn:thm:core-fractions,dsn:cor:core-normal}; birthday
estimates and cutoffs are classical inputs; transitive $M$.\\
Bounded birthday escape and noncollection & 07 (04) & Proposed
consequences & \Cref{dsn:thm:bdescape,dsn:thm:noncollection}.\\
First non-HOD birthday & 04 & Duplicate, printed by citation &
\lbl{univ:prop:beta} with $M=\HOD$; also \lbl{hset:thm:firstbirthday}; the
compactness restriction is \cite{BCHGP2026}.\\
Bounded-complexity escape & 04 & Standard-method refinement &
Fixed-fragment truth and Replacement; no all-level truth predicate and no
asserted one-step complexity bound.\\
Forcing without new reals & 04, 07 & Application & Standard homogeneous
closed forcing over $L$, combined with the independent-family theorem.\\
Density iff $V=\HOD$ & 07 & Hamkins 2012 & Credited in
\cref{dsn:thm:hamkins}.\\
Old-field closed nowhere denseness & 07 & Specialization &
\lbl{univ:thm:nowheredense}.\\
Simplicity defines $\omega$ & 07 & Elementary & The formula is that of
\lbl{hset:prop:nondef}.\\
Pair $(\No,\Oz)$ & merge & From the collection & \lbl{opa:thm:fixed},
\lbl{opa:thm:parameters}, \lbl{odg:def:thm:realrecovery}.\\
Finite quotients, constant term, non-Noetherian, equations, size boundary &
07 & Restricted duplicates & \lbl{odg:thm:finitequotients},
\lbl{odg:prop:notnormal}, \lbl{odg:thm:transfer}, \lbl{osq:thm:universal},
\lbl{osq:sub:foundations}.\\
Finite verification & 04, 07 & Regression tests only & Finite symbolic codes,
decoders, formal support-coset and floor checks. They prove no infinite or
set-theoretic theorem.\\
\bottomrule
\end{longtable}

\subsection{Stale statements and corrections}\label{dsn:app:stale}

\begin{enumerate}
\item \emph{``All 51 reports.''} 04's inspection note said that it was not a
review of ``all 51 reports''. At the pin the collection had 51 reports; at
the placement commit it has 53 written reports, and this report was placed
together with two others.
\item \emph{First new birthday.} 04 credited only an ``older repository
discussion'' in the birthday-cutoff report. The exact theorem, with the same
proof, is \lbl{univ:prop:beta}; the birthday-cutoff report has it for same-ordinal
models as \lbl{hset:thm:firstbirthday}. 04's theorem is printed by citation.
\item \emph{Hamkins 2012.} 04 did not cite Hamkins's density/$V=\HOD$
answer, which is close to its equivalence of (i) and (ii) in
\cref{dsn:thm:globalOD}; the credit is added.
\item \emph{Size boundary.} 07 described its size-boundary warning as
``supplied here''. It is in the quotient report
(\lbl{osq:sub:foundations}).
\item \emph{OrdSign.} 07's formula defining the ordinals from order and
prefix is the one in the proof of \lbl{hset:prop:nondef}; credit added.
\item \emph{Cutoff theorem.} 07 cited van den Dries--Ehrlich without the
erratum; the erratum and Bournez--Guilmant are added.
\item \emph{Universes report.} 07's ``Theorem 14.3 in the report'' is
\lbl{univ:thm:nowheredense}. 07's statement that the universes report
``develops this compatibility explicitly'' (for normal forms and the
exponential) is only partly accurate: \lbl{univ:prop:absolute} covers the
arithmetic and cuts.
\item \emph{Restricted duplicates.} 07's arithmetic results are restricted
versions of \lbl{odg:thm:finitequotients}, \lbl{odg:prop:notnormal},
\lbl{odg:thm:transfer}, \lbl{odg:def:thm:ideal} and the open-induction
discussion \lbl{odg:rem:foursquares}; the labels are now named.
\item \emph{Omnific-predicate questions.} Neither source knew the
automorphisms report or the reconstruction section of the omnific
Diophantine report; their language questions are re-scoped
(\cref{dsn:prop:pair,dsn:q:languages}).
\item \emph{Proof gap.} 04's proof that a proper $\KH$ is not intrinsically
definable took the interval endpoints in the subfield
(\cref{dsn:prop:notdefinable}).
\item \emph{Bibliography.} 07 gave Gonshor's book as Lecture Note Series 124;
it is number 110, as 04 has it.
\item \emph{Delivery files.} Both sources described their own archives
(\texttt{.tex}, PDF, README). Those files are not shipped here; the shipped
files are listed in the README.
\end{enumerate}

\subsection{What each source inspected}\label{dsn:app:inspected}

\paragraph{Source 04.} The repository was accessed through its GitHub
connector at the pin. Inspection covered its README, the report catalogue,
the introduction and provenance section of the birthday-cutoff/hereditary-set
report, and the introduction and convention section of the
omnific-Diophantine report. It was not a line-by-line audit of all reports
or their preserved source manuscripts. No repository files were edited, no
commit was made, and no Lean build was performed. The literature comparison
used the supplied definable-real reference, primary works and author
announcements on pointwise definability, $\HOD$, surreal arithmetic,
omnific normal forms, and quantifier elimination; the recent
first-$\HOD$-disagreement preprint was checked as well. The cited works and
the catalogue did not reveal the fixed-interval code or the support-subfield
amplification. This limited negative search evidence cannot establish
publication priority.

\paragraph{Source 07.} The repository was inspected at the pin through its
connected GitHub interface: the README, the catalogue, and the detailed
guides for the universes, birthday and omnific Diophantine reports, whose
cross-references identified the quotient report. These guides identify the
statements, provenance and proof-review boundaries of the earlier
manuscripts; they are not an exhaustive independent review of every proof
or file. The literature search, on 23 September 2026, combined ``definable
surreal'', ``ordinal definable surreal'', ``surreal HOD'', ``definable
omnific'', ``first undefinable ordinal surreal'' and ``definable surreal
initial ordinal'' with the named classical papers. The most relevant
precedent located was Hamkins's 2012 density answer; the epsilon-cutoff
statement was checked against Corollary 5.5 of the original paper. No
matching statement of the coding/core package was located; that supports
describing these as proposed new formulations, and does not certify
priority or exclude an equivalent result under different terminology.

\paragraph{The merge.} Every cross-report label cited here was checked in
the collection at the placement commit. The Bournez--Guilmant abstract was
checked for the epsilon-number characterization.

\subsection{Files and checks}

The source is self-contained apart from ordinary LaTeX packages: its
bibliography is included below, and no external graphics or bibliography
database is required. The two finite scripts shipped with the report are
exact regression checks. Neither decides definability, approximates
$\HOD$, or replaces $\omega$ by a large real number: that numerical
substitution would destroy the exact non-Archimedean ordering relevant to
the proofs.

\paragraph{Source 04's checks.} The script represents the nested exponents
$e_n$ by formal tags. Its recorded run passed 8,191 finite code round trips
through sign length 12, 11 malformed-code rejections, 3,121 exact rational
normalizations of the two-term map $h$ with 3,120 strict order comparisons
and 2 rejected endpoints, 2,157 formal infinitesimal floor cases, and 500
finite support-coset polynomial checks with 55,105 noncancelled formal
support terms. A deliberately dependent exponent example correctly produced
a vanishing polynomial, checking that the independence hypothesis is not
silently omitted.

\paragraph{Source 07's checks.} The script checks finite normal-form bit
encodings, their decoders and marker bounds, rational instances of the
compression/inverse identities, and floor-correction examples. Its recorded
run passed 8,191 bit codes (exhaustive for lengths 0 through 12), 4
malformed-code rejections, 555 rational compression/inverse pairs, 554
adjacent monotonicity comparisons, 1,161 floor-correction cases and 24
support-shift cases. It verifies no transfinite, set-theoretic
satisfaction, $\HOD$, epsilon-cutoff, elementarity, or core statement.

Both scripts were rerun for this merge on copies and reproduced their
recorded results up to line endings.

\section{Non-claims}\label{dsn:app:nonclaims}

Every limitation stated by either source is kept here; most also appear at
the relevant statement.

\subsection{Source 04}
\begin{enumerate}[label=N4.\arabic*]
\item The report is not refereed, not formalized in Lean, and makes no
mapping to Lean declarations; no Lean build was performed.
\item The coding and transcendence results are proposed as original on the
strength of a limited, non-exhaustive search; priority is not established,
and no theorem settles a named open problem in the published literature.
\item $\Code$ is not a ring homomorphism, is not definable in the pure
ordered field, and has no arithmetic compatibility (the code of $0$ is not
$0$).
\item ``$\Ord$-indexed independent'' means only that finite subfamilies are
independent. No transcendence basis, no Global Choice, no class sum, and no
ordinary cardinal of a proper class is asserted; in a set model the
assertions have their internal meaning.
\item $D^V_\varnothing$ is not an internal set or class; there is no truth
predicate for $V$.
\item For nonstandard models, codes are internal objects; the
finite-ambiguity argument concerns genuinely finite sets.
\item The definable closure of $(\No;+,-,\cdot,<,0,1,\Oz)$ is not classified
(see \cref{dsn:prop:pair} for what the collection now proves).
\item Definability is not identified with computability or with the
computable-surreal representations; normal-form uniqueness supplies no
algorithm for infinite codes and no topological convergence of partial
sums.
\item $\Frac(\Oz)=\No$ is not new; the contribution is the uniformity of the
denominator; no priority is claimed for the general closure principle.
\item Termwise $\HOD$ membership is insufficient; $\KH$ is not the Hahn field
of $\HOD$ entries.
\item The widths in the localization are infinite monomial widths; a real
width below $1$ contains at most one omnific integer; no density is
claimed.
\item The $B_n$ escape asserts no $m=n+1$, no closed form, no algorithm, and
no cofinality of the $d_n$; the sequence of the $B_n$ is not uniformly
definable.
\item The BCGHP corollary is a translation of a cited theorem, not a new
surreal theorem; $\lambda_H$ is not claimed to be regular.
\item Forcing: old definitions need not keep their truth or uniqueness;
$V=\HOD$ is not preserved by arbitrary forcing; a generic filter over the
actual universe is not assumed to be a set in it.
\item The surcomplex results use canonical coordinates; the abstract
algebraically closed-field language behaves differently.
\item The questions are not published open problems and not omitted
hypotheses.
\item The finite checks are regression tests only; $\omega$ is never
replaced by a real, and no infinite or set-theoretic theorem is checked.
\item The Chen--Hamkins--Yang announcement is not a black-box dependency; the
set-coding mechanism is not new.
\item No nonclassical Diophantine, quotient, automorphism or
bi-interpretation claim of the repository is a proof dependency.
\item The collector has a nondefinable coefficient only under the stated
hypothesis that some member of $\R\cap\HOD$ is not parameter-free
definable.
\item No birthday estimate for codes is claimed or used.
\item The existence of pointwise definable models is prior work (HLR).
\item No claim to have introduced $\HOD$-valued surreal fields; the
first-birthday threshold is the standard first subset disagreement.
\item The two-term code is an elementary observation, not a major novelty
claim.
\end{enumerate}

\subsection{Source 07}
\begin{enumerate}[label=N7.\arabic*]
\item The main setting is an externally fixed transitive set model; its
existence is not a ZFC theorem, and theorems quantified over it are
conditional.
\item The core theorem and its corollaries are not claimed for nonstandard
models.
\item $D^M_A$ is not claimed to be elementary in $M$; $K^M_A$ need not be
represented by a set of $M$.
\item Provable uniqueness does not give absoluteness (the CH example).
\item There is no algorithm deciding definability; definability is not
identified with computability.
\item The $\Lexp$ results do not cover the restricted-analytic language or
other uncountable languages.
\item $\Mc$ is not onto an interval, is not ordered-field-definable, and uses
$\Omega$, not $\exp$.
\item The birthday escape does not give density.
\item The set-coding lemma gives no canonical definable injection of sets
into $\Oz$; its use of Choice is existential.
\item Noncollection concerns sets of $M$, not classes.
\item The reverse gap is conditional; an internally uncountable ordinal
need not have an undefinable predecessor.
\item $\rho$ is external, not a parameter (the Berry distinction).
\item An initial real closed subfield need not be the fraction field of its
omnific part (real algebraic numbers with $\Z$).
\item The finite-parameter discontinuity needs an externally uncountable
ordinal.
\item Hamkins 2012, the old-field topology result, pointwise-model
existence, birthday bi-interpretations, full-class denominator clearing and
full-class finite quotients are credited, not claimed.
\item There is no exact definable closure for the real exponential field and
no claim about particular transcendental reals.
\item The complexification depends on the chosen orientation and is not
transferred to the pure algebraically closed-field language.
\item The full-class quotient theorem does not transfer to fragments; no
Peano arithmetic and no stronger induction is inferred from finite-quotient
agreement.
\item The questions are not published open problems.
\item The finite checks verify no transfinite, $\HOD$, epsilon-cutoff,
elementarity or core statement.
\item No Lean formalization; priority is not certified; novelty and
correctness are separate, and many arguments are short combinations of
established theorems.
\item Interdefinability by formula substitution need not preserve
quantifier complexity; no optimal description-length theorem is asserted.
\item The Cohen example is relative to the existence of the ground model; the
homogeneity argument does not assert that no new $\OD$ sets of arbitrary
kind are added, nor that every forcing extension has the same $\HOD$.
\item The $\Lexp$ elementarity is proved formula by formula, with no truth
predicate in $M$.
\item The self-similar value group is not a claim that $K_A$ contains every
Hahn series with entries in $K_A$.
\end{enumerate}

\subsection{The merge}
\begin{enumerate}[label=NM.\arabic*]
\item \Cref{dsn:prop:pair} does not identify $\dcl_{\mathfrak S}(\varnothing)$;
the omega-map expansion is not addressed.
\item \lbl{odg:q:size} is not answered; \cref{dsn:rem:size} gives one instance.
\item Support-subfield amplification is not claimed for general same-ordinal
inner models (\cref{dsn:q:support}).
\item The agreement of $\HOD$'s internal omega map, exponential and normal
forms with the ambient ones rests on the sources' recursion argument, not on
the universes report, which covers arithmetic and cuts.
\end{enumerate}

\section{A small gallery of examples \src{07}}

\begin{center}
\renewcommand{\arraystretch}{1.25}
\begin{tabular}{@{}lll@{}}
\toprule
Input or datum & Constructed number & What it illustrates\\
\midrule
$x=0$ & $\Mc(0)=\omega^{1/2}$ & A definable bounded omnific code\\
$x=3/4$ & $\Mc(3/4)=\omega^{4/5}$ & Exact algebraic compression\\
$x=-3/4$ & $\Mc(-3/4)=\omega^{1/5}$ & Order preservation\\
$\alpha=0$, $X=\varnothing$ & $E_0(X)=\omega^{3/4}$ & Marker and empty decoder\\
$\alpha=2$, $X=\{1\}$ & $\omega^{3/4}+\omega^{1/2}+2\omega^{1/3}$ & Finite bit assignment\\
$x=\omega-\omega^{-1}$ & $\lfloor x\rfloor_{\Oz}=\omega-1$ & Negative infinitesimal correction\\
$(a,b)=(\omega,2\omega)$ & $\Mc_{a,b}(x)=\omega+\omega^{\tau(x)}$ & Coding in an infinite-width interval\\
$s=1$ & $\Code(1)=\omega+2\omega^{\omega^{-1}}+3\omega^{\omega^{-2}}$ & The sign code (04)\\
\bottomrule
\end{tabular}
\end{center}

For the second row, $\sqrt{1+(3/4)^2}=5/4$, so $\tau(3/4)=(1+3/5)/2=4/5$.
Every identity is an exact field identity, not a numerical approximation.
The examples involving $\omega$ are normal forms with symbolic exponents
and are not evaluated by substituting a large real number for $\omega$.

\section{Notation and reading guide}

\begin{longtable}{@{}P{0.24\textwidth}P{0.70\textwidth}@{}}
\caption{Notation.}\label{dsn:tab:notation}\\
\toprule
Notation & Meaning\\
\midrule
\endfirsthead
\toprule
Notation & Meaning\\
\midrule
\endhead
$\No$, $\Oz$, $\Pi$ & Canonical surreal field, classical omnific-integer
ring, purely infinite ideal (07's $\mathcal J$).\\[3pt]
$\bd(s)$, $\ps$ & Ordinal length of the canonical sign sequence; proper
prefix.\\[3pt]
$\omega^x=\Omega(x)$ & Conway monomial map; not $\exp(x\log\omega)$.\\[3pt]
$\underline\alpha$ & The surreal all-plus ordinal $\alpha$.\\[3pt]
$\NF(x)$, $\supp(x)$ & Whole indexed normal form; its exponent support.\\[3pt]
$M$, $A$ & A set model of ZFC (transitive where stated) and an external
allowed parameter family.\\[3pt]
$D_A^M$ & External definable closure of $A$ in $(M,\in)$ (07's
$\Def_M(A)$).\\[3pt]
$K_A$, $I_A$, $O_A$ & Ambiently $A$-definable surreals, omnific integers,
ordinals (07's $D_M(A)$, $I_M(A)$, $O_M(A)$).\\[3pt]
$\Pi_A$, $R_A$ & $K_A\cap\Pi$ and $K_A\cap\R$.\\[3pt]
$\KH$, $\IH$ & Ordinal-definable surreals and omnific integers;
$\KH=\No^{\HOD}$.\\[3pt]
$\Cell$, $\Cell_2$ & $\{z\in\Oz:\omega<z<\omega+\omega^{1/2},\ \ct(z)=0\}$
and its two-term part.\\[3pt]
$e_\alpha$ & $\omega^{-(\alpha+1)}$, a positive infinitesimal surreal.\\[3pt]
$\iota_\beta$ & $1/(\underline\beta+2)$ (07's $e_\beta$).\\[3pt]
$\Code$, $\Decode$ & 04's sign code and its partial inverse.\\[3pt]
$\Two$, $h$ & 04's two-term code and its rational compression.\\[3pt]
$\tau$, $\Mc$, $\Mc_{a,b}$ & 07's algebraic compression, monomial code and
local code (07's $C$, $C_{a,b}$).\\[3pt]
$\esc(\alpha)$ & 07's bounded code with birthday above $\alpha$ (07's
$T$).\\[3pt]
$E_\alpha(X)$ & 07's bit code of $X\subseteq\alpha$.\\[3pt]
$\SetDec$, $X_a$ & 04's decoder of sets from signs, and the ordinal code of
$a$ (04's $E$, $B$).\\[3pt]
$\zeta$, $q_*$ & 04's collector map and collector.\\[3pt]
$\rho^M_A$, $\Core^M_A$ & First ordinal of $M$ not $A$-definable (external),
and the maximal initial core (07's $K_M(A)$).\\[3pt]
$\epsilon$, $\varepsilon$ & An epsilon number; an infinitesimal.\\[3pt]
$\lambda_H$ & First ordinal with a non-$\HOD$ subset; the first birthday of
a non-$\OD$ surreal.\\[3pt]
$B_n$, $d_n$ & Fixed-$\Sigma_n$ uniquely definable surreals, and the
definable ordinal surreal escaping that level.\\[3pt]
$\No_{\le\alpha}$ & Surreals of birthday at most $\alpha$ (07's
$B_\alpha$).\\[3pt]
$\mathfrak S$ & $(\No;+,-,\cdot,<,0,1,\Oz)$.\\[3pt]
$\ct$, $\st$, $v$ & Constant term, finite standard part, normal-form
valuation.\\
\bottomrule
\end{longtable}

For the main new algebraic argument alone, read
\cref{dsn:sec:foundations,dsn:sec:HOD,dsn:sec:transcendence}. For the
definability concept and its model-theoretic limitations, read
\cref{dsn:sec:definitions,dsn:sec:global,dsn:sec:complexity}. For the
initial core, read \cref{dsn:sec:hulls,dsn:sec:core}. The codes in
\cref{dsn:sec:code} connect the routes.

% ===================== Part II (batch 93, manuscripts 02 and 05) =====================
\clearpage
\part{Real parameters, countable assembly, and strong symmetries}
\label{dsn:part:rp}

% Part II continues the section numbers of Part I after its appendices.
\setcounter{section}{18}
\renewcommand{\thesection}{\arabic{section}}
\renewcommand{\theHsection}{rp.\arabic{section}}
\crefalias{section}{section}
\crefalias{subsection}{subsection}
\makeatletter\gdef\Hy@chapapp{\Hy@chapterstring}\makeatother
\fancyhead[L]{\small Definable surreals, Part~II: real parameters}

{\large Real parameters, HOD, countable assembly, Hahn summation, forcing, and
the strong stabilizer of the reals and the ordinals\par}
\smallskip
{Merged from two research manuscripts prepared for Vladimir Reshetnikov,
both dated 4 October 2026\par}

\paragraph{Abstract.}
Definability of a surreal number depends both on its language and on its
allowed parameters. This Part distinguishes external parameter-free
definability, real-only definability, ordinal definability, definability
from a fixed real and ordinals, and definability from an arbitrarily
varying real and ordinals. The last regime yields the initial real closed
field $\KR=\bigcup_{r\in\R}\No^{\HOD(r)}$ and its omnific integer part
$\IR$. This union is captured by one real parameter exactly when $\HR$
satisfies Choice, and every surreal is definable from a real and ordinals
only if one real suffices for the whole universe. Countable assembly of
ordinal information, closure of $\KR$ under countable Hahn sums,
interpolation of countable cuts, closure under countable sequences of
definitions, and $\ROD=\COD$ are all equivalent; the test can be made with
one explicit coefficient-$1$ omnific integer below $\omega$, and the
surreals definable from countable ordinal sequences are the least
ambiently stable, countably Hahn-closed enlargement. Cohen forcing gives a
proper such field with an exactly $\omega_1$-long cofinal parameter
hierarchy; a Namba extension over $L$ makes countable closure fail with no
large-cardinal hypothesis, so it is independent of ZFC if ZFC is
consistent; a Prikry extension gives a cardinal-preserving failure from a
measurable cardinal. Finally, the common fixed field of the strong
automorphisms fixing every real and every ordinal is exactly the field of
normal forms with exponents in the rational span of the ordinals, which
answers the strong form of a question of the self-embeddings report.

\begin{writenote}[Added 4 October 2026, batch~93]
This Part merges manuscripts~02 and~05 of batch~93 of ProveIt's
incoming-reports intake (sources and pins in
\cref{dsn:rp:sub:provenance}). Like Part~\ref{dsn:part:one} it is
AI-assisted, unrefereed and not formalized. Manuscript~02, \emph{Definable
Surreal Numbers and Omnific Integers: Real Parameters, Countable Closure, and
Strong Symmetries}, is the base text and is printed in full: \textbf{its
Section~$k$ is Section~$k+18$ here and its statement or equation $k.m$ is
$(k+18).m$}, so its Theorem~5.1 is \cref{dsn:rp:thm:countable}, its
Corollary~4.5 is \cref{dsn:rp:cor:global-single-real} and its Theorem~7.6 is
\cref{dsn:rp:thm:namba}. Manuscript~05, \emph{Definable Surreal Numbers and
Omnific Integers: Real parameters, ordinal information, and countable
assembly}, proves the same central theorems by its own routes and adds the
canonical completion theorem, parameter spectra and a covering criterion.
Its material follows source~02's in the section it belongs to, marked
\src{05}, so source~02's numbers are undisturbed; a theorem proved by both is
printed once, with both sources named and the different proof kept as a
marked second route. \Cref{dsn:rp:tab:corr05} locates every numbered item of
source~05. Source numbers 02, 05 and 06 in Parts~II and~III are batch-93
manuscript numbers; Part~\ref{dsn:part:one}'s sources 04 and~07 are batch-27
manuscripts. Notation changed from the sources is listed in
\cref{dsn:rp:sub:notation}; apart from those changes, the dated
\textbf{[write]} notes and the editorial changes listed in
\cref{dsn:rp:sub:provenance}, no statement, proof, number or non-claim of
source~02 was changed. Throughout, ``this paper'' in source~02's text and
``this article'' in source~05's text mean the respective manuscript.
\end{writenote}

% ===== Begin sections/01_foundations.tex =====
\section{Scope, contributions, and relationship to the repository}
\label{dsn:rp:sec:scope}

There is no single satisfactory meaning of ``a definable surreal number''
until three choices have been made: the ambient universe, the language,
and the parameter policy. For example, $\omega$ has an elementary
parameter-free description in set theory, but cannot be defined from any
real parameters in the pure ordered-field structure of the surreals.
Allowing ordinal parameters avoids one obstruction caused by Tarski's
undefinability theorem, but it does not make every surreal definable in
every universe. Allowing arbitrary real parameters includes every real,
yet can still leave a proper class of surreal numbers outside the
resulting field.

The present paper is an extension of the existing research collection,
not a replacement attribution for it. The ProveIt snapshot used here is
\texttt{3c25fbb5515748040decb853139ed39543885b90}, observed on 4 October 2026. The directly relevant predecessor is
\emph{Definable Surreal Numbers and Omnific Integers}
(Part~\ref{dsn:part:one}). It already develops fixed-parameter definable fields,
their omnific integer parts, the identification with $\HOD$, reversible
bounded coding, transcendence amplification, forcing examples, and
several intrinsic-language obstructions. The reports on universes,
automorphisms, and self-embeddings provide additional antecedents
\cite{RepoUniv,RepoOpa,RepoSse}. The foundational facts used below
also have classical sources \cite{Conway,Gonshor,vdDE,DoraisHamkins}.

\subsection{What is proved here}

The principal extensions established in this paper are as follows.

\begin{enumerate}
\item A classification of the fields $K_r$ and rings $I_r$ obtained from
one real parameter and arbitrary finite ordinal information. Their
inclusion order is exactly ordinal-definability reducibility on real
parameters. The field obtained by allowing the real parameter to vary is
captured by one real precisely when $\HR$ satisfies Choice.
\item An equivalence between countable ordinal information, external
countable-sequence closure, countable Hahn summation, and countable-cut
interpolation. This is stronger than merely placing countably many
individual numbers in a common fixed-real field.
\item An explicit Cohen model in which the parameter fields have an
exactly $\omega_1$-long cofinal chain, all countable cuts are filled, and a
cut with $\aleph_1$ endpoints is omitted. A Namba model gives the opposite countable-closure behavior without
a large-cardinal assumption; a Prikry model adds a cardinal-preserving
example under the stated measurable-cardinal assumption.
\item An exact answer to the \emph{strong} version of the question
\texttt{sse:sr:q:stabilizers} in the self-embeddings report
\cite{RepoSse}: the common fixed field of the strong automorphisms
fixing every real and every ordinal is the full normal-form field over
the rational span of the ordinals. The answer remains the same if
preservation of $\Oz$ is imposed as well.
\end{enumerate}

These are proved mathematical statements with their hypotheses displayed.
The claim of extension is relative to the inspected repository snapshot;
it is not a certification of priority over the entire literature.
In particular, the general Hahn-lifting and support-orbit techniques are
antecedents, not discoveries of this paper. The unconstrained, possibly
nonstrong stabilizer problem remains open here.

\subsection{Proof and formalization status}

Every new theorem has an ordinary mathematical proof below. The main
arguments were independently checked for parameter, model, and support
issues during preparation. No Lean or Rocq verification of the new
theorems is claimed. The source repository itself distinguishes
unformalized research documents from its trusted theorem boundary.
Compilation of this paper checks its presentation, not its mathematics.
The small validation script shipped with the paper checks only explicit
rational inequalities used in one coding construction.


\subsection{What source 05 asks, and what it adds \src{05}}
\label{dsn:ca:sec:intro}

Source~05 organizes the subject around one concrete question:
\begin{quote}
If every term of a countable summable surreal family is definable from
some real and finitely many ordinals, must its sum admit such a definition?
\end{quote}
The answer is not an unconditional yes. Source~05 proves an exact
criterion and exhibits both a proper field satisfying it and an explicit
relative-model counterexample; the counterexample uses only
coefficient-$1$ monomials and produces an omnific integer smaller than
$\omega$. It also stresses the same three choices as source~02, with the
example of $\omega$, which is parameter-free definable in set theory and
hence as a canonical sign sequence, but is not parameter-free definable in
the pure ordered field, whose parameter-free definable elements are the
real algebraic numbers. Naming the set $\R$ as one ambient parameter adds
nothing, since $\R$ is parameter-free definable; allowing an arbitrary
member of $\R$ is very different; allowing ordinal parameters as well is a
third, stronger policy; and allowing a \emph{countable sequence of
ordinals} as one parameter is not the same as allowing a finite tuple of
ordinals.

\paragraph{What source 05 found in the repository.}
Source~05 read Part~\ref{dsn:part:one} at its pin as already establishing
the definable real closed field and its omnific fraction ring, the
identification with $\No^{\HOD}$, fixed-set-of-ordinals relativizations,
reversible bounded omnific codes, whole-family versus termwise
definability gaps, and noncollection results, and it presents none of
these as new. The real-parameter special case follows by taking the fixed
set of ordinals of \cref{dsn:prop:relative} to be a real; what needs
additional analysis is the \emph{union over all real parameters}, above
all its external countable assembly properties. Source~05 used Part~I as a
research predecessor, not as a substitute for proofs or as an
independently certified library. Its description of this report's README
as ``a merged, 55-page report from two manuscripts dated 23~September
2026'' whose formalization mapping is ``marked pending'' was true at its
pin; \textbf{[write]} the report had 56 pages from batch~90 on and now has
the three Parts listed in \cref{dsn:rp:sub:provenance}; the formalization
mapping is still pending.

\paragraph{Relation to the published literature.}
Canonical sign sequences, Conway normal forms and the omnific integer part
are classical; useful modern accounts are \cite{EhrlichKaplanJSL,LM},
alongside \cite{Conway,Gonshor}. Definability from countable ordinal
sequences is classical, not a new notion: Solovay uses it explicitly in
\cite[III.2.2--III.2.6]{Solovay}. In the Levy--Solovay extension, countable
sequences of real--ordinal definable sets are real--ordinal definable; a
recent explicit formulation is \cite[Proposition~2.2(i), version~1]{KL}.
These facts prevent the sources from claiming that countable assembly
itself was previously unnoticed. Source~05's contribution is the
\emph{surreal and omnific equivalence}, its explicit separated-band code,
the precise minimal-closure statement, and the detailed Cohen/Prikry
comparison. The forcing ingredients are standard, but their role is proved
rather than hidden behind a general appeal to homogeneity. Source~05's
literature search was selective; priority for these particular
combinations remains unverified.

\paragraph{Source 05's status statement.}
Complete mathematical arguments are supplied for the stated new bridges
and applications. These are proposed research results, not refereed or
Lean-verified theorems. Classical set-theoretic ingredients and earlier
repository results are credited separately. No claim of historical
priority, a resolved named classical conjecture, or an optimal consistency
bound is made.

\paragraph{Source 05's guide, in this Part's numbering.}
The multiplexer \cref{dsn:ca:thm:code} and the further forms of the
assembly theorem in \cref{dsn:ca:thm:assembly} identify the exact
obstruction to countable sums; the completion
\cref{dsn:ca:thm:completion} supplies the canonical enlargement when the
obstruction is present; \cref{dsn:rp:thm:cohen-atlas} with source~05's
ring form \eqref{dsn:ca:eq:cohenring} describes the real-parameter field
in a Cohen product, and \cref{dsn:ca:thm:cohenconsequences} makes its
properness, countable closure and first omitted birthday explicit;
\cref{dsn:ca:thm:prikry} produces source~05's negative example without
adding reals.

\subsection{Relation to Part I and to neighbouring reports \writetag}
\label{dsn:rp:sub:relation}

The following antecedents in the collection were checked at the write
(4~October 2026) and are credited here; the sources credit some of them
themselves, and the others were written before or after their pins.
\begin{itemize}
\item \Cref{dsn:rp:prop:hereditary-model} combines the classical inner
model $\HOD(\R)$ with the relativized theory of
\cref{dsn:prop:relative} (Part~\ref{dsn:part:one}) for $p$ a real; the
fixed-real fields $K_r$ are that proposition's $K_p$.
\item \Cref{dsn:rp:sec:prikry} strengthens to real-ordinal definability
the ordinal-definable gaps of Part~\ref{dsn:part:one}:
\cref{dsn:thm:termwise}, \cref{dsn:cor:entries-insufficient} and the Cohen
real of \cref{dsn:ex:cohen} already separate termwise from uniform
definability for $\OD$; the $\ROD$ versions are new here.
\item The Prikry gap of \cref{dsn:rp:sec:prikry} is related to the
omitted $(\omega,\omega)$ gap of the old field in Prikry extensions,
\lbl{univ:cor:prikry}~(iii),~(v) of
\repopath{foundations-and-computation/surreal-fields-across-universes}
\cite{RepoUniv}; Namba forcing under CH is already used there
(\lbl{univ:gs:thm:nambainput}, imported from \cite[Section~7.6]{FKW} in
its classical tree presentation, and \lbl{univ:gs:cor:namba}).
\Cref{dsn:rp:thm:namba} uses the Laver-style presentation and is new.
\item In \cref{dsn:rp:sec:strong}, \cref{dsn:rp:cor:dcl-bound} without
the predicate for $\Ord$, and the nondefinability assertions after
\cref{dsn:rp:thm:independence} without that predicate, follow already
(per finite tuple, and more sharply) from
\lbl{opa:sc:cor:innerbound}, \lbl{opa:thm:parameters} and
\lbl{opa:sc:thm:undef}~(ii),~(iv) of
\repopath{surreal/omnific-preserving-automorphisms} \cite{RepoOpa}; see the
dated notes in \cref{dsn:rp:sec:strong}. With a predicate for $\Ord$ they are new.
\cref{dsn:rp:thm:strong-fixed} answers the strong form of
\lbl{sse:sr:q:stabilizers} of \repopath{surreal/surreal-self-embeddings}
\cite{RepoSse}; the nonstrong remainder is \cref{dsn:rp:q:nonstrong}.
\item \Cref{dsn:rp:prop:finite-arithmetic} restricts
\cref{dsn:thm:finite-quotients} and \cref{dsn:prop:equations}
(Part~\ref{dsn:part:one}, themselves restrictions of results of the
omnific Diophantine and set-sized-quotient reports
\cite{RepoOdg,RepoOsq}); the topological proposition of
\cref{dsn:rp:sec:consequences} re-proves the mechanism of
\lbl{univ:thm:nowheredense} for the ZF inner model $\HR$.
\item Part~\ref{dsn:part:op} (source~06) proves
\cref{dsn:rp:cor:global-single-real} independently (its Theorem~4.2,
printed here once, with source~06's second proof and its further
clauses in \cref{dsn:rp:rem:consolidation06}), and its interpretation
theorem \cref{dsn:op:thm:V-interpretation} goes far beyond
\cref{dsn:rp:prop:omega-code}.
\end{itemize}
No statement of this Part is formalized. Its placement beside the
collection's formal developments confers no formal status; no Lean or Rocq
declaration of ProveIt states any of its theorems.

\subsection{Provenance of this Part \writetag}
\label{dsn:rp:sub:provenance}

This Part was built from two manuscripts of batch~93, delivered in the
archives \fn{definable_surreals.zip} and
\fn{definable_surreals_real_parameters.zip} (arrival commit
\texttt{2faa3b37a}) and placed in commit \texttt{47a77daba}.

\begin{center}
\small
\begin{tabular}{@{}P{0.07\linewidth}P{0.33\linewidth}P{0.24\linewidth}P{0.27\linewidth}@{}}
\toprule
Source & Title and date & Pin & Printed here\\
\midrule
02 & \emph{Definable Surreal Numbers and Omnific Integers: Real Parameters,
Countable Closure, and Strong Symmetries}, 4~October 2026 &
\texttt{3c25fbb55} &
Sections~19--29 in full (base); labels \texttt{dsn:rp:}\\
05 & \emph{Definable Surreal Numbers and Omnific Integers: Real
parameters, ordinal information, and countable assembly}, 4~October 2026 &
\texttt{a9ab9a698} &
inserted by subject, \cref{dsn:rp:tab:corr05}; labels \texttt{dsn:ca:}\\
\bottomrule
\end{tabular}
\end{center}
The full pins, both commits of 4~October 2026 and ancestors of the placement
commit, are
\begin{center}\small
\texttt{3c25fbb5515748040decb853139ed39543885b90} (source~02; ``Verify
batch-90 file placement and route five later research arrivals''),\
\texttt{a9ab9a69872c2fb7464097af32fd7b76c75cb030} (source~05; ``Audit eight
arithmetic compiler reports'').
\end{center}

The two manuscripts are not editions of one text: their pins differ, they
share under one per cent of their eight-word sequences, and neither cites
the other. They prove the same central theorems independently. The merge
made the following choices.
\begin{enumerate}
\item \emph{Base.} Source~02 is the base: its negative direction needs only
the consistency of ZFC (Namba forcing over $L$, \cref{dsn:rp:thm:namba})
where source~05 needs a measurable cardinal, and it proves the larger set of
statements on the shared spine.
\item \emph{One copy of each shared theorem.} The countable-assembly
equivalence (source~02's Theorem~5.1, source~05's Theorem~7.2) is printed
as \cref{dsn:rp:thm:countable}; source~05's additional equivalent clauses
are \cref{dsn:ca:thm:assembly} with source~05's proof as a second route.
The Cohen classification (02's Theorem~6.1, 05's Theorem~9.3) is
\cref{dsn:rp:thm:cohen-atlas}, with 05's ring form and proof as a second
route; 05's Theorem~9.4 is printed as \cref{dsn:ca:thm:cohenconsequences},
naming which clauses are source~02's. The global test (05's
Proposition~5.4) is printed for its new invariant $\lambda_\R$ only.
\item \emph{Two codes.} Source~02's rational-slot code $Z_f$ (in the proof of
\cref{dsn:rp:thm:countable}) and source~05's multiplexer $\Zc(f)$
(\cref{dsn:ca:thm:code}) are the same construction with different band
constants, $[q_n,\tfrac54q_n)$ and $[q_n,\tfrac32q_n)$ for
$q_n=2^{-n-1}$; both are kept.
\item \emph{Prikry.} Source~02's \cref{dsn:rp:thm:prikry-sum} and source~05's
Theorem~10.2 prove the same phenomenon with different codes; source~05's
version is \cref{dsn:ca:thm:prikry}, with its graph-coding proof that the
Prikry sequence is not real-ordinal definable as a second route to
source~02's order-type argument.
\item \emph{Background.} Source~05's re-proofs of Part~\ref{dsn:part:one}
background (its Sections~3--4) are printed once: as pointers to Part~I and
to source~02, with source~05's proofs kept where they take a different
route (\cref{dsn:ca:sec:algebra}).
\item \emph{Superseded statements.} Source~05's Section~10.3 and its
Question~13.1 present a measurable cardinal as the known upper bound for
failure of countable assembly; source~02's \cref{dsn:rp:thm:namba} needs
only $\mathrm{Con}(\ZFC)$. Both texts are kept with dated notes. Source~05's
Question~13.7 is answered negatively by \cref{dsn:rp:cor:global-single-real}.
\item \emph{Questions.} Questions asked by both sources are printed once,
with both discussions; source~05's other questions follow source~02's in
\cref{dsn:rp:sec:agenda}, with the questions added at the write.
\end{enumerate}

Editorial changes, besides the notation of \cref{dsn:rp:sub:notation}: the
two abstracts were merged into the abstract above; citations of the
predecessor report became references into Part~\ref{dsn:part:one}, and
citations of other reports of the collection name their labels; the two
bibliographies were merged into the report's (\cite{vdDE},
\cite{DoraisHamkins}, \cite{HamkinsDensity}, \cite{vddTame} and \cite{LM}
are the sources' \texttt{VDE}, \texttt{DoraisHamkins},
\texttt{HamkinsDense}, \texttt{vdD} and \texttt{LM}); source~02's two labels
on its Section~8 became one; source~05's cover page, contents and build
appendix were dropped, their content kept in this subsection and in the
README.

\paragraph{Files.} The report directory ships, from source~02,
\fn{code/08-real-parameters-finite_code_check.py} (exact rational checks of
the slot inequalities and the inverse formula of $Z_f$; no recorded output
was delivered) and \fn{code/08-real-parameters-build.sh}, which builds the
undelivered \fn{definable_surreals.tex} and is kept as delivered, not for
use; from source~05, \fn{code/09-countable-assembly-verify_code.py} with its
recorded output \fn{data/09-countable-assembly-verification.json} (20,508
inverse checks, 206 sequence checks, 1,000 band checks and four
invalid-input rejections), \fn{code/09-countable-assembly-Makefile} and the
audit \fn{09-countable-assembly-PROOF_AUDIT.md}. The delivered files are
byte-identical to the delivery and still use delivery names; the README
gives rerun instructions on a copy. Both scripts check finite rational
specializations only; they do not implement infinite ordinals,
satisfaction, $\HOD$ or forcing, and they are not proofs of any theorem
here. The manuscripts' sources, PDFs and delivery READMEs and source~05's
checksum manifest are not shipped and survive in the arrival commit's
archives.

\subsection{Notation in this Part \writetag}
\label{dsn:rp:sub:notation}

Symbols of Part~\ref{dsn:part:one} keep their meaning; in particular
$\Cell$ is Part~I's code interval \eqref{dsn:eq:maincell}, $\Code$ is
Part~I's sign code, $\mathfrak S=(\No;+,-,\cdot,<,0,1,\Oz)$, and
$e_\beta=\omega^{-(\beta+1)}$ as in \eqref{dsn:eq:Khyp}. Symbols defined
inside a proof are local to it. The sources' symbols were changed as
follows.
{\small
\begin{longtable}{@{}P{0.27\linewidth}P{0.22\linewidth}P{0.42\linewidth}@{}}
\toprule
Source symbol & Here & Reason\\
\midrule
\endhead
02: $\mathcal U=\bigcup_r\HOD(r)$ & $\Hcup$ & $\mathcal U$ is Part~I's code
interval $\Cell$, in which source~02's own examples land\\
02: $\mathsf C_\omega^\R$ (assembly principle); 05: $\mathsf A_\omega$ &
$\Asm$ & $\mathsf C$ is Part~I's sign code; source~02's principle
\eqref{dsn:rp:eq:countable-principle} (countable sets of ordinals) and
source~05's (countable ordinal sequences) are equivalent\\
02: $\operatorname{OD}_\R$; 05: $\mathrm{ROD}$ & $\ROD$ & one name for one
class\\
05: $K_{\mathrm{ro}}$, $I_{\mathrm{ro}}$, $H_{\mathrm{ro}}$,
$\lambda_{\mathrm{ro}}$ & $\KR$, $\IR$, $\HR$, $\lambda_\R$ & source~02's
names; the classes are the same\\
05: $K_{\mathrm{od}}=\No\cap\OD$, $I_{\mathrm{od}}$ & $\KH$, $\IH$ & Part~I's
names; equal by \cref{dsn:thm:HOD}\\
05: $\mathcal Z(f)$ (macro \texttt{\textbackslash Code}) & $\Zc(f)$ & the
macro name \texttt{\textbackslash Code} is Part~I's sign code $\Code$; the
glyph $\Zc$ is unchanged; it differs from source~02's $Z_f$\\
05: $e_n(a)$ (band exponents) & $b_n(a)$ & $e_\beta$ is Part~I's
$\omega^{-(\beta+1)}$\\
02, \cref{dsn:rp:sec:countable}: $e_n$ (slot exponents) & $g_n$ & the
same\\
02, \cref{dsn:rp:sec:prikry}: $e_n=\Omega(-(\kappa_n+1))$ & $e_{\kappa_n}$ &
this is Part~I's $e_\beta$ at $\beta=\kappa_n$; no change of meaning\\
02, \cref{dsn:rp:sec:strong}: $H$, $G$, $K$ & $\Gamma_\Z$, $\Gamma_\Q$,
$\R(\Ord)$ & $H_r$, Part~I's $K$ (\cref{dsn:thm:amplify}) and $G$
(\cref{dsn:lem:cosets}) mean other things\\
02: $\mathcal F_G$, $I_G$ & $\mathcal F_\Gamma$, $I_\Gamma$ & follows the
previous line\\
02: $\ell(x)$ (leading exponent) & $\operatorname{lead}(x)$ & Part~I's
$\ell$ is a width\\
02: $E=\sum_n\omega^{-n}/n!$ & $\mathcal E$ & Part~I's $E_\alpha(X)$ is a bit
code\\
02, \cref{dsn:rp:sec:strong}: $\tau$, $\tau_s$, $S_\tau$ (additive
automorphisms) & $\psi$, $\psi_s$, $S_\psi$ & $\tau$ is the bounded-code
exponent of \eqref{dsn:eq:tauinv}, also in source~02's
\cref{dsn:rp:sec:real-parameters}\\
02: $e_\alpha=\omega^{-\alpha}$ (\cref{dsn:rp:thm:independence}) &
$m_\alpha$ & differs from Part~I's $e_\alpha=\omega^{-(\alpha+1)}$ by the
index shift\\
02: $B=\{\Omega(t):0<t<1\}$; 06: $U$ & $\Qrange$ & the range of the bounded
code $\Mc$ (\cref{dsn:thm:bounded-code}); $B$ and $U$ are used otherwise\\
02: $\mathfrak P=(\No;+,\cdot,<,\Oz)$ & $\mathfrak S$ & Part~I's name\\
02: $\operatorname{bd}$, $\operatorname{tc}$, $\operatorname{rcl}$,
$\sqsubset_{\neq}$; 05: $\mathrm b$ & $\bd$, $\TC$, $\RC$, $\ps$ & Part~I's
typography\\
\bottomrule
\end{longtable}}
Tempting false readings: $\Hcup$ is a union of inner models, not an inner
model in general (\cref{dsn:rp:thm:stabilization}); $\KR$ is the union of the
fixed-real fields $K_r$, not $\No\cap\OD(\R)$ for the single set $\R$, which
is just $\KH$ (\cref{dsn:ca:prop:realset}); the field $\mathcal F_\Gamma$ is
a fixed field of automorphisms, not a definable closure
(\cref{dsn:rp:prop:fixed-not-definable}); and $\omega^a$ is the Conway map
$\Omega(a)$, not $\exp(a\log\omega)$.


\subsection{Where source 05 is printed \writetag}
\label{dsn:rp:sub:corr05}

{\small
\begin{longtable}{@{}P{0.30\linewidth}P{0.62\linewidth}@{}}
\caption{Source~05's numbered items and their places in this report.}
\label{dsn:rp:tab:corr05}\\
\toprule
Source 05 & Here\\
\midrule
\endfirsthead
\toprule
Source 05 & Here\\
\midrule
\endhead
Definitions 2.1--2.2 & \cref{dsn:ca:sec:foundations} (\cref{dsn:ca:def:summable})\\
Proposition 3.1, Lemma 3.2 & \cref{dsn:ca:prop:realset}, \cref{dsn:ca:lem:cert}\\
Proposition 3.3 & \eqref{dsn:rp:eq:fielddcl} and \cref{dsn:thm:purefield}\\
Lemma 4.1, Theorem 4.2, Lemma 4.4, Remark 4.5, Theorem 4.6, Proposition 4.7, Section 4.3 &
second routes in \cref{dsn:ca:sec:algebra} to \cref{dsn:lem:closure},
\cref{dsn:thm:hull}, \cref{dsn:prop:relative}, \cref{dsn:lem:HODsupport} and
\cref{dsn:thm:valuation} of Part~\ref{dsn:part:one} and to
\cref{dsn:rp:prop:basic-closure,dsn:rp:prop:hereditary-model}\\
Proposition 4.3 & \eqref{dsn:rp:eq:floor} and \cref{dsn:prop:floor}, with source~05's proof in \cref{dsn:ca:sec:algebra}\\
Proposition 5.1 & \cref{dsn:rp:prop:hereditary-model} and the paragraph after it\\
Theorem 5.2 & \cref{dsn:ca:thm:countbirthday}\\
Proposition 5.3 & \cref{dsn:rp:thm:parameter-order}, last clause\\
Section 5.1, Proposition 6.1 & \cref{dsn:ca:sec:realunion} (spectra; the bounded code \cref{dsn:thm:bounded-code})\\
Proposition 5.4 & \cref{dsn:ca:prop:globaltest}, with \cref{dsn:rp:cor:global-single-real}\\
Theorem 6.2, Remark 6.3, Example 6.4 & \cref{dsn:ca:thm:code}, the remark after it, \cref{dsn:ca:ex:finite}\\
Definition 7.1, Theorem 7.2 & \cref{dsn:ca:sec:assembly}: \cref{dsn:rp:thm:countable} with \cref{dsn:ca:thm:assembly}\\
Corollary 7.3, Remark 7.4, Proposition 7.5 & \cref{dsn:ca:cor:ringassembly}, the remark after it, \cref{dsn:ca:prop:cover}\\
Lemma 8.1, Definition 8.2, Theorem 8.3, Corollary 8.4, Remark 8.5 &
\cref{dsn:ca:lem:codclosure}, \cref{dsn:ca:def:stable}, \cref{dsn:ca:thm:completion} and the two statements after it\\
Lemmas 9.1--9.2 & \cref{dsn:ca:lem:ccc,dsn:ca:lem:homogeneity}\\
Theorem 9.3 & \cref{dsn:rp:thm:cohen-atlas}; source~05's form and proof in \cref{dsn:ca:sec:cohen}\\
Theorem 9.4, Remark 9.5 & \cref{dsn:ca:thm:cohenconsequences} and the remark after it\\
Lemma 10.1 & \cref{dsn:rp:lem:prikry-cones,dsn:rp:lem:prikry-rod-ground}; source~05's proof in \cref{dsn:ca:sec:prikry}\\
Theorem 10.2, Remark 10.3 & \cref{dsn:ca:thm:prikry} and the remark after it\\
Sections 10.3, 11, 12, 14 & \cref{dsn:ca:sec:prikry} (with a dated note), \cref{dsn:ca:sec:consequences}, \cref{dsn:ca:sec:formalization}, \cref{dsn:ca:sec:consequences}\\
Questions 13.1--13.10 & \cref{dsn:ca:sec:questions}\\
Appendices A and B & \cref{dsn:ca:app:certificates}; \cref{dsn:ca:sec:formalization}\\
\bottomrule
\end{longtable}}

\section{Foundations: sets, signs, series, and integers}
\label{dsn:rp:sec:foundations}

\subsection{The ambient convention}

Unless otherwise stated, work in an ambient universe $V$ satisfying ZFC.
Each surreal is a \emph{set}, although the collection $\No$ is a proper
class. We use the canonical sign-sequence representation
\[
 x:\alpha\longrightarrow\{-,+\},\qquad \alpha\in\Ord.
\]
The ordinal $\alpha=\bd(x)$ is its birthday. Prefix inclusion is
simplicity. Numerical order is lexicographic with
$-<\text{termination}<+$. The canonical ordinal $\alpha$ is the all-plus
sequence of length $\alpha$. We identify an ordinal with this canonical
surreal when performing surreal field arithmetic; ordinal addition and
surreal addition must otherwise be distinguished.

An inner model in this paper is a transitive definable class containing
all ordinals and satisfying the indicated set theory. For an inner model
$M$, $\No^M$ means its canonical sign sequences. In particular
\[
 \No^M=\No\cap M
\]
as carriers, and their field operations agree with the ambient ones.
For set models we explicitly say when satisfaction and cardinality are
external. Forcing constructions can be read over transitive models or
through the usual forcing-theoretic relative-consistency interpretation;
we do not infer the existence of a transitive model from ZFC.

Statements involving a fixed first-order formula on $\No$ are translated
into set theory by relativizing quantifiers to its definable carrier.
Statements about all formulas are metatheorems or schemes unless a
set-sized satisfaction predicate is supplied. Statements about all class
automorphisms may be read in a class metatheory. The separating
automorphisms constructed below are definable from set parameters, so
their use does not require a class-sized choice function or a
proper-class basis.

\subsection{Normal forms and summability}

Write $\Omega(x)=\omega^x$ for Conway's canonical omega map. Every
surreal has a unique normal form
\begin{equation}\label{dsn:rp:eq:nf}
 x=\sum_{\xi<\lambda}r_\xi\omega^{g_\xi},
 \qquad r_\xi\in\R\setminus\{0\},\qquad
 g_\xi>g_\eta\quad(\xi<\eta<\lambda).
\end{equation}
The exponent support is a set and is reverse well ordered: every
nonempty subset has a greatest element. Conversely, every such set-sized
normal form represents a surreal. The map satisfies
\[
 \Omega(a+b)=\Omega(a)\Omega(b),\qquad
 a<b\Longrightarrow\Omega(a)\ll\Omega(b),
\]
where $u\ll v$ for positive quantities means $nu<v$ for every positive
ordinary integer $n$. We also write $u\prec v$ if $|u|\ll|v|$.
The omega map is not the ordinary surreal exponential. In particular,
$\Omega(1)=\omega$, whereas $\exp(1)$ is the ordinary real $e$.

A family $(x_i)_{i\in J}$ is \emph{Hahn summable} (also called strongly
summable) if the union of its exponent supports is reverse well ordered
and each exponent occurs in only finitely many members. Its sum is
defined coefficient by coefficient. This is an algebraic summation
operation, not a claim of sequential convergence in the order topology.
All uses of summability below have set-sized index sets.

The canonical normal form, the omega map and its inverse on monomials,
and the simplest-number operation on set-sized cuts are uniquely
definable set-theoretic constructions. These facts follow from the
standard sign and normal-form theory \cite{Gonshor}; the relevant
absoluteness and definability issues are discussed in the predecessor
(Part~\ref{dsn:part:one}). We give the hereditary membership argument whenever
normal forms are used to transfer information between parameter regimes.

\subsection{Omnific integers and the sign of the infinitesimal tail}

An omnific integer is a surreal whose normal form has no negative
exponents and whose coefficient at exponent zero is an ordinary integer:
\begin{equation}\label{dsn:rp:eq:omnific}
 \Oz=\left\{p+n:\supp(p)\subseteq\No_{>0},\ n\in\Z\right\}.
\end{equation}
The coefficients of positive-exponent terms may be arbitrary reals. Thus
$\sqrt2\,\omega$, $\omega^{1/2}$, and $\omega^{1/\omega}$ are omnific
integers. This does not make them ordinary integers or ordinal numbers.
The ring $\Oz$ is discretely ordered with least positive element $1$.

Decompose a surreal uniquely as
$x=p+r+\varepsilon$, where $p$ is purely infinite, $r\in\R$, and
$\varepsilon$ is infinitesimal. Its unique omnific floor is
\begin{equation}\label{dsn:rp:eq:floor}
 \lfloor x\rfloor_{\Oz}=
 \begin{cases}
 p+r-1,&r\in\Z\text{ and }\varepsilon<0,\\
 p+\lfloor r\rfloor,&\text{otherwise}.
 \end{cases}
\end{equation}
It is characterized by $z\le x<z+1$. The first case matters: for example,
$\lfloor\omega-\omega^{-1}\rfloor_{\Oz}=\omega-1$, not $\omega$.
Uniqueness also makes the floor definable in the ordered field with the
predicate $\Oz$.


\subsection{Source 05's conventions \src{05}}
\label{dsn:ca:sec:foundations}

Source~05 works with the same conventions; its formulations add the
following points.

\paragraph{The universe convention.}

We reason in ZFC about canonical set codes. A surreal is a sign function
$s:\alpha\to\{-,+\}$ for an ordinal $\alpha$; its birthday is
$\bd(s)=\alpha$. The class of all such functions is $\No$. No surreal is
itself a proper class. The notation $\No$ is a definable proper class,
not an extra set to which unrestricted power-set operations may be
applied.

A claim about arbitrary ambient formulas is read as a scheme, one claim
for each fixed standard formula. Equivalently, one may fix externally a
transitive set model $M\models\mathrm{ZFC}$ and use genuine external
satisfaction for $(M,\in)$. Then $\No^M$ is a set in the metatheory, while
$M$ regards its surreal numbers as a proper class. The existence of such
an $M$ is not asserted by ZFC; model-based results are conditional on the
specified model. All forcing statements use this usual relative-model
interpretation.

Unless marked otherwise, ordinals, reals, countability and definability
are those of the ambient universe currently under discussion. In a
forcing section that universe is the extension, not the ground model.
We restrict the main model comparisons to transitive models; we make no
claim about a least externally omitted ordinal in an ill-founded model.

\paragraph{Order, simplicity, and two exponentials.}

The order on signs is lexicographic, with a missing sign between $-$ and
$+$. An ordinal $\alpha$ is represented by the all-plus sequence of length
$\alpha$. The simplicity relation is proper initial-segment inclusion
of sign sequences. A subclass is \emph{initial} when it contains every
proper sign prefix of every one of its members.

The Conway map
\[
\Omega(a)=\omega^a
\]
is the canonical monomial map. It satisfies
$\Omega(a+b)=\Omega(a)\Omega(b)$ and is strictly increasing. In particular,
$0<a<1$ gives $1<\omega^a<\omega$. Throughout, $\omega^a$ means this
map; it is \emph{not} shorthand for
$\exp(a\log\omega)$ using Gonshor's exponential. Both constructions are
ambiently definable, but they play different roles \cite{Gonshor,EhrlichKaplanJSL}.

When ordinals occur in a displayed \emph{field} expression such as
$\alpha/(1+\alpha)$, they are first embedded into $\No$ and the operations
are surreal field operations. Ordinal-index manipulations, such as a
well-order type or the domain of a sign sequence, remain ordinary ordinal
operations. This convention is important in the code below.

\paragraph{Normal forms and omnific integers.}

Every surreal has a unique Conway normal form
\begin{equation}\label{dsn:ca:eq:nf}
 x=\sum_{i<\lambda} r_i\omega^{a_i},
 \qquad r_i\in\R\setminus\{0\},\qquad a_i>a_j\quad(i<j),
\end{equation}
where $\lambda$ is an ordinal and the support is a set. Conversely, a
set-sized reverse-well-ordered support with nonzero real coefficients
gives a surreal. We use $\NF(x)$ for the entire indexed graph
$\seq{(a_i,r_i):i<\lambda}$, not for a collection of individually
specified entries. These are standard normal-form facts
\cite{Conway,Gonshor,LM}.

\begin{definition}
The class $\Oz$ of omnific integers consists of the normal forms with no
negative exponent and with integer coefficient at exponent $0$. A
missing constant term has coefficient $0$. For a parameter policy
$\mathcal P$, a \emph{$\mathcal P$-definable omnific integer} is an
omnific integer whose canonical surreal code is $\mathcal P$-definable.
\end{definition}

Thus $\sqrt2\omega$, $\omega^{1/2}$ and
$\omega^{1/2}+\omega^{1/4}+\cdots$ are omnific integers; $1/2$ and
$\omega^{-1}$ are not. An omnific integer need not have ordinary integer
coefficients at positive exponents. Nor does the interval $(0,\omega)$
contain only the ordinary positive integers. The integer-part theorem
for $\Oz$ is classical \cite{LM}; we give its useful explicit formula in
\cref{dsn:prop:floor}.

\paragraph{What a countable sum means.}

\begin{definition}\label{dsn:ca:def:summable}
A family $(x_j)_{j\in J}$, with $J$ a set, is \emph{Hahn-summable} if
$\bigcup_{j\in J}\supp(x_j)$ is reverse well ordered and each exponent
occurs in the support of only finitely many $x_j$. Its sum is obtained by
adding the finitely many coefficients at each exponent and removing
zeros.
\end{definition}

This is a formal, canonical summation operation, not convergence of
partial sums in the order topology. In particular a reverse-well-ordered
sequence of distinct monomials is summable. Countable support does not
imply countable birthday: its exponents may encode very large ordinals.
All sums in the main theorems are set-sized, and no sum over all ordinals
is taken.

The normal-form and summation constructions are fixed definable
operations on canonical set codes. Consequently a definable \emph{whole
summable family} always has a definable sum with the same parameters.
The difficult issue is whether pointwise definability yields such a
whole-family definition.

\section{A parameter policy is part of a definition}
\label{dsn:rp:sec:definitions}

\subsection{External definitions without ordinal parameters}

For a set model $M\models\mathrm{ZFC}$ and an external family $A\subseteq M$,
let $D_A^M$ consist of objects uniquely defined in $(M,\in)$ by a standard
first-order formula using a finite tuple from $A$. Write
\[
 D_{A,\No}^M=D_A^M\cap\No^M,\qquad
 D_{A,\Oz}^M=D_A^M\cap\Oz^M.
\]
Taking $A=\varnothing$ gives parameter-free definability. Taking a fixed
real $r$ as the only parameter, or taking $A=\R^M$, gives two different
real-only notions. There are at most
$\max(\aleph_0,|A|)$ such objects externally. The external enumeration of
definitions need not belong to $M$ and does not provide an internal
enumeration or an internal birthday bound. Pointwise definable models
are compatible with this observation \cite{HLR}.

In the intended universe, the corresponding phrase ``definable without
parameters in $V$'' is a semantic notion. It must not be replaced by an
unrestricted internally definable truth predicate. Tarski's obstruction
applies here, as it also does to real-only definability without ordinal
parameters \cite{DoraisHamkins}. A fixed complexity level or a fixed
set structure $V_\theta$ has a satisfaction predicate; the entire
universe does not have one in its own language.

\subsection{Ordinal and real-plus-ordinal definitions}

For a real $r$, define $\OD(r)$ internally by rank-local satisfaction:
an object $x$ belongs to $\OD(r)$ when there are an ordinal $\theta$, a
formula code, and finitely many ordinals below $\theta$ such that
$x,r\in V_\theta$ and $x$ is the unique object selected by that formula
in $(V_\theta,\in)$ using those parameters and $r$.
Reflection identifies this with semantic definability in $V$ from $r$
and finitely many ordinals. Conversely, a rank-local definition is an
ambient definition because its rank and formula code may be included
among the ordinal parameters. This construction uses only satisfaction
for a set structure.

Let $\OD$ denote the case with no additional real, and define
\begin{equation}\label{dsn:rp:eq:rod}
 \ROD=\bigcup_{r\in\R}\OD(r),\qquad
 \HR=\{x:\TC(\{x\})\subseteq\ROD\}.
\end{equation}
We use $\HR$ for what is often written $\HOD(\R)$ in the
real-plus-ordinal sense. A finite tuple of real parameters can be joined
into a single real. The predicates in \eqref{dsn:rp:eq:rod} are first-order
definable in set theory. The standard general theory of $\OD(X)$ provides
this framework \cite{DoraisHamkins}; no uniform truth predicate for $V$
is being asserted.

For a fixed real let
\[
 H_r=\HOD(r)=\{x:\TC(\{x\})\subseteq\OD(r)\},\quad
 K_r=\No\cap\OD(r),\quad I_r=\Oz\cap\OD(r).
\]
Also set $\KR=\No\cap\ROD$ and $\IR=\Oz\cap\ROD$. These definitions
make explicit the difference between \emph{a fixed real} and \emph{a
real that may depend on the number being defined}.

Giving the whole set $\R$ as one additional ambient parameter does not
have the same effect as allowing its individual members as parameters.
The set $\R$ is already parameter-free definable in $V$, so giving that
set alone adds no ambient definability. Similarly, naming all ordinals
as separate constants allows only finitely many of them in any one
formula; it does not supply a countable sequence of ordinal parameters.

\begin{center}
\small
\begin{tabular}{>{\raggedright\arraybackslash}p{0.25\linewidth}>{\raggedright\arraybackslash}p{0.32\linewidth}>{\raggedright\arraybackslash}p{0.32\linewidth}}
\toprule
Regime & Allowed information in one definition & Formal status\\
\midrule
Parameter-free & One finite formula & External semantics, or a fixed-model notion\\
Real-only & A finite tuple of reals & External semantics without a rank parameter\\
Ordinal-definable & Finitely many ordinals & Internal rank-satisfaction predicate $\OD$\\
Fixed-real ordinal-definable & One specified $r$ and finitely many ordinals & $\OD(r)$ and $H_r$\\
Varying-real ordinal-definable & Some real $r$, depending on the object, and finitely many ordinals & $\ROD$ and $\HR$\\
Countable ordinal information & A countable ordinal sequence and finitely many ordinals & $\OD(\Ord^\omega)$; usually stronger\\
\bottomrule
\end{tabular}
\end{center}

\subsection{Basic field and integer-part closure}
\label{dsn:rp:subsec:closure-background}

The next proposition is background from the predecessor, with proof
included to make the later arguments independent of informal closure
claims (Part~\ref{dsn:part:one}).

\begin{proposition}[Same-parameter closure]\label{dsn:rp:prop:basic-closure}
In any of the ambient finite-parameter regimes, the definable surreals
form an ordered real closed subfield. They are closed under every
uniquely specified parameter-free ambient operation whose inputs belong
to the regime. In particular this includes the omega map, its inverse
on monomials, normal forms, the omnific floor, and summation when the
\emph{whole family} is available as a definable set. Their omnific
integers form an integer part, and the field is their fraction field.
\end{proposition}

\begin{proof}
Combine the finitely many input definitions by substitution into the
graph of the desired operation. This gives addition, multiplication,
negation, reciprocal, and the other uniquely specified outputs without
new parameters. A polynomial over the resulting field has its real roots
in a finite numerical order. Its $j$th root, if it exists, is therefore
uniquely definable from its coefficients and the ordinary integer $j$.
Positive square roots and roots of odd-degree polynomials follow, giving
real closure. The floor characterization gives the integer part.

For the fraction-field assertion take $x\ne0$ with exponent support $S$.
Let $\gamma$ be the simplest surreal strictly greater than
$\{0\}\cup\{-g:g\in S\}$. It is definable from $x$. Put
$b=\omega^\gamma$ and $a=xb$. Every exponent of $a$ and $b$ is positive,
so $a,b$ are purely infinite omnific integers, $b\ne0$, and $x=a/b$.
The case $x=0$ uses $0/1$. This construction preserves the original
parameter information.
\end{proof}

The words ``whole family'' cannot be dropped. The main countable
theorem below measures exactly when that additional information can be
assembled from individual real-plus-ordinal definitions.

\subsection{Intrinsic ordered-field definability}

Fix the language $L_{\mathrm{of}}=\{0,1,+,-,\cdot,<\}$. Quantifier
elimination for real closed fields gives
\begin{equation}\label{dsn:rp:eq:fielddcl}
 \dcl_{L_{\mathrm{of}}}^{\No}(A)
   =\{x\in\No:x\text{ is algebraic over }\Q(A)\}
   =\RC_{\No}\Q(A).
\end{equation}
For a proper parameter class the right side uses only finite parameter
subsets; this is the same convention as on the left.
To see the nontrivial inclusion directly, a formula reduces to finitely
many polynomial sign conditions. At a point transcendental over its
coefficient field all nonzero polynomials in that finite list are
nonzero, and their signs persist on an interval. The formula cannot
isolate that point. Conversely, a root is isolated by its numerical
position among the finitely many roots.

Consequently
\[
 \dcl_{L_{\mathrm{of}}}^{\No}(\varnothing)=\R_{\mathrm{alg}},\qquad
 \dcl_{L_{\mathrm{of}}}^{\No}(\R)=\R.
\]
With no parameters, the omnific intersection is $\Z$. These intrinsic
conclusions do not conflict with ambient definitions of infinite
surreals. Later we strengthen the intrinsic obstruction even after every
real and ordinal is named and the omnific predicate is added.

% ===== End sections/01_foundations.tex =====


% ===== Begin sections/02_real_parameters.tex =====

\subsection{Source 05's parameter dictionary \src{05}}
\label{dsn:ca:sec:definitions}

For a fixed structure $(M,\in)$ and a family $A\subseteq M$ of allowed
parameters, let $D_A^M$ be the external class of elements uniquely
specified in $M$ by a standard first-order formula and a finite tuple
from $A$. Define
\[
 K_A^M=D_A^M\cap\No^M,\qquad I_A^M=D_A^M\cap\Oz^M.
\]
For the intended universe the same notation is interpreted schematically.
Parameter-free definability means $A=\varnothing$. A fixed parameter
$r$ does not authorize arbitrary elements of $r$ as unrelated parameters;
any extraction must itself be definable. For real parameters, natural
indices are harmless because each natural number is parameter-free
definable.

\begin{center}
\small
\begin{tabular}{@{}P{0.18\textwidth}P{0.47\textwidth}P{0.27\textwidth}@{}}
\toprule
Notation & Allowed information in an ambient definition & Surreal field\\
\midrule
$D_\varnothing$ & No parameters & $K_\varnothing$\\
$\RD$ & Finitely many real parameters, no ordinal parameters & $K_{\mathrm{rd}}$\\
$\OD$ & Finitely many ordinals & $\KH$\\
$\OD(r)$ & One fixed real $r$, and finitely many ordinals & $K_r$\\
$\ROD$ & An arbitrary real, and finitely many ordinals & $\KR$\\
$\COD$ & One countable sequence of ordinals & $\Kco$\\
\bottomrule
\end{tabular}
\end{center}

A finite real tuple, or a given countable sequence of reals, can be encoded
by one real using a fixed pairing function on $\omega$. We freely replace
real numbers by their standard subset-of-$\omega$ codes, which preserves
the corresponding definability policies. Thus
\begin{equation}\label{dsn:ca:eq:rodunion}
 \ROD=\bigcup_{r\in\R}\OD(r),\qquad
 \KR=\bigcup_{r\in\R}K_r,\qquad
 \IR=\bigcup_{r\in\R}I_r.
\end{equation}
Only finitely many ordinals are permitted in each $\ROD$ definition.
A countable sequence of ordinal parameters cannot silently be substituted
for that finite tuple.

\begin{center}\small
Source~05's table writes $K_{\mathrm{od}}$, $K_{\mathrm{ro}}$ for $\KH$, $\KR$
(\cref{dsn:rp:sub:notation}). The class $\COD$ and its fields $\Kco$, $\Ico$
are source~02's $\OD(\Ord^\omega)$ of \cref{dsn:rp:cor:countable-variants}.
\end{center}

\begin{proposition}\label{dsn:ca:prop:realset}
Adding $\R$ itself as a single set parameter to the ambient language
$(V,\in)$ does not enlarge parameter-free definability. Adding $\R$ to
the permitted ordinal parameters does not enlarge $\OD$. In contrast,
allowing arbitrary individual real parameters includes every real.
\end{proposition}
\begin{proof}
Replace each occurrence of the parameter $\R$ by its parameter-free
set-theoretic definition. Conversely a given real $r$ is uniquely defined
from the parameter $r$ by $x=r$.
\end{proof}

This proposition is about ambient set theory. Naming a real predicate in
a weaker structure need not be a definitional extension of that weaker
structure. The language must not be changed unnoticed.
\paragraph{Rank-local certificates, without a truth predicate for $V$.}
For ordinal-allowing policies there is a uniform internal definition.
A certificate for $x\in\OD(r)$ consists of an ordinal $\theta$, a finite
formula code $e$, and a finite ordinal tuple $\bar\alpha$, with all
parameters and $x$ in $V_\theta$, such that
\begin{equation}\label{dsn:ca:eq:certificate}
 (V_\theta,\in)\models
 \text{``$x$ is the unique $y$ satisfying $\varphi_e(y,r,\bar\alpha)$''}.
\end{equation}
Satisfaction for \emph{set structures} is a single definable relation.
The certificate therefore makes sense uniformly as $e$ varies.

\begin{lemma}[Certificate equivalence]\label{dsn:ca:lem:cert}
A set is globally ordinal definable from $r$ if and only if it has a
certificate of the form \eqref{dsn:ca:eq:certificate}. Consequently $\OD(r)$,
$\ROD$, and $\COD$ are internally definable classes, with the displayed
parameter policy respected.
\end{lemma}
\begin{proof}
Given a fixed global definition with a unique output, reflection for that
formula and its uniqueness assertion supplies a sufficiently large
$V_\theta$ containing the parameters and the output. Conversely the
unique output in \eqref{dsn:ca:eq:certificate} is globally definable from
$r,\theta,\bar\alpha$ and the finite numeral coding $e$. The extra ordinal
$\theta$ is permitted. For $\COD$, place $\theta$, the syntax code, and
any finite auxiliary ordinals at the beginning of the countable ordinal
parameter sequence. Existentially quantifying over these certificates
gives the uniform internal predicates.
\end{proof}

The same argument does \emph{not} supply a uniform predicate for global
parameter-free definability, or for global real-only definability:
the rank parameter $\theta$ cannot simply be introduced under those
policies. External counting of formulas likewise does not provide an
internal enumeration of all definable objects. Pointwise definable
models give an important warning against a naive ``least undefinable
object'' argument \cite{HLR}. This distinction will be enforced in the
countable-assembly proof.

\paragraph{Intrinsic definability in the field reduct.}
Source~05's Proposition~3.3 is \eqref{dsn:rp:eq:fielddcl} above
(\cref{dsn:thm:purefield} in Part~\ref{dsn:part:one}), proved by the same
quantifier-elimination argument \cite{vddTame}.

For ordinal parameters in the field language, the answer is similarly
the relative real closure of the field generated by those parameters.
It is not, merely by virtue of those parameter names, the ambient
$\HOD$-surreal field. Surreal sign coding and set-theoretic membership

\subsection{Source 05's algebraic baseline: pointers and second routes \src{05}}
\label{dsn:ca:sec:algebra}

The principal statements of source~05's Section~4 are already in
Part~\ref{dsn:part:one} and in \cref{dsn:rp:prop:basic-closure} above, and are
printed there. Source~05 includes proofs to isolate exactly which parameter
and collection rules they use; where its route differs, the statement is
recalled without a number and its proof is kept.

\begin{dupstatement}{Source 05's Lemma 4.1 (= \cref{dsn:lem:closure})}
If $x_1,\ldots,x_m$ are definable under a parameter policy closed under
finite combination of parameters, and $F$ is a fixed parameter-free
ambiently definable partial function defined on this tuple, then
$F(x_1,\ldots,x_m)$ is definable under the same policy.
\end{dupstatement}
\begin{proof}[Second proof (source 05)]
Conjoin the finitely many definitions of the inputs, quantify over those
inputs, and apply the fixed graph formula for $F$. Input uniqueness and
output uniqueness give a unique result. This is a formula-by-formula
argument, not an internal quantification over global truth.
\end{proof}
\begin{dupstatement}{Source 05's Theorem 4.2 (\cref{dsn:thm:hull} with \cref{dsn:rp:prop:basic-closure}; the closure under $\exp$ and $\log$ is part of \cref{dsn:thm:hull})}
Each field associated with a policy in the table of
\cref{dsn:ca:sec:definitions} is a real closed subfield of $\No$. It is closed
under the Conway map, its inverse on monomials, and the canonical surreal
exponential and logarithm on their domains. If $K$ is such a field and
$I=K\cap\Oz$, then $I$ is an integer part of $K$ and
\begin{equation}\label{dsn:ca:eq:frac}
 K=\Frac(I).
\end{equation}
These assertions also apply to a fixed arbitrary allowed parameter
family $A$, interpreted as in \cref{dsn:ca:sec:definitions}.
\end{dupstatement}
\begin{proof}[Second proof (source 05)]
The field operations and their domain restrictions are canonical
ambiently definable maps, so \cref{dsn:lem:closure} applies. Every positive
element has a uniquely specified positive square root. Every odd-degree
polynomial has a real root in $\No$, and one can specify its least real
root, or more generally its $k$th real root in increasing order. Hence
$K$ is real closed. The named maps are canonical, so unique-output
closure applies to them as well. The floor and fraction assertions are
proved explicitly below. Only finite input tuples are used here.
\end{proof}
Source~05's Proposition~4.3 is the floor formula \eqref{dsn:rp:eq:floor}
(\cref{dsn:prop:floor}); its proof:
\begin{proof}
If $c$ is not an integer, its distances from the adjacent integers are
positive reals and dominate the infinitesimal $\varepsilon$. If $c$ is
an integer, a negative infinitesimal places $x$ just below
$x_{>0}+c$, accounting for the first case. The stated inequalities follow
in all cases. A nonzero omnific integer has absolute value at least $1$:
its leading positive exponent makes it infinite, or its normal form is
a nonzero ordinary integer. Thus two candidates for $z$ cannot differ
by a nonzero omnific integer of absolute value less than $1$.
\end{proof}
\begin{dupstatement}{Source 05's Lemma 4.4 (a denominator; compare \eqref{dsn:eq:denominators} and the proof of \cref{dsn:rp:prop:basic-closure})}
For every nonzero surreal $x$ there is a parameter-free definable choice
$q(x)\in\Oz_{>0}$ such that $xq(x)\in\Oz$. In particular if $x\in K$,
both $q(x)$ and $xq(x)$ belong to $I$.
\end{dupstatement}
\begin{proof}[Second proof (source 05)]
Let $S=\supp(x)$. Choose the least ordinal $\delta>0$ such that
$\delta>-a$ for every $a\in S$, with the comparison made in $\No$.
This ordinal exists: each surreal is below some ordinal; Replacement
on the set $S$ and an ordinal supremum give one common upper bound.
The least such ordinal is definable from $x$, without adding $\delta$
as an independent parameter. Set $q(x)=\omega^\delta$.
Multiplication shifts every exponent of $x$ to $a+\delta>0$, so
$xq(x)$ has zero constant term and no negative exponent. Both factors
are omnific. This proves \eqref{dsn:ca:eq:frac}, with $x=0$ handled by
$0/1$.
\end{proof}
\begin{remark}
For a proper-class ring, \eqref{dsn:ca:eq:frac} means elementwise representation
as a quotient with nonzero denominator; it does not require forming a
set of all equivalence classes. In an external set-model interpretation,
it is the usual fraction-field assertion for the resulting set-sized
ring. Likewise the real-closed-field assertion is a scheme of finite
algebraic conditions. Quantifier elimination also yields ordered-field
elementarity of these inclusions in the corresponding semantic setting.
\end{remark}

For a fixed real $r$, use $\HOD(r)$ for the sets hereditarily in
$\OD(r)$, including the object itself in its transitive closure.
Because $r$ is a set of natural numbers, $r\in\HOD(r)$; this is the
standard relative HOD inner model of ZFC \cite{Jech}.
The superscript $\No^N$ will mean the canonical sign sequences belonging
to a transitive inner model $N$, not a quotient by arbitrary game forms.
\begin{dupstatement}{Source 05's Theorem 4.6 (= \cref{dsn:prop:relative} for a real; compare \cref{dsn:thm:HOD} and \cref{dsn:rp:prop:hereditary-model})}
For a fixed real $r$,
\begin{equation}\label{dsn:ca:eq:hod}
 K_r=\No\cap\OD(r)=\No\cap\HOD(r)=\No^{\HOD(r)},
 \qquad I_r=\Oz\cap K_r.
\end{equation}
The same holds with no real parameter, replacing $\OD(r)$ by $\OD$.
Moreover $K_r$ and $\KH$ are initial in the sign-sequence order.
\end{dupstatement}
\begin{proof}[Second proof (source 05)]
A canonical sign sequence is a set of ordered pairs whose first
coordinates are ordinals and whose second coordinates are fixed finite
tags. Every member of the transitive closure of such a graph other than
the graph itself is hereditarily ordinal definable: it is built from
ordinals and finite sets by canonical finite operations. If the graph
is $\OD(r)$, it is therefore hereditarily $\OD(r)$. The reverse inclusion
is immediate. Membership in the canonical surreal class is absolute for
transitive inner models, since a specified domain ordinal and a given
function graph have the same meaning there. This proves the set-code
identification.

For any $\beta<\bd(x)$, the prefix $x\res\beta$ is definable from $x$
and $\beta$. The ordinal $\beta$ is allowed. Thus every prefix of an
$\OD(r)$ surreal is again $\OD(r)$, proving initiality.
\end{proof}
The equality here is first an equality of canonical sign-code classes,
with ambient order and operations on the left. The later arguments do
not require an unproved blanket absoluteness theorem for all analytic
operations between arbitrary inner models with different real lines.
Ambient definability and the preceding hereditary argument suffice.
\begin{dupstatement}{Source 05's Proposition 4.7 (= the whole-normal-form paragraph after \cref{dsn:rp:prop:hereditary-model}; compare \cref{dsn:lem:HODsupport})}
A surreal $x$ belongs to $K_r$ if and only if its entire normal-form graph
belongs to $\HOD(r)$. In particular every coefficient and every exponent
of a member of $K_r$ is $\OD(r)$.
\end{dupstatement}
\begin{proof}[Second proof (source 05)]
From $x$ one defines $\NF(x)$ and, using the allowed ordinal index $i$,
each entry $(a_i,r_i)$. Each exponent is a sign sequence and each
coefficient is a real, so the preceding hereditary argument applies
to them. The graph, its indexed pairs, and all their transitive members
are consequently $\OD(r)$. Conversely the graph is itself $\OD(r)$ if
it is in $\HOD(r)$, and its canonical evaluation uniquely defines $x$.
\end{proof}
This proposition concerns the \emph{whole graph}. Its converse cannot
be weakened to ``every entry is individually $\OD(r)$'': the infinite
assignment of entries to indices can carry information absent from
$r$ and every finite tuple of ordinal parameters. The strongest example
in this article will have every entry already $\OD$.
\paragraph{The natural valuation.}
For $x\ne0$ with leading exponent $a_0$, put $v(x)=-a_0$. On any of the
fields $K$ above the value group is $(K,+)$: the leading exponent is
definable from $x$, and every value $a\in K$ is realized by
$\omega^{-a}$. The residue field is $K\cap\R$. Indeed, for a finite
surreal the real constant term is its unique real standard part, and
the kernel is the infinitesimal ideal. These are direct normal-form
arguments, also recorded in (Part~\ref{dsn:part:one}).
For $\KR$ and $\Kco$ the residue field is the full ambient $\R$.
Having all real residues will not imply that the whole field is $\No$.

This is \cref{dsn:thm:valuation} of Part~\ref{dsn:part:one} for the fields
of this Part.

\section{Varying real parameters and the associated inner models}
\label{dsn:rp:sec:real-parameters}

Throughout this section the ambient universe satisfies ZFC.  We represent
real parameters by subsets of $\omega$; the standard conversions to
ordinary real numbers and to their canonical surreal representatives are
definable.  Allowing an arbitrary finite tuple of real parameters is
equivalent to allowing one real, by interleaving binary codes.  Define
\[
 \ROD=\bigcup_{r\in\mathbb R}\operatorname{OD}(r),
 \qquad
 H_r=\operatorname{HOD}(r),
\]
and distinguish the two classes
\begin{equation}\label{dsn:rp:eq:two-unions}
 \HR
  =\{x:\operatorname{tc}(\{x\})\subseteq\ROD\},
 \qquad
 \Hcup=\bigcup_{r\in\mathbb R}H_r.
\end{equation}
The first is hereditary definability using individual real and ordinal
parameters.  The second is the union of the classes obtained by fixing
one real parameter at a time.  The distinction matters even though
their canonical surreal numbers coincide.  These definitions use the
standard rank-satisfaction presentation of ordinal definability, so all
four displayed classes are first-order definable.  They do not require
a truth predicate for $V$; compare the general $\operatorname{OD}(X)$
framework in \cite[Section~4]{DoraisHamkins}.

\begin{proposition}[The hereditary model and its surreal numbers]
\label{dsn:rp:prop:hereditary-model}
$\HR$ is a transitive inner model of ZF containing all ordinals,
all ambient reals, and the ambient set $\mathbb R$ itself.  If
\[
 K_r=\No\cap H_r,\qquad I_r=\Oz\cap H_r,
 \qquad
 \KR=\No\cap\ROD,\qquad
 \IR=\Oz\cap\ROD,
\]
then
\begin{align}
 \KR
   &=\No\cap \HR
     =\No\cap \Hcup
     =\bigcup_{r\in\mathbb R}K_r
     =\No^{\HR},\label{dsn:rp:eq:surreal-union}\\
 \IR
   &=\Oz\cap \HR
     =\Oz\cap \Hcup
     =\bigcup_{r\in\mathbb R}I_r.\label{dsn:rp:eq:integer-union}
\end{align}
Here the last term of \eqref{dsn:rp:eq:surreal-union} uses the canonical
sign-sequence construction inside $\HR$.
\end{proposition}
\begin{proof}
Hereditary membership immediately makes $\HR$ transitive.
Every ordinal is ordinal definable; every real is definable from itself.
The set $\mathbb R$ is parameter-free definable, and its hereditary
members are reals and the elementary sets used to represent them.
Thus all these objects belong to $\HR$.

We give the inner-model argument to identify where choice is, and is
not, used.  A finite tuple of ROD parameters can be replaced by one real
and finitely many ordinals.  Consequently a uniquely defined output
from such parameters is ROD.  Pairing and union therefore preserve
$\HR$.  Given $a,\bar p\in \HR$, a set obtained by
separation using a formula relativized to $\HR$ is definable
in $V$ from $a,\bar p$, since $\HR$ itself is definable.  It is
ROD, and its hereditary members already lie in $\HR$.
For replacement, apply replacement in $V$ to the relativized functional
formula.  The resulting range is again ROD and has hereditary members
in $\HR$.  Its internal powerset of $a$ is
\[
 \{b\subseteq a:b\in \HR\}.
\]
This set exists by separation in $V$, is ROD from $a$, and is hereditary
ROD.  Infinity is supplied by $\omega$, and extensionality and foundation
are inherited by the transitive class.  Thus the class satisfies ZF.
This proof does not establish choice in $\HR$.

If a canonical sign sequence $s$ is OD$(r)$, every other member of
$\operatorname{tc}(\{s\})$ is ordinal definable: it is an ordinal, a
finite sign tag, or one of the finite sets used to encode an ordered
pair.  Hence $s\in H_r$.  Conversely, membership in $\HR$
implies that $s$ itself is ROD.  This proves the asserted intersections
and unions; the omnific assertion is their restriction to
$\Oz$.  Finally the property of being a function from an ordinal
to the fixed two-element sign set is absolute for the transitive
same-ordinal model $\HR$.  Its canonical surreal class is
therefore precisely $\No\cap \HR$.
\end{proof}

By \cref{dsn:rp:prop:basic-closure} and the fixed-real theory, each $K_r$ is a real closed field and $I_r$ is
its canonical omnific integer part.  The unions in
\eqref{dsn:rp:eq:surreal-union}--\eqref{dsn:rp:eq:integer-union} are directed:
any finitely many real parameters can be combined into one real.
Consequently $\KR$ is a real closed field containing all
ambient reals, and $\IR$ is its omnific integer part.  For
example, the coefficients of a given polynomial all lie in one $K_r$,
which supplies the needed positive square roots and odd-degree roots.
This reasoning does not assume that $\HR$ satisfies choice.

Every $K_r$ and $\KR$ is initial in the sign tree: if $x$ is in the
field, its prefix of any ordinal length is definable from $x$ and that
ordinal. All these fields contain every canonical ordinal. By
quantifier elimination for real closed fields they are elementary
subfields of $\No$ in the ordered-field language, interpreted
formula by formula. They also satisfy the same-parameter fraction-field
property of \cref{dsn:rp:prop:basic-closure}.

There is a useful whole-normal-form version. If $x\in K_r$, its
complete normal form is OD$(r)$; each coefficient and exponent is
definable from $x$ and its ordinal index, hence belongs to $H_r$.
The hereditary graph coding then puts the entire normal form in
$H_r$. Conversely a whole normal form in $H_r$ has an OD$(r)$ value.
Therefore $x\in K_r$ if and only if its whole normal-form code belongs
to $H_r$. The corresponding statement for $\KR$ is obtained by
allowing $r$ to vary. Entrywise membership alone is not sufficient.

\subsection{The real-parameter fields record exactly the parameter degrees}

Define $r\leq_{\mathrm{OD}}s$ to mean $r\in\operatorname{OD}(s)$.
All occurrences of relative HOD in the next theorem refer to the same
ambient universe, rather than to iterated computations inside its
various inner models.

\begin{theorem}[Exact inclusion classification]
\label{dsn:rp:thm:parameter-order}
For real parameters $r,s$, the following are equivalent:
\begin{equation}\label{dsn:rp:eq:parameter-order}
 r\leq_{\mathrm{OD}}s
 \quad\Longleftrightarrow\quad H_r\subseteq H_s
 \quad\Longleftrightarrow\quad K_r\subseteq K_s
 \quad\Longleftrightarrow\quad I_r\subseteq I_s.
\end{equation}
The field $K_{r\oplus s}$ is the least member of this family containing
both $K_r$ and $K_s$, and similarly for $H_r$ and $I_r$.
Every countable collection of members of $\Hcup$ is contained
in some $H_t$.  In particular, every countable collection of members of
$\KR$ is contained in some $K_t$.
\end{theorem}
\begin{proof}
If $r$ is OD$(s)$, substitute a definition of $r$ into every definition
from $r$ and ordinal parameters.  This gives
$\operatorname{OD}(r)\subseteq\operatorname{OD}(s)$, and the hereditary
condition gives $H_r\subseteq H_s$.  Restriction gives the field and
integer-ring inclusions.

Conversely, the canonical surreal representative of $r$ lies in $K_r$.
If $K_r\subseteq K_s$, its inverse real code is OD$(s)$.  For the ring
inclusion use the injective, parameter-free definable map
\[
 \mathbb R\longrightarrow\Oz,\qquad a\longmapsto a\omega.
\]
For $a\ne0$ its image has one positive normal-form exponent and zero
constant term; for $a=0$ the assertion is immediate.  The inverse on
the image is division by $\omega$.  Thus $I_r\subseteq I_s$ also
implies that $r$ is OD$(s)$.

The two coordinates of $r\oplus s$ are definable from their interleaved
code, so its class contains the two given classes.  If $t$ bounds both
parameters in the preorder, the pair and its interleaving are OD$(t)$.
This proves the least-upper-bound statement.  It concerns least upper
bounds within this family; no equality with a field-theoretic
compositum is asserted.

Given $x_n\in \Hcup$, ordinary choice selects reals $r_n$ such
that $x_n\in H_{r_n}$.  Interleave the sequence of real codes into a
single real $t$.  Each $r_n$ is definable from $t$ and the integer $n$,
so $H_{r_n}\subseteq H_t$.  Hence $x_n\in H_t$ for every $n$.
Apply \eqref{dsn:rp:eq:surreal-union} for the surreal version.
\end{proof}

The countable conclusion concerns the individual values.  It does not
say that the sequence $n\mapsto x_n$, or its range as a set, belongs to
$H_t$.  The ordinal information selecting those values can fail to be
recoverable from a real.  The Prikry example below realizes this
failure even when all the values are ordinal definable.

\subsection{Choice is equivalent to stabilization at one real}

\begin{theorem}[Stabilization theorem]
\label{dsn:rp:thm:stabilization}
The following conditions are equivalent:
\begin{enumerate}
\item $\KR=K_r$ for some real $r$;
\item $\IR=I_r$ for some real $r$;
\item $\mathbb R\subseteq H_r$ for some real $r$;
\item $\HR=H_r$ for some real $r$;
\item $\HR\models\mathrm{AC}$;
\item $\mathbb R\in \Hcup$;
\item $\Hcup=\HR$;
\item $\Hcup$ has an internal powerset of $\omega$;
\item $\Hcup$ is an inner model of ZF.
\end{enumerate}
If these conditions fail, the family of fields $\{K_r:r\in\mathbb R\}$
has no countable cofinal subfamily under inclusion.
\end{theorem}
\begin{proof}
Suppose first that $\HR$ satisfies choice.  Since it contains
the actual set $\mathbb R$, it contains a wellordering $W$ of that set.
Its internal isomorphism to an ordinal is still an isomorphism in $V$,
so $W$ really wellorders all ambient reals.  There is a real $s$ such
that $W$ is OD$(s)$.  For each real $a$, its ordinal index in $W$
uniquely defines $a$ from $W$.  Substitution shows that $a$ is OD$(s)$.
Since a real is hereditary OD$(s)$ whenever it is OD$(s)$, we obtain
$\mathbb R\subseteq H_s$.  This proves (5)$\Rightarrow$(3).

If (3) holds, every allowed real parameter can be replaced by a
definition from $r$ and ordinals.  Thus
$\ROD=\operatorname{OD}(r)$, and taking hereditary
membership yields (4).  Condition (4) implies (5), since the fixed-real
class $H_r$ is an inner model of ZFC.  It also implies (1), (2), (7)
and (9).  Conversely (1) implies (3), since $\KR$ contains
all reals; (2) implies it by the coding $a\mapsto a\omega$ from
Theorem~\ref{dsn:rp:thm:parameter-order}.

The class $\Hcup$ is transitive and contains every ambient
subset of $\omega$.  If it has an internal powerset of $\omega$, that
set must therefore be the actual $\mathbb R$.  Hence
(9)$\Rightarrow$(8)$\Rightarrow$(6).  If
$\mathbb R\in H_r\subseteq \Hcup$, transitivity of $H_r$
implies (3).  Finally (7) implies (6), since $\HR$ contains
$\mathbb R$.  These implications prove all equivalences.

For the last assertion, suppose a countable subfamily were cofinal.
Theorem~\ref{dsn:rp:thm:parameter-order} supplies a single real $t$ whose
class contains every member of that subfamily.  Cofinality then implies
that it contains every $K_r$, giving $K_t=\KR$, contrary to
failure of stabilization.
\end{proof}

Thus the failure of choice in $\HR$ is visible in a concrete
field question: there is no single real parameter whose definable
surreal field includes all the others.  In that case the union
$\Hcup$ contains every real individually but omits their
collection as a set and fails even powerset at $\omega$.

\begin{proposition}[The internally hereditarily wellorderable part]
\label{dsn:rp:prop:hereditary-wellorder}
One has
\[
 \Hcup
 =\{x\in \HR:
     \HR\models
       ``\operatorname{tc}(\{x\})\text{ is wellorderable}''\}.
\]
\end{proposition}
\begin{proof}
If $x\in H_r$, its transitive closure belongs to $H_r$, and choice in
$H_r$ supplies a wellordering there.  The same witness belongs to
$\HR$.

Conversely suppose $\HR$ contains a bijection
$e:\lambda\to T$, where $T=\operatorname{tc}(\{x\})$ and $\lambda$
is an ordinal.  Inside $\HR$ form
\[
 E=\{(\xi,\eta)\in\lambda^2:e(\xi)\in e(\eta)\},
\]
and choose the ordinal $\eta_x$ with $e(\eta_x)=x$.  The relation $E$
is ROD, so it is OD$(s)$ for some real $s$.  All the other members of
its transitive closure are ordinal coded; hence $E\in H_s$.  Its
well-foundedness and extensionality hold in $H_s$, which therefore
computes its Mostowski collapse.  By uniqueness of well-founded
collapse, this is the same collapse as in $V$, and its value at
$\eta_x$ is $x$.  Thus $x\in H_s\subseteq \Hcup$.
\end{proof}

\begin{corollary}[A single real suffices if every surreal is ROD]
\label{dsn:rp:cor:global-single-real}
Over ambient ZFC,
\[
 \KR=\No
 \quad\Longleftrightarrow\quad \IR=\Oz
 \quad\Longleftrightarrow\quad V=\HR
 \quad\Longleftrightarrow\quad
 \exists r\in\mathbb R\;V=\operatorname{HOD}(r).
\]
\end{corollary}
\begin{proof}
If every surreal is ROD, the characteristic sign sequence of each
subset of an ordinal puts that subset in $\HR$.  For an
arbitrary set $x$, choice in $V$ supplies an ordinal code for the
membership graph of $\operatorname{tc}(\{x\})$.  This code belongs
to $\HR$, and its well-founded collapse there recovers $x$.
Thus $V=\HR$.  The stabilization theorem then applies,
because this inner model has the ambient axiom of choice, and yields
one real $r$ with $V=H_r$.  The reverse implications are immediate.

To recover all surreals from all omnific integers, use the
parameter-free injective map $x\mapsto\Omega(\Omega(x))$.  Its outer
exponent is positive, so its image consists of purely infinite
omnific integers.  Its inverse on its image is uniquely definable.
Alternatively the bounded reversible coding from the existing
framework gives the same implication.
\end{proof}


\begin{corollary}[A bounded omnific test for the entire universe]
The equivalent global conditions of
\cref{dsn:rp:cor:global-single-real} hold if and only if every omnific
integer $z$ with $1<z<\omega$ belongs to $\ROD$.
\end{corollary}
\begin{proof}
Use the predecessor's reversible bounded monomial code, whose
parameter behavior we spell out:
\[
 \tau(x)=\frac12\left(1+\frac{x}{\sqrt{1+x^2}}\right),
 \qquad Q(x)=\Omega(\tau(x)).
\]
The function $\tau$ bijects $\No$ with $(0,1)$, and $Q(x)$ is a
purely infinite omnific integer between $1$ and $\omega$. If
$t=\Omega^{-1}(Q(x))$, then
\[
 x=\frac{2t-1}{2\sqrt{t(1-t)}}.
\]
Both directions are uniquely definable without new parameters. Thus
ROD membership of every such omnific integer implies ROD membership
of every surreal, and the global corollary applies. The bounded code
itself is antecedent material from (Part~\ref{dsn:part:one}).
\end{proof}

% ===== End sections/02_real_parameters.tex =====


% ===== Begin sections/03_countable_closure.tex =====

\begin{remark}[Source 06's form of \cref{dsn:rp:cor:global-single-real} \src{06}]
\label{dsn:rp:rem:consolidation06}
Source~06 (Part~\ref{dsn:part:op}) proves \cref{dsn:rp:cor:global-single-real}
independently, as its Theorem~4.2, which is printed only here
(\cref{dsn:op:thm:real-consolidation} in Part~\ref{dsn:part:op} points back to
this remark). Its formulation: over ZFC the following are equivalent:
(1)~$V=\HR$; (2)~every set is $\ROD$; (3)~there is a real $r$ with $V=H_r$;
(4)~every surreal is $\ROD$; (5)~every omnific integer is $\ROD$;
(6)~every monomial $\omega^a$ with $0<a<1$ is $\ROD$. Consequently, if every
surreal has some real--ordinal definition, then one real $r$ works for all
surreals and all omnific integers simultaneously. Clause~(6) strengthens
the preceding corollary, since these monomials form the subclass $\Qrange$ of
$\Oz\cap(1,\omega)$.
\end{remark}
\begin{proof}[Second proof (source 06)]
The equivalence of (1) and (2) follows directly from hereditary
membership.  To prove (2)$\Rightarrow$(3), choose in $V$ a well-order $W$
of $\R$.  By (2), $W$ is definable from some real $r$ and ordinal
parameters.  Every real $s$ is then $\OD(r)$: use, in addition, its
ordinal position in $W$.  Given any set $x$, a ROD definition of $x$
uses a real $s$, which can now be replaced by its $\OD(r)$ definition.
Thus every set is $\OD(r)$ and hence hereditarily so.  This proves
$V=H_r$.  Conversely (3) implies (2).

Clearly (2) implies (4), (4) implies (5), and (5) implies (6).
For (6)$\Rightarrow$(4), use the parameter-free reversible monomial
coding $x\mapsto\omega^{\tau(x)}$, where
\[
 \tau(x)=\frac12\left(1+\frac{x}{\sqrt{1+x^2}}\right)\in(0,1).
\]
The canonical omega-map inverse and
$\tau^{-1}(t)=(2t-1)/\sqrt{1-(2t-1)^2}$ recover $x$.
Finally, any set $x$ has, by Choice, a code consisting of a pointed
well-founded extensional relation on an ordinal.  Encode that relation
as a set of ordinals and then as a surreal sign sequence.  If that
surreal is ROD, the unique Mostowski collapse and its marked element
are ROD, so $x$ is ROD.  This proves (4)$\Rightarrow$(2).
\end{proof}
The same proof works relative to a fixed set of ordinals $p$: if every
set is definable from $p$, some real depending on that set, and ordinal
parameters, then $V=\HOD(p,r)$ for one real $r$.  It is not enough merely
to know that every real is available as a parameter; the hypothesis
concerns all sets, or equivalently the full surreal or bounded omnific
coding class.  Nor does the theorem identify $\HOD_{\R}$ with a fixed
$H_r$ when $\HOD_{\R}$ is a proper inner model.

\subsection{Countable birthdays, spectra, and the first omitted birthday \src{05}}
\label{dsn:ca:sec:realunion}

Source~05's Proposition~5.1 (the basic structure of $\KR$) is the paragraph
after \cref{dsn:rp:prop:hereditary-model}; its proof joins finitely many real
parameters into one and applies the fixed-real results, initiality following
either from the same argument or by taking a union of the initial classes
$K_r$.

\begin{theorem}[All countable-birthday surreals need only a real]\label{dsn:ca:thm:countbirthday}
\begin{equation}\label{dsn:ca:eq:countbirthday}
 \No_{<\omega_1}\subseteq K_{\mathrm{rd}}\subseteq\KR.
\end{equation}
The first inclusion needs no ordinal parameters.
\end{theorem}
\begin{proof}
If $x:\alpha\to\{-,+\}$ and $\alpha$ is countable, choose a well-order
on a subset of $\omega$ of type $\alpha$, together with a label on each
point giving the sign of its image under the unique order isomorphism.
One real codes the domain, relation, and labels. From that real alone,
well-founded collapse recovers the order type, its order isomorphism,
and hence the canonical sign sequence $x$. Finite and empty domains
are covered by the same tagged coding convention.
\end{proof}

This theorem concerns countable \emph{sign length}, not countable
normal-form support. The distinction is decisive in
\cref{dsn:ca:thm:prikry}.
This is also source~06's \cref{dsn:op:cor:countable-birthday-real} (Part~\ref{dsn:part:op}),
where it follows from the surjectivity of the real decoder, and source~02
states it in \cref{dsn:rp:sec:consequences}.

Source~05's Proposition~5.3 is the last clause of
\cref{dsn:rp:thm:parameter-order}, with the same proof; source~05 adds that the
witness selection can be performed in ZFC using the internally definable
certificate relation of \cref{dsn:ca:lem:cert}, Collection, and Choice.

It is tempting, but wrong, to conclude that a countable summable family
in $\KR$ has its sum in this common $K_r$. Individual definitions may
require a non-$\OD(r)$ sequence of ordinal parameters. Joining the real
parameters has not joined that ordinal information.

\paragraph{Parameter spectra.}
For a set $x$, define its real-parameter spectrum by
\[
 \Spec(x)=\{r\in\R:x\in\OD(r)\}.
\]
This is a definable subset of $\R$. On reals, put
$r\preceq s$ when $r\in\OD(s)$, equivalently $r\in\HOD(s)$. It is a
preorder. Formula substitution shows that $r\preceq s$ implies
$K_r\subseteq K_s$, so $\Spec(x)$ is upward closed in this preorder.
The join of a countable sequence of reals is an upper bound for all
its entries. No assertion about the existence of a \emph{least}
real in a spectrum is implicit in this observation.

Finite algebraic operations may lose information. For instance,
$\Spec(x)\cap\Spec(y)\subseteq\Spec(x+y)$, but equality need not hold:
$x+(-x)=0$. Reversible coding, in contrast, will preserve the entire
spectrum exactly.

\begin{dupstatement}{Source 05's Proposition 6.1 (the bounded code $\Mc$ of \cref{dsn:thm:bounded-code}, with spectra)}
Put
\begin{equation}\label{dsn:ca:eq:bounded}
 \tau(x)=\frac12\left(1+\frac{x}{\sqrt{1+x^2}}\right),
 \qquad Q(x)=\omega^{\tau(x)}.
\end{equation}
Then $Q$ is a parameter-free ambiently definable bijection from $\No$
onto the monomials $\omega^a$ with $0<a<1$, a subclass of
$\Oz\cap(0,\omega)$. Its inverse is parameter-free definable. It
preserves every parameter policy in this article, and
$\Spec(Q(x))=\Spec(x)$.
\end{dupstatement}
\begin{proof}[Second proof (source 05)]
The positive square root of $1+x^2$ is greater than $|x|$, so
$0<\tau(x)<1$. Given a monomial $z=\omega^a$ with $0<a<1$, recover $a$
by the inverse Conway map. Put $y=2a-1$, so $-1<y<1$, and recover
\[
 x=\frac{y}{\sqrt{1-y^2}}.
\]
Direct field algebra verifies that these formulas are inverse.
Unique-output closure gives parameter preservation in both directions.
\end{proof}
This is the monomial code in (Part~\ref{dsn:part:one}), not a new code
invented here. It is not a ring homomorphism and is not definable in the
pure ordered-field language. It shows, for example, that the global
conditions in \cref{dsn:ca:prop:globaltest} are equivalent to the much smaller
looking assertion that every omnific integer in $(0,\omega)$ is $\ROD$.
A bounded omnific interval can retain arbitrary ambient definability
information.

\paragraph{Bounded testing and the first omitted birthday.}
The hereditary real--ordinal class will be denoted
\[
 \HR=\{x:\TC(\{x\})\subseteq\ROD\}.
\]
We use it as a definable class, without assuming that it has a canonical
well-order or satisfies Choice. For the next statement no inner-model
axiom for $\HR$ is needed.
\begin{proposition}[First omitted real--ordinal birthday \src{05}]
\label{dsn:ca:prop:globaltest}
If $\KR\ne\No$, the least birthday of a missing surreal is
\begin{equation}\label{dsn:ca:eq:lambda}
 \lambda_{\R}
 =\min\{\alpha:\exists X\subseteq\alpha\ (X\notin\ROD)\}
 \ \ge\ \omega_1.
\end{equation}
Source~05 proves this together with the equivalence of: every surreal is
$\ROD$; every subset of every ordinal is $\ROD$; every set is $\ROD$; and
$V=\HR$, which are part of \cref{dsn:rp:cor:global-single-real}.
\end{proposition}
\begin{proof}
A subset of an ordinal and its labelled sign sequence are interdefinable
using that ordinal. This gives the surreal/subset equivalence. By Choice,
any set $a$ has a pointed well-founded extensional membership code on an
ordinal domain: well-order its transitive closure and transport
membership to that domain. A fixed ordinal coding of pairs turns this
into a subset of an ordinal, with the distinguished point and domain
included as finite ordinal parameters. If the code is $\ROD$, its
Mostowski collapse defines $a$ under the same policy. No canonical
choice of the code is needed, since the hypothesis covers all such
codes. The hereditary equivalence is immediate. The same sign coding
identifies the two least ordinals in \eqref{dsn:ca:eq:lambda}; the lower bound
is \cref{dsn:ca:thm:countbirthday}.
\end{proof}
The predicate in \eqref{dsn:ca:eq:lambda} is internally definable, so the least
ordinal exists whenever there is a witness. This is not a ``least
undefinable object'' paradox: the least \emph{birthday} is definable,
while it need not pick a unique missing sign sequence. We do not claim
that $\lambda_{\R}$ must be a cardinal in every universe;
\cref{dsn:ca:thm:cohenconsequences} computes it in a specific model.

\begin{remark}[The real-parameter field is a support field \writetag]
\label{dsn:rp:rem:support}
The field $\KR$ satisfies the hypothesis \eqref{dsn:eq:Khyp} of the
support-subfield amplification theorem \cref{dsn:thm:amplify}, so if
$V\ne\HR$, then for every $t\in(0,1)\setminus\KR$ the family
$z_\alpha=\omega+\omega^{te_\alpha}$ $(\alpha\in\Ord)$ of
\eqref{dsn:eq:indfam} lies in $\Cell$ and is algebraically independent over
$\KR$. This was observed at the write (4~October 2026); neither source
states it.
\end{remark}
\begin{proof}
$\KR$ is a subfield of $\No$ (\cref{dsn:rp:prop:hereditary-model} and the
paragraph after it). The surreal $\omega$ is parameter-free definable, and
$e_\alpha=\omega^{-(\alpha+1)}$ is definable from the ordinal $\alpha$, so
both lie in $\KH\subseteq\KR$. If $0\ne a\in\KR$, say $a\in K_r$, then every
exponent of $a$ is definable from $a$ and its ordinal index in the normal
form (the whole-normal-form paragraph after
\cref{dsn:rp:prop:hereditary-model}), hence lies in $K_r\subseteq\KR$; so
$\supp(a)\subseteq\KR$. Finally $\KR\ne\No$ exactly when $V\ne\HR$, by
\cref{dsn:rp:cor:global-single-real}, and then a $t\in(0,1)\setminus\KR$
exists: if $x\notin\KR$ then $\tau(x)\notin\KR$, because $\KR$ is real closed
and $x=\tau^{-1}(\tau(x))$ is given by the algebraic formula
\eqref{dsn:eq:tauinv}. \Cref{dsn:thm:amplify} applies.
\end{proof}

\section{Countable information, Hahn sums, and countable cuts}
\label{dsn:rp:sec:countable}

The preceding pointwise amalgamation theorem does not supply the sequence
of definitions as one object. We now characterize exactly the missing
condition. Define the first-order principle
\begin{equation}\label{dsn:rp:eq:countable-principle}
 \Asm:\quad\text{every countable set of ordinals belongs to }\ROD.
\end{equation}
Here and throughout this section, countability is measured in $V$.
A countable sequence $f$ of ordinals can be coded by the countable set
$\{\lambda\cdot n+f(n):n<\omega\}$ for an ordinal $\lambda$ exceeding all
$f(n)$. Conversely, a countable set is the range of a countable ordinal
sequence. Thus $\Asm$ is equivalent to condition (1) below. Ordinal
multiplication in this coding sentence is set-theoretic ordinal
multiplication, not surreal multiplication.

Work in an ambient universe of ZFC.  Write $\operatorname{OD}(r)$ for
ordinal definability with one additional real parameter $r$, and put
\[
 K_r=\No\cap\operatorname{OD}(r),\qquad
 \KR=\bigcup_{r\in\mathbb R}K_r,
 \qquad
 \IR=\Oz\cap\KR.
\]
Here $\HR$ means hereditary definability
allowing an arbitrary real and finitely many ordinals at each object.
It is not shorthand for one fixed-real HOD.

\begin{theorem}[Countable assembly criterion]\label{dsn:rp:thm:countable}
The following conditions are equivalent.
\begin{enumerate}
\item Every function $f:\omega\longrightarrow\operatorname{Ord}$ is
ordinal definable from a real.
\item $\HR$ is closed under the ambient
countable sequences of its elements.
\item Every ambient countable sequence of elements of
$\KR$ has an entire sequence code ordinal definable
from a real.
\item $\KR$ is closed under arbitrary countable
strongly summable families of its elements.
\item Every countable strongly summable family of ordinal-definable
purely infinite omnific integers has its sum in
$\IR$.
\item Whenever $A,B\subseteq\KR$ are countable
sets and $A<B$, some $y\in\KR$ satisfies $A<y<B$.
\end{enumerate}
In (6), countability is computed in the ambient universe.  Thus (6)
is the ambient $\eta_1$ interpolation property.
\end{theorem}

\begin{proof}
Assume (1), and let $(x_n)_{n<\omega}$ be a sequence of elements of
$\HR$.  For each $n$, choose a rank-local
definition of $x_n$: a real $r_n$, a rank $\theta_n$, a finite tuple
of ordinal parameters, and a formula code that singles out $x_n$
inside $V_{\theta_n}$.  Collection and ambient Choice permit these
countably many choices.  Join the reals and the natural-number formula codes into one
real.  The ordinal witnesses form a single countable ordinal sequence,
which, by (1), is itself ordinal definable from a real.  Join the two
real parameters.  Set-sized satisfaction then gives a uniform definition
of the whole sequence $(x_n)_{n<\omega}$.  Its transitive closure contains
only the $x_n$, their hereditary descendants, and the elementary sequence
coding apparatus.  Consequently the sequence is in
$\HR$.  This proves (2).

Condition (2) implies (3), since a surreal sign sequence is in
$\KR$ precisely when it belongs to
$\HR$.  Condition (3) implies (4), because
the strongly summable family is now available as one set-coded input
and Hahn summation is a parameter-free, uniquely determined operation.
Condition (4) implies (5), because the sum of a strongly summable
family of purely infinite surreals is purely infinite (or zero).

We prove (5)$\Rightarrow$(1) by an explicit code.  Given
$f:\omega\to\operatorname{Ord}$, regard each ordinal $f(n)$ as its
canonical surreal, and set
\[
 q_n=2^{-(n+1)},\qquad
 g_n=q_n+\frac{q_n}{4}\frac{f(n)}{1+f(n)},\qquad
 t_n=\omega^{g_n}.
\]
Since $0\le f(n)/(1+f(n))<1$,
\[
 q_n\le g_n<\frac54q_n,\qquad
 g_{n+1}<\frac58q_n<g_n.
\]
The family $(t_n)_{n<\omega}$ is strongly summable; each $t_n$ is an
ordinal-definable purely infinite omnific integer.  Its sum
\begin{equation}\label{dsn:rp:eq:ordinal-sequence-code}
 Z_f=\sum_{n<\omega}\omega^{g_n}
\end{equation}
belongs to $(0,\omega)\cap\Oz$, with the displayed expression
as its exact normal form.  From its $n$th exponent recover
\[
 u_n=4(g_n/q_n-1),\qquad f(n)=\frac{u_n}{1-u_n}.
\]
The canonical ordinal embedding has a definable inverse on its range.
Thus $Z_f$ and $f$ are interdefinable without additional parameters.
Condition (5) gives (1).

Condition (3) implies (6): enumerate the two countable sets, join their
whole sequence codes, and take the unique simplest surreal between
them.  This construction is definable from that set-coded cut.

For (6)$\Rightarrow$(1), use the same $t_n$ and put
$s_n=\sum_{k<n}t_k$.  Set
\[
 A_f=\{s_{n+1}:n<\omega\},\qquad
 B_f=\{s_n+2t_n:n<\omega\}.
\]
Every endpoint is ordinal definable using only finitely many entries
of $f$.  The full sum $Z_f$ witnesses $A_f<B_f$.  By (6), choose
$y\in\KR$ between these sets.  For every $n$,
\[
 0<y-s_{n+1}<2t_{n+1}\prec t_n.
\]
Hence the first $n+1$ normal-form terms of $y$ are exactly
$t_0,\ldots,t_n$.  Its first $\omega$ terms therefore recover all the
$g_n$ and then $f$.  Any additional normal-form tail occurs below
every $g_n$ and does not affect decoding.  Thus $f$ is ordinal
definable from a real, proving (1).
\end{proof}

\begin{remark}
Even without the equivalent conditions, every countable set of real
parameters has a join $r$, so every countable collection of individual
$\KR$ elements is contained pointwise in one
$K_r$.  This does not in general put the whole chosen
sequence in $\operatorname{HOD}(r)$: the ordinal witnesses may carry
additional countable information.  The theorem identifies precisely
the extra assembly principle.
\end{remark}


\subsection{Two further forms of the same principle}

\begin{corollary}\label{dsn:rp:cor:countable-variants}
Each condition of \cref{dsn:rp:thm:countable} is also equivalent to each of:
\begin{enumerate}
\item The transitive union $\Hcup=\bigcup_r H_r$ is closed under all ambient
$\omega$-sequences of its elements.
\item $\HR=\HOD(\Ord^\omega)$, where the right side permits a countable
sequence of ordinals and finitely many ordinals as parameters.
\end{enumerate}
Under these conditions, $\HR$ satisfies Dependent Choice.
\end{corollary}
\begin{proof}
Assume $\Asm$ and $x_n\in\Hcup$. Choose reals $r_n$ with $x_n\in H_{r_n}$
and join them into $r$. Every $x_n$ and all its hereditary descendants
then belong to $H_r$. As in the theorem, use $\Asm$ to assemble the
ordinal records of the definitions and join their real witness to $r$.
The whole graph is now $\OD(t)$ for one real $t$, and its descendants
belong to $H_t$. Hence the graph belongs to $H_t\subseteq\Hcup$.
Conversely, ordinal sequences are sequences of elements of $\Hcup$,
so closure implies their real-plus-ordinal definability.

For the second equivalence, reals are countable ordinal sequences, so
$\ROD\subseteq\OD(\Ord^\omega)$. Under $\Asm$ every allowed ordinal-sequence
parameter on the right is itself in $\ROD$, and substitution gives
$\OD(\Ord^\omega)=\ROD$ and the hereditary equality. Conversely each
ordinal-sequence graph is hereditarily definable from itself, so that
equality makes every such graph $\ROD$.

Finally, let a nonempty set and a serial relation on it belong to
$\HR$. In the ambient universe choose a dependent sequence, starting
at any specified member. All its entries lie in $\HR$, and the
countable-sequence closure just proved puts the whole sequence in
$\HR$. This proves Dependent Choice there. We assert no converse from
Dependent Choice alone.
\end{proof}

\begin{remark}[No automatic countable least join]
Interleaving a specified countable sequence of real parameters gives an
upper bound for their ordinal-definability degrees, but need not give
their least upper bound. In a universe with a non-OD real, its sequence
of bits consists of individually parameter-free definable bit values;
their interleaving can nevertheless encode that non-OD real. Finite
joins and countable sequence assembly have different logical content.
\end{remark}

\subsection{What a countable sum can carry}
The code $Z_f$ in \eqref{dsn:rp:eq:ordinal-sequence-code} is uniformly
interdefinable with $f$, without any additional parameters. Its
individual monomials require only one ordinal each. This is a precise
explanation for why a countable sum can leave a definability class:
the sum can recover the entire ordering and association of infinitely
many ordinal choices. The failure is not caused by divergent series;
the family is Hahn summable with an explicitly verified support.

The interpolation implication is also exact. It does not require the
interpolant to equal the sum. Every possible interpolant has the same
first $\omega$ normal-form terms and therefore the same coded
information. Allowing an arbitrary smaller tail cannot hide the code.

% ===== End sections/03_countable_closure.tex =====


% ===== Begin sections/04_cohen.tex =====

\subsection{Source 05's separated-band multiplexer \src{05}}
\label{dsn:ca:sec:codes}

Source~05 builds the code of \cref{dsn:rp:thm:countable} as a separate
theorem, with dyadic bands $[2^{-n-1},3\cdot2^{-n-2})$ in place of source~02's
slots $[q_n,\tfrac54q_n)$; it also records the parameter spectra of
\cref{dsn:ca:sec:realunion}.

The general code does not by itself expose a sum of individually
ordinal-definable terms. We now give a version tailored to that task.
For a nonnegative surreal $a$ set
\begin{equation}\label{dsn:ca:eq:band}
 u(a)=\frac{a}{1+a},\qquad
 b_n(a)=2^{-n-1}+2^{-n-2}u(a)\quad(n<\omega).
\end{equation}
For $f:\omega\to\Ord$, interpreting $f(n)$ as a surreal ordinal, put
\begin{equation}\label{dsn:ca:eq:code}
 \boxed{\displaystyle
 \Zc(f)=\sum_{n<\omega}\omega^{b_n(f(n))}.}
\end{equation}
No monotonicity assumption is made on $f$.

\begin{theorem}[Countable ordinal multiplexer]\label{dsn:ca:thm:code}
For every $f:\omega\to\Ord$ the expression \eqref{dsn:ca:eq:code} is a
well-defined omnific integer in $(0,\omega)$. Its normal form has
support of order type $\omega$ in decreasing order and all coefficients
are $1$. Every individual monomial in the display is $\OD$.
The whole set $f$ and the whole surreal $\Zc(f)$ are interdefinable
by fixed parameter-free ambient formulas. In particular
\begin{equation}\label{dsn:ca:eq:codeiff}
 \Zc(f)\in\KR\ \Longleftrightarrow\ f\in\ROD,
 \qquad \Spec(\Zc(f))=\Spec(f).
\end{equation}
\end{theorem}
\begin{proof}
For $a\ge0$ one has $0\le u(a)<1$, whence
\begin{equation}\label{dsn:ca:eq:bounds}
 2^{-n-1}\le b_n(a)<3\cdot2^{-n-2}.
\end{equation}
The entire band for index $n+1$ lies below the band for $n$, since
\[
 3\cdot2^{-n-3}<2^{-n-1}.
\]
Thus the sequence $b_n(f(n))$ is strictly decreasing regardless of the
values or repetitions in $f$. Any nonempty subset of this support has
a largest element, namely the one with the least natural index. The
normal form is therefore valid and Hahn-summable. All its exponents
are positive, so it is omnific with zero constant term. Its leading
exponent is less than $3/4<1$; comparison of leading terms gives
$0<\Zc(f)<\omega$.

Fixing $n$, the term $\omega^{b_n(f(n))}$ is definable from the ordinal
$f(n)$ and the natural numeral $n$. This proves individual $\OD$
definability without asserting a uniform $\OD$ assignment of these
ordinals as $n$ varies.

Conversely, the unique normal form of $z=\Zc(f)$ supplies its $n$th
exponent $a_n$. Set
\begin{equation}\label{dsn:ca:eq:decode}
 v_n=2^{n+2}\bigl(a_n-2^{-n-1}\bigr),\qquad
 f(n)=\frac{v_n}{1-v_n}.
\end{equation}
The result is the all-plus surreal corresponding to the original
ordinal, whose domain recovers that ordinal as a set. Normal-form
uniqueness and the disjoint bands make this decoder unambiguous.
Replacement over $n<\omega$ recovers the whole graph of $f$.
Encoding and decoding are fixed definable operations, so each direction
preserves all permitted parameters.
\end{proof}

\begin{remark}[Why this is stronger than small coefficients]
The construction requires no nondefinable real coefficients, no
uncountable support, and no freedom to permute the support after it is
chosen. The entire possible failure of $\ROD$ definability is the
countable sequence of ordinal information hidden \emph{inside the
exponents}. The real leading parts of their bands are fixed in advance.
Even a repeated or oscillating ordinal sequence is faithfully encoded.
\end{remark}

\begin{example}[A finite check of the formulas]\label{dsn:ca:ex:finite}
For the initial values $f(n)=n$, the first four exponents are
\[
 b_0(0)=\tfrac12,\quad b_1(1)=\tfrac5{16},\quad
 b_2(2)=\tfrac16,\quad b_3(3)=\tfrac{11}{128}.
\]
Substituting each one into \eqref{dsn:ca:eq:decode} recovers $0,1,2,3$ exactly.
The supplied verification script tests the rational formula and band
separation on many finite examples. It does not implement infinite
ordinal arithmetic, decide definability, or prove the forcing results.
\end{example}

\subsection{Further equivalent forms of countable assembly \src{05}}
\label{dsn:ca:sec:assembly}

\begin{definition}
The \emph{real--ordinal assembly principle} $\Asm$ is the assertion
\[
 \text{every function }f:\omega\to\Ord\text{ belongs to }\ROD.
\]
Quantification is over set functions: every such function has an ordinal
bound on its range. This is an ordinary first-order assertion of ZFC,
using the certificate predicate for $\ROD$.
\end{definition}

The principle does not say that every ordinal is $\ROD$; that is automatic.
It says that the entire \emph{assignment of ordinals to natural indices}
can be described using one real and only finitely many ordinal parameters.
The principle $\Asm$ is equivalent to source~02's
\eqref{dsn:rp:eq:countable-principle}, as explained there; the two sources
state it for sequences and for sets respectively.

\begin{theorem}[Countable-assembly equivalence: source 05's further forms \src{05}]
\label{dsn:ca:thm:assembly}
Over ZFC the conditions of \cref{dsn:rp:thm:countable} are equivalent to
each of the following.
\begin{enumerate}[label=\textup{(\roman*)}]
\setcounter{enumi}{2}
\item Whenever $(x_n)_{n<\omega}$ is a sequence of sets with each
$x_n\in\ROD$, the whole sequence is $\ROD$.
\setcounter{enumi}{4}
\item For every $f:\omega\to\Ord$, the particular omnific integer
$\Zc(f)$ in \eqref{dsn:ca:eq:code} belongs to $\IR$.
\item $\ROD=\COD$.
\item $\KR=\Kco$.
\end{enumerate}
Source~05's clauses \textup{(i)} $\Asm$, \textup{(ii)} every countable set of
ordinals is $\ROD$, and \textup{(iv)} every countable Hahn-summable sequence
of elements of $\KR$ has its sum in $\KR$, are \cref{dsn:rp:thm:countable}~(1),
\eqref{dsn:rp:eq:countable-principle} and \cref{dsn:rp:thm:countable}~(4). In
particular, testing countable sums of the very special individually $\OD$
monomials in \eqref{dsn:ca:eq:code} is already sufficient to test arbitrary
countable Hahn-sum closure of the whole field $\KR$. Clause \textup{(vi)} is
also source~02's \cref{dsn:rp:cor:countable-variants}~(2).
\end{theorem}
\begin{proof}[Second proof of \cref{dsn:rp:thm:countable} with these clauses (source 05)]
We prove the implications in several explicit steps.

\emph{(i)$\Rightarrow$(iii): assembling definitions, not merely real
parameters.} Let $x_n\in\ROD$. For each $n$ select a rank-local
certificate consisting of a real $r_n$, a rank ordinal $\theta_n$, a
formula code $e_n$, and a finite ordinal tuple $\bar\alpha_n$ uniquely
specifying $x_n$ in $V_{\theta_n}$. Selection is legitimate: the
certificate relation is internally definable, Collection first bounds
a set of witnesses for the countably many inputs, and Choice selects
one for each index.

Encode the sequence $(r_n)$, the natural formula codes, and the lengths
of the finite tuples by a real $r_*$. Flatten all rank ordinals
$\theta_n$ and all entries of $\bar\alpha_n$ into one sequence
$g:\omega\to\Ord$, using the tuple lengths to delimit the blocks and
padding unused positions by $0$. The exact pairing convention is fixed
once and for all. By (i), $g$ is $\OD(s)$ for some real $s$ and a finite
ordinal tuple $\bar\beta$. Join $s$ and $r_*$ into a real $t$.

From $t,\bar\beta$ recover $g$, then all the rank-local certificates.
Uniform satisfaction for the set structures $V_{\theta_n}$ recovers
each unique $x_n$ using one formula, uniformly in $n$. Replacement
therefore yields the whole sequence, which is definable from
$t,\bar\beta$. This uses no satisfaction predicate for the full
universe. The rank sequence is precisely the information that (i)
allows us to compress.

\emph{(iii)$\Rightarrow$(iv).} Apply (iii) to the sequence of canonical
surreal codes. The whole family is $\ROD$, and its Hahn sum is a fixed
uniquely definable operation on that family, by
\cref{dsn:ca:def:summable,dsn:lem:closure}.

\emph{(iv)$\Rightarrow$(v)$\Rightarrow$(i).} Each term in
$\Zc(f)$ is $\OD$ and hence in $\KR$, and the family is
Hahn-summable. Thus (iv) implies (v). The decoder
\eqref{dsn:ca:eq:decode} then recovers $f$ from $\Zc(f)$, giving (i).
This is the crucial reduction from unrestricted sums to a very narrow
omnific test family.

\emph{(i)$\Leftrightarrow$(ii).} Under (i), a countable set of ordinals
has an enumeration whose graph is $\ROD$, so its range is $\ROD$;
the empty set is immediate. Conversely assume (ii), and let
$B=\operatorname{ran}(f)$. The increasing enumeration of $B$ has a
countable ordinal domain $\delta$. A real can code a well-order on a
subset of $\omega$ of type $\delta$, together with the index in that
well-order corresponding to each $f(n)$. Decoding this real and the
$\ROD$ set $B$ recovers $f$ using one real and the finite ordinal
parameters in a definition of $B$. This gives (i).

\emph{(i)$\Leftrightarrow$(vi).} A real is a countable sequence of bits,
so a real and any finite ordinal tuple can be joined into one countable
ordinal sequence. Hence $\ROD\subseteq\COD$ always. Under (i), any
countable ordinal parameter sequence is itself $\ROD$, and substituting
its definition shows $\COD\subseteq\ROD$. Conversely any
$f:\omega\to\Ord$ is $\COD$ from the parameter $f$ itself, so (vi)
implies (i).

Finally (vi) implies (vii). If (vii) holds, then for every $f$ its code
$\Zc(f)$ belongs to $\Kco$ by definition and therefore to $\KR$;
being omnific, it belongs to $\IR$. Thus (vii) implies (v), completing
the equivalence.
\end{proof}

\begin{corollary}[The integer ring detects the same obstruction]\label{dsn:ca:cor:ringassembly}
The conditions of \cref{dsn:ca:thm:assembly} are also equivalent to
$\IR=\Ico$, and to closure of $\IR$ under countable Hahn-summable
families of its elements.
\end{corollary}
\begin{proof}
A Hahn-summable family of omnific integers has no negative exponents,
and only finitely many nonzero contributions at exponent $0$. Its
constant coefficient is therefore an integer, so its sum is omnific.
Field closure implies ring closure. Ring closure implies condition
(v), since the test monomials are omnific. Equality of the integer
rings similarly implies (v), and field equality implies equality of
the rings by intersection with $\Oz$.
\end{proof}

\begin{remark}[Whole-family closure remains unconditional]
Even when $\Asm$ fails, a summable family whose \emph{whole graph} is
$\ROD$ has a $\ROD$ sum. Failure concerns pointwise membership only.
Thus the theorem does not contradict standard closure under a specified
definable transfinite construction. It identifies when an unspecified
countable choice of definitions can be turned into one specified
construction.
\end{remark}

\paragraph{A useful covering criterion.}
\begin{proposition}[Countable ordinal covering]\label{dsn:ca:prop:cover}
Suppose every countable set of ordinals is contained in a countable
ordinal-definable set of ordinals. Then $\Asm$ holds, hence $\KR$ and
$\IR$ are countably Hahn-sum closed.
\end{proposition}
\begin{proof}
Let $f:\omega\to\Ord$, and choose a countable $\OD$ set $B$ containing
its range. As in the proof of (ii)$\Rightarrow$(i), a real codes the
countable order type of $B$ together with the indices of the values
$f(n)$. The increasing enumeration of $B$ is definable from $B$.
Combining the real code with an ordinal definition of $B$ gives an
$\ROD$ definition of $f$.
\end{proof}

The hypothesis is a sufficient \emph{covering} principle, not the
assertion that a single real codes all countable ordinal sequences.
Different input sequences may require different reals, ordinal
parameters, and ordinal-definable covers.

\subsection{A canonical countable-assembly completion \src{05}}
\label{dsn:ca:sec:completion}


The class $\COD$ is the classical class of sets definable from countable
ordinal sequences; compare \cite[III.2.2]{Solovay}. It is useful here
because countably many such parameter sequences can themselves be
flattened into one. Write
\[
 \Kco=\No\cap\COD,\qquad \Ico=\Oz\cap\COD.
\]
This is an enlargement of $\KR$, not a change in the ambient arithmetic.

\begin{lemma}[Countable-ordinal assembly]\label{dsn:ca:lem:codclosure}
A countable sequence of $\COD$ sets is itself $\COD$.
\end{lemma}
\begin{proof}
Choose rank-local definition certificates for its entries, with
countable ordinal parameter sequences. Flatten the countable matrix
of parameter ordinals, the rank ordinals, and the natural syntax codes
into one countable ordinal sequence. This single parameter recovers
all the certificates. Uniform set satisfaction and Replacement recover
the whole input sequence. The witness selection is justified by
Collection and Choice exactly as in \cref{dsn:ca:thm:assembly}. This is the
familiar ordinal-sequence assembly mechanism underlying
\cite[III.2.6]{Solovay}.
\end{proof}

\begin{definition}[Ambient stability]\label{dsn:ca:def:stable}
A subclass $H\subseteq\No$ is \emph{ambiently stable} if, for every
fixed parameter-free set-theoretic formula defining a partial
surreal-valued function on a finite tuple, applying that function to
inputs from $H$ gives an element of $H$ whenever the output is unique
and exists. This is a scheme, not a putative uniform truth predicate.
\end{definition}

This is deliberately stronger than closure under the field operations
alone. Without this qualification the minimality statement below would
not follow: an arbitrary ambient definition cannot be recovered by
field arithmetic merely because its ordinal parameter has been encoded
in a monomial.

\begin{theorem}[Canonical ambient-stable Hahn completion]\label{dsn:ca:thm:completion}
The class $\Kco$ is the least subclass of $\No$ which contains $\KR$,
is ambiently stable, and is closed under countable Hahn sums. More
explicitly, any subclass $H$ with these three properties contains
$\Kco$. Moreover $\Kco$ is initial and real closed,
$\Kco=\Frac(\Ico)$, and $\Ico$ is countably Hahn-sum closed.
\end{theorem}
\begin{proof}
The finite-parameter combination property for $\COD$,
\cref{dsn:lem:closure,dsn:ca:lem:codclosure}, and definability of summation prove
ambient stability and countable Hahn-sum closure. The algebraic
properties follow from \cref{dsn:thm:hull}; initiality follows by
adjoining an ordinal prefix length to the countable ordinal parameter.
The integer-ring assertion follows as in \cref{dsn:ca:cor:ringassembly}.

For minimality, let $x\in\Kco$ and choose
$f:\omega\to\Ord$ and a fixed formula $\varphi(y,f)$ which uniquely
defines the surreal $x$. Every term of $\Zc(f)$ is in $\KH$ and
hence in $\KR\subseteq H$. Countable Hahn-sum closure places
$\Zc(f)$ in $H$.

There is a fixed parameter-free partial function on surreals which
first checks whether its input is in the range of the code, decodes
$f$ by \eqref{dsn:ca:eq:decode}, and returns the unique surreal satisfying
$\varphi(y,f)$, when there is such a unique surreal. Its graph is a
single set-theoretic formula with $\varphi$ substituted literally.
It may equivalently be totalized by returning $0$ on other inputs.
Ambient stability now gives $x\in H$. No additional parameter and no
uniform evaluator for all formulas has been introduced.
\end{proof}

\begin{corollary}[The completion has no further countable-assembly stages]
Applying the same ambient-stability and countable-Hahn-closure operation
to $\Kco$ produces $\Kco$ again. In a universe satisfying $\Asm$,
$\KR$ is already this completion; in the Prikry model below the
inclusion $\KR\subsetneq\Kco$ is strict.
\end{corollary}

\begin{remark}[What the word completion does not mean]
This is a minimality statement for a specified closure scheme. It is
not a Dedekind completion, an order-topological completion, a completion
of a metric space, or a claim that $\Kco$ is closed under every
uncountable summable family. It also does not assert that the field
operations and countable sums \emph{alone} generate $\Kco$.
\end{remark}

\begin{remark}[The completion needs only the ordinals \writetag]
\label{dsn:ca:rem:ordsuffice}
In \cref{dsn:ca:thm:completion} the hypothesis ``contains $\KR$'' can be
weakened to ``contains every ordinal'': $\Kco$ is the least subclass of $\No$
that contains every ordinal, is ambiently stable, and is closed under
countable Hahn sums. Hence such a subclass contains all of $\KR$, and in
particular every real. This was observed at the write (4~October 2026); the
proof is source~05's.
\end{remark}
\begin{proof}
$\Kco$ contains every ordinal, so by \cref{dsn:ca:thm:completion} it has the
three properties. Conversely let $H$ have them, and let $x\in\Kco$ be the
unique $y$ with $\varphi(y,f)$, where $f:\omega\to\Ord$. For each fixed $n$
the map $\alpha\mapsto\omega^{b_n(\alpha)}$ is given by one parameter-free
formula (the numeral $n$ is written into it), so ambient stability applied to
the input $f(n)\in\Ord\subseteq H$ puts the $n$th term of $\Zc(f)$ in $H$.
Countable Hahn-sum closure puts $\Zc(f)$ in $H$, and the fixed decoding map
of the proof of \cref{dsn:ca:thm:completion} puts $x$ in $H$. Only the
ordinals, not all of $\KR$, were used.
\end{proof}

\subsection{Certificate coding details \src{05}}
\label{dsn:ca:app:certificates}

This is source~05's Appendix~A.

Here is one completely definite version of the bookkeeping used in
\cref{dsn:ca:thm:assembly}. Fix a recursive pairing bijection
$\pi:\omega^2\to\omega$. A real $r_*$ has disjoint tagged columns
coding the bits of each $r_n$, the formula code $e_n$, and the tuple
length $\ell_n$. Define an ordinal sequence $g$ by
\[
 g(\pi(n,0))=\theta_n,\qquad
 g(\pi(n,j+1))=
 \begin{cases}
 \alpha_{n,j},&j<\ell_n,\\
 0,&j\ge\ell_n.
 \end{cases}
\]
The natural-number information in $r_*$ tells the decoder which zeroes
are genuine entries and which are padding. Once $g$ and $r_*$ are
available, every local formula, rank, and finite ordinal tuple is
uniformly available.

The statement that $x_n$ is the unique output is checked using
$\Sat(V_{\theta_n},e_n,x_n,r_n,\bar\alpha_n)$. The outer operation is
one fixed formula concerning satisfaction in a set, not an evaluation
of $\varphi_{e_n}$ in $V$. The only global formula that is substituted
in the proof is the single selected definition of $g$ from a real and
finitely many ordinals. This is why no uniform global truth predicate
is needed.

For $\COD$, use an additional pairing layer to store the countable
ordinal sequence parameter for each $x_n$ and interleave the ranks,
syntax codes, and tuple lengths as further ordinal entries. All
finite syntactic data are natural numbers and therefore ordinals.
One $\omega$-sequence still suffices. This makes explicit that the
flattening does not require an ordinal-indexed family of real
parameters or a class choice principle.

\section{A proper real-parameter field with exact countable closure}
\label{dsn:rp:sec:cohen}

The following construction uses standard Cohen forcing and its
homogeneity. Its purpose is to compute the full real-parameter field
hierarchy and its countable behavior, rather than to claim a new forcing
notion. For background on the forcing machinery see
\cite{Kunen,KaragilaForcing}.

Let $W\models\mathrm{ZFC}+V=L$, let $\lambda$ be an uncountable
cardinal of $W$, and force with
\[
 \mathbb P=\operatorname{Fn}(\lambda\times\omega,2,{<}\omega).
\]
Write $V=W[G]$, $c_\xi(n)=(\bigcup G)(\xi,n)$, and for ground-model
countable $S\subseteq\lambda$ let
$W_S=W[G\upharpoonright(S\times\omega)]$.
Using a canonical constructible enumeration of $S$, code this countable
coordinate generic by a real $r_S$, so $W_S=L[r_S]$.

\begin{theorem}[Exact atlas and its inclusion order]\label{dsn:rp:thm:cohen-atlas}
In $V$,
\[
 K_{r_S}=\No^{W_S},\qquad
 \KR
 =\bigcup_{S\in([\lambda]^{\le\omega})^W}\No^{W_S}.
\]
Furthermore,
\[
 K_{r_S}\subseteq K_{r_T}
 \quad\Longleftrightarrow\quad S\subseteq T.
\]
All equalities concern canonical sign sequences and ambiently absolute
surreal operations.
\end{theorem}

\begin{proof}
The forcing is ccc.  A nice name for a real uses countably many maximal
antichains, each countable, and thus only a ground-model countable set
of coordinates.  Consequently every real $r\in V$ belongs to some
$W_S$.

Factor the forcing over $W_S$ into the remaining Cohen coordinates.
This tail is weakly homogeneous.  If a set of ordinals is definable
in $V$ from $r\in W_S$ and ordinals, homogeneity decides each
membership assertion at the top condition: compatible images of
conditions prevent two contrary decisions about check parameters.
The forcing relation in $W_S$ therefore reconstructs that set there.
The sign code of a surreal is a set of ordinals up to fixed definable
coding.  It follows that
$K_r\subseteq\No^{W_S}$.

For $r=r_S$ this gives one inclusion in the first equality.  Conversely,
the canonical $L[r_S]$ well-order shows that every element of $W_S$
is ordinal definable in $V$ from $r_S$.  The first equality follows,
and the real-name support argument proves the union formula.

If $S\subseteq T$, inclusion of the intermediate models gives field
inclusion.  If $\xi\in S\setminus T$, the real $c_\xi$ belongs to
$\No^{W_S}$ but is Cohen generic over $W_T$ and is not in
$W_T$.  Thus field inclusion fails.
\end{proof}

\begin{corollary}[The $\omega_1$ chain]\label{dsn:rp:cor:cohen-chain}
For $\lambda=\omega_1^W=\omega_1^V$, let $r_\delta$ code the
coordinates below $\delta$.  Then
\[
 K_{r_\delta}=\No^{L[G_{<\delta}]},\qquad
 K_{r_\delta}\subsetneq K_{r_{\delta+1}},\qquad
 \KR=\bigcup_{\delta<\omega_1}K_{r_\delta}.
\]
The preorder of real-parameter fields, ordered by inclusion, has
cofinality exactly $\omega_1$.  It has no largest field.
\end{corollary}

\begin{proof}
Every ground countable set of coordinates is bounded in $\omega_1$,
giving the cofinal chain, and $c_\delta$ witnesses strictness.  A
countable cofinal family of real parameters would have a real join;
its field would then contain every field in the family and so be
largest, contradicting strict extension by a fresh Cohen coordinate.
\end{proof}

\begin{proposition}[The general cofinality invariant]
Let $\mathscr D$ be the preorder of the fields $K_r$ in
the Cohen extension above.  Then
\[
 \operatorname{cof}(\mathscr D)
 =\operatorname{cof}\bigl(([\lambda]^{\le\omega})^W,\subseteq\bigr).
\]
The right-hand cofinality has the same value in $W$ and $V$.
In particular it is $\lambda$ when $W=L$ and $\lambda$ is regular
uncountable.
\end{proposition}

\begin{proof}
The canonical fields form a cofinal copy of the displayed support
poset.  Conversely, for each member $K_r$ of a cofinal
family choose a countable ground support $S_r$ containing a name for
$r$.  Any canonical field $K_{r_T}$ is contained in some
$K_r$, hence in $K_{r_{S_r}}$, and therefore
$T\subseteq S_r$.  Thus the chosen supports are cofinal.

To see preservation, let a name enumerate a new cofinal family in the
ground support poset using $\mu$ indices.  A ccc maximal antichain
deciding each value yields only countably many possible ground values
per index.  The ground union of all these possibilities has size at
most $\mu\cdot\aleph_0$ and is cofinal; otherwise a condition forcing
cofinality would contradict a missed ground support.  The reverse
inequality holds because an old cofinal family remains cofinal in
the unchanged ground poset.  Finally, in $L$, regular uncountable
$\lambda$ satisfies $\lambda^{\aleph_0}=\lambda$ by GCH.  The
support poset has cardinality $\lambda$, and any cofinal family must
cover $\lambda$ by its countable members, so its cofinality is
$\lambda$.
\end{proof}

\begin{theorem}[Countable closure and its exact first failure]\label{dsn:rp:thm:cohen-threshold}
In the Cohen extension with $\lambda=\omega_1$,
\begin{enumerate}
\item Every countable sequence of ordinals is ordinal definable from a
real, and all the countable assembly conditions above hold.
\item Every surreal of countable birthday belongs to
$\KR$.
\item The least birthday of a surreal outside
$\KR$ is exactly $\omega_1$.
\item The order $\KR$ fills every external countable cut but omits a cut with $\aleph_1$ endpoints (the $\eta_1$ but not $\eta_2$ property).
\end{enumerate}
\end{theorem}

\begin{proof}
For a name $\dot f$ of an ordinal sequence, choose for each $n$ a
maximal antichain deciding $\dot f(n)$.  The ccc makes their total
coordinate support countable, so $f\in L[G_S]=L[r_S]$ and hence
$f\in\operatorname{OD}(r_S)$.  This proves (1).

For (2), any countable sign sequence can be encoded by a real coding
a countable well-order of its domain and the sign assignment.  The
Mostowski collapse decodes the actual ordinal domain.  Such a surreal
is definable from that real, even without additional ordinal parameters.

For (3), let
\[
 A=\{\xi<\omega_1:c_\xi(0)=1\}.
\]
This set is in no $W_S$ with $S$ ground countable.  Indeed, for any
$B\subseteq\omega_1$ in $W_S$, the tail forcing has dense conditions
disagreeing with $B$ at a fresh coordinate outside $S$.  Hence $A$ is
not real-ordinal definable by the atlas theorem.  Its sign sequence
$s_A:\omega_1\to\{-,+\}$ is likewise outside
$\KR$, and has birthday $\omega_1$.

The $\eta_1$ assertion follows from countable assembly.  To disprove
$\eta_2$, take the canonical cut of $s_A$: a proper prefix
$s_A\upharpoonright\alpha$ goes on the left when $s_A(\alpha)=+$
and on the right when $s_A(\alpha)=-$.  Every prefix has countable
birthday and lies in $\KR$.  Any interpolant
must extend the entire sign sequence $s_A$; a first disagreement,
or an early end, violates the corresponding cut inequality.  But
$\KR$ is initial: a prefix of a real-ordinal
definable sign sequence is definable using the same real and the
prefix length as an additional ordinal.  Thus an interpolant there
would put $s_A$ there, a contradiction.  This is an unfilled cut of
size $\omega_1$.
\end{proof}

\begin{remark}[Do not extend the atlas equality to all hereditary sets]
The full set $\mathbb R^V$ belongs to
$\HR^V$: it is definable without parameters
and each of its elements is definable using itself.  Yet no countable
coordinate model $W_S$ contains $\mathbb R^V$ as a set, since it would
then contain every real in that set by transitivity.  Thus the surreal
atlas does not imply
$\HR=\bigcup_S W_S$.
\end{remark}

\begin{remark}[Which cofinality is being computed]
The cofinality in \cref{dsn:rp:cor:cohen-chain} concerns the \emph{inclusion
preorder of parameter fields}. It is not the order cofinality of the
surreal field: every ordinal already belongs to every $K_r$.
The strict chain and the absence of a largest parameter field imply,
by the Choice criterion, that $\HR$ fails Choice in this example.
Nevertheless $\Asm$ holds, so $\HR$ is closed under ambient countable
sequences and satisfies Dependent Choice. Thus countable coherence and
capture by one real are independent features of this construction.
\end{remark}

% ===== End sections/04_cohen.tex =====


% ===== Begin sections/05_prikry.tex =====

\subsection{Source 05's treatment of the Cohen model \src{05}}
\label{dsn:ca:sec:cohen}

Source~05 works in the same extension, written with $\kappa$ for
source~02's $\lambda$ and $J$ for source~02's $S$:
We now construct a nontrivial positive model. Let $W$ be a transitive
model of ZFC with $W\models V=L$. Let $\kappa$ be an uncountable cardinal
of $W$, let
\[
 \mathbb P=\operatorname{Add}(\omega,\kappa)^W
 =\operatorname{Fn}(\kappa\times\omega,2,{<}\omega)^W,
\]
and let $G$ be $W$-generic. Set $V=W[G]$. For a ground-model countable
$J\subseteq\kappa$, put
\[
 G_J=G\res(J\times\omega),\qquad W_J=W[G_J].
\]
We identify the generic filter with its union when restricting its
coordinate function. All $\OD(r)$ and $\ROD$ classes in this section
are computed in $V$.

\begin{lemma}[Countable support and ordinal covering]\label{dsn:ca:lem:ccc}
The Cohen product is ccc. Every real of $V$ belongs to some $W_J$ with
$J\in W$ countable in $W$. Every countable set of ordinals of $V$ is
contained in a countable set of ordinals belonging to $W$.
\end{lemma}
\begin{proof}
For ccc, apply the delta-system lemma to the finite domains of an
uncountable family of conditions, and thin further so all conditions
agree on their common finite root. Any two are compatible, so the
family is not an antichain.

For a name for a real, choose for each natural index a maximal antichain
deciding its bit. Each antichain is countable in $W$. The union of the
finite coordinate supports appearing in these countably many
antichains is a countable $J\in W$. The decisions, and hence the real,
are determined by $G_J$.

For a name for $f:\omega\to\Ord$, work below a condition forcing that
it is such a function and choose a maximal antichain deciding each
$f(n)$. The union $B$ of all ordinal values decided by these antichains
is a countable set in $W$, and $\operatorname{ran}(f)\subseteq B$ in
the generic extension. Enumerating an arbitrary countable set gives
the final assertion. This proof also explains the usual preservation
of $\omega_1$ by ccc forcing \cite{Jech}.
\end{proof}

\begin{lemma}[Homogeneity and ordinal subsets]\label{dsn:ca:lem:homogeneity}
Let $N$ be a transitive ground model and $P\in N$ a forcing such that
any two conditions have stronger conditions with isomorphic forcing
cones, with the induced name isomorphism fixing canonical names for
ground-model sets. If $X\subseteq\theta$ in an extension $N[H]$ is
definable there from finitely many ground-model parameters, then
$X\in N$.
\end{lemma}
\begin{proof}
For each $\xi<\theta$, write membership as a sentence
$\psi(\check\xi,\check{\bar a})$ using the given ground parameters.
If one condition forced this sentence and another forced its negation,
strengthen to isomorphic cones. Transport of names preserves the ground
parameters, giving opposite decisions at corresponding tops, a
contradiction. Hence the top condition decides each membership sentence.
The definability theorem for forcing gives the ground-model set
\[
 X=\{\xi<\theta:1_P\Vdash\psi(\check\xi,\check{\bar a})\}.
\]
This is the standard homogeneity argument, stated only for subsets of
a specified ground ordinal. It does not assert that no new definable
sets of arbitrary kind can appear.
\end{proof}

For a Cohen product, finite bit flips make any two conditions compatible
after an automorphism, and fix all ground check names. Thus the lemma
applies also to the remaining-coordinate product over an intermediate
model $W_J$.


\begin{dupstatement}{Source 05's Theorem 9.3 (= \cref{dsn:rp:thm:cohen-atlas}, with the ring equality)}
In the extension just described,
\begin{align}
 \KR^V
 &=\bigcup_{J\in[\kappa]^{\le\omega}\cap W}\No^{W_J},
 \label{dsn:ca:eq:cohenfield}\\
 \IR^V
 &=\Oz^V\cap
   \bigcup_{J\in[\kappa]^{\le\omega}\cap W}\No^{W_J}.
 \label{dsn:ca:eq:cohenring}
\end{align}
Here $[\kappa]^{\le\omega}\cap W$ means the sets countable \emph{in
$W$}. For each such $J$, choose a real $g_J$ coding $G_J$ using a
constructible enumeration of $J$. Then
\begin{equation}\label{dsn:ca:eq:fixedcohen}
 K_{g_J}^{V}=\No^{W_J}.
\end{equation}
The equalities are equalities of canonical sign-code subclasses, with
ambient surreal order and arithmetic. In particular their right-hand
sides are subfields under that arithmetic.
\end{dupstatement}
\begin{proof}[Second proof (source 05)]
Let $x\in\KR^V$ and choose a real $r$ with $x\in\OD^V(r)$.
By \cref{dsn:ca:lem:ccc}, $r\in W_J$ for some ground-model countable $J$.
Factor the forcing into its $J$ and $\kappa\setminus J$ coordinates.
Over $W_J$ the remaining factor is a weakly homogeneous Cohen product.
Let $\alpha=\bd(x)$ and let $A\subseteq\alpha$ be the positive-sign
positions of $x$. The set $A$ is definable in $V$ from $r$ and finitely
many ordinals. These parameters all belong to $W_J$. By
\cref{dsn:ca:lem:homogeneity}, $A\in W_J$, and hence its sign sequence $x$
belongs to $W_J$. This proves the forward inclusion.

Conversely, choose a constructible enumeration of the ground-model
countable set $J$. With this enumeration, a real $g_J$ codes the whole
array $G_J$, and
\[
 W_J=L[g_J]
\]
with the constructible hierarchy taken up to the common ordinals of
the models. The constructible enumeration and $J$ itself are already
in $L$. Every member of $L[g_J]$ is ordinal definable from $g_J$ in
$V$: specify a stage of the relative constructible hierarchy and a
position in its canonical $g_J$-definable well-order. The hierarchy
$L[g_J]$ has the same meaning in the intermediate and final models.
Consequently every canonical surreal in $W_J$ is in $K_{g_J}^{V}$,
proving the reverse inclusion and \eqref{dsn:ca:eq:fixedcohen}. Intersecting
with the ambient omnific predicate gives \eqref{dsn:ca:eq:cohenring}.
\end{proof}
The formulation of the ring equality avoids using additional
normal-form absoluteness as a hidden premise. The field property of
each displayed subclass follows directly from its equality with the
ambient definable field $K_{g_J}^{V}$.

\begin{theorem}[Properness, closure, and the first omitted birthday]\label{dsn:ca:thm:cohenconsequences}
In this Cohen extension the following all hold:
\begin{enumerate}[label=\textup{(\roman*)}]
\item $\KH^V=\No^W$ and
$\KH^V\subsetneq\KR^V\subsetneq\No^V$.
\item $\Asm$ holds. Thus $\KR^V=\Kco^V$ and both the field and its
omnific integer part are countably Hahn-sum closed.
\item No single real $r$ satisfies $K_r^V=\KR^V$.
\item $\lambda_{\R}^V=\omega_1^V$.
\item There is a non-$\ROD$ omnific integer in $(0,\omega)$.
\end{enumerate}
\end{theorem}
\begin{proof}
Homogeneity with ordinal parameters gives
$\KH^V\subseteq\No^W$, and every constructible set is ordinal
definable in $V$, giving equality. A coordinate Cohen real $c_i$ is
not in $W$ but is real--ordinal definable from itself. This makes the
first inclusion in (i) strict.

For (ii), the countable ordinal cover from \cref{dsn:ca:lem:ccc} lies in
$W=L$, and every constructible set is ordinal definable in $V$.
Thus \cref{dsn:ca:prop:cover} applies.

For properness, let $\eta=\omega_1^W=\omega_1^V$ and encode the generic
array on $\eta\times\omega$ as a subset $A$ of $\eta$, using the
ordinal bijection $(\xi,n)\mapsto\omega\cdot\xi+n$. The multiplication
in this sentence is ordinal multiplication. If $A$ were $\ROD$, the
support argument in \cref{dsn:rp:thm:cohen-atlas} would put it in some $W_J$ for
countable $J\in W$. Choose $i<\eta$ outside $J$. The coordinate real
$c_i$ is definable from $A$ and $i$, so it would be in $W_J$. But $c_i$
is Cohen generic over $W_J$ and is not in $W_J$, a contradiction.
Thus the sign sequence of length $\eta$ with positive set $A$ is a
surreal outside $\KR^V$, proving the second strict inclusion in (i).
Its birthday is exactly $\eta$, and \cref{dsn:ca:thm:countbirthday} supplies
the matching lower bound, proving (iv).

For (iii), given $r$, choose countable $J$ with $r\in W_J$.
The forward support argument gives $K_r^V\subseteq\No^{W_J}$.
Choose any $i\in\kappa\setminus J$; then $c_i\in\KR^V$ but
$c_i\notin K_r^V$. Finally apply the reversible bounded code $Q$ of
\cref{dsn:thm:bounded-code} to the missing sign sequence to obtain (v).
\end{proof}
Clauses (ii), (iii) and (iv) are also source~02's: they are
\cref{dsn:rp:thm:cohen-threshold}~(1), \cref{dsn:rp:cor:cohen-chain} with the
remark on which cofinality is computed, and
\cref{dsn:rp:thm:cohen-threshold}~(3); source~05's proofs above are second
routes, through the covering criterion \cref{dsn:ca:prop:cover} for (ii).
Clauses (i) and (v) are source~05's alone.

\begin{remark}[The countable-coordinate boundary]
If only countably many Cohen coordinates are added over $L$, the entire
generic array is coded by one real $g$, and $V=L[g]$. In that case every
set is $\OD(g)$, so $K_g=\KR=\No$. Uncountably many coordinates are
what make the positive example both countably closed and proper.
In the uncountable product, \cref{dsn:rp:thm:parameter-order} is sharp in the sense
that countable families of real parameters have a common real bound,
while the full family of coordinate reals does not.
\end{remark}

\section{A countable Hahn sum can defeat every real parameter}
\label{dsn:rp:sec:prikry}

We now construct a precise obstruction to arbitrary countable summation.
The large-cardinal assumption in this example is explicit: let $W$ be
a transitive ground model of ZFC with a measurable cardinal $\kappa$
and a normal nonprincipal $\kappa$-complete ultrafilter $D$ on $\kappa$.
Force over $W$ with ordinary Prikry forcing $\mathbb P_D$, obtaining
$V=W[G]$.  Its conditions are pairs $(s,A)$, where $s$ is a finite
increasing sequence of ordinals below $\kappa$, $A\in D$, and every
element of $A$ lies above $s$.  An extension appends finitely many
elements from the old upper set and shrinks that upper set.  A direct
extension, written $\leq^*$, leaves the stem $s$ fixed.

The forcing ingredients are classical results of Prikry
\cite{Prikry}; we include their proofs in the form needed here.
For a detailed author-written exposition see
\cite[Section~9.2]{KaragilaForcing}.

\begin{lemma}[The Prikry property and preservation of reals]
\label{dsn:rp:lem:prikry-property}
Every forcing statement is decided by a direct extension of any given
condition.  Consequently $\mathbb P_D$ adds no bounded subset of
$\kappa$, and $W$ and $V$ have the same reals.
\end{lemma}
\begin{proof}
First normality gives the following finite-set homogeneity fact: if
$c:[\kappa]^{<\omega}\to\mu$ with $\mu<\kappa$, some $H\in D$ makes
$c$ constant on $[H]^n$ separately for each $n<\omega$.
For completeness prove the fixed-length statement by induction on $n$.
The case $n=1$ follows from $\kappa$-completeness of the ultrafilter.
For the successor step, for each $a\in[\kappa]^n$ choose a color $i_a$
whose set $B_a$ of possible last coordinates above $a$ belongs to $D$.
For each $\eta<\kappa$ intersect the fewer than $\kappa$ sets $B_a$
with $a\subseteq\eta+1$, and take the diagonal intersection of these
intersections.  Normality makes the result $D$-large.  If $a$ lies
in it and $\beta$ is a later point of it, then $\beta\in B_a$.
The induction hypothesis applied to $a\mapsto i_a$ now gives
homogeneity for length $n+1$.  Intersecting the countably many
fixed-length homogeneous sets proves the stated fact.

Fix $p=(s,A)$ and a statement $\varphi$.  Color a finite tuple
$t\in[A]^{<\omega}$ by $1$ if some $B\in D$ makes
$(s^\frown t,B)$ a condition forcing $\varphi$, and by $0$ otherwise.
Choose $H\subseteq A$ in $D$ homogeneous at every length, and put
$q=(s,H)$.  If $q$ did not decide $\varphi$, it would have extensions
$q_0,q_1$ forcing opposite answers.  Extend the shorter stem until
their stems have equal length.  Their appended tuples $t_0,t_1$
belong to the same $[H]^n$, and hence have the same color.  The tuple
of the extension forcing $\varphi$ has color $1$.  The other tuple
therefore has an upper-set witness forcing $\varphi$ as well.
That witness and the extension forcing $\neg\varphi$ have the same
stem and compatible upper sets, a contradiction.  Thus $q$ decides
$\varphi$.

Suppose a name $\dot X$ is forced to be a subset of $\lambda<\kappa$.
For each $\xi<\lambda$ use the property just proved to find a direct
extension, with a common fixed stem, deciding $\xi\in\dot X$.
Intersect their upper sets.  By $\kappa$-completeness the result is
a condition deciding every bit, hence forcing $\dot X$ to equal a
ground-model set.  Such conditions are dense, so no bounded subset
of $\kappa$ is added.  In particular no subset of $\omega$ is added.
\end{proof}

\begin{lemma}[An explicit isomorphism between tail cones]
\label{dsn:rp:lem:prikry-cones}
Prikry forcing is cone homogeneous.  As a result, its largest
condition decides every sentence whose parameters are canonical
names for ground-model sets.
\end{lemma}
\begin{proof}
Given $p=(s,A)$ and $q=(t,B)$, remove from $A\cap B$ a bounded initial
segment containing both stems, leaving $D_0\in D$.  Set
$p'=(s,D_0)$ and $q'=(t,D_0)$.  The map
\[
 (s^\frown u,E)\longmapsto(t^\frown u,E)
\]
is an order isomorphism from $\mathbb P_D/p'$ onto
$\mathbb P_D/q'$.  Both cones allow exactly the same additional
finite tuples $u$ from $D_0$ and exactly the same remaining upper
sets $E$.  The inverse replaces $t$ by $s$.

An isomorphism of forcing notions preserves forcing and takes each
canonical name for a ground-model set to the same canonical name
over the other cone.  If two conditions forced opposite answers to
a sentence with only such parameters, strengthen them to the two
isomorphic cones above.  Invariance of forcing would give the
opposite answers at corresponding largest conditions, which is
impossible.  The density of deciding conditions now shows that the
largest condition decides the sentence.
\end{proof}

\begin{lemma}[Real-ordinal-definable ordinal sets remain in the ground]
\label{dsn:rp:lem:prikry-rod-ground}
Every ROD subset of an ordinal in $V=W[G]$ belongs to $W$.
\end{lemma}
\begin{proof}
Suppose $X\subseteq\eta$ is uniquely defined in $V$ by
$\varphi(X,r,\bar\alpha)$, where $r$ is a real and $\bar\alpha$
is a finite tuple of ordinals.  By
Lemma~\ref{dsn:rp:lem:prikry-property}, $r\in W$; all the ordinal
parameters belong to $W$ as well.  For each $\xi<\eta$, let
$\theta(\xi,r,\bar\alpha)$ assert that the uniquely
$\varphi$-defined object contains $\xi$.  By
Lemma~\ref{dsn:rp:lem:prikry-cones} the largest condition decides this
statement.  The forcing theorem and separation in $W$ produce
\[
 X_0=\{\xi<\eta:
   1_{\mathbb P_D}\Vdash\theta(\check\xi,\check r,
                                      \check{\bar\alpha})\}\in W.
\]
The truth lemma gives $X=X_0$.  This argument does not require
$\mathbb P_D$ itself to be ordinal definable: it is an ordinary
ground-model parameter in the construction of $X_0$.
\end{proof}

\begin{theorem}[Failure of countable Hahn closure in a Prikry extension]
\label{dsn:rp:thm:prikry-sum}
In $V=W[G]$ there is a countable summable sequence $(y_n)_{n<\omega}$
of ordinal-definable omnific integers such that
\[
 z=\omega+\sum_{n<\omega}y_n
\]
is not ROD.  Moreover $z$ is an omnific integer with zero constant
term and
\[
 \omega<z<\omega+\Omega(1/2).
\]
Thus all summands belong to the single field
$\No\cap\operatorname{HOD}^{V}$, while their Hahn sum does
not belong to $\KR^{V}$.
\end{theorem}
\begin{proof}
The union of the generic stems is an increasing cofinal sequence
$C=\{\kappa_n:n<\omega\}$ in $\kappa$.  The sets of conditions
requiring any specified finite stem length, and requiring a point
above any specified $\beta<\kappa$, are dense.  This proves both
countability and cofinality.  Since $\kappa$ is regular in $W$,
$C\notin W$.  Lemma~\ref{dsn:rp:lem:prikry-rod-ground} therefore gives
$C\notin\ROD^{V}$.

Put
\[
 e_{\kappa_n}=\Omega(-(\kappa_n+1)),\qquad y_n=\Omega(e_{\kappa_n}).
\]
Each $y_n$ is definable from the ordinal $\kappa_n$ and is a monomial
with positive exponent, so it is an ordinal-definable omnific integer.
The exponents $e_{\kappa_n}$ are strictly decreasing, making the family
Hahn-summable.  They all satisfy
$0<e_{\kappa_n}\leq\Omega(-1)<1/2$.  The sum $q=\sum_n y_n$ is therefore
purely infinite and positive, with leading exponent less than $1/2$.
Normal-form comparison gives $0<q<\Omega(1/2)$, so $z=\omega+q$
has the asserted location and constant term.

Its normal form recovers the generic set by the fixed formula
\[
 C=\{\alpha\in\operatorname{Ord}:
       \Omega(-(\alpha+1))\in\operatorname{supp}_{\rm exp}(z-\omega)\}.
\]
The support and this inverse decoding are uniquely definable set
operations.  If $z$ were ROD, then so would be $C$, a contradiction.
All the decoding takes place in $V$; no additional absoluteness
theorem about normal forms between $W$ and $V$ is needed.
\end{proof}

The theorem is a relative-consistency construction from a measurable
cardinal.  It does not assert failure of countable closure in every
universe.  In particular, $V=L$ makes every set ordinal definable.
Its role is to show exactly why a countable family of values can fit
inside one definable surreal field without its complete description,
or its Hahn sum, fitting inside any field obtained from real and
ordinal parameters.


\begin{corollary}[A missing countable cut]
In the same Prikry extension there are countable sets of
ordinal-definable surreal endpoints $A<B$ with no interpolant in
$\KR$. Thus the pointwise real-parameter amalgamation theorem does
not imply even external countable-cut interpolation.
\end{corollary}
\begin{proof}
Apply the finite-prefix cut construction in \cref{dsn:rp:thm:countable} to
the increasing enumeration of the new cofinal set $C$. Every
interpolant would recover that ordinal sequence, contradicting
\cref{dsn:rp:lem:prikry-rod-ground}.
\end{proof}

\subsection{A negative example without a large-cardinal hypothesis}
\label{dsn:rp:sec:namba}

The measurable cardinal in the Prikry construction is useful for its
cardinal preservation, but is unnecessary for the failure of countable
assembly.  A suitable Namba extension gives that failure from the
consistency of ZFC alone.  We specify the forcing presentation because
two versions commonly called Namba forcing are not forcing equivalent.

Let $W\models\mathrm{ZFC}+V=L$ and let
$\kappa=(\omega_2)^W$.  Use the version $\mathbb N$ whose conditions
are trees $T\subseteq\kappa^{<\omega}$ with a stem $s_T$, with a
single path below that stem, such that \emph{every} node at or above
the stem has $\kappa$ immediate successors in $T$.  The order is
inclusion.  This is the Laver-style presentation.  The classical
no-new-reals theorem under CH applies to this version:
\cite{LietzNamba2024} explicitly distinguishes the two presentations
and states that the preservation proof applies to both;
\cite[Theorem~23 and the comparison table on p.~4]{MillerForcing2009}
also records the result for this presentation.  We import this
classical preservation theorem and prove the additional properties
used here.

\begin{theorem}[Stabilization with failure of countable assembly]
\label{dsn:rp:thm:namba}
Let $G\subseteq\mathbb N$ be $W$-generic and $V=W[G]$.  Then
\[
 \HR^{V}=\operatorname{HOD}^{V}=L^{V}=W,
 \qquad
 \KR^{V}=\No^{W}.
\]
Every real-parameter field $K_r^V$ equals this same field.  Nevertheless
there is a countable set of ordinals that is not ROD in $V$.
Consequently countable Hahn closure, and the equivalent countable
assembly and ambient $\eta_1$ properties, fail.  In particular,
choice in $\HR$ and stabilization of the real-parameter
fields do not imply these countable-closure properties.
\end{theorem}
\begin{proof}
We first verify cone homogeneity for exactly the specified forcing.
Given two conditions $T,S$, consider only their nodes at or above
their respective stems.  Each is an abstract rooted tree of height
$\omega$ in which every node has exactly $\kappa$ children.  By
choosing a bijection between the successor sets at each pair of
corresponding nodes and recursing through the finite levels, construct
a rooted-tree isomorphism $F$ between them in $W$.

For a condition $R\subseteq T$, map its nodes above $s_T$ by $F$ and
restore the fixed proper prefixes of $s_S$.  The result is a condition
below $S$: its stem is the image of $s_R$, and the number of children
of every node above that stem remains $\kappa$.  This construction
preserves inclusion and has the analogous inverse induced by
$F^{-1}$.  Thus it is an isomorphism
\[
 \mathbb N/T\cong\mathbb N/S.
\]
In particular the forcing is cone homogeneous, with the canonical
names for all ground-model sets fixed under the induced name maps.

CH holds in $W$, so the cited Namba theorem gives
$\mathbb R^V=\mathbb R^W$.  The argument of
Lemma~\ref{dsn:rp:lem:prikry-rod-ground}, whose only inputs were cone
homogeneity and preservation of reals, now shows that every ROD
subset of an ordinal in $V$ belongs to $W$.

All reals of $V$ are constructible and hence ordinal definable in $V$.
It follows by substitution that
$\operatorname{ROD}^{V}=\operatorname{OD}^{V}$ and
$\HR^{V}=\operatorname{HOD}^{V}$.  To see that this class
is contained in $W$, take any $x\in\operatorname{HOD}^{V}$.
Choice inside that inner model supplies an ordinal code for the
membership graph of $\operatorname{tc}(\{x\})$.  The code is an
ordinal-definable set of ordinals, hence belongs to $W$ by the preceding
homogeneity argument.  Its well-founded collapse in $W$ recovers $x$.
Conversely every constructible set is ordinal definable in $V$ by its
position in the canonical constructible wellorder.  Thus
$\operatorname{HOD}^{V}=W=L^V$.  The canonical sign-sequence
identification gives the equality of surreal fields; because each
real is already ordinal definable, adding any one real parameter
does not enlarge the field.

Namba forcing adds a cofinal function $g:\omega\to\kappa$.  The
dense sets requiring longer stems and a stem value above a specified
ordinal prove this directly.  Its range $A$ is a countable unbounded
subset of $\kappa$ in $V$.  It cannot belong to $W$: a fixed subset
of an ordinal has the same order type in transitive universes, while
any countable ordinal remains below $\omega_1^W=\omega_1^V$; a
ground-model countable cofinal subset would contradict regularity
of $\kappa$ in $W$.  Equivalently, take the increasing cofinal
sequence of successive new maxima extracted from $g$; its order
type is $\omega$, making this contradiction immediate.
Therefore $A$ is not ROD, by the proved ground-model containment.

Apply the countable assembly criterion.  An explicit failed Hahn sum
is obtained, as before, from the monomials
\[
 \Omega\bigl(\Omega(-(\alpha+1))\bigr)\qquad(\alpha\in A).
\]
They are individually ordinal-definable omnific integers, and their
sum recovers $A$ from its exponent support.  Adding $\omega$ places
the missing omnific integer in the same bounded cell as in
Theorem~\ref{dsn:rp:thm:prikry-sum}.
\end{proof}

Since $V=L$ itself satisfies the countable assembly criterion and
the Namba extension just constructed violates it, that criterion is
independent of ZFC provided ZFC is consistent.  The same is therefore
true of arbitrary countable Hahn closure for the real-ordinal-definable
surreal field.  The Prikry construction supplies the complementary
illustration in which the relevant ground-model cardinal is preserved.


% ===== End sections/05_prikry.tex =====


% ===== Begin sections/06_strong_stabilizer.tex =====

\begin{writenote}[On the imported Namba input, 4 October 2026]
Source~02's only imported input in \cref{dsn:rp:thm:namba} is that this
presentation of Namba forcing adds no reals under CH. The write could not
confirm a textbook statement for this presentation; the references and what is
missing are recorded in \cref{dsn:rp:q:nambainput}. The cone homogeneity, the
computation $\HR^V=\HOD^V=L^V=W$ and the nondefinability argument are proved
above.
\end{writenote}

\subsection{Source 05's Prikry example \src{05}}
\label{dsn:ca:sec:prikry}

The positive Cohen model does not establish a theorem of ZFC that
$\Asm$ always holds. A forcing that adds a new countable ordinal
sequence while preserving all reals supplies the opposite behavior,
provided its homogeneity prevents the new sequence from becoming
ordinal definable from an old real.

Let $W\models\mathrm{ZFC}$ be transitive and let $U\in W$ be a normal,
nonprincipal, $\kappa$-complete ultrafilter on a measurable cardinal
$\kappa$. Ordinary Prikry forcing $P_U$ has conditions $(s,A)$, where
$s$ is a finite increasing sequence from $\kappa$, $A\in U$, and all
members of $A$ lie above the last member of $s$. An extension appends
finitely many members of $A$ to the stem and shrinks the remaining
tail. A \emph{direct extension} does not change the stem.
Let $C=\seq{\kappa_n:n<\omega}$ be the generic Prikry sequence and
$V=W[C]$.

\paragraph{The precise imported forcing input.}
The classical Prikry property says that every statement can be decided
by a direct extension. For a normal measure, ordinary Prikry forcing
adds an increasing cofinal $\omega$-sequence in $\kappa$ while adding
no bounded subsets of $\kappa$. A modern proof and discussion appear
in \cite{Benhamou}. We use this standard theorem, not a general claim
about all Prikry-type or tree-Prikry forcings.

For clarity, the no-new-bounded-subsets consequence follows directly
from the Prikry property and $\kappa$-completeness. Given a name for a
subset of $\gamma<\kappa$, decide its membership statements by direct
extensions with a fixed stem and intersect their $U$-tails. The resulting
tail still belongs to $U$ and gives a condition deciding the entire
subset to be a ground-model set. This argument works below every
condition. In particular
\begin{equation}\label{dsn:ca:eq:samereals}
 \R^{W[C]}=\R^W.
\end{equation}
The cofinality of $C$ follows from the dense requirements to extend the
stem beyond each fixed $\xi<\kappa$.
Source~02 proves this input instead (\cref{dsn:rp:lem:prikry-property}).

\begin{dupstatement}{Source 05's Lemma 10.1 (= \cref{dsn:rp:lem:prikry-cones,dsn:rp:lem:prikry-rod-ground}; source~05's ultrafilter is written $U$)}
The forcing $P_U$ satisfies the cone-isomorphism hypothesis of
\cref{dsn:ca:lem:homogeneity}. Consequently, every $\ROD$ subset of an ordinal
in $V=W[C]$ belongs to $W$.
\end{dupstatement}
\begin{proof}[Second proof (source 05)]
Given $p=(s,A)$ and $q=(t,B)$, choose $\gamma<\kappa$ at least as
large as every entry of either finite stem, and put
\[
 D=(A\cap B)\setminus(\gamma+1).
\]
Then $D\in U$, and $(s,D)$ and
$(t,D)$ strengthen the original conditions. Every condition in the
first cone has the form $(s\mathbin{{}^\frown}v,E)$ with the appended
finite sequence $v$ drawn from $D$, and the second cone has the
corresponding condition $(t\mathbin{{}^\frown}v,E)$. Replacing $s$ by
$t$ and leaving $v,E$ unchanged is an order isomorphism. When $v$ is
empty, every allowed tail is above both stems by the choice of $D$;
when it is nonempty, both stems have the same last appended point.
Thus the construction works even if the original stem lengths differ.

The recursively induced isomorphism of names takes each canonical
ground-model name to its counterpart. Apply
\cref{dsn:ca:lem:homogeneity}. A $\ROD$ definition in $V$ uses a real from
$W$ by \eqref{dsn:ca:eq:samereals} and finitely many ordinals, also from $W$.
All its parameters are therefore ground-model parameters, proving the
last assertion. No assumption that $U$ itself is ordinal definable
is required for this containment argument.
\end{proof}

Source~05's version of \cref{dsn:rp:thm:prikry-sum} uses its multiplexer and
proves that the Prikry sequence is not real-ordinal definable by coding its
graph, a second route to source~02's argument through order types.

\begin{theorem}[A non-$\ROD$ countable omnific sum]\label{dsn:ca:thm:prikry}
In the ordinary Prikry extension above, let $f(n)=\kappa_n$. Then
\begin{equation}\label{dsn:ca:eq:prikrycode}
 z_C=\sum_{n<\omega}
 \omega^{\displaystyle 2^{-n-1}
       +2^{-n-2}\frac{\kappa_n}{1+\kappa_n}}
 \quad\in\quad\Oz^V\cap(0,\omega)
\end{equation}
is not $\ROD$ in $V$. Every individual summand is $\OD$ in $V$,
the family is Hahn-summable, and every coefficient in its normal form
is $1$. Therefore
\[
 \neg\Asm,\qquad \KR^V\subsetneq\Kco^V,
 \qquad \IR^V\subsetneq\Ico^V.
\]
In particular $\KR^V$ and $\IR^V$ need not be closed under countable
Hahn sums, even of elements of the common ordinal-definable subfield
and subring.
\end{theorem}
\begin{proof}
The function $f$ cannot belong to $W$: if it did, its domain $\omega$
and cofinality in $\kappa$ would be absolute, contradicting regularity
of $\kappa$ in $W$. Code its graph by
\[
 X_f=\{\kappa\cdot n+f(n):n<\omega\}
       \subseteq\kappa\cdot\omega,
\]
using ordinal multiplication and addition. The coding and decoding are
definable from the ordinal $\kappa$. If $f$ were $\ROD$, so would be
$X_f$. By \cref{dsn:rp:lem:prikry-rod-ground}, $X_f\in W$, and decoding inside $W$
would give $f\in W$, a contradiction. Thus $f\notin\ROD$. This
argument uses the graph, not an assumption that every ground-model set
which is countable in the extension was countable in the ground.

Apply \cref{dsn:ca:thm:code} to this $f$. It proves every assertion about the
normal form, summability, bounds, and individual definability in
\eqref{dsn:ca:eq:prikrycode}. If $z_C$ were $\ROD$, its fixed decoder would
recover the whole $f$ as a $\ROD$ set, contrary to the preceding
paragraph. On the other hand $f$ is itself an allowed countable ordinal
parameter, so $z_C\in\Ico$. The strict inclusions and failure of
$\Asm$ follow.
\end{proof}

\begin{remark}[What the negative result really separates]
The failure is not caused by a new undefinable real: there are no new
reals at all. It is not caused by choosing different real parameters
for different summands: every summand is already $\OD$. Nor is it
caused by unsummability or a noncanonical limit convention. What is
missing is a real--ordinal definition of the sequence
$\seq{\kappa_n:n<\omega}$. A countable sequence of very large ordinals
need not be encodable by one real and finitely many ordinal parameters.
\end{remark}

\paragraph{The exact relative-consistency conclusion.}
The Cohen construction is available over a constructible transitive
model of ZFC and gives a proper $\KR$ satisfying countable Hahn-sum
closure. The Prikry construction is available over a transitive model
with a measurable cardinal and gives failure of that closure. Thus,
under the stated model-existence hypotheses, the two behaviors occur
in models of ZFC. In the usual forcing-relative-consistency reading,
ZFC supplies the positive direction, while a measurable cardinal is
an upper bound for the consistency strength of the negative direction.

No lower bound matching measurability is proved. In particular we do
not claim equiconsistency with a measurable cardinal, and we do not
claim that a measurable cardinal is necessary for failure of $\Asm$.
Ordinary Prikry forcing was chosen because both no-new-reals and the
needed cone isomorphisms have transparent proofs. Replacing it by a
weaker-consistency construction is a separate research problem.

\begin{writenote}[Superseded, 4 October 2026, batch~93]
The last two paragraphs are source~05's, written without knowledge of
source~02. \Cref{dsn:rp:thm:namba} is exactly the weaker-consistency
construction they call a separate research problem: from a transitive model
of $\ZFC+V=L$, which exists if any transitive model of ZFC does, Namba
forcing gives failure of $\Asm$ with no new reals and with the needed cone
homogeneity proved. In the forcing relative-consistency reading, failure of
$\Asm$ is therefore consistent relative to ZFC alone, and since it is a
statement of ZFC its consistency strength is exactly that of ZFC. The
measurable cardinal remains what source~05 says it is: a sufficient
hypothesis for this particular, cardinal-preserving construction. Nothing in
source~05's paragraphs is wrong; they claim no lower bound. See also
\cref{dsn:ca:q:strength}.
\end{writenote}

\section{The strong stabilizer of the reals and all ordinals}
\label{dsn:rp:sec:strong}

Naming all real and ordinal parameters is still much weaker than allowing
arbitrary set-theoretic definitions. This section computes an exact
invariant that makes the difference visible. Arithmetic on surreal
exponents is surreal arithmetic; addition of ordinal surreals is their
natural sum. We work in ZFC, using canonical sign sequences. Statements
about class automorphisms can equivalently be read in a conservative class
theory with set choice.

For a subclass \(B\subseteq\No\), put
\[
 \operatorname{Ser}(B)=\{x\in\No:\supp(x)\subseteq B\},
\]
where \(\supp(x)\) is the set of exponents in the Conway normal form of
\(x\). The empty support represents zero. We use the classical facts that
each support is a set, reverse well ordered by the surreal order, and that
an ordinal's Conway normal form is its finite Cantor normal form
\cite{Gonshor}. Define
\begin{equation}\label{dsn:rp:eq:fields}
 \Gamma_{\Z}=\operatorname{span}_{\mathbb Z}(\Ord),\qquad
 \Gamma_{\Q}=\operatorname{span}_{\Q}(\Ord),\qquad
 \mathcal F_\Gamma=\operatorname{Ser}(\Gamma_{\Q}).
\end{equation}
The spans use finite sums, so no class sum occurs.

A field automorphism is called \emph{strong} here if it and its inverse
transport every set-indexed summable family to a summable family and
commute with its Conway sum. Every Hahn automorphism constructed below has
this property. The exact theorem concerns this explicitly specified
class; it does not assume that every field automorphism is strong.

\subsection{The rational span of the ordinals}

\begin{lemma}\label{dsn:rp:lem:rational-span}
The elements of \(\Gamma_{\Q}\) are precisely the finite normal forms
\[
 g=\sum_{j=1}^{m}q_j\omega^{\gamma_j},\qquad
 q_j\in\Q,\quad \gamma_1>\cdots>\gamma_m,\quad\gamma_j\in\Ord.
\]
For \(\Gamma_{\Z}\), replace rational coefficients by integer coefficients. Hence
\[
 \Gamma_{\Q}\cap\R=\Q,\qquad
 \Gamma_{\Q}\cap\{\text{infinitesimal surreals}\}=\{0\}.
\]
Moreover, \(\operatorname{Ser}(\Ord)\) consists of finite normal forms with
real coefficients and ordinal exponents.
\end{lemma}

\begin{proof}
Expand each ordinal in a finite rational linear combination into Cantor
normal form and collect equal exponents. Conversely, \(\omega^\gamma\)
is an ordinal when \(\gamma\) is an ordinal, so every displayed expression
belongs to \(\Gamma_{\Q}\). The integer assertion is identical.

A nonzero such normal form with leading exponent zero is a nonzero
rational. If its leading exponent is positive, it is infinite in absolute
value. These are the only possibilities, proving the intersections.
Finally, a reverse well-ordered subset of the ordinals is finite: an
infinite set of ordinals has a strictly increasing sequence of its first
\(\omega\) elements. This proves the last assertion.
\end{proof}

\subsection{Two automorphism tools}

The Hahn-lift construction and invariant-support method below already
occur in the repository's automorphism and self-embedding reports
\cite{RepoOpa,RepoSse}. We give the details needed
for the new fixed-field calculation. In particular, the set-parameter
span-of-supports result \texttt{opa:thm:hull} and elementary shears in
\cite{RepoOpa} are antecedents. The issue settled here is the stabilizer
of the proper class of all canonical ordinals, which remains explicitly
open in \cite{RepoSse} at the pinned snapshot.

\begin{lemma}[Hahn lifts]\label{dsn:rp:lem:hahn-lift}
Let \(\psi:(\No,+,<)\to(\No,+,<)\) be an additive order automorphism.
Then
\begin{equation}\label{dsn:rp:eq:hahn-lift}
 S_\psi\left(\sum_{g\in S}r_g\omega^g\right)
       =\sum_{g\in S}r_g\omega^{\psi(g)}
\end{equation}
is a strong ordered-field automorphism of \(\No\), fixes \(\R\)
pointwise, and preserves \(\Oz\) in both directions. If \(\psi\)
fixes every ordinal, then so does \(S_\psi\).
\end{lemma}

\begin{proof}
An increasing bijection transports a reverse well-ordered set to a
reverse well-ordered set. Reindexing therefore defines the map.
It preserves addition, the largest exponent where normal forms differ,
and the coefficient there; hence it preserves order. Inverses reindex
by the inverse exponent maps.

For \(S_\psi\), additivity of \(\psi\) gives
\(\omega^{\psi(g+h)}=\omega^{\psi(g)}\omega^{\psi(h)}\), so normal-form
convolution for multiplication is preserved. Reindexing preserves the
reverse-well-order and local finiteness conditions for summability in
both directions; coefficientwise addition proves strongness.

An additive order automorphism fixes zero and preserves positive
exponents. Thus \(S_\psi\) fixes real constant series and preserves
the normal-form conditions for \(\Oz\): all nonconstant exponents
positive, constant term an integer. Finally, an ordinal has finite
normal form with ordinal exponents and natural coefficients, so an
exponent map fixing ordinals fixes this normal form.
\end{proof}

\begin{lemma}[A moving support point prevents a fixed series]
\label{dsn:rp:lem:orbit}
If an increasing bijection \(h\) of a linearly ordered class satisfies
\(h(S)=S\) for a reverse well-ordered set \(S\), it fixes every member
of \(S\). Consequently
\[
 S_\psi(x)=x\ \Longrightarrow\
 \supp(x)\subseteq\operatorname{Fix}(\psi).
\]
\end{lemma}

\begin{proof}
If \(h(s)>s\), then \(s,h(s),h^2(s),\ldots\) is a strictly increasing
sequence in \(S\). If \(h(s)<s\), use inverse iterates. Both contradict
reverse well ordering. Equality of reindexed normal forms implies
invariance of their supports, proving the consequences.
\end{proof}

\subsection{Moving every exponent outside the rational ordinal span}

\begin{theorem}[The additive stabilizer]\label{dsn:rp:thm:additive-stabilizer}
The common fixed subgroup of all additive order automorphisms of \(\No\)
that fix \(\Ord\) pointwise is exactly \(\Gamma_{\Q}\). For every \(e\notin \Gamma_{\Q}\)
there is such an automorphism \(\psi\), definable from set parameters,
that moves \(e\).
\end{theorem}

\begin{proof}
An additive automorphism of a divisible torsion-free group is
\(\Q\)-linear, so every automorphism in question fixes \(\Gamma_{\Q}\).

Choose once and for all a \(\Q\)-linear map \(d:\R\to\R\) whose
kernel is exactly \(\Q\). Such a map is obtained by extending
\(\{1\}\) to a Hamel basis, writing \(\R=\Q\oplus W\), and taking
the projection onto \(W\). This is the explicit use of set choice;
the single map \(d\) is a set parameter.

For each \(s\in\No\), let \(c_s(x)\) be the coefficient at exponent
\(s\), zero if absent, and define
\[
 d_s=\begin{cases}d,&s\in\Ord,\\
                   \operatorname{id}_{\R},&s\notin\Ord.
      \end{cases}
\]
Put
\begin{equation}\label{dsn:rp:eq:shear}
 N_s(x)=d_s(c_s(x))\omega^{s-1},\qquad
 \psi_s(x)=x+N_s(x).
\end{equation}
Coefficient extraction is additive. Since \(s-1<s\), the
coefficient at \(s\) in \(N_s(x)\) vanishes. Thus \(N_s^2=0\) and
\(\psi_s^{-1}=1-N_s\).

If \(N_s(x)\ne0\), then \(x\) has leading exponent at least
\(s\), whereas the added term has exponent \(s-1<s\).
The leading sign is unchanged. Thus \(\psi_s\) preserves signs
and, by additivity, order.
An ordinal has coefficient zero at a nonordinal exponent and a
natural-number coefficient at an ordinal exponent. Both are killed
by the corresponding \(d_s\), so every \(\psi_s\) fixes all ordinals.

Finally, an element fixed by every \(\psi_s\) has no nonordinal
support exponent and has rational coefficients at every ordinal
support exponent. Its support is finite by
Lemma~\ref{dsn:rp:lem:rational-span}, so it belongs to \(\Gamma_{\Q}\).
Thus some \(\psi_s\) moves each \(e\notin \Gamma_{\Q}\).
The entire family is given by one class formula with the one set
parameter \(d\); no basis of the proper-class group \(\No\) is used.
\end{proof}

\begin{theorem}[Exact strong fixed field]\label{dsn:rp:thm:strong-fixed}
The common fixed field of all strong field automorphisms of \(\No\)
fixing \(\R\cup\Ord\) pointwise is
\begin{equation}\label{dsn:rp:eq:strong-fixed}
 \operatorname{Fix}_{\mathrm{strong}}(\R\cup\Ord)
    =\mathcal F_\Gamma
    =\{x\in\No:\supp(x)\subseteq\operatorname{span}_{\Q}(\Ord)\}.
\end{equation}
The same equality holds if these automorphisms are additionally
required to preserve \(\Oz\).
\end{theorem}

\begin{proof}
Let \(\sigma\) be any strong field automorphism fixing the reals and
ordinals. A field automorphism of a real closed field preserves its
order, since positive elements are the nonzero squares. For
\(g=\sum_{j=1}^{m}q_j\alpha_j\in \Gamma_{\Q}\), one has
\[
 \omega^g=\prod_{j=1}^{m}(\omega^{\alpha_j})^{q_j}.
\]
Each \(\omega^{\alpha_j}\) is an ordinal, hence fixed. Positive
rational powers are unique in an ordered field, so \(\omega^g\)
is fixed. Real coefficients are fixed as well. Strongness then fixes
every normal-form sum supported on \(\Gamma_{\Q}\). This proves the lower
inclusion even without an assumption about \(\Oz\).

If \(x\notin\mathcal F_\Gamma\), choose \(e\in\supp(x)\setminus \Gamma_{\Q}\).
The additive stabilizer theorem gives \(\psi\) fixing every ordinal
but moving \(e\). Its Hahn lift \(S_\psi\) is strong, fixes
\(\R\cup\Ord\), and preserves \(\Oz\). If it fixed \(x\), the
support-orbit lemma would force \(\psi(e)=e\), a contradiction.
Thus the upper inclusion already follows from the explicitly
constructed \(\Oz\)-preserving automorphisms.
\end{proof}

This exactly answers the \emph{strong} version of the
\(\R\cup\Ord\)-stabilizer question recorded as
\texttt{sse:sr:q:stabilizers} in the pinned self-embedding report
\cite{RepoSse}. The lower inclusion uses preservation of
arbitrary summable families; the theorem does not compute the common
fixed field without that hypothesis.

No collection of classes is formed as a set. The lower inclusion is
a theorem about each class automorphism with the stated properties.
For the upper inclusion, each \(x\notin\mathcal F_\Gamma\) has an explicit
automorphism defined from set parameters. These statements suffice
for applications to each fixed first-order formula.

\subsection{Consequences for intrinsic definability}

Let \(\mathfrak S\) be the surreal ordered field with predicate \(\Oz\).
We may also add predicates for \(\R\) and \(\Ord\); the next assertion
holds for either language. A permitted parameter class means that each
formula may use a \emph{finite tuple} from that class.

\begin{corollary}\label{dsn:rp:cor:dcl-bound}
For every \(A\subseteq\R\cup\Ord\),
\[
 \operatorname{dcl}_{\mathfrak S}(A)\subseteq\mathcal F_\Gamma.
\]
This includes permitting all real and ordinal parameters simultaneously.
\end{corollary}

\begin{proof}
The automorphisms in the upper-inclusion proof fix \(A\), preserve
\(\Oz\), and fix the real and ordinal subclasses pointwise. They are
automorphisms of either expansion. A uniquely defined element is
fixed by every such automorphism. For a proper-class structure,
invariance is proved by induction on the subformulas of the particular
formula; no uniform truth predicate is needed.
\end{proof}
\begin{writenote}[Antecedent, 4 October 2026]
For a finite tuple of real and ordinal parameters, and in the language with
$\Oz$ and a constant for every real but without predicates for $\Ord$, this
corollary follows already, and more sharply, from
\lbl{opa:sc:cor:innerbound} of
\repopath{surreal/omnific-preserving-automorphisms}: an element definable over
a set $A$ there is supported in the rational span of the exponents of the
members of $A$, and for $A\subseteq\R\cup\Ord$ that span is a
finite-dimensional rational subspace of $\Gamma_\Q$. Source~02 does not cite it. The versions
with a predicate for $\Ord$, which elementary shears need not preserve, are
new here.
\end{writenote}


\begin{proposition}\label{dsn:rp:prop:fixed-field-arithmetic}
The class \(\mathcal F_\Gamma\) is a real closed subfield of \(\No\).
Its intersection \(I_\Gamma=\mathcal F_\Gamma\cap\Oz\) is an integer part, and
\(\mathcal F_\Gamma=\operatorname{Frac}(I_\Gamma)\).
\end{proposition}

\begin{proof}
It is a subfield by its fixed-field characterization. Every
automorphism fixing the coefficients of a polynomial preserves its
finite ordered list of real roots, hence each individual root.
Positive square roots and roots of odd-degree polynomials exist in
\(\No\), and therefore belong to the fixed field. This gives real
closedness.

The omnific floor is preserved by every order automorphism preserving
\(\Oz\), since it is the unique \(z\in\Oz\) with \(z\le x<z+1\).
The strong \(\Oz\)-preserving fixed-field characterization gives the
integer-part assertion.

For \(0\ne x\in\mathcal F_\Gamma\), choose an ordinal \(\delta\) with
\(\delta+g>0\) for every \(g\in\supp(x)\). Such an ordinal exists
because the support is a set and ordinals are cofinal in \(\No\).
Then \(b=\omega^\delta\) is an ordinal, and \(a=bx\) has only positive
exponents. Thus \(a,b\in I_\Gamma\) and \(x=a/b\).
\end{proof}

\subsection{The smaller field obtained from pure-field definitions}

Write \(F_{\mathrm{alg}}=\operatorname{rcl}_{\No}(\R(\Ord))\).
Quantifier elimination for real closed fields gives the classical
identity
\[
 \operatorname{dcl}_{(\No;+,\cdot,<)}(\R\cup\Ord)=F_{\mathrm{alg}},
\]
already recorded in the definability report (Part~\ref{dsn:part:one}).
It uses finite parameters even when the permitted parameter class
is proper. The next calculation makes its growth restrictions exact.

For \(x\ne0\), write \(\operatorname{lead}(x)=\max\supp(x)\) for its leading exponent.
One has \(\operatorname{lead}(xy)=\operatorname{lead}(x)+\operatorname{lead}(y)\), and a finite sum with a unique
largest leading exponent is nonzero.

\begin{proposition}[Exact leading-exponent group]\label{dsn:rp:prop:value-group}
One has
\[
 \operatorname{lead}(\R(\Ord)^\times)=\Gamma_{\Z},\qquad
 \operatorname{lead}(F_{\mathrm{alg}}^\times)=\Gamma_{\Q}.
\]
Every positive infinite \(x\in F_{\mathrm{alg}}\) therefore satisfies
\(x>\omega^{1/n}\) for some positive integer \(n\).
\end{proposition}

\begin{proof}
The ring \(\R[\Ord]\) consists of finite normal forms with real
coefficients and ordinal exponents. Finite products add ordinal
exponents by natural addition, which still gives ordinals;
conversely each \(\omega^\alpha\), \(\alpha\in\Ord\), is a generator.
A nonzero quotient in \(\R(\Ord)\) consequently has leading exponent
\(\alpha-\beta\in \Gamma_{\Z}\). Every member of \(\Gamma_{\Z}\) is a difference of two
ordinals, after collecting positive and negative summands, and
\(\omega^{\alpha-\beta}=\omega^\alpha/\omega^\beta\in \R(\Ord)\).

For \(0\ne x\in F_{\mathrm{alg}}\), choose a polynomial relation
\(\sum_i a_i x^i=0\) over \(\R(\Ord)\), with nonzero displayed coefficients.
The largest of \(\operatorname{lead}(a_i)+i\operatorname{lead}(x)\) must occur at least twice.
For suitable \(i\ne j\),
\[
 \operatorname{lead}(x)=\frac{\operatorname{lead}(a_j)-\operatorname{lead}(a_i)}{i-j}\in \Gamma_{\Q}.
\]
Conversely, for \(g\in \Gamma_{\Q}\), choose a positive integer \(n\) with
\(ng\in \Gamma_{\Z}\). The positive monomial \(\omega^g\) is an \(n\)-th root
of \(\omega^{ng}\in \R(\Ord)\), so belongs to \(F_{\mathrm{alg}}\).

For positive infinite \(x\), one has \(\operatorname{lead}(x)>0\).
By Lemma~\ref{dsn:rp:lem:rational-span} this exponent is either a positive
rational or positive infinite. Choose \(n\) with \(1/n<\operatorname{lead}(x)\).
The quotient \(x/\omega^{1/n}\) has positive leading exponent and
positive leading coefficient, so is infinite.
\end{proof}

\begin{proposition}[A strict distinction]\label{dsn:rp:prop:strict}
The inclusion \(F_{\mathrm{alg}}\subseteq\mathcal F_\Gamma\) is strict.
An explicit witness is
\begin{equation}\label{dsn:rp:eq:formal-exp}
 \mathcal E=\sum_{n=0}^{\infty}\frac{\omega^{-n}}{n!},
\end{equation}
which lies in \(\mathcal F_\Gamma\) and is transcendental over \(\R(\Ord)\).
\end{proposition}

\begin{proof}
The inclusion follows because \(\mathcal F_\Gamma\) is real closed and
contains \(\R(\Ord)\). The displayed series is summable and supported on
negative integers, so belongs to \(\mathcal F_\Gamma\).

If it were algebraic over \(\R(\Ord)\), clear denominators of a nonzero
polynomial relation to obtain coefficients in \(\R[\Ord]\). Their
combined support is a finite subset of \(\Gamma_{\Z}\). Partition it into
cosets of \(\mathbb Z\) in \(\Gamma_{\Z}\), with representatives \(h_c\), to write
\[
 P(Y)=\sum_c\omega^{h_c}P_c(\omega,Y),\qquad
 P_c(T,Y)\in\R[T,T^{-1},Y],
\]
where at least one \(P_c\ne0\). After substituting \(\mathcal E\), each term
has support in \(h_c+\mathbb Z\). Disjointness of the cosets forces
\(P_c(\omega,\mathcal E)=0\) for every \(c\).
A nonzero Laurent-polynomial coefficient stays nonzero upon
substituting \(\omega\), by uniqueness of finite normal forms.
Thus one nonzero \(P_c\) gives an algebraic relation over \(\R(\omega)\).

Put \(t=\omega^{-1}\). The integer-exponent subfield is the usual
Laurent-series field \(\R((t))\), and \(\mathcal E\) becomes
\(f(t)=\sum_{n\ge0}t^n/n!\), satisfying \(f'=f\).
For completeness, \(f\) is transcendental over \(\R(t)\).
If it had monic minimal polynomial
\[
 Q(t,Y)=Y^d+\sum_{i<d}a_i(t)Y^i,\qquad a_i(t)\in\R(t),
\]
differentiation would give
\((\partial_tQ+Y\partial_YQ)(t,f(t))=0\).
Minimality makes \(Q\) divide \(\partial_tQ+Y\partial_YQ\).
The latter has degree \(d\) and leading coefficient \(d\), so equals
\(dQ\). Hence \(a_i'=(d-i)a_i\) for \(i<d\).
A nonzero rational function \(a=p/q\) cannot satisfy \(a'=ca\) for a
nonzero constant \(c\): the equation
\(p'q-pq'=cpq\) has left-hand degree at most
\(\deg p+\deg q-1\) and right-hand degree \(\deg p+\deg q\).
Thus every \(a_i=0\), giving \(Q=Y^d\), impossible at \(f\ne0\).
\end{proof}

This does not assert that \(\mathcal E\) is undefinable after adding the
omnific predicate. It separates pure-field definable closure from
the strong fixed field, which are distinct invariants.

\subsection{An independent family in one omnific interval}

\begin{theorem}[Bounded ordinal-definable independence]
\label{dsn:rp:thm:independence}
For \(0<\alpha\in\Ord\), put
\[
 m_\alpha=\omega^{-\alpha},\qquad
 y_\alpha=\omega^{m_\alpha},\qquad z_\alpha=\omega+y_\alpha.
\]
The families \((y_\alpha)\) and \((z_\alpha)\) are algebraically
independent over \(\mathcal F_\Gamma\).
Every \(z_\alpha\) is an ordinal-definable omnific integer with
\[
 \omega<z_\alpha<\omega+\omega^{1/2},\qquad
 \operatorname{ct}(z_\alpha)=0.
\]
None is intrinsically definable in \(\mathfrak S\) from real and
ordinal parameters, even if every such parameter is permitted.
These assertions hold in every universe of ZFC, including
\(V=\mathrm{HOD}\).
\end{theorem}

\begin{proof}
Every \(m_\alpha\) is a positive infinitesimal. For distinct
\(\alpha_1,\ldots,\alpha_m\) and integers \(k_j\), not all zero,
\[
 \sum_{j=1}^{m}k_jm_{\alpha_j}\ne0
\]
by uniqueness of normal form. Its exponents are negative, so it
is infinitesimal and lies outside \(\Gamma_{\Q}\).

For a nonzero polynomial in finitely many variables, combine equal
monomials to write
\[
 P(Y_1,\ldots,Y_m)=\sum_{\nu\in F}a_\nu Y^\nu,\qquad
 0\ne a_\nu\in\mathcal F_\Gamma.
\]
After substituting \(Y_j=y_{\alpha_j}\), the term at \(\nu\) has
support in
\[
 \Gamma_{\Q}+\sum_{j=1}^{m}\nu_jm_{\alpha_j}.
\]
Distinct multi-indices give disjoint cosets. No coefficient of one
nonzero term can cancel with another term, so the sum is nonzero.
This applies the repository's disjoint-support method
(Part~\ref{dsn:part:one}) to the newly identified field \(\mathcal F_\Gamma\).
Since \(\omega\in\mathcal F_\Gamma\), translating each variable by
\(\omega\) is an invertible polynomial change of variables, so
the \(z_\alpha\) are independent as well.

The inequalities \(0<m_\alpha<1/2\) give
\(1<y_\alpha<\omega^{1/2}\). The normal form of \(z_\alpha\)
has exactly the two positive exponents \(1,m_\alpha\), coefficients
one, and zero constant term. This proves the omnific and interval
assertions.

Canonical arithmetic and \(\Omega\) are parameter-free definable
in set theory, so \(z_\alpha\) is ambiently definable from \(\alpha\),
and \(z_1\) is ambiently parameter-free definable.
But \(m_\alpha\notin \Gamma_{\Q}\), so neither \(y_\alpha\) nor \(z_\alpha\)
lies in \(\mathcal F_\Gamma\). The intrinsic nondefinability follows from
Corollary~\ref{dsn:rp:cor:dcl-bound}.

Only finite subfamilies and set-sized normal forms were used.
Ordinal-indexed independence means that each finite subfamily is
independent, not that a proper-class polynomial or sum is formed.
\end{proof}

The first omitted monomial has a familiar analytic name:
\[
 y_1=\omega^{1/\omega}=\log\omega
\]
for the canonical Gonshor logarithm. This identity is classical:
van den Dries--Ehrlich, proof of Lemma~5.3, p.~186
\cite{vdDE}. Thus \(\log\omega\) is an ambiently
parameter-free definable omnific integer, but is not definable in
the field with omnific, real, and ordinal predicates from any finite
tuple of real and ordinal parameters.

Neither the canonical omega map nor the canonical exponential is
definable in this reduct from such parameters. A definition of the
omega map would define its value at \(1/\omega\); a definition of
the exponential would define the unique solution of
\(\exp(x)=\omega\). Both would define \(y_1\), a contradiction.
These statements concern the specified reduct, not structures where
either function has already been named.
\begin{writenote}[Antecedent, 4 October 2026]
Without the predicate for $\Ord$, these nondefinability statements also follow
from the omnific-preserving automorphisms report: $\omega^{1/\omega}$ has the
exponent $\omega^{-1}$ outside the rational span of the exponents of any tuple
from $\R\cup\Ord$, so \lbl{opa:sc:cor:innerbound} excludes it; and
\lbl{opa:thm:parameters} and \lbl{opa:sc:thm:undef}~(ii),~(iv) show that the
graph of $x\mapsto\omega^x$ and the Gonshor exponential are not definable over
any set of parameters in the field with $\Oz$ and all real constants.
Source~02 does not cite them; with the predicate for $\Ord$ the statements are
new.
\end{writenote}


\subsection{Naming the omega map}

In the ordered-field language expanded by \(\Omega\),
\(\omega=\Omega(1)\) is a closed term, and
\[
 y_1=\Omega\bigl(1/\Omega(1)\bigr)
\]
is intrinsically parameter-free definable. Every \(z_\alpha\) is
definable from \(\alpha\).

\begin{proposition}[Intrinsic bounded coding]\label{dsn:rp:prop:omega-code}
In the ordered-field language with \(\Omega\), the function
\[
 q(x)=\Omega\left(
    \frac12\left(1+\frac{x}{\sqrt{1+x^2}}\right)\right)
\]
is a parameter-free definable bijection from \(\No\) onto
\[
 \Qrange=\{\Omega(t):0<t<1\}\subseteq\Oz\cap(1,\omega),
\]
with parameter-free definable inverse. For every parameter family
\(A\) and every \(x\),
\[
 \operatorname{dcl}_{\Omega}(A\cup\{x\})
   =\operatorname{dcl}_{\Omega}(A\cup\{q(x)\}).
\]
\end{proposition}

\begin{proof}
In every real closed field,
\[
 t=\frac12\left(1+\frac{x}{\sqrt{1+x^2}}\right)
\]
is a bijection onto \((0,1)\), with inverse
\[
 x=\frac{2t-1}{2\sqrt{t(1-t)}}.
\]
Positive square roots are uniquely definable, and direct substitution
verifies the inverse identities. The omega map is increasing and
injective, so \(q\) has the stated inverse: recover the unique \(t\)
with \(\Omega(t)=q(x)\), then apply the algebraic formula.
For \(0<t<1\), the monomial \(\omega^t\) is an omnific integer
between \(1\) and \(\omega\). Mutual definability gives the closure
identity.
\end{proof}

The monomial code already occurs as an ambient code in the repository
(Part~\ref{dsn:part:one}). This proposition answers a specific part of
its question about intermediate languages: the code and its inverse
become intrinsic when the omega map is named. It does not classify
the entire definable closure in that expansion.
\begin{writenote}[4 October 2026]
With $\Oz$ named as well, Part~\ref{dsn:part:op} goes much further in this
language: every coefficient and truncation is definable
(\cref{dsn:op:thm:support}), a bounded lexicographic pairing is definable
(\cref{dsn:op:lem:pair}), and $(V,\in)$ has a parameter-free quotient
interpretation with representatives in $\Qrange$
(\cref{dsn:op:thm:V-interpretation}). The definable closure is still not
classified there.
\end{writenote}


\subsection{Remaining stabilizer and language questions}

Let \(F_{\mathrm{all}}\) be the common fixed field of all
\(\Oz\)-preserving field automorphisms fixing \(\R\cup\Ord\), with no
strongness assumption, and put
\(D_{\Oz}=\operatorname{dcl}_{(\No;+,\cdot,<,\Oz)}(\R\cup\Ord)\).
The proven bounds are
\begin{equation}\label{dsn:rp:eq:open-bounds}
 F_{\mathrm{alg}}\subseteq D_{\Oz}
 \subseteq F_{\mathrm{all}}\subseteq\mathcal F_\Gamma,\qquad
 F_{\mathrm{alg}}\subsetneq\mathcal F_\Gamma.
\end{equation}
They follow respectively from pure-field definitions, automorphism
invariance, and the inclusion of the explicit strong automorphisms
among all automorphisms considered.

The unresolved part of the repository's stabilizer question is to
identify \(F_{\mathrm{all}}\), and independently to identify
\(D_{\Oz}\). Definable closure need not automatically equal an
automorphism fixed field in a proper-class structure.
The series \(\mathcal E\) in \eqref{dsn:rp:eq:formal-exp} is a concrete test:
every strong automorphism fixes it, while pure-field definability
has been ruled out. Can an \(\Oz\)-preserving automorphism fixing
all real and ordinal parameters move it?

Further questions include intrinsic definable closure in the
omega-map expansion, comparisons between real constants and a
predicate for the real subclass, and versions without set choice.
The irrational-coefficient shear uses a \(\Q\)-linear functional on
\(\R\), so eliminating choice requires a separate argument.

\begin{proposition}[The fixed field need not consist of definable elements]
\label{dsn:rp:prop:fixed-not-definable}
If $V\ne\HR$, then $\mathcal F_\Gamma$ contains an element outside $\ROD$.
In particular the bound on intrinsic definable closure by
$\mathcal F_\Gamma$ is strict in such a universe.
\end{proposition}
\begin{proof}
The global coding theorem yields an ordinal $\kappa$ and
$A\subseteq\kappa$ with $A\notin\ROD$. Form the normal form
\[
 X_A=\sum_{\alpha<\kappa}(1+\chi_A(\alpha))\omega^{-\alpha}.
\]
Its support is reverse well ordered, and $-\alpha\in \Gamma_{\Q}$ for every
ordinal $\alpha$. Thus $X_A\in\mathcal F_\Gamma$. Its support order type
recovers $\kappa$ and its coefficients recover $A$. If $X_A$ were
ROD, so would be $A$. Finally every element intrinsically definable
from finite real and ordinal parameters in the stated predicate
languages is ambiently ROD: translate the defining formula into set
theory using the definable carriers and predicates. Hence $X_A$ is
outside those intrinsic closures.
\end{proof}

% ===== End sections/06_strong_stabilizer.tex =====


% ===== Begin sections/07_synthesis.tex =====
\section{Arithmetic consequences and distinctions that remain}
\label{dsn:rp:sec:consequences}

\subsection{The omnific ring remembers the parameter hierarchy}

The equality $\KR=\operatorname{Frac}(\IR)$ does not mean that
omitting nondefinable surreals is an exclusively field-theoretic
phenomenon. The entire real-parameter inclusion hierarchy is already
visible in the omnific rings $I_r$. Likewise the countable closure
criterion can be tested using only positive, purely infinite omnific
summands. Neither fractional parts nor arbitrary real coefficients in
the test series are necessary to detect the missing countable ordinal
information.

There are nevertheless many arithmetic invariants that cannot detect
this hierarchy. We record elementary instances, already established in
the fixed-parameter predecessor and now applicable to $\IR$
(Part~\ref{dsn:part:one}).

\begin{proposition}[Constant term and finite arithmetic]\label{dsn:rp:prop:finite-arithmetic}
Let $\Pi_{\R}$ be the purely infinite part of $\IR$. Then
\[
 \IR=\Pi_{\R}\oplus\Z,
 \qquad \ct:\IR\longrightarrow\Z
\]
is a split unital ring retraction, and $\Pi_{\R}$ is divisible as an
additive group. For every positive ordinary integer $n$,
\[
 n\IR=\Pi_{\R}\oplus n\Z,\qquad
 \IR/n\IR\cong\Z/n\Z.
\]
Every unital ring homomorphism from $\IR$ to a finite ring factors
through $\ct$. A finite system of polynomial equations with ordinary
integer coefficients has a solution in $\IR$ if and only if it has a
solution in $\Z$.
\end{proposition}
\begin{proof}
Normal-form extraction preserves $\ROD$, giving the decomposition.
Products of positive-exponent terms still have positive exponents, so
constant term is a ring homomorphism. Dividing a purely infinite
omnific integer by $n$ preserves its positive exponent support and its
definability. This proves divisibility and the displayed residue-ring
formula. The latter can be formalized using the ordinary integer
representatives $0,\ldots,n-1$, without treating a proper class as a
set-valued equivalence class.

If the additive group of a finite target ring has exponent $m$, any
$p\in\Pi_{\R}$ has the form $mq$ and maps to zero. Thus the homomorphism
factors through $\ct$. Applying the retraction coordinatewise to a
polynomial solution proves the forward implication in the last
assertion; ordinary integers supply the reverse implication.
\end{proof}

Nonvanishing constraints change this conclusion. For example,
$X^2=2Y^2$ has the nonzero solution
$(X,Y)=(\sqrt2\,\omega,\omega)$ in $\IR$, but no nonzero ordinary integer
solution. Also $\IR$ is not Noetherian, since
\[
 (\omega)\subsetneq(\omega^{1/2})\subsetneq
 (\omega^{1/3})\subsetneq\cdots.
\]
The inclusion follows by monomial division; the reverse inclusion
would require a nonzero infinitesimal to be omnific. These examples
show why agreement of finite quotients with $\Z$ should not be confused
with ordinary integer arithmetic in all logical languages.

\subsection{A real parameter can encode long countable signs}

Every countable-birthday surreal is definable from one real alone, in
every ambient ZFC universe. A real can encode a well-order on a subset
of $\omega$ and a sign on each of its elements. The well-founded
collapse supplies the actual ordinal domain, not just a numerical
approximation to it. This gives an exact definition of the surreal.

This statement does not say that every ordinal sequence is real-only
definable. A countable sequence of ordinals may use ordinals far above
$\omega_1$; its entries need not have countable codes. The principle
$\Asm$ concerns exactly this extra information. Prikry forcing exploits
that difference: it changes a high ordinal's cofinality without adding
a real.

\subsection{Initiality and order topology}

\begin{proposition}[Topological behavior of the varying-real field]
If $\KR\ne\No$, then $\KR$ is closed and nowhere dense in the
surreal order topology, understood through intervals and definable
classes. It can nevertheless fill every external countable cut.
\end{proposition}
\begin{proof}
If $x\notin\KR$, let $x^-$ and $x^+$ be the signs of $x$ extended by
one minus or one plus. Every $y\in(x^-,x^+)$ extends $x$ as a prefix.
If such a $y$ belonged to $\KR$, its prefix of length $\bd(x)$ would
belong to $\KR$, using the same real and that ordinal. This is
impossible. Thus the complement is open.

No proper subfield of an ordered field contains an open interval.
Indeed, if it did, choose $c<d$ in such an interval and in the subfield.
For any positive $u$ of the ambient field,
$c+(d-c)/(u+1)$ lies strictly between $c$ and $d$, so belongs to the
subfield; solving for $u$ puts $u$ there as well. Negation gives the
whole field. Hence $\KR$ has empty interior. The final assertion is
realized by \cref{dsn:rp:thm:cohen-atlas,dsn:rp:thm:cohen-threshold}.
\end{proof}

The prefix-gap mechanism is an antecedent of the repository and of
Hamkins's analysis of ordinal-definable dense maps
\cite{HamkinsDensity,RepoUniv}. Its application here emphasizes that
countable interpolation is much weaker than ambient density. The
omitted gaps can require uncountable information to specify from
inside the field.

\subsection{Examples by language}

The following table separates representative notions. The negative
claims about the enriched pair allow all real and ordinal parameters,
but do not add the omega map or simplicity.

\begin{center}
\small
\begin{tabular}{>{\raggedright\arraybackslash}p{0.25\linewidth}>{\raggedright\arraybackslash}p{0.30\linewidth}>{\raggedright\arraybackslash}p{0.33\linewidth}}
\toprule
Object or class & Ambient set-theoretic information & Intrinsic comparison\\
\midrule
$\omega$ & Parameter-free & Not definable from real parameters in the pure ordered field\\
$\omega^{1/\omega}=\log\omega$ & Parameter-free, purely infinite omnific & Not definable in the field/omnific/real/ordinal predicate language even from all real and ordinal constants\\
$\omega+\omega^{\omega^{-\alpha}}$, $0<\alpha\in\Ord$ & Ordinal-definable from $\alpha$ & Outside the computed strong fixed field, hence outside the same intrinsic definable closure\\
A countable-birthday sign sequence & Definable from one coding real & This does not imply field definability from that real\\
$Z_f$ from an ordinal sequence $f$ & Exactly the information of $f$ & Each individual summand is OD; the whole sum need not be $\ROD$\\
All of $\KR$ & Varying real plus finite ordinals & One real suffices exactly when $\HR$ satisfies Choice\\
\bottomrule
\end{tabular}
\end{center}


\subsection{Consequences and common incorrect inferences \src{05}}
\label{dsn:ca:sec:consequences}

\paragraph{Seven conclusions about definable omnific integers.}
The results can now be read directly as a theory of definable integers,
not merely as a theory of surreal fields with an extra predicate.
First, for each policy the integer part determines the field by
fractions with a denominator chosen from the same parameters.
Second, the real-parameter spectrum of an arbitrary surreal can be
transported intact to a single omnific monomial below $\omega$.
Third, a countable sequence of ordinals can be transported intact to
a coefficient-$1$ omnific normal form of length $\omega$.

Fourth, the countable-summation behavior of $\IR$ detects exactly the
same assembly principle as the entire field $\KR$.
Fifth, replacing finite ordinal information by one countable ordinal
parameter sequence gives $\Ico$, a countably sum-closed integer part
of the canonical ambient-stable completion.
Sixth, in an uncountable Cohen product there are non-$\ROD$ omnific
integers below $\omega$, although every countable Hahn sum of $\ROD$
omnific integers is $\ROD$.
Finally, in the Prikry example even a countable Hahn sum of $\OD$
omnific monomials can fail to be $\ROD$.

The last two conclusions are compatible. In the Cohen model, a missing
bounded integer cannot have a Hahn-summable countable presentation whose
terms all lie in $\KR$. In the Prikry model, the formula
\eqref{dsn:ca:eq:prikrycode} gives exactly such a presentation.
\paragraph{Definable entries, definable support, and definable graphs.}
These are different kinds of information. If a whole normal form is
$\ROD$, every indexed exponent and coefficient is $\ROD$ because its
ordinal index may be used as an additional parameter. Conversely, even
when all coefficients are fixed to $1$ and every exponent is $\OD$,
the support may fail to be $\ROD$: \eqref{dsn:ca:eq:prikrycode} proves this.

The separated bands make the support and the ordinal sequence
interdefinable. Since the coefficients are constant, a $\ROD$ support
would give a $\ROD$ normal form and hence a $\ROD$ sum. In the negative
example the support itself therefore cannot be $\ROD$. This pinpoints
the failure instead of attributing it to an unspecified difficulty with
infinite series.

A definable \emph{set} of possible coefficients or possible exponents
is likewise insufficient. One needs the actual assignment, including
which values occur and where. A collection of individually available
objects is not automatically an available sequence of those objects.
\paragraph{All reals and real closedness do not give all surreals.}
The Cohen model has
\[
 \R\subseteq\KR\subsetneq\No,
 \qquad \KR\text{ initial, real closed, and countably Hahn-sum closed}.
\]
It also contains every ordinal. Thus neither possessing every real and
every ordinal, nor adjoining the natural countable summation operation,
forces equality with the full surreal class. Information distributed
across uncountably many generic coordinates is still absent.
Every element of $\No\setminus\KR$ is transcendental over $\KR$,
since a real closed field has no proper ordered algebraic extension.
This is an ordinary algebraic observation; no claim about a
proper-class transcendence basis is needed.
\paragraph{Why there is no definability paradox.}
Nothing here forms a set of all surreal numbers, a set of all classes,
or a uniform truth predicate for $V$. The code $Q$ maps a proper class
into a bounded \emph{proper-class} interval, so there is no conflict
with cardinal arithmetic for sets. The code $\Zc(f)$ sums only an
$\omega$-indexed set of terms. The assembly theorem manipulates
rank-local certificates, not a list of all true global formulas.
The completion theorem uses a scheme of fixed definable maps, not an
operation evaluating arbitrary global formula codes.

The internal definitions of $\ROD$ and $\COD$ legitimately absorb
rank ordinals into their permitted parameters. The parameter-free
and real-only notions are instead treated externally or as schemes.
This is why the least missing $\ROD$ birthday is meaningful while an
unqualified ``least parameter-free undefinable surreal'' construction
is not justified by the same reasoning.

\paragraph{Source 05's conclusion.}
Definability is not an intrinsic adjective attached to a surreal
independently of its language, parameters, and universe. With those
choices fixed, the algebra is stable: ambient definable surreals form
real closed fields, and the corresponding definable omnific integers
recover those fields by fractions. Ordinal parameters make the
hereditary sign-code principle available; arbitrary real parameters
produce a countably directed union of the fixed-real fields.

The central additional phenomenon is assembly. Real parameters absorb
countable real data, but not automatically countable ordinal data.
The explicit code $\Zc(f)$ translates that distinction into a
countable coefficient-$1$ omnific normal form in $(0,\omega)$.
It gives an exact criterion for countable Hahn-sum closure, rather
than a vague warning about nonuniform definitions. Allowing a
countable ordinal sequence as a parameter yields the corresponding
minimal ambient-stable countably summable extension.

The Cohen and Prikry models then separate what the definitions can do.
A field can be initial, real closed, contain all reals and all ordinals,
and be countably Hahn-sum closed, while still omitting surreals as early
as birthday $\omega_1$. Conversely, without adding any new reals, a
countable sum of ordinal-definable omnific monomials can fail to be
real--ordinal definable. Both conclusions are obtained with explicit
parameter accounting and without a proper-class sum or a truth predicate
for the universe.

\section{A precise agenda for further research}
\label{dsn:rp:sec:agenda}

The questions below distinguish what has been resolved from what is
still missing. They are proposed problems arising from this paper and
the cited repository, not claims that every formulation is new to the
literature.

\subsection{The remaining stabilizer and language questions}

\begin{question}[Remove the strong-sum hypothesis]\label{dsn:rp:q:nonstrong}
Compute the common fixed field of \emph{all} ordered-field
automorphisms fixing $\R\cup\Ord$ pointwise, both with and without
preservation of $\Oz$.
\end{question}
The theorem in \cref{dsn:rp:sec:strong} computes the strong case exactly.
Without strong summation the lower inclusion for an arbitrary normal
form is unavailable. The relative real closure of $\R(\Ord)$ is still
fixed, while the full common fixed field is contained in the strong
one. Determine whether either bound is exact, or whether a proper
intermediate field occurs. A proof must control automorphisms that do
not commute with an infinite Hahn sum; the coefficient-by-coefficient
argument cannot simply be reused.

\begin{question}[Definable closure versus fixed elements]
Determine the exact intrinsic definable closure of $\R\cup\Ord$ in
$(\No;+,-,\cdot,<,0,1,\Oz)$ and in the expansion also naming the
predicates $\R$ and $\Ord$.
\end{question}
The strong fixed field is an upper bound, not an identification of
definable closure. In a universe with $V\ne\HR$, it contains elements
outside $\ROD$, so it is strictly larger than these intrinsic closures.
Even when $V=\HR$, invariance under automorphisms need not supply a
finite isolating formula in this proper-class structure. The
parameter-free real trace of the omnific expansion is a separate
remaining question from (Part~\ref{dsn:part:one}).

\begin{question}[Name the omega map]
Classify definable closure in the omega-map expansion, and compare it
with the exponential and simplicity expansions.
\end{question}
Naming $\Omega$ immediately defines $\omega=\Omega(1)$ and
$\omega^{1/\omega}=\Omega(\Omega(-1))$. It also makes bounded
monomial coding intrinsic. The previous automorphisms do not preserve
this function and therefore give no nondefinability theorem in that
larger language. Determine which additional set-theoretic coding is
actually expressible, rather than inferring it from the ambient graph
of the map.

\subsection{The real-parameter hierarchy and countable information}

\begin{question}[Forcing preservation of countable coherence]
Characterize forcing notions that preserve $\Asm$, and determine its
behavior under iterations and under forcing that makes new objects
ordinal definable.
\end{question}
The Namba example in \cref{dsn:rp:thm:namba} already proves failure without
a large-cardinal assumption, so the basic independence question is
settled here. The Cohen and Namba examples also show that $\Asm$ and
Choice in $\HR$ imply neither one another. What remains is a structural
preservation theorem: the property concerns definability in the final
universe, not merely whether ordinal sequences were added. Coding
forcing can change that definability, so distributivity alone is not
a complete account of its behavior.

\begin{question}[Higher cardinal versions]
For an uncountable regular cardinal $\kappa$, identify the correct
parameter class and coding principle corresponding to Hahn sums of
size less than $\kappa$ and to interpolation of cuts of size less than
$\kappa$.
\end{question}
At the countable level, countably many reals can be joined into one
real. That is the special fact powering the theorem. For uncountable
families it fails, so merely replacing $\omega$ by $\kappa$ in the
statement is unjustified. Parameters from $\R^{<\kappa}$ or
$\Ord^{<\kappa}$ are natural alternatives, and their hereditary models
should be compared explicitly.

\begin{question}[Cofinality and order types of real-parameter fields]
Which inclusion preorders can the fields $K_r$ realize? Which
uncountable cardinals can be their cofinality? How do these invariants
interact with $\Asm$ and with Choice or Dependent Choice in $\HR$?
\end{question}
The Cohen computation gives the precise ground countable-support
poset, and in particular cofinality $\lambda$ for regular uncountable
$\lambda$ over $L$. More complicated forcing may produce degree
patterns that are not coordinate-inclusion orders. Finite joins are
known exactly; countable least upper bounds require separate analysis.

\begin{question}[Possible first missing birthdays]
If $\KR\ne\No$, which ordinals can be the first birthday outside
$\KR$, and what restrictions follow from $\Asm$?
\end{question}
The lower bound is always $\omega_1$, and the Cohen model attains it.
For a general transitive inner class, the first missing sign sequence
is equivalent to the first missing subset of an ordinal. Here the
relevant class can fail Choice, so arguments based on an internal
cardinal well-order must be reconsidered rather than imported without
checking. The interaction with cofinality-changing forcing is a
particularly concrete target.

\begin{question}[Uniformity and complexity]
Find sharp syntactic complexity bounds for the maps from definitions
of ordinal sequences to definitions of their omnific codes, and for
the reverse normal-form decoding.
\end{question}
The proof gives fixed formula substitutions but does not bound their
optimal position in the Levy hierarchy. It also does not identify
semantic definability with computability. A useful formalized result
would be an explicit transformation on formula codes with a proved
parameter and complexity budget.

\subsection{Normal forms, parameter recovery, and integer arithmetic}

\begin{question}[Minimal parameter fields for individual numbers]
For which $x\in\KR$ does the collection
$\{K_r:x\in K_r\}$ have a least member? Can its failure be witnessed
by a purely infinite omnific integer in a fixed bounded interval?
\end{question}
The finite-join theorem does not give such a least field: it provides
upper bounds for parameters already selected. Reversible omnific
codes transfer any answer between general surreals and omnific
integers, so the integer restriction does not remove the issue.

\begin{question}[Exact birthday cost of countable ordinal codes]
Estimate sharply $\bd(Z_f)$ from information about $f$, and compare
the rational-slot code with nested-monomial codes.
\end{question}
The construction controls support order type and numerical size, not
the birthday of its output. These are very different invariants in
$\No$. Any optimality claim must specify the permitted coding maps and
parameter policy.

\begin{question}[Arithmetic recovery of countable coherence]
Can $\Asm$ be characterized by a natural first-order or infinitary
property of the ring $\IR$ alone, without explicitly naming canonical
normal forms or the omega map?
\end{question}
The current characterization uses external countable families and
Hahn summation. Finite residue rings and finite integer-coefficient
equations cannot distinguish the relevant models by
\cref{dsn:rp:prop:finite-arithmetic}. A successful intrinsic characterization
would need additional arithmetic structure or a controlled infinitary
language.


\begin{writenote}[Status of source 02's questions, 4 October 2026, batch~93]
All questions above stay open as printed, with these re-scopings.
\Cref{dsn:rp:q:nonstrong} is the nonstrong remainder of
\lbl{sse:sr:q:stabilizers}. The question on naming the omega map continues
\cref{dsn:q:languages} of Part~\ref{dsn:part:one}; Part~\ref{dsn:part:op}
answers its ``which additional set-theoretic coding is actually
expressible'' part for the omega map together with $\Oz$, at the level of an
interpretation (\cref{dsn:op:thm:V-interpretation}), while the classification
of definable closure stays open there (its \cref{dsn:op:q:dcl}) and in the
language with the omega map alone. The question on uniformity and complexity
meets \cref{dsn:q:cost}; Part~\ref{dsn:part:op} gives $\Delta_1$ graphs for
the canonical operations and level preservation for the bounded code
(\cref{dsn:op:thm:canonical-delta-one,dsn:op:cor:code-levels}), not for the
codes $Z_f$. The question on the birthday cost of countable codes refines
\cref{dsn:q:cost-birthday}.
\end{writenote}

\subsection{Source 05's questions \src{05}}
\label{dsn:ca:sec:questions}

Source~05 proposes ten problems. Except where an earlier source is
explicitly named, they are not asserted to be previously published open
problems, and none is an unstated hypothesis of the proved theorems. Two are
answered by source~02 and are printed with their answers; four are source~02's
questions above and are printed once, with source~05's discussion here; four
are new.

\begin{question}[The consistency strength of failed assembly \src{05}; answered by \cref{dsn:rp:thm:namba}]\label{dsn:ca:q:strength}
What is the exact consistency strength of $\neg\Asm$? Can a transitive
model of ZFC alone support a no-new-reals, sufficiently homogeneous
extension adding a non-$\ROD$ countable ordinal sequence?
\end{question}
The Prikry construction gives a measurable-cardinal upper bound, not
a lower bound. A proposed replacement by Namba-type or other
cofinality-changing forcing must separately verify no new reals and
the relevant homogeneity or nondefinability argument. It is not enough
merely to add an $\omega$-sequence of ordinals.

\begin{writenote}[Answered, 4 October 2026, batch~93]
\Cref{dsn:rp:thm:namba} answers both parts: a transitive model of ZFC alone
supports such an extension (Namba forcing over its $L$), and the consistency
strength of $\neg\Asm$ is exactly that of ZFC; see the note at the end of
\cref{dsn:ca:sec:prikry}. Source~02 carried out exactly the checks the
discussion asks for: no new reals (imported, under CH; see
\cref{dsn:rp:q:nambainput}), and cone homogeneity with the resulting
nondefinability (proved).
\end{writenote}

\begin{question}[Weakening ambient stability in the completion theorem \src{05}]
Which explicitly listed surreal operations, together with countable
Hahn sums, generate $\Kco$ from $\KR$? Is there a useful natural
intermediate field obtained using only field operations, the Conway
map and its inverse, normal-form extraction, and countable summation?
\end{question}
The current minimality theorem uses all fixed ambient definable maps.
Removing that hypothesis is a substantive algebraic/model-theoretic
problem, not a cosmetic sharpening. A negative answer should identify
an invariant preserved by the restricted operation list.


\paragraph{Source 05's Question 13.3} (longer parameter sequences and longer
sums) is the question on higher cardinal versions above. Source~05 asks: for
an infinite cardinal $\mu$, compare definability from a $\mu$-sequence of
ordinals with closure under $\mu$-indexed Hahn sums; can one obtain a
reversible normal-form code with a prescribed support order type and an exact
analogue of \cref{dsn:ca:thm:assembly}?
The present dyadic-band construction uses disjoint bands indexed by
$\omega$. Uncountably many real bands are unavailable. A longer code
would need transfinite surreal scales, with careful control of reverse
well ordering and of the parameters needed to identify those scales.

\paragraph{Source 05's Question 13.4} (real-parameter spectra and least
information) contains the question on minimal parameter fields above. Source~05
asks: which upward-closed subsets of the real-ordinal reducibility preorder
occur as $\Spec(x)$ for a surreal $x$; when does a nonempty spectrum have a
least degree; and which spectra already occur for single monomials in
$(0,\omega)$?
The bounded code answers the last occurrence question relatively:
every spectrum occurring for a surreal also occurs for such a monomial.
The remaining task is the classification of spectra themselves,
not another coding argument. The Cohen support theorem gives concrete
necessary exclusions for coordinate generics.

\begin{question}[Minimal coordinate supports in Cohen products \src{05}]
Refine \cref{dsn:rp:thm:cohen-atlas} from existence of a countable ground-model support
to a classification of possible minimal supports and intersection
behavior. Which surreal elements lie in every $W_J$ from a given
family of countable coordinate sets?
\end{question}
A useful first step is to distinguish pairwise intersections of
mutually generic subextensions from arbitrary intersections of a
countable decreasing family. The distinction should be proved before
using it to assert existence of a least support.

\paragraph{Source 05's Question 13.6} (the first omitted real--ordinal
birthday) is the question on possible first missing birthdays above, asked
also for cofinalities: what values can $\lambda_\R$ take, must it be a cardinal
of the ambient universe, what cofinalities are possible, and can one prescribe
it while independently controlling $\Asm$?
The Cohen product gives exactly $\omega_1$. For $\OD$, the repository
has a stronger first-birthday analysis through a single ZFC inner
model $\HOD$. The union over all reals lacks the same immediate
canonical well-order argument, so those conclusions should not be
transferred without proof.

\begin{question}[Global real parameters versus one fixed real \src{05}; answered negatively by \cref{dsn:rp:cor:global-single-real}]\label{dsn:ca:q:global}
Can $\ROD=V$ hold while $V\ne\HOD(r)$ for every real $r$? How do any
such examples appear through the bounded omnific test of
\cref{dsn:ca:prop:globaltest,dsn:thm:bounded-code}?
\end{question}
The Cohen example proves that a proper $\KR$ need not equal any
single $K_r$. It does not settle the stronger global situation posed
here. The quantifiers $\forall x\,\exists r$ and
$\exists r\,\forall x$ must remain distinct.
\begin{writenote}[Answered negatively, 4 October 2026, batch~93]
No. Over ZFC, \cref{dsn:rp:cor:global-single-real} (and source~06's
\cref{dsn:rp:rem:consolidation06}) shows that $\ROD=V$ implies $V=\HOD(r)$ for
one real $r$: by Choice the reals have a well-order, which is $\OD(r)$ for some
real $r$, and then every real, hence every set, is $\OD(r)$. The quantifiers
$\forall x\,\exists r$ and $\exists r\,\forall x$ are distinct in general
(the Cohen example), but they coincide when the $x$ range over all sets.
Through the bounded omnific test the answer reads: if every omnific integer
of $(0,\omega)$, or even every monomial $\omega^a$ with $0<a<1$, is $\ROD$,
then one real suffices for all of them.
\end{writenote}

\paragraph{Source 05's Question 13.8} (birthday and syntax cost of the
multiplexer) is the question on the birthday cost of countable codes above,
together with the question on uniformity and complexity. Source~05 asks for
sharp bounds on $\bd(\Zc(f))$ from the birthdays or ordinal sizes of the
entries of $f$, a comparison with alternative reversible countable omnific
codes, and effective translations of definition complexity through the
encoder and decoder.
The construction fixes the support order type and coefficient alphabet,
not the exact birthday. This refines the earlier repository question
about the birthday cost of bounded coding. Fixed-formula definability
preservation does not automatically preserve a specific Levy-hierarchy
level or an optimal description length.

\begin{question}[Real-only definability without ordinal parameters \src{05}]
Study the external first omitted ordinal and the maximal initial
fragment of $K_{\mathrm{rd}}$ in specified transitive models. How do
these depend on the supply of reals and on the model's height?
\end{question}
Every countable-birthday surreal is already in $K_{\mathrm{rd}}$,
but the rank-local-certificate argument cannot silently add rank
ordinals under this policy. A solution should specify its external
truth semantics and not confuse $K_{\mathrm{rd}}$ with $\KR$.

\paragraph{Source 05's Question 13.10} (intrinsic languages between the field
and set theory) is the question on naming the omega map above and
\cref{dsn:q:languages}: which parts of the real-parameter hierarchy are
captured by the ordered field expanded by the omnific predicate, the Conway
map, or simplicity, and can any of the countable-assembly phenomena be
detected in a precisely specified set-sized first-order structure?
This continues the language question in (Part~\ref{dsn:part:one}). The pure field
answer is \cref{dsn:thm:purefield}, but the present ambient codes use
normal forms and set-theoretic reconstruction. They do not solve the
intrinsic definability problem merely by existing as ambient functions.
See the status note above for what Part~\ref{dsn:part:op} adds.

\subsection{Further questions and research added at the write \writetag}
\label{dsn:rp:sub:further}

The following item is recorded under Vladimir's standing rule of 4~October
2026 that unproved claims are kept as open questions with their source,
argument sketch and what is missing.

\begin{question}[A reference for the Namba input \writetag]
\label{dsn:rp:q:nambainput}
Give a complete proof, or a textbook reference that states it for the
presentation used, that Namba forcing in the Laver-style presentation of
\cref{dsn:rp:sec:namba} (every node above the stem has $\omega_2$ immediate
successors) adds no reals over a ground model of CH.
\end{question}
Source~02 imports this from Lietz's exposition \cite{LietzNamba2024}, which
distinguishes the two tree presentations and states that the classical
no-new-reals proof applies to both, and from Miller's lecture notes
\cite[Theorem~23 and the table on p.~4]{MillerForcing2009}. The standard
textbook treatments of Namba forcing under CH are \cite[Chapter~28]{Jech} and
\cite[Chapter~XI]{ShelahPIF}; the latter also shows that the two presentations
are not forcing equivalent (Magidor and Shelah, its Claim~XI.4.2, as reported
in \cite{LietzNamba2024}). The collection imports the theorem for the
classical presentation from \cite[Section~7.6]{FKW}
(\lbl{univ:gs:thm:nambainput}). The write consulted neither textbook and did
not confirm for which presentation they prove it. The argument for either
presentation is a fusion over $\aleph_2$-splitting trees that uses CH to
handle the $\aleph_1$ possible values of a name for a real; what is missing
here is a written check of that fusion for trees in which every node above the
stem splits. Everything else in \cref{dsn:rp:thm:namba}, in particular cone
homogeneity and the computation of $\HR$, is proved in the text. Should the
input fail for this presentation, the theorem would still hold for any
presentation that adds no reals under CH, is cone homogeneous, and adds a
cofinal $\omega$-sequence in $\omega_2^W$.

\section{Formalization plan and reproducibility ledger}
\label{dsn:rp:sec:formalization}

\subsection{A safe sequence of formal developments}

The following order keeps set-theoretic semantics separate from
surreal calculations.

\begin{enumerate}
\item Define canonical signs, birthdays, and the translation from
sets of ordinals to sign sequences. Prove that all hereditary
components of a sign graph except the graph itself are OD.
\item Formalize satisfaction for set structures, rank-local
$\OD(r)$, real joining, and substitution of definitions. Prove the
inner-model axioms for $H_r$ and $\HR$ with Choice kept as a separate
conclusion for the latter.
\item Derive the parameter inclusion and capture theorems. This
stage needs no sophisticated surreal arithmetic: real injection
and the canonical sign representation are enough.
\item Formalize the rational-slot exponents and their inequalities,
then normal-form decoding and the finite-prefix cut argument.
These are the central algebraic lemmas of the countable theorem.
\item Formalize additive index relabeling and Hahn lifting, the
reverse-well-ordered orbit lemma, and the rational-span description.
For the coefficient shear, isolate the ordinary set-sized Hamel-basis
choice as an explicit hypothesis or a ZFC construction.
\item Add the Cohen, Namba, and Prikry models through a set-theoretic
forcing library. Until that is available, keep their exact forcing
inputs as named assumptions and formally prove the conditional
surreal consequences.
\end{enumerate}

For an implementation in a type theory with universes, $\No$ at one
universe level is a set or type only from a larger level. Moving to the
larger level also changes which parameter families and sequences are
available. The external countable closure statements must retain their
ambient-universe annotation. Likewise, a finite formula with constants
from a proper parameter class is represented by its finite list of
constants, not by a set-sized language assumed to contain the entire
class at the same level.

\subsection{What was checked and what was imported}

\begin{longtable}{>{\raggedright\arraybackslash}p{0.30\linewidth}>{\raggedright\arraybackslash}p{0.57\linewidth}}
\toprule
Item & Status and dependency\\
\midrule
\endhead
Canonical signs, real closure, omega map and normal forms & Classical surreal theory; sources \cite{Conway,Gonshor,vdDE}.\\
Internal ordinal definability and real-plus-ordinal parameters & Classical rank-satisfaction framework; \cite{DoraisHamkins}.\\
Fixed-parameter closure, floor, denominators and bounded coding & Already in the predecessor (Part~\ref{dsn:part:one}); short proofs included where needed.\\
Varying-real inclusion and Choice criterion & Proved here by real joining, hereditary coding, and a well-order of the actual real set.\\
Countable assembly and interpolation equivalence & Proved here with a uniform reversible ordinal-sequence code and explicit prefix inequalities.\\
Cohen field hierarchy & Standard forcing inputs, with the exact parameter-field computation and closure consequences proved here.\\
Negative countable-assembly models & Namba gives failure relative to ZFC consistency; Prikry gives the displayed measurable-cardinal example. The exact Namba no-new-reals theorem is imported; cone homogeneity and definability consequences are proved here.\\
Hahn lifts and support-orbit mechanisms & Antecedents in \cite{RepoOpa,RepoSse}; rederived for the specific maps used.\\
Strong stabilizer of all real and ordinal constants & Exact rational-span fixed field proved here; strong version of the repository question resolved.\\
Full nonstrong stabilizer or full intrinsic definable closure & Not resolved.\\
Machine verification & No new theorem is claimed to be checked in Lean or Rocq. The finite code-inequality script is an auxiliary check only.\\
\bottomrule
\end{longtable}

\subsection{Package contents}

The supplied source compiles with a standard LaTeX installation using
\texttt{latexmk -pdf definable\_surreals.tex}, or repeated runs of
\texttt{pdflatex}. The package includes a self-contained TeX source, the
compiled PDF, a short build guide, and an exact-arithmetic check of the
coding inequalities. There are no downloaded third-party articles in
the package. Bibliographic links identify the primary sources and the
immutable repository snapshot instead.


\begin{writenote}[On the package contents, 4 October 2026]
The preceding subsection describes source~02's delivery. In this report only
the checking script and the build script of source~02 are shipped, as
\fn{code/08-real-parameters-finite_code_check.py} and
\fn{code/08-real-parameters-build.sh}; the latter builds the undelivered
\fn{definable_surreals.tex} and is kept only as delivered. The manuscript
source and PDF survive in the arrival commit's archive. See
\cref{dsn:rp:sub:provenance} and the README.
\end{writenote}

\subsection{Source 05's formalization plan and proof ledger \src{05}}
\label{dsn:ca:sec:formalization}

\paragraph{Separate the mathematical interfaces.}
A realistic formalization should not begin by treating the entire
proper class $\No$ as a small type in its own universe. The first
component is a universe-polymorphic sign-sequence representation,
with the ordinal domain recorded in a larger universe when necessary.
A set-model semantics can alternatively represent exactly the
surreal codes belonging to a specified model.

The second component is a normal-form interface: uniqueness of the
support/coefficient graph, definable construction of the Conway map,
and the formal Hahn summation operation for set-indexed families.
The separated-band theorem only needs ordered-field identities,
monomial comparison, the normal-form existence theorem, and an
ordinal-code decoder. It should be formalized independently of HOD
or forcing.

The third component is first-order syntax and semantics for set
structures, including satisfaction for $V_\theta$ and reflection
as an explicitly stated theorem. It must distinguish external standard
formulas from any internal nonstandard syntax. The proof of
\cref{dsn:ca:thm:assembly} then becomes a certificate-coding theorem plus the
normal-form code, rather than an appeal to an informal truth oracle.

The fourth component is forcing: ccc countable supports, the
cone-homogeneity lemma, and the Prikry property. These can initially
be isolated as named hypotheses so that the new transfer theorem is
checked without pretending that an existing library already contains
a complete development of Prikry forcing.
\paragraph{What has actually been checked here.}
The accompanying program uses exact rational arithmetic to check the
finite specialization of \eqref{dsn:ca:eq:band} and \eqref{dsn:ca:eq:decode}, including
nonmonotone inputs and repetitions. A successful run is useful for
catching an incorrect factor of $2$ or band endpoint. It is not a
formal proof of any statement involving infinite ordinals, real
parameters, satisfaction, or generic extensions.

The LaTeX source is compiled to PDF and inspected for layout defects.
No Lean source is presented as verified, and no Lean build was run.
The repository's earlier definability report is treated according to
its own pending formalization status (Part~\ref{dsn:part:one}).

\begin{longtable}{@{}P{0.22\textwidth}P{0.43\textwidth}P{0.27\textwidth}@{}}
\caption{Proof and provenance ledger.}\label{dsn:ca:tab:ledger}\\
\toprule
Result & Mathematical input & Status in this manuscript\\
\midrule
\endfirsthead
\toprule Result & Mathematical input & Status in this manuscript\\\midrule
\endhead
Parameter taxonomy and certificates & Set satisfaction, reflection, finite parameter coding & Standard foundations, made explicit\\
Field, floor, fractions, fixed-real HOD & Classical surreal theory and unique-output closure & Earlier repository results; proofs reproduced in a self-contained form\\
Countable-birthday real coding & Countable well-order codes & Elementary consequence, proved here\\
Bounded monomial code $Q$ & Real closed field algebra and $\Omega$ & Explicitly credited to the repository\\
Countable code $\Zc$ & Separated rational bands and normal forms & Proposed tailored construction; full proof\\
\Cref{dsn:ca:thm:assembly} & Certificates, Choice, $\Zc$ & Main proposed equivalence; full proof\\
\Cref{dsn:ca:thm:completion} & Classical countable-ordinal coding and $\Zc$ & Proposed surreal minimality bridge; closure policy explicit\\
\Cref{dsn:rp:thm:cohen-atlas,dsn:ca:thm:cohenconsequences} & ccc, product factorization, homogeneity, $L[r]$ & Proposed classification/application; detailed proof\\
\Cref{dsn:ca:thm:prikry} & Classical normal-measure Prikry property; cone isomorphisms & Proposed omnific application; exact imported premise identified\\
Finite verification & Python standard library, exact fractions & Executed regression checks only\\
Historical priority & Selected repository and primary-source search & Not certified\\
\bottomrule
\end{longtable}

Source~05's build notes (its Appendix~B) describe a delivery of its source,
its PDF, a README, a separate proof audit and the regression script with its
recorded JSON output, built with
\verb|latexmk -pdf -interaction=nonstopmode -halt-on-error article.tex| and
checked with \verb|python3 code/verify_code.py --output verification.json|.
The regression script tests only the finite rational specializations of the
countable code; the mathematical proofs, rather than that program, justify
its transfinite use. No third-party papers, repository source copies or font
files were redistributed, and the delivery modified nothing in the
repository. \textbf{[write]} Of these, the script, its recorded output, the
audit and the \fn{Makefile} are shipped under the names given in
\cref{dsn:rp:sub:provenance}; the audit still writes the band exponents
$e_n(a)$ (here $b_n(a)$). Rerun only on a copy, as the README explains: the
delivered \fn{Makefile} and command line write over the recorded output.

% ===== End sections/07_synthesis.tex =====




% ===================== Part III (batch 93, manuscript 06) =====================
\clearpage
\part{Definable operations: parameter hierarchies, absolute normal forms, uniform coding, and an interpretation of set theory}
\label{dsn:part:op}

% Part III continues the section numbers of Part II.
\setcounter{section}{29}
\renewcommand{\theHsection}{op.\arabic{section}}
\fancyhead[L]{\small Definable surreals, Part~III: definable operations}

{\large Definable Surreals and Omnific Arithmetic\par}
\smallskip
{Research report prepared for Vladimir Reshetnikov, 4 October 2026\par}

\paragraph{Abstract.}
We distinguish definability in the surrounding universe of sets from
definability in a specified surreal structure, and develop separate notions
using no parameters, fixed parameters, ordinal parameters, real parameters,
and combinations of these.  The principal construction is a parameter-free
first-order interpretation of the set-theoretic universe in
\(\Somega=(\No;+,-,\cdot,<,0,1,\Oz,\Omega)\), where
\(\Omega(x)=\omega^x\) is Conway's omega map.  The interpretation uses
equivalence classes of codes lying in a definable subclass of the bounded
omnific monomial interval \(\{\omega^a:0<a<1\}\).  Its mechanism is explicit:
recover the real coefficients, define normal-form supports, encode all
subsets and relations on a support, and collapse pointed well-founded
extensional graphs.  We also define Hahn summability and summation for
uniformly coded set-indexed families in this same first-order language.

Further results concern parameters and operations.  Allowing a different
real and finitely many ordinals for each surreal yields every surreal
exactly when one real \(r\) already satisfies \(V=\HOD(r)\).  A canonical
decoder of real codes for countable sign sequences has a parameter-free
right inverse exactly when \(\omega_1^{\HOD}=\omega_1\), with a corresponding
criterion at every infinite cardinal.  We prove absoluteness of canonical
normal forms between transitive inner models without requiring them to
have the same reals, and analyze the syntactic cost of reversible coding.
The article builds on, and explicitly separates its extensions from, the
existing ProveIt reports.  It establishes mathematical claims by written
proofs; it does not claim a formal verification, a complete determination
of intrinsic definable closure, or a certified historical priority.

\begin{writenote}[Added 4 October 2026, batch~93]
This Part is manuscript~06 of batch~93 of ProveIt's incoming-reports intake,
\emph{Definable Surreals and Omnific Arithmetic: Parameter hierarchies,
absolute normal forms, uniform coding, and an interpretation of set theory},
printed in full after Part~\ref{dsn:part:rp} (provenance, notation and the
relation to Parts~\ref{dsn:part:one}--\ref{dsn:part:rp}:
\cref{dsn:op:sub:provenance,dsn:op:sub:notation,dsn:op:sub:relation}). Like
the other Parts it is AI-assisted, unrefereed and not formalized. To keep the
earlier numbers, its sections continue Part~\ref{dsn:part:rp}'s: \textbf{the
manuscript's Section~$k$ is Section~$k+29$ here, and its statement or equation
$k.m$ is $(k+29).m$}. Thus its Theorem~5.1 is
\cref{dsn:op:thm:all-reals-free-absoluteness}, its Theorem~8.1 is
\cref{dsn:op:thm:V-interpretation} and its Theorem~11.4 is
\cref{dsn:op:thm:theory-faithful}. Its Theorem~4.2 is the same theorem as
source~02's Corollary~4.5 and is printed once, in Part~\ref{dsn:part:rp}
(\cref{dsn:rp:cor:global-single-real} with \cref{dsn:rp:rem:consolidation06});
its Corollary~4.3 is printed once as source~05's Theorem~5.2
(\cref{dsn:ca:thm:countbirthday}); both keep their numbers here as pointer
remarks. Apart from this, the notation changes of \cref{dsn:op:sub:notation},
the dated \textbf{[write]} notes and the editorial changes listed in
\cref{dsn:op:sub:provenance}, no statement, proof, number or non-claim of the
manuscript was changed. ``This article'' in its text means the manuscript.
\end{writenote}

\section{Scope, conclusions, and relation to the repository}
\label{dsn:op:sec:scope}

The expression ``a definable surreal number'' has no unique meaning until
three choices have been made: the ambient model, the language, and the
allowed parameters.  The first choice matters because forcing may add new
sign sequences or change which old objects have definitions.  The second
matters because an ordered field can express much less than its canonical
set-theoretic construction.  The third matters because an ordinal can
specify a transfinite position, whereas a real can carry an arbitrary
countable amount of information.

This distinction becomes especially consequential for omnific integers.
An omnific integer is discrete in the numerical order, but it may have a
transfinite normal form with arbitrary real coefficients at positive
exponents.  Discreteness therefore places almost no bound on its
information content.  A single bounded omnific monomial can already encode
an arbitrary surreal when Conway's omega map is available.

\subsection{The existing starting point}

The repository snapshot used here is commit
\texttt{1404038dfb1cad09eb8a42e38b95f193e5fa54ba}, dated 4 October 2026.
The exact-topic report (Part~\ref{dsn:part:one}) already proves real closedness and
elementarity of ambient definability hulls, same-parameter fraction
representations, the identification with HOD, bounded omnific coding,
several independence results, and a maximal initial definable core.
Its \cref{dsn:q:languages} leaves the omega-map expansion open; \cref{dsn:q:support}
asks for further support fields, and \cref{dsn:q:cost} asks about syntactic
costs.  The companion normal-form report supplies an inner-model instance
under a same-reals hypothesis \cite{RepoLce}.  The reconstruction and
automorphism reports distinguish the omnific predicate from a named omega
map \cite{RepoOdg,RepoOpa}.

Consequently, elementary closure statements and the bounded unary code
are background in the present article.  The main additions are the
support-based interpretation, its operation calculus, the real-parameter
uniformization and section criteria, and the removal of the same-reals
restriction from normal-form absoluteness.  The syntactic analysis gives
explicit preservation statements rather than an assertion that
interdefinability automatically preserves every complexity measure.

The novelty audit was targeted, not exhaustive.  In particular,
Chen--Hamkins--Yang have announced a bi-interpretation for a
\emph{birthday-enriched} surreal structure \cite{CHYKobe}; the general
well-founded-graph technique is not new.  Camacho's work shows the logical
strength of truncation in Hahn fields \cite{Camacho}.  Here the specific
signature, the coefficient-recovery route, and the full universe coded
through arbitrary reverse-well-ordered supports are the point of the
construction.  We prove a mutual interpretation, with one composite
identified in set theory.  We do not infer a bi-interpretation of the
omega-map structure from this.

\subsection{Principal results in operational form}

The following list also indicates how to read the article.
\begin{enumerate}
\item \textbf{Parameter classification.} Ambient definability hulls form
real closed fields and their omnific parts are integer parts with the same
fraction field.  Ordinal-definable sign sequences are precisely the
surreals belonging to HOD.  Individual real parameters and the single set
\(\R\) are different permissions.
\item \textbf{Global real consolidation.} If each surreal is definable
from some real and ordinals, then one real works for the entire universe:
\(V=\HOD(r)\).  This is stronger than a separate statement about each
number, and uses Choice.
\item \textbf{Uniform code selection.} Every countable-birthday surreal
has a real code.  A uniform parameter-free selection of such codes exists
exactly when HOD computes \(\omega_1\) correctly.  The analogous obstruction
at \(\kappa\) is correctness of \(\kappa^+\).
\item \textbf{Absolute operations.} Canonical arithmetic, the omega map,
and normal-form evaluation agree in transitive inner models on common
inputs.  The real coefficients need not exhaust the ambient reals.
\item \textbf{Intrinsic support calculus.} The omnific predicate recovers
the canonical real field and the zero-exponent coefficient.  Adding
\(\Omega\) recovers every coefficient and every support, and gives a
definable pairing compatible with support order.
\item \textbf{Set theory in bounded integer codes.} Pointed well-founded
extensional graphs have first-order codes in \(\Somega\); their quotient
is the ambient \((V,\in)\).  Codes can be compressed into omnific
monomials strictly between \(1\) and \(\omega\).
\item \textbf{Infinite operations with explicit inputs.} The same
language defines summability and the Hahn sum of a coded set-indexed
family.  This does not introduce a proper-class sum or a universal truth
evaluator.
\end{enumerate}


\subsection{Provenance of this Part \writetag}
\label{dsn:op:sub:provenance}

This Part is one manuscript: batch-93 manuscript~06, \emph{Definable Surreals
and Omnific Arithmetic: Parameter hierarchies, absolute normal forms, uniform
coding, and an interpretation of set theory}, ``research report prepared for
Vladimir Reshetnikov'', dated 4~October 2026, pinned to commit
\texttt{1404038dfb1cad09eb8a42e38b95f193e5fa54ba} (the batch-90 write of
Part~VIII of the birthday-cutoff report). It arrived as
\fn{definable_surreal_operations.zip} in commit \texttt{de37a66d1} and was
placed in commit \texttt{47a77daba}; its labels carry the prefix
\texttt{dsn:op:}. At the pin, Part~\ref{dsn:part:one} read as now, except
for the batch-90 note before \cref{dsn:thm:amplify} and the dated notes of
batch~93. It is a separate Part rather than a third member of
Part~\ref{dsn:part:rp} because its spine is different: intrinsic definability
in $\Somega$, absoluteness without a same-reals hypothesis, and uniform code
sections. It shares one theorem with Part~\ref{dsn:part:rp} (its Theorem~4.2,
source~02's Corollary~4.5) and one elementary fact (its Corollary~4.3, source~05's
Theorem~5.2); each is printed once, in Part~\ref{dsn:part:rp}.

The only delivered file shipped besides the text is
\fn{10-definable-operations-PROOF_STATUS.txt}, byte-identical; its theorem
numbers are the manuscript's (add~29 to the section number). The manuscript
ships no code, data or checksum manifest; its source, PDF (41~pages) and
delivery README are not shipped and survive in the arrival commit's archive.
Its preparation checks (a clean build, rendering of all pages, an independent
AI review) are recorded in that file and could not be re-checked at intake.

Editorial changes: citations of the predecessor report became references into
Part~\ref{dsn:part:one}, with its numbered questions and sections cited by
label; the bibliography was merged into the report's (\cite{DoraisHamkins} is
the manuscript's \texttt{DoraisHamkins2019}, \cite{CHYKobe} its \texttt{CHY},
\cite{RepoLce}, \cite{RepoOdg}, \cite{RepoOpa} and \cite{RepoRvs} its
\texttt{RepoLC}, \texttt{RepoDio}, \texttt{RepoAut} and \texttt{RepoVec}); its
Theorem~4.2 with its proof and the paragraph after it, and the statement of
its Corollary~4.3, were moved to Part~\ref{dsn:part:rp} as described above;
its \texttt{\textbackslash sq} ($\lhd$) for proper sign prefix is written
$\ps$.

\subsection{Notation in this Part \writetag}
\label{dsn:op:sub:notation}

Part~\ref{dsn:part:one}'s symbols keep their meaning, and so do
Part~\ref{dsn:part:rp}'s $\ROD$, $\HR$ and $\Qrange$. The manuscript's
$U=\{\Omega(a):0<a<1\}$ is written $\Qrange$, as in Part~\ref{dsn:part:rp},
because $\mathcal U$ is Part~I's code interval $\Cell$; and its proper-prefix
symbol $\lhd$ is written $\ps$. Its $\Szero$ is Part~I's $\mathfrak S$; its
$\tau$ and $Q$ are Part~I's $\tau$ and $\Mc$; its $e_\alpha=\omega^{-(\alpha+1)}$
in \cref{dsn:op:cor:inner-independent} is Part~I's; and for a real $p$ its
$H_p=\HOD(p)$ is Part~II's $H_r$. The following letters have several local
meanings in the manuscript and are kept, because each is introduced where it
is used: $D_A^M$ (definability hulls, as in Part~I), $D$ and $D_\kappa$ (the
real and $\kappa$-code decoders of \cref{dsn:op:sec:real-parameter-extension}),
$D_b$ (the derivations of \cref{dsn:op:thm:derivations}) and $D$ for a support
set in \cref{dsn:op:sec:interpretation}; $B(a,b)$ and $p(a,b)$ (the pairing of
\eqref{dsn:op:eq:B}, not Part~I's $B_n$ nor a set of ordinals $p$); $K_b$, $J_b$
(the operators of \cref{dsn:op:prop:partial-integration}, not hulls $K_A$);
$H_D$ (a product ring); $T_{>a}$ (truncations, not Part~I's $\Two$); and $U$ in
\cref{dsn:op:sec:consequences}, a putative universal function in a diagonal
argument.

\subsection{Relation to Parts I and II and to neighbouring reports \writetag}
\label{dsn:op:sub:relation}

The manuscript credits Part~\ref{dsn:part:one}, the large-cardinal, omnific
Diophantine, automorphism and real-vector-space reports. The following
antecedents in the collection, checked at the write, are not cited by it.
\begin{itemize}
\item Parts (1)--(2) of \cref{dsn:op:thm:all-reals-free-absoluteness}
(old cuts, arithmetic and the omega map) are \lbl{sse:pf:lem:absolute} of
\repopath{surreal/surreal-self-embeddings} \cite{RepoSse}, whose note already
observes that the same-reals hypothesis is not needed for them, and items
(i)--(iii) of \lbl{univ:prop:absolute} \cite{RepoUniv}; its
\cref{dsn:op:cor:inner-rcf-elementarity} is \lbl{univ:prop:absolute}~(iv).
Parts (3)--(5), absoluteness of complete normal forms and sums without the
same-reals hypothesis of \lbl{lce:lem:absolute} \cite{RepoLce}, are new; so is
\cref{dsn:op:cor:inner-support-fields}, which supplies the inner-model instance
that \cref{dsn:rem:largecardinal} (Part~\ref{dsn:part:one}) and
\lbl{lce:rem:dsnsupport} record as not covered.
\item \Cref{dsn:op:prop:native-omnific} largely re-proves
\lbl{odg:def:thm:realrecovery} and \lbl{odg:def:cor:arithmetic} of
\repopath{surreal/omnific-diophantine-geometry} \cite{RepoOdg} and
\cref{dsn:prop:pair}~(i) of Part~\ref{dsn:part:one}, with the same dilation
$S_2$ as \lbl{odg:def:ex:dilation}; its inclusion
$\dcl_{\Szero}(\R)\subseteq\R$ also follows from
\lbl{opa:sc:cor:innerbound} \cite{RepoOpa}.
\item The quotient interpretation \cref{dsn:op:thm:V-interpretation} is for
the signature $(\No;+,\cdot,<,\Oz,\Omega)$; the birthday-cutoff report proves
uniform bi-interpretations for the birthday-enriched structures
(\lbl{hset:thm:biinterpretation}, \lbl{hset:thm:main-rounding}) and records
the transfer $M\equiv N\iff\mathfrak B^M\equiv\mathfrak B^N$
(\lbl{hset:rem:internal}, \lbl{hset:cor:equiv}) and the detectability of an
outer universe, $\mathfrak B^M\prec\mathfrak B^N\iff M=N$
(\lbl{hset:thm:outer}) \cite{RepoHset}. \Cref{dsn:op:thm:theory-faithful} is
the analogue for $\Somega$; see the note after it.
\item \Cref{dsn:op:sec:first-omissions} re-proves for inner models the epsilon
lemma of \cref{dsn:sec:core} and the cardinal threshold of
\lbl{univ:prop:beta} \cite{RepoUniv} (general $M\subseteq N$), of which
\cref{dsn:thm:firstbirthday} is the case $M=\HOD$.
\item Part~\ref{dsn:part:rp}'s \cref{dsn:rp:prop:omega-code} (source~02) is the
special case of the bounded code in the omega-map language; Part~III's support
calculus and interpretation go far beyond it.
\end{itemize}
No statement of this Part is formalized, and no Lean or Rocq declaration of
ProveIt states any of its theorems.

\section{Foundations and the objects being defined}
\label{dsn:op:sec:foundations}

\subsection{Set theory, models, and class notation}

Unless stated otherwise, the background is ZFC.  A surreal is represented
by a function \(s:\alpha\to\{-,+\}\), where \(\alpha\) is an ordinal.
Thus every surreal is a set, while \(\No\) is a definable proper class.
Its birthday is \(\bd(s)=\alpha\).  We use standard set encodings of
functions and ordered pairs and fixed finite tags for the two signs.
The numerical order is lexicographic, treating an ended sequence as a
temporary symbol lying strictly between minus and plus.  Proper initial
segment is the simplicity relation \(s\ps t\).

Statements about a fixed formula on a definable class are translated into
set theory by restricting its quantifiers to that class.  For example,
\(\Somega\models\varphi\) means the corresponding translation of this
particular formula into the ambient universe.  It does not invoke a
set-theoretic truth predicate that accepts arbitrary formula codes for
\(V\).

When discussing all definitions at once, we may instead fix a set model
\(M\models\mathrm{ZFC}\) in an external metatheory.  Then its satisfaction
relation is an ordinary set externally, and its surreal class gives an
external set structure \(\No^M\).  For absoluteness results we explicitly
require transitivity.  The interpretation theorem itself also has an
internal reading in nonstandard models: subsets, well-foundedness, and
Mostowski collapse are then those of the model.  No assumption that ZFC
proves the existence of a transitive set model is made.

No argument forms \(\mathcal P(\No)\), sums over all ordinals, or treats
all class functions as a set.  A quotient interpretation will be used
through a domain formula and an equivalence-relation formula; its
equivalence classes need not be sets.

\subsection{Normal forms, monomials, and omnific integers}

We use the classical Conway--Gonshor normal-form theorem
\cite{Conway,Gonshor,PK}.  Every surreal has a unique expression
\begin{equation}\label{dsn:op:eq:nf}
 x=\sum_{\xi<\lambda}r_\xi\omega^{a_\xi},
 \qquad r_\xi\in\R\setminus\{0\},\qquad
 a_\xi>a_\eta\quad(\xi<\eta<\lambda),
\end{equation}
with a set ordinal \(\lambda\).  Conversely every such family has a
surreal value.  The support is the reverse-well-ordered set
\(\supp(x)=\{a_\xi:\xi<\lambda\}\).  Here and throughout,
``reverse well ordered'' means that every nonempty subset has a greatest
element in the numerical order.  The zero series has empty support.

The map \(\Omega(a)=\omega^a\) is an increasing isomorphism from the
additive group of \(\No\) onto the multiplicative group of Conway
monomials.  In particular,
\begin{equation}\label{dsn:op:eq:omega-gap}
 a<b\quad\Longrightarrow\quad
 n\Omega(a)<\Omega(b)\quad\text{for every }n\in\N.
\end{equation}
Its notation must not be confused with the Gonshor exponential
\(\exp\).  In general \(\Omega(x)\ne\exp(x\log\omega)\).
The distinction between omega maps and ordinary exponentials is also
central in \cite{BKMMe}.

Write \(x=x_{>0}+r+x_{<0}\) by restricting its normal form to positive,
zero, and negative exponents.  The middle term is an ordinary real,
possibly zero.  The negative part is infinitesimal: its absolute value
is less than every positive real.  Define
\begin{equation}\label{dsn:op:eq:oz}
 \Pi=\{x:\supp(x)\subseteq\No_{>0}\},\qquad
 \Oz=\Pi\oplus\Z.
\end{equation}
These are the purely infinite additive subgroup and the omnific integer
ring, respectively.  A negative element of \(\Pi\) is allowed;
``positive support'' describes exponents, not the sign of the number.

The ring \(\Oz\) is discretely ordered, with least positive element
\(1\).  If an omnific integer has a nonempty positive support, its
leading term has absolute value larger than every ordinary integer;
otherwise it is an ordinary integer.  This proves the discreteness
directly.  The coefficients at positive exponents can be arbitrary reals:
for example \(\pi\omega\in\Oz\), while \(\pi\notin\Oz\).

\begin{proposition}[The omnific floor]\label[proposition]{dsn:op:prop:floor-background}
For every \(x\in\No\) there is a unique \(z\in\Oz\) such that
\(z\le x<z+1\).  If \(x=p+r+\eps\) is the preceding decomposition, then
\[
 \lfloor x\rfloor_{\Oz}=
 p+\begin{cases}
 \lfloor r\rfloor-1,&r\in\Z\text{ and }\eps<0,\\
 \lfloor r\rfloor,&\text{otherwise}.
 \end{cases}
\]
\end{proposition}
\begin{proof}
When \(r\notin\Z\), both distances to its adjacent integers are positive
reals, so adding an infinitesimal changes neither inequality.  When
\(r\in\Z\), a negative infinitesimal puts the number just below \(p+r\),
requiring the displayed correction.  The remaining cases are immediate.
Two different answers would be omnific integers at distance less than
\(1\), contradicting discreteness.
\end{proof}
For instance \(\lfloor\omega-\omega^{-1}\rfloor_{\Oz}=\omega-1\).
Simply discarding negative exponents gives the wrong answer in this case.
This floor formula and the surrounding facts are already developed in
(Part~\ref{dsn:part:one}; \cite{RepoOdg}).

\section{A precise hierarchy of definable numbers}
\label{dsn:op:sec:hierarchy}

\subsection{Ambient parameters and language parameters}

Fix externally a model \(M\) and an allowed collection \(A\subseteq M\).
An element \(x\in M\) is \emph{ambiently definable from \(A\)} if a
standard finite first-order formula \(\varphi(v,\bar a)\), using a finite
tuple \(\bar a\) from \(A\), uniquely defines \(x\) in \((M,\in)\).
Write
\[
 D_A^M=\Def^M(A),\qquad K_A^M=D_A^M\cap\No^M,
 \qquad I_A^M=D_A^M\cap\Oz^M.
\]
The entire collection \(A\) is not itself a parameter unless explicitly
allowed.  In particular, permitting every real as an individual parameter
differs from naming the single set \(\R^M\).

By contrast, for a specified structure \(\mathfrak S\) on surreals,
\(\dcl_{\mathfrak S}(A)\) consists of unique solutions of formulas in
the \emph{language of that structure}.  Such a formula has no access to
arbitrary membership questions about a surreal's sign representation
unless the structure can interpret them.  Every intrinsic definition in
one of the canonical structures considered here translates to an ambient
definition, with the same parameters.  The converse can fail radically.

\begin{center}
\begin{tabular}{P{0.25\textwidth}P{0.26\textwidth}P{0.39\textwidth}}
\toprule
Permission & Notation & Meaning \\
\midrule
No parameters & \(D_\varnothing\) & Unique ambient definition using no named input.\\
A fixed real \(r\) & \(D_{\{r\}}\) & That real is available; arbitrary ordinals are not.\\
Arbitrary individual reals & \(D_\R\) & A finite real tuple, equivalently one real, may vary with the object.\\
Ordinal parameters & \(\OD\) & A finite ordinal tuple may vary with the object.\\
One real and ordinals & \(\OD(r)\) & The fixed real and any finite ordinal tuple.\\
Reals and ordinals & \(\ROD\) & Some real and some finite ordinal tuple.\\
The single set \(\R\) & \(D_{\{\R\}}\) & Equals \(D_\varnothing\) in the ambient language of sets.\\
\bottomrule
\end{tabular}
\end{center}

\subsection{Why ordinal definability is internally expressible}

For a fixed set of ordinals \(p\), define \(x\in\OD(p)\) by the existence
of a rank segment \(V_\theta\), a formula code, and finitely many ordinal
parameters such that \(x\) is uniquely defined there using \(p\) and those
parameters.  Satisfaction for a \emph{set} structure is definable.  The
reflection theorem identifies this with the usual external notion of
definability from \(p\) and ordinals.  The extra rank \(\theta\) can be
included in the ordinal parameter tuple.  The same argument cannot
silently be used to define parameter-free definability: adding \(\theta\)
would change the permitted parameters.  See \cite{DoraisHamkins,HLR}
for this distinction.

Set
\[
 \HOD(p)=\{x:\TC(\{x\})\subseteq\OD(p)\},\qquad
 \HOD=\HOD(\varnothing).
\]
For a set of ordinals \(p\), this is a transitive inner model of ZFC and
has its usual canonical well-order definable from \(p\).  These statements
are about \(\HOD(p)^V\).  They do not replace it by
\(\HOD^{\HOD(p)}\), which is a separate computation.

\begin{proposition}[Hereditary definability of sign codes]
\label[proposition]{dsn:op:prop:sign-hod}
For a surreal in its canonical sign representation and a set of ordinals
\(p\),
\[
 x\in\OD(p)\quad\Longleftrightarrow\quad x\in\HOD(p).
\]
Consequently \(\No\cap\OD=\No\cap\HOD\), and the analogous statement
holds for \(\Oz\).
\end{proposition}
\begin{proof}
All hereditary components of the graph of
\(x:\alpha\to\{-,+\}\), except possibly the full graph itself, are
ordinals or fixed finite constructions from ordinals and sign tags.
They are ordinal definable.  Thus definability of \(x\) from \(p\) and
ordinals already gives hereditary such definability.  The reverse
implication is part of the definition of HOD.
\end{proof}
This identification is representation-sensitive in its formulation: an
arbitrary noncanonical game name could contain extraneous nondefinable
sets.  The canonical number and that particular game name need not have
the same definability degree.

\subsection{Elementary closure and the role of whole inputs}

\begin{lemma}[Closure under a uniquely specified operation]
\label[lemma]{dsn:op:lem:unique-closure}
If a class function \(F\) is ambiently definable from a finite tuple
\(\bar b\), and its input tuple \(\bar x\) is definable from \(A\), then
\(F(\bar x)\) is definable from \(A\cup\{\bar b\}\).
\end{lemma}
\begin{proof}
Conjoin definitions of the inputs with the graph definition of \(F\),
and existentially quantify the input variables.  Uniqueness of the
inputs and of the function value proves uniqueness of the output.
Only finitely many formulas and parameters are combined.
\end{proof}

\begin{theorem}[Background structure of ambient definability hulls]
\label[theorem]{dsn:op:thm:ambient-hull}
Whenever the hulls above are interpreted in a fixed transitive model,
\(K_A\) is a real closed field, \(I_A=K_A\cap\Oz\) is its integer part,
and
\[
 K_A=\Frac(I_A),\qquad K_A\preccurlyeq\No
 \quad\text{in the ordered-field language}.
\]
They are closed under every parameter-free, uniquely defined ambient
operation whose inputs and outputs have the indicated types.  In
particular this includes \(\Omega\), its inverse on monomials, the
Gonshor exponential and logarithm, and the omnific floor.
\end{theorem}
\begin{proof}
Arithmetic and the named canonical functions have parameter-free
set-theoretic definitions, so \cref{dsn:op:lem:unique-closure} applies.  The
positive square root is unique.  For an odd-degree polynomial choose its
least root in the real closed surreal field.  Its
finite root set is ordered, so this is a uniquely specified operation.
These two root properties characterize real closedness.  Quantifier
elimination for real closed ordered fields gives elementarity.
The floor provides the integer part.

For completeness, a denominator can be selected without adding a fresh
ordinal parameter.  If \(S=\supp(x)\), let \(\delta(x)\) be the least
ordinal such that \(a+\delta(x)>0\) for every \(a\in S\), where addition
is surreal addition.  Such an ordinal exists because \(S\) is a set and
the ordinal surreals are cofinal in \(\No\).  More explicitly,
\(|a|\le\bd(a)\), and the ordinal supremum of the values
\(\bd(a)+1\) suffices.  For empty support take \(\delta(0)=0\).
Both \(q=\omega^{\delta(x)}\) and \(p=xq\) are omnific integers, with
\(q>0\), and both are definable from \(x\).  Thus \(x=p/q\) belongs to
\(\Frac(I_A)\).  The other inclusion follows from field closure.
\end{proof}
These are results of the existing reports (Part~\ref{dsn:part:one}; \cite{RepoOdg}); the
proof is included to specify exactly which closure principles are used.
The denominator is a definable \emph{function of the input}.  No new
permission to use arbitrary ordinal parameters has been introduced.

The phrase ``closed under normal forms'' needs an input discipline.  If
an entire ordinal-indexed family is definable, its canonical sum is
definable.  Knowing that each term separately has some definition does
not supply a single definable family or a definable collector of those
definitions.  The repository already supplies counterexamples to this
mistake (\cref{dsn:sec:uniformity} of Part~\ref{dsn:part:one}).  The intrinsic summation theorem later
in this article therefore includes an explicit code for the entire
family among its inputs.

\begin{proposition}[Elementary arithmetic of the definable integer part]
\label[proposition]{dsn:op:prop:integer-part-algebra}
Let \(\Pi_A=K_A\cap\Pi\).  Then
\[
 I_A=\Pi_A\oplus\Z,
 \qquad nI_A=\Pi_A\oplus n\Z\quad(1\le n\in\N),
 \qquad I_A/nI_A\cong\Z/n\Z.
\]
\end{proposition}
\begin{proof}
Constant-term extraction is a canonical operation, so the zero exponent
coefficient of an element of \(I_A\) is an ordinary integer in the same
hull.  The difference lies in \(\Pi_A\).  Division of a purely infinite
element by a positive integer preserves its support and belongs to the
same hull, so \(n\Pi_A=\Pi_A\).  The formulas and quotient identification
follow.  In particular the small finite quotients do not measure the
transfinite information in an omnific normal form.
\end{proof}

\subsection{A reusable bounded code}

Define parameter-free maps
\begin{equation}\label{dsn:op:eq:tau-q}
 \tau(x)=\frac12\left(1+\frac{x}{\sqrt{1+x^2}}\right),
 \qquad Q(x)=\Omega(\tau(x)).
\end{equation}
The positive square root is understood.  The map \(\tau\) bijects
\(\No\) increasingly with \((0,1)\), with inverse
\[
 \tau^{-1}(t)=\frac{2t-1}{\sqrt{1-(2t-1)^2}}.
\]
Hence \(Q\) bijects \(\No\) with
\begin{equation}\label{dsn:op:eq:bounded-U}
 \Qrange=\{\Omega(a):0<a<1\}\subseteq\Oz\cap(1,\omega).
\end{equation}
Both directions are definable in ambient set theory, and already in
the ordered field expanded by \(\Omega\).  Thus
\begin{equation}\label{dsn:op:eq:degree-preservation}
 \Def(A,x)=\Def(A,Q(x))
\end{equation}
for any ambient allowed parameters, with the corresponding intrinsic
statement in a language containing \(\Omega\).
This bounded code is prior repository work (\cref{dsn:thm:bounded-code} of Part~\ref{dsn:part:one}).
It will serve as a compression step for the new interpretation.

\begin{corollary}[The ordinal-definability test, recalled]
\label[corollary]{dsn:op:cor:od-test}
Over ZFC, each of the following is equivalent to \(V=\HOD\): every
surreal is OD; every omnific integer is OD; every member of \(\Qrange\) is OD.
\end{corollary}
\begin{proof}
The bounded code recovers every surreal from a member of \(\Qrange\).
Every subset of an ordinal has a canonical characteristic sign
sequence.  If all surreals are OD, every subset of every ordinal is in
HOD.  By Choice, any set has a pointed well-founded extensional relation
code on an ordinal; that relation and its distinguished point can be
encoded by a subset of an ordinal.  Mostowski collapse then puts the
original set in HOD.  The converse is immediate.
\end{proof}


\section{Real parameters: individual descriptions and uniform coding}
\label{dsn:op:sec:real-parameter-extension}

Throughout this section the ambient universe satisfies ZFC, and all
definability is in its language of set theory.  Real parameters may
equivalently be taken from $2^\omega$: there are fixed definable coding
maps between ordinary reals and subsets of $\omega$, and finitely many
such codes combine into one.  We keep four notions separate:
\[
 \begin{array}{c|l}
 \mathrm{D}_{\R} & \text{definability from one real, without ordinal parameters},\\
 \OD(r) & \text{definability from the fixed real $r$ and ordinal parameters},\\
 \mathrm{ROD} & \displaystyle\bigcup_{r\in\R}\OD(r),\\
 \HOD_{\R} & \{x:\operatorname{TC}(\{x\})\subseteq\mathrm{ROD}\}.
 \end{array}
\]
The first line is an external semantic notion, or a scheme over a fixed
model.  The last three have first-order definitions by reflection and
satisfaction in rank segments.  Our $\HOD_{\R}$ is the convention often
written $\HOD(\R)$; that notation must not be confused with permitting
only the single set $\R$ as a parameter.  That single parameter adds no
definability here, since the set of reals is itself parameter-free
definable.  Dorais and Hamkins discuss the reflection construction and
the real-only distinction explicitly in \cite[Section 4]{DoraisHamkins}.

For a fixed set of ordinals $p$, write $H_p=\HOD(p)$, with definitions
allowed to use $p$ and ordinal parameters.  This is the usual transitive
inner model of ZFC with its canonical well-order definable from $p$.
The requirement that $p$ be a set of ordinals ensures that it is itself
hereditarily $\OD(p)$.  A real is a special case.  All occurrences of
$H_p$ in what follows are computed in the same ambient universe $V$.

\begin{proposition}[Exactly when arbitrary real parameters add information]
\label[proposition]{dsn:op:prop:real-collapse-od}
The following are equivalent:
\begin{enumerate}
\item every real is ordinal definable;
\item $\mathrm{ROD}=\OD$;
\item $\No\cap\mathrm{ROD}=\No\cap\OD$;
\item $\Oz\cap\mathrm{ROD}=\Oz\cap\OD$.
\end{enumerate}
\end{proposition}
\begin{proof}
If every real is OD, substitute an ordinal definition of the real
parameter into any ROD definition; this proves (1)$\Rightarrow$(2),
which implies (3) and (4).  Condition (3) implies (1), since the
canonical embedding of a real in $\No$ and its inverse are definable.
For (4), use the omnific integer $r\omega$ associated to a real $r$.
It is ROD, and division by the parameter-free definable $\omega$ recovers
$r$.  Thus (4) also implies (1).
\end{proof}

The hereditary sign-code argument gives
\[
 \No\cap\mathrm{ROD}=\No\cap\HOD_{\R}.
\]
Indeed, every element in the transitive closure of a sign sequence other
than the sequence itself is constructed from ordinals and fixed finite
tags.  This assertion concerns hereditary membership and does not assume
that $\HOD_{\R}$ satisfies Choice.

\subsection{A single real suffices when all surreals are ROD}

The next statement strengthens a pointwise permission to choose a new
real for each object into one common real for the entire universe.  It
uses ambient Choice essentially.

\begin{remark}[Global real-parameter consolidation; printed in Part~\ref{dsn:part:rp}]
\label[remark]{dsn:op:thm:real-consolidation}
The manuscript's Theorem~4.2, with its proof and the paragraph after it on
relativization to a set of ordinals $p$, is printed once, in
Part~\ref{dsn:part:rp}: it is the same theorem as source~02's
\cref{dsn:rp:cor:global-single-real}, and its formulation with six clauses,
its proof and the relativization are \cref{dsn:rp:rem:consolidation06}. Over
ZFC it says: $V=\HR$ if and only if every set, or every surreal, or every
omnific integer, or every monomial $\omega^a$ with $0<a<1$ is $\ROD$, if and
only if $V=H_r$ for one real $r$. Its proof uses ambient Choice essentially:
a well-order of $\R$ is $\OD(r)$ for some real $r$.
\end{remark}

\subsection{A uniform decoder of countable sign sequences}

Fix a definable coding of triples
\[
 (d,E,\varepsilon),\qquad
 d\subseteq\omega,\quad E\subseteq d\times d,\quad
 \varepsilon:d\longrightarrow\{-,+\},
\]
by subsets of $\omega$.  Such a code is \emph{valid} when $E$ is a
strict well-order.  Let $\alpha$ be its order type and let
$\pi:(d,E)\cong(\alpha,\in)$ be the unique order isomorphism.
The decoded surreal is the sign sequence $s$ of length $\alpha$ given by
$s(\pi(n))=\varepsilon(n)$.  Send invalid codes to $0$.  We obtain one
parameter-free function
\begin{equation}\label{dsn:op:eq:real-decoder}
 D:\mathcal P(\omega)\longrightarrow\No_{<\omega_1}.
\end{equation}
Every countable-birthday surreal has a valid code, by taking any
enumeration of its domain and pulling back the ordinal order and its
signs.  Hence $D$ is onto.  In particular:
\begin{remark}[One-real descriptions below $\omega_1$; printed in Part~\ref{dsn:part:rp}]
\label[remark]{dsn:op:cor:countable-birthday-real}
The manuscript's Corollary~4.3 says that every surreal of countable birthday,
in particular every omnific integer of countable birthday, is definable from
one real without ordinal parameters:
$\No_{<\omega_1}\subseteq\mathrm D_{\R}\cap\No\subseteq\mathrm{ROD}\cap\No$.
It is printed once, as source~05's \cref{dsn:ca:thm:countbirthday}; here it
follows at once from the surjectivity of $D$.
\end{remark}

This construction selects a code only existentially.  Reversing its
direction by one definable operation is a genuinely stronger demand.

\begin{theorem}[Uniform real coding detects correctness of $\omega_1$]
\label[theorem]{dsn:op:thm:real-code-section}
Let $p$ be a set of ordinals.  The following are equivalent:
\begin{enumerate}
\item there is an $\OD(p)$ function
\[
 C:\No_{<\omega_1^V}\longrightarrow\mathcal P(\omega)
 \quad\text{such that}\quad D(C(x))=x;
\]
\item there is such a function definable from $p$ alone;
\item $(\omega_1)^{H_p}=(\omega_1)^V$.
\end{enumerate}
Thus allowing ordinal parameters in the coding operation does not
improve the criterion.  When $p$ is omitted, a parameter-free section
exists exactly when HOD computes $\omega_1$ correctly.
\end{theorem}
\begin{proof}
For (1)$\Rightarrow$(3), let $\alpha<\omega_1^V$ and regard $\alpha$ as
its all-plus surreal.  If $\alpha>0$, $C(\alpha)$ must be a valid real
code for a well-order of type $\alpha$, since invalid codes decode to
$0$.  The case $\alpha=0$ is already countable in every inner model.
The code $C(\alpha)$ is $\OD(p)$: a definition of $C$ and the ordinal
parameter $\alpha$ define its value.  As a subset of $\omega$, it
belongs to $H_p$.  Its well-order and order-type collapse are absolute
between $H_p$ and $V$.  Therefore $H_p$ regards $\alpha$ as countable.
Every $\alpha<\omega_1^V$ is consequently below $\omega_1^{H_p}$.
Together with the general inequality
$\omega_1^{H_p}\leq\omega_1^V$, this gives equality.

Assume (3).  For each $\alpha<\omega_1^V$, the inner model $H_p$ has
a real code for an enumeration of $\alpha$.  Select the least such
enumeration under the canonical $H_p$ well-order.  More explicitly,
select a triple $(d_\alpha,E_\alpha,\pi_\alpha)$ with
$(d_\alpha,E_\alpha)\cong(\alpha,\in)$ and $d_\alpha\subseteq\omega$;
it is enough to select the least real coding $(d_\alpha,E_\alpha)$,
as the collapse $\pi_\alpha$ is then unique.  This selection is uniform
in $\alpha$ and definable from $p$.

For an arbitrary input $x$, of birthday $\alpha$, keep
$d_\alpha,E_\alpha$ and add the labels
\[
 \varepsilon_x(n)=x(\pi_\alpha(n)).
\]
The resulting real is $C(x)$.  The input $x$ need not belong to $H_p$:
only the chosen enumeration of its ordinal domain does.  The labels
are added in $V$, by a uniquely defined operation on $p,x$.
By construction $D(C(x))=x$.  This proves (3)$\Rightarrow$(2), and
(2)$\Rightarrow$(1) is immediate.
\end{proof}

\begin{remark}[Formula complexity and absoluteness]
The decoder $D$ can be given a $\Delta_1$ graph over ZFC, in the
Levy hierarchy of set theory.  A valid
code has a set-sized ordinal collapse witnessing its output; failure
of the strict order axioms has a bounded witness, and failure of
well-foundedness has a nonempty-subset witness with no least element.
Both the graph and its complement therefore have existential set
witnesses.  For the encoder, use the standard $\Delta_2(p)$ surjection
from $\Ord$ onto $\OD(p)$, as in
\cite[Section 4]{DoraisHamkins}, and take the least ordinal code
that produces a real well-order of the required type.  The tested
predicate is $\Delta_2(p)$ and leastness adds only bounded ordinal
quantification.  The resulting uniquely determined output has a
$\Sigma_2(p)$ graph; on its total domain, a different uniquely
determined output gives a $\Sigma_2(p)$ definition of the complement.
Under the equivalent hypotheses of the theorem, extending the
function by a fixed real off $\No_{<\omega_1}$ gives a total
$\Delta_2(p)$ operation.  The intended input domain is $\Sigma_1$:
the birthday admits an injection into $\omega$, or into the fixed
cardinal $\kappa$ in the version below.  Its complement is $\Pi_1$,
so adjoining the fixed default off that domain preserves the
$\Sigma_2$ graph bound.  These are upper bounds, not optimality claims.
\end{remark}

\subsection{The corresponding theorem at every infinite cardinal}

The real case is part of an exact family of statements.  Fix an infinite
$V$-cardinal $\kappa$ and use a definable pairing of $\kappa$ to encode
a signed well-order on a subset of $\kappa$ by a member of
$\mathcal P(\kappa)$.  Decoding valid codes and sending invalid ones to
$0$ gives a function definable from $\kappa$,
\[
 D_\kappa:\mathcal P(\kappa)\longrightarrow
            \No_{<(\kappa^+)^V}.
\]
It is onto: the domain of a surreal on the right has cardinality at most
$\kappa$.  The ordinal pairing may be chosen by ordering pairs first by
their maximum and then lexicographically; on an infinite initial
ordinal its order type is $\kappa$.  It is therefore available with the
same interpretation in $H_p$ and $V$.

\begin{theorem}[Uniform cardinal-sized coding]
\label[theorem]{dsn:op:thm:kappa-code-section}
For a set of ordinals $p$ and an infinite $V$-cardinal $\kappa$, the
decoder $D_\kappa$ has an $\OD(p)$ right inverse if and only if
\begin{equation}\label{dsn:op:eq:successor-cardinal-correct}
 (\kappa^+)^{H_p}=(\kappa^+)^V.
\end{equation}
When this holds, a right inverse is definable from $p,\kappa$ alone,
and has a $\Delta_2(p,\kappa)$ total extension.
\end{theorem}
\begin{proof}
The ordinal $\kappa$ is also a cardinal in $H_p$, because a bijection
with a smaller ordinal there would be one in $V$.  Evaluate a proposed
right inverse at each ordinal $0<\alpha<(\kappa^+)^V$.  Its
$\OD(p)$ output is a subset of $\kappa$, hence belongs to $H_p$,
and codes there a well-order of cardinality at most $\kappa$ and
order type $\alpha$.  Thus $\alpha<(\kappa^+)^{H_p}$.  The reverse
inequality between the two successor cardinals always holds, proving
necessity.  For sufficiency, choose in $H_p$ the least $\kappa$-code
of an enumeration of each such $\alpha$, and add the input's signs
just as in Theorem~\ref{dsn:op:thm:real-code-section}.  The complexity argument
is unchanged, with $\kappa$ added as a parameter.
\end{proof}

\begin{corollary}[A uniform family detects all cardinal correctness]
\label[corollary]{dsn:op:cor:all-cardinal-coding}
For a fixed set of ordinals $p$, the following are equivalent:
\begin{enumerate}
\item $H_p$ and $V$ have exactly the same cardinals;
\item for every infinite $V$-cardinal $\kappa$, $D_\kappa$ has an
$\OD(p)$ right inverse;
\item one class operation definable from $p$ assigns to each pair
$(\kappa,x)$, where $\kappa$ is an infinite $V$-cardinal and
$\bd(x)<(\kappa^+)^V$, a $\kappa$-code for $x$.
\end{enumerate}
\end{corollary}
\begin{proof}
Cardinal agreement implies successor-cardinal agreement, and the
least-code construction in Theorem~\ref{dsn:op:thm:kappa-code-section} is one
formula uniform in $\kappa$, proving (1)$\Rightarrow$(3).  Clearly
(3) implies (2).  If (2) holds, every $V$-cardinal has the same successor
in $H_p$ and $V$.  All $V$-cardinals are already $H_p$-cardinals.
For any additional infinite $H_p$-cardinal $\lambda$, take
$\kappa=|\lambda|^V$.  Then $\kappa<\lambda<(\kappa^+)^V$, contradicting
$(\kappa^+)^{H_p}=(\kappa^+)^V$.  The finite cardinals and $\omega$
are absolute, so (1) follows.
\end{proof}

Cardinal correctness is substantially weaker than $V=H_p$.
For example, adjoining a Cohen real to $L$ preserves cardinals while
HOD remains $L$; hence the parameter-free real section exists even
though there are non-OD real surreals.  In the other direction, a
homogeneous collapse making $\omega_1^L$ countable has HOD equal to
$L$ and fails the parameter-free real-section criterion.  These
applications use the ordinary homogeneity argument: an ordinal-definable
set of ordinals in an ordinal-definable homogeneous forcing extension
comes from the ground model.  The constructible sets remain OD.

\begin{corollary}[An exact all-real-parameter test for a coding operation]
There is a ROD right inverse of the real decoder $D$ if and only if
there is a real $r$ such that
\[
 \omega_1^{\HOD(r)}=\omega_1^V.
\]
\end{corollary}
\begin{proof}
A ROD definition uses some single real $r$ and ordinal parameters.
Apply Theorem~\ref{dsn:op:thm:real-code-section} with $p=r$.
\end{proof}

The coding theorem is an operation-level obstruction, stronger than
the observation that a particular set lacks a canonical real code.
Every countable-birthday surreal has a real description, and ZFC
supplies some set-theoretic right inverse of $D$.  The theorem determines
exactly when such a right inverse can be specified with a fixed
permitted parameter budget.  The cardinal-sized version gives the
same answer at every successor cardinal.



\section{Canonical surreal operations are absolute without a same-reals hypothesis}
\label{dsn:op:sec:absolute-extension}

The repository's question about inner-model support fields leaves open the
case in which the inner model has fewer reals than the ambient universe.
The following argument removes that restriction.  Its key point is that
evaluating an \emph{existing} real coefficient never requires enumerating
all reals.  At limit stages, a normal form can be evaluated using only the
sign sequences in one existing, set-sized sequence.

Throughout this section, a surreal is represented canonically by a sign
sequence, equivalently by a pair $(\alpha,A)$ where $\alpha$ is an ordinal
and $A\subseteq\alpha$ records the positions with sign $+$.  Write
$u\preceq_s v$ when $u$ is an initial sign segment of $v$, including
equality, and write $u\prec_s v$ for a proper segment.  ``Simplest'' means
having the least birthday in a nonempty convex class; that element is
unique and is an initial sign segment of every element of the class.

\begin{theorem}[Absoluteness of the canonical structure]
\label[theorem]{dsn:op:thm:all-reals-free-absoluteness}
Let $M\subseteq W$ be transitive models of $\mathsf{ZFC}$.  The models may
be set models in a larger metatheory or transitive inner classes.  For
inner classes the assertions below are understood formula by formula.
They need not have the same reals.  Then:
\begin{enumerate}
\item The canonical surreals of $M$ are exactly the canonical surreal
sign sequences of $W$ that belong to $M$.  Their order, birthdays and
simplicity relation agree in the two models.
\item The inclusion preserves every set cut whose two option sets
belong to $M$, the ordered-field operations, and Conway's omega map
$a\mapsto\omega^a$.
\item If $x\in\No^M$, its complete canonical normal-form graph
computed in $W$ belongs to $M$ and is exactly the graph computed in $M$.
In particular its coefficients lie in $\mathbb R^M$ and its exponents in
$\No^M$.
\item Canonical normal-form evaluation agrees for every normal-form
graph belonging to $M$.  More generally, summability and the sum agree
for every set-indexed family of surreals that belongs to $M$.
\item The canonical real and omnific-integer predicates agree on old
surreals, and so does the omnific floor.  Consequently
\[
 \mathbb R^M=\mathbb R^W\cap\No^M,
 \qquad
 \Oz^M=\Oz^W\cap\No^M.
\]
\end{enumerate}
For these local assertions, equality of the ordinal heights is not
needed.  In the usual inner-model application both models contain the
same ordinals.
\end{theorem}

\begin{lemma}[A bounded test for the simplest separator]
\label[lemma]{dsn:op:lem:bounded-cut-test}
Let $L,R$ be sets of canonical surreal sign sequences with $L<R$.  A
surreal $z$ equals $\{L\mid R\}$ if and only if
\[
 L<z<R
 \quad\hbox{and}\quad
 \text{no proper sign prefix of $z$ lies strictly between $L$ and $R$}.
\]
The test uses only comparisons with the members of $L,R$ and the
ordinal-indexed prefixes of $z$.
\end{lemma}
\begin{proof}
The forward implication follows from simplicity.  Conversely suppose
$z$ satisfies the displayed test and that $y$ is another separator with
a smaller birthday.  If $y$ is a prefix of $z$, it is a forbidden proper
prefix.  Otherwise $y,z$ first differ at some position $\gamma$; their
common prefix of length $\gamma$ lies strictly between them in the
surreal order.  Since the class of separators is convex, that common
prefix is again a separator, also a contradiction.  Thus $z$ has the
least birthday.  The same common-prefix argument rules out two distinct
least-birthday separators.
\end{proof}

\begin{proof}[Proof of Theorem~\ref{dsn:op:thm:all-reals-free-absoluteness},
parts (1) and (2)]
Being a function from an ordinal to the two fixed sign symbols is
absolute for transitive models.  The order is determined by the first
differing sign, with an absent sign placed between $-$ and $+$, so it is
absolute as well.  Prefixes of an old sign sequence, and the sequence of
all its prefixes, belong to $M$ by Separation and Replacement.

If $L,R\in M$ and $M$ computes $z=\{L\mid R\}$, every test in
Lemma~\ref{dsn:op:lem:bounded-cut-test} is absolute.  Hence $z$ is also the
simplest separator in $W$.  This argument does not assume that all
surreals with birthday below that of $z$ belong to $M$.

For addition, let $x_\beta=x\restriction\beta$ for
$\beta\leq\bd(x)$, and similarly define $y_\gamma$.
All additions needed to compute $x+y$ occur in the set-sized table
\[
 T(\beta,\gamma)=x_\beta+y_\gamma.
\]
The recursive definition of each table entry uses entries with at least
one strictly smaller coordinate and the other coordinate unchanged.
The coordinatewise predecessor relation is well founded; for example
the natural ordinal sum $\beta\mathbin{\oplus}\gamma$ strictly decreases
along it.  Corresponding entries agree by well-founded induction, since
the cut at each step agrees.  Negation is just sign reversal.

The multiplication table on the same prefix pairs is treated by the
same induction.  Each of its cut options is a finite sum or difference
of products with smaller prefix pairs.  Addition has already been
shown absolute, so the options and their simplest separator agree.
Reciprocals then agree by uniqueness of $y$ with $xy=1$; existence is
supplied in $M$ by its field theorem.  This avoids any assumption about
an external class recursion or about the correctness of a choice of
noncanonical game representatives.

For the omega map use its usual uniform cut in the form
\[
 \omega^a=
 \left\{\,0,\ n\omega^{a^L}
 \ \middle|\ 2^{-n}\omega^{a^R}\,\right\},
 \qquad n=1,2,3,\ldots,
\]
where $a^L,a^R$ range over the canonical left and right prefixes of
$a$.  Integer multiples are cofinal among positive real multiples on
the left, and the positive dyadic multiples are coinitial on the right;
the displayed cut therefore gives the standard omega map.  Its entire
dependency table is indexed by the prefixes of $a$.  Induction on
their lengths, using the arithmetic and cut results, proves agreement.
There is no quantification over a possibly different set of positive
reals in this recursion.
\end{proof}

\subsection{The limit step for complete normal forms}

Here is the exact sign-limit operation needed in the proof.  If
$f=\langle x_\xi:\xi<\lambda\rangle$ is an ordinal-indexed sequence and
$\lambda$ is a nonzero limit ordinal, put
\[
 \operatorname{slim}_{\xi<\lambda}x_\xi
 =\bigcup\left\{u:\begin{array}{l}
       u\preceq_s x_\xi\text{ for some }\xi<\lambda,\\[-2pt]
       (\exists\eta<\lambda)(\forall\xi\in[\eta,\lambda))
                    \,u\preceq_s x_\xi
                    \end{array}\right\}.
\]
Two eventually persistent prefixes are comparable, since a sufficiently
late member of the sequence extends both.  The union is therefore a
sign sequence.  Its length is at most
$\sup_{\xi<\lambda}\bd(x_\xi)$, and its construction uses
only a subset of the set of prefixes of the members of $f$.

Equivalently a sign at position $\gamma$ survives if there is a
\emph{single} tail of the sequence on which the entire prefix through
$\gamma$ is defined and constant.  Coordinatewise eventual agreement,
with different tails and no common tail for a whole transfinite prefix,
is insufficient.  This distinction is essential at infinite positions.
For a sequence of successor length its sign limit is its last member;
for the empty sequence set its sign limit equal to $0$.

\begin{lemma}[Absoluteness of the sign limit]
\label[lemma]{dsn:op:lem:slim-absolute}
If the sequence $f$ belongs to $M$, its sign limit computed in $M$ and
in $W$ is the same sign sequence.
\end{lemma}
\begin{proof}
The collection of all prefixes of members of $f$ is the same set in
both models.  Eventual persistence of any such prefix is tested by
quantifiers bounded by $\lambda$ and the existing sign graphs.  The
selected collection and its union thus agree.
\end{proof}

We use the following standard normal-form fact in a particularly useful
form.  Lipparini and Mezo state it as Corollary~3.2 of
\emph{Surreal limits}, citing Conway's normal-form theory and Gonshor's
Theorem~5.12 \cite{LippariniMezo}.  If
$\langle(r_\xi,a_\xi):\xi<\lambda\rangle$ is a legal Conway normal
form, evaluate it by
\begin{align*}
 S_0&=0,\\
 S_{\xi+1}&=S_\xi+r_\xi\omega^{a_\xi},\\
 S_\delta&=\operatorname{slim}_{\xi<\delta}S_\xi
               &&(\delta\text{ a nonzero limit}).
\end{align*}
Then $S_\lambda$ is its canonical Conway sum.  The fact applies to
arbitrary set ordinal lengths and to arbitrary real coefficients,
provided the exponents occur in strictly decreasing order and zero
coefficients have been omitted.  It does \emph{not} assert that sign
limits or sign sums agree with arbitrary analytic limits or arbitrary
rearrangements of series.
\begin{writenote}[A check of the imported normal-form fact, 4 October 2026]
The fact is also a consequence of Gonshor's sign-expansion theorem
\cite[Theorems~5.11--5.12]{Gonshor}, in the formulation of Bournez and
Guilmant \cite[Definition~2.20 and Theorem~2.21]{BG}, which the write
consulted (it did not consult Gonshor's book itself). That theorem gives the
sign expansion of a legal normal form $\sum_{\xi<\lambda}r_\xi\omega^{a_\xi}$
as the juxtaposition, in normal-form order, of one block for each term: the
sign expansion of $\omega^{a'_\xi}r_\xi$, where the reduced exponent $a'_\xi$
is obtained from $a_\xi$ by deleting the minus signs already treated by the
earlier exponents $a_\eta$ $(\eta<\xi)$ and, at a successor $\xi$, one more
minus when $r_{\xi-1}$ is not dyadic. The block of the $\xi$th term thus
depends only on the terms of index at most $\xi$. Hence for $\xi\le\eta$ the
sign expansion of the partial sum $S_\xi$ is an initial segment of that of
$S_\eta$, and at a limit $\delta$ the sign expansion of $S_\delta$ is the union
of those of the $S_\xi$, $\xi<\delta$. Along a $\preceq_s$-increasing sequence
a prefix is eventually persistent exactly when it is a prefix of some member,
so that union is the sign limit $\operatorname{slim}_{\xi<\delta}S_\xi$, as
the displayed recursion requires. The large-cardinal report's proof of
\lbl{lce:lem:absolute} rests on the same sign-expansion theorem of Gonshor's
Chapter~5. Thus the dependence of parts (3)--(4) of
\cref{dsn:op:thm:all-reals-free-absoluteness}, and of item~(4) of
\cref{dsn:op:thm:canonical-delta-one}, on the Lipparini--Mezo corollary
reduces to Gonshor's classical theorem.
\end{writenote}


\begin{proof}[Proof of Theorem~\ref{dsn:op:thm:all-reals-free-absoluteness},
parts (3) and (4)]
An old canonical real remains a real in the outer model: equivalently,
its rational Dedekind cut is unchanged.  In sign notation a canonical
real has finite length, or has length $\omega$ with both signs occurring
arbitrarily far along the sequence.  This characterization is absolute.

Let
$N=\langle(r_\xi,a_\xi):\xi<\lambda\rangle\in M$ be a normal-form
graph computed in $M$.  It is still a legal normal-form graph in $W$.
Its coefficients are the same nonzero canonical reals; its exponents
are the same surreals in the same decreasing order; and its ordinal
indexing is unchanged.  The ordinal indexing itself witnesses the
well-order type of the support and cannot be destroyed by passing to
$W$.

Construct the displayed evaluation table $\langle S_\xi:
\xi\leq\lambda\rangle$ in each model.  The tables agree by transfinite
induction: successor steps use the already absolute field operations
and omega map, while limit steps use
Lemma~\ref{dsn:op:lem:slim-absolute}.  The normal-form fact identifies the last
entry with the canonical sum in both models.  This proves the
evaluation assertion without presupposing the absoluteness of normal
forms.

Now take $x\in\No^M$ and set $N=\operatorname{NF}^M(x)$.
The graph $N$ belongs to $M$ by the internal normal-form theorem.  What
we have just proved says that its canonical value in $W$ is also $x$.
Uniqueness of normal form in $W$ yields
\[
 \operatorname{NF}^W(x)=N=\operatorname{NF}^M(x).
\]
This step proves that $\operatorname{NF}^W(x)$ belongs to $M$; it does
not assume that fact.

For the more general assertion let $f:I\to\No$ belong to $M$.
Its family of support graphs and the union $B$ of their monomials are
sets in $M$.  Summability says that $B$ is reverse well ordered and
that each monomial occurs in only finitely many members of the family.
The latter condition is absolute for transitive models.  If $M$ judges
$B$ reverse well ordered, its order-isomorphism to an ordinal belongs
to $M$ and witnesses the same assertion in $W$.  The converse follows
by restricting the well-order assertion to subsets belonging to $M$.
Thus summability agrees.  The coefficients of the sum are finite sums
of old coefficients, hence agree.  The resulting normal-form graph
belongs to $M$, and the evaluation assertion proves agreement of its
sum.
\end{proof}

\begin{proof}[Proof of Theorem~\ref{dsn:op:thm:all-reals-free-absoluteness},
part (5)]
The real assertion was established above.  A normal form represents an
omnific integer exactly when all its exponents are nonnegative and its
coefficient at exponent zero, if present, is an ordinary integer.
These conditions are absolute once the normal form is.  If $M$
computes $z=\lfloor x\rfloor_{\Oz}$, then in $W$ the same $z$
is an omnific integer and satisfies $z\leq x<z+1$.  Uniqueness of the
integer part in $W$ proves that the floors agree.
\end{proof}

\begin{corollary}[Support fields from arbitrary inner models]
\label[corollary]{dsn:op:cor:inner-support-fields}
For every transitive inner model $M$ of $\mathsf{ZFC}$, the class
$K=\No^M$ is an initial real closed subfield of
$\No^V$.  It contains all ordinals, is closed under the omega
map and its inverse on monomials, and contains every coefficient,
exponent and truncation of the normal form of each of its elements.
Its normal forms agree with the ambient normal forms.  All sums of
summable families \emph{coded in $M$} are preserved.

Thus the hypotheses in the repository's support-field construction
hold for $\No^M$ without the additional assumption
$\mathbb R^M=\mathbb R^V$.  In particular they apply to $L$ or to
$\mathrm{HOD}$ even when those models omit reals.
\end{corollary}
\begin{proof}
Initiality follows because all prefixes of an element of $M$ belong to
$M$.  The internal real-closed field theorem supplies roots of every
finite old polynomial; its assertions remain true externally since
the arithmetic and order agree.  The omega-map assertion is the
theorem.  If an old element is an ambient monomial, its absolute normal
form has a single term with coefficient one; its exponent is therefore
old, which proves closure under the inverse omega map.  The entire
normal form belongs to $M$, hence so do its entries and its
ordinal-indexed truncations and their values.
\end{proof}
\begin{writenote}[Effect on Part~\ref{dsn:part:one}, 4 October 2026]
This corollary answers the instance of \cref{dsn:q:support} that
\cref{dsn:rem:largecardinal} and \lbl{lce:rem:dsnsupport} leave open: for
every transitive inner model $M\ne V$ of ZFC, including $L$ and $\HOD$ when
they lack reals, $\No^M$ satisfies \eqref{dsn:eq:Khyp} and is proper (a
transitive inner model containing every set of ordinals is $V$), so
\cref{dsn:thm:amplify} applies to it; \cref{dsn:op:cor:inner-independent}
carries this out. The classification asked in \cref{dsn:q:support} stays
open (\cref{dsn:op:q:support}). Part~\ref{dsn:part:rp}'s
\cref{dsn:rp:rem:support} adds the field $\KR$, which need not be of the form
$\No^M$.
\end{writenote}


\begin{corollary}[Ordered-field elementarity]
\label[corollary]{dsn:op:cor:inner-rcf-elementarity}
The inclusion $\No^M\subseteq\No^W$ is elementary in
the language of ordered fields.  For class-sized fields this is a
scheme, one assertion for each standard first-order formula.
\end{corollary}
\begin{proof}
Both structures are real closed fields and the inclusion is an
ordered-field embedding.  Quantifier elimination for real closed
fields reduces each formula to a quantifier-free formula, whose truth
is preserved by the inclusion.  No assertion of elementary inclusion
in the language of set theory is involved.
\end{proof}

\subsection{A Levy-hierarchy refinement}

Fix standard effective set codes for ordinals, sign sequences, ordered
pairs and finite syntax.  A class relation is called $\Delta_1$ here
when, over $\mathsf{ZFC}$, it has equivalent $\Sigma_1$ and $\Pi_1$
definitions on its stated domain.  Using the sign-pair code
$(\alpha,A)$, comparison, restriction tests and the test in
Lemma~\ref{dsn:op:lem:bounded-cut-test} are bounded formulas in the standard
rudimentary coding expansion.  This is sufficient for the
$\Delta_1$ conclusions in the pure membership language as well.

\begin{lemma}[Set-sized computation certificates]
\label[lemma]{dsn:op:lem:computation-certificates}
Suppose a total, uniquely specified operation is evaluated on a
set-sized well-founded dependency diagram.  Assume its dependency
diagram is specified by bounded conditions from its inputs and that
each local transition has a $\Sigma_1$ graph.  If the required
evaluation tables exist in $\mathsf{ZFC}$, then its graph is
$\Sigma_1$.  Totality and uniqueness make its graph $\Delta_1$ on
the stated domain.
\end{lemma}
\begin{proof}
A certificate consists of the dependency diagram, an ordinal-valued
rank function when one is needed to verify its well-foundedness, the
table of values, and certificates for the local transitions.  For the
prefix and ordinal-indexed diagrams used here their well-foundedness
already follows from their specified ordinal coordinates.  Write the latter graphs
as $\exists w\,\vartheta$, with $\vartheta$ bounded.  There are only
set many local transitions.  Collection supplies one set of their
witnesses, so all local checks become bounded checks on a single
existentially quantified certificate.  The output graph therefore has
the form $\exists C\,\Psi(C,\text{input},\text{output})$, with
$\Psi$ bounded.  On an input in the domain, failure of the graph is
equivalent to existence of a certificate for a different output,
again a $\Sigma_1$ assertion.  This gives the complementary
$\Pi_1$ definition.

One must actually specify a set-sized dependency diagram.  Merely
saying that an operation is ``defined by a class recursion'' does not
establish the certificate claim.
\end{proof}

\begin{theorem}[$\Delta_1$ graphs of canonical operations]
\label[theorem]{dsn:op:thm:canonical-delta-one}
On their natural domains, the graphs of the following canonical
operations are $\Delta_1$ over $\mathsf{ZFC}$:
\begin{enumerate}
\item simplest set cuts, negation, addition, multiplication,
nonzero reciprocals and nonnegative square roots;
\item the omega map, and its inverse graph on monomials;
\item the sign limit of an ordinal-indexed set sequence;
\item evaluation of a legal canonical normal-form graph and the map
$x\mapsto\operatorname{NF}(x)$;
\item the omnific-integer predicate and the omnific floor.
\end{enumerate}
\end{theorem}
\begin{proof}
For simplest cuts use the bounded prefix test.  The addition and
multiplication dependency diagrams are precisely the prefix-pair
tables described above.  Every use of an already constructed
arithmetic operation is accompanied by its certificate, and
Collection packages the set many certificates into one.  This proves
the $\Sigma_1$ assertions for the arithmetic graphs without a hidden
proper-class object.  Their standard existence and uniqueness
theorems give $\Delta_1$ graphs.  Alternatively, the reciprocal graph
is the already $\Delta_1$ relation $xy=1$, restricted to $x\ne0$.
The nonnegative square-root graph is the $\Delta_1$ relation
$y\geq0$ and $y^2=x$, restricted to $x\geq0$; its existence and
uniqueness follow from real closedness.

The omega-map dependency diagram is the sequence of prefixes of its
input.  At a stage $\beta$, the option sets have an explicit indexing
by $\omega$ times a subset of $\beta$, and their values use earlier
omega values and the previously certified field operations.  The
cut test is bounded.  The certificate lemma applies.  Transposing the
graph gives the inverse graph; after the normal-form result below,
its monomial domain is also $\Delta_1$.

For the sign limit, take the set of prefixes of the members of the
input sequence and apply the bounded persistence test.  A certificate
for that prefix set and its union is a set-sized certificate of the
required form.  Equivalently one can directly test that every proper
prefix through a position of the proposed output is eventually
persistent, and that the output extends every eventually persistent
prefix of an input member.  These are bounded tests on the input
sequence and its sign graphs.

Legality of a normal-form graph is a bounded condition in the sign-code
expansion: its index is an ordinal, coefficients are nonzero canonical
reals, and its exponents strictly decrease.  The sign characterization
of a canonical real uses only $\omega$ and its sign graph.  The
evaluation table has domain $\lambda+1$.  At successor stages it
uses certified arithmetic and omega values, and at limit stages it
uses the sign-limit test.  The certificate lemma therefore proves
that evaluation has a $\Delta_1$ graph.

Finally the assertion that $N=\operatorname{NF}(x)$ says that $N$ is
legal and its certified value is $x$.  This is $\Sigma_1$; the
existence and uniqueness of the canonical normal form make it
$\Pi_1$ as well.  Testing nonnegative exponents and an integral
constant coefficient in this graph gives the $\Delta_1$
omnific-integer predicate.  Its complement is verified using the
same unique graph, so no unbounded search through all normal forms
is needed.  The floor graph is then the $\Delta_1$ assertion
$z\in\Oz$ and $z\leq x<z+1$.
\end{proof}

\begin{remark}[What the certificate theorem does not say]
The theorem is a complexity classification after a fixed standard
coding and theory are chosen.  It is not an optimal symbol-count or
proof-length bound, and it does not claim that the operations are
computable on finite descriptions of arbitrary surreal inputs.
The normal-form theorem is an explicit mathematical dependency, not
something silently replaced by an algorithm for all class-sized
objects.
\end{remark}

\begin{proposition}[Transfer of definition complexity]
\label[proposition]{dsn:op:prop:definition-complexity-transfer}
Let $n\geq1$ be a fixed standard integer.
\begin{enumerate}
\item A total operation with $\Delta_1$ graph sends a uniquely
$\Sigma_n$-definable input tuple to a uniquely
$\Sigma_n$-definable output, with the same permitted parameters.
It also preserves uniquely $\Delta_n$-definable objects.
\item Let $E:X\to Y$ be a fixed injective coding whose graph,
domain, range and inverse graph are $\Delta_1$.  Then $x$ and
$E(x)$ have the same $\Sigma_n$-, $\Pi_n$- and
$\Delta_n$-singleton definability status, with the same parameters.
\end{enumerate}
\end{proposition}
\begin{proof}
For (1), if $\varphi(u)$ is a unique $\Sigma_n$ definition of the
input, the output is defined by
$\exists u\,[\varphi(u)\wedge F(u)=y]$.  Substitution of a
$\Delta_1$ graph preserves the $\Sigma_n$ class.  When the input is
$\Delta_n$, the complementary output predicate is given by
$\exists u\,[\varphi(u)\wedge F(u)\ne y]$ and is also
$\Sigma_n$, by uniqueness and totality.

For (2), a $\Pi_n$ definition $\varphi(x)$ transfers to
\[
 y\in\operatorname{ran}(E)\ \wedge\
 \forall x\,[E(x)=y\ \longrightarrow\ \varphi(x)].
\]
Use the $\Sigma_1$ version of the graph inside the antecedent; its
negation is $\Pi_1$, so this remains $\Pi_n$.  The
$\Sigma_n$ transfer uses the existential form, and the inverse
gives the reverse implications.  The $\Delta_n$ claim follows.
\end{proof}

The injectivity and inverse/range conditions matter in part (2).
The image of an arbitrary uniquely $\Pi_n$-definable input under a
total $\Delta_1$ operation is not asserted to retain a $\Pi_n$
definition.  This is the appropriate distinction when applying the
result to a reversible omnific code instead of to an arbitrary
definable operation.

\subsection{Limits of the absoluteness theorem}

The theorem concerns the canonical operations just specified.  A
parameter-free ambient set-theoretic definition need not be absolute:
for example, the constant surreal-valued function that is zero if
the continuum hypothesis holds and one otherwise can change in a
forcing extension.  Similarly, an old set of surreals whose members
each belong to $M$, but whose \emph{whole family} is not a set in
$M$, is outside the summation assertion.  No claim is made that
$\No^M$ contains every ambient Hahn sum of old coefficients
and old exponents.  Those are precisely the distinctions needed for
the termwise-versus-uniform-definability phenomenon.


\section{What the choice of language permits}
\label{dsn:op:sec:languages}

\subsection{The ordered field alone}

Let \(\mathfrak F=(\No;+,-,\cdot,<,0,1)\).  Quantifier elimination for
real closed fields gives an exact answer in this language
\cite{Marker}:
\begin{equation}\label{dsn:op:eq:pure-dcl}
 \dcl_{\mathfrak F}(A)
 =\{x\in\No:x\text{ is algebraic over }\Q(A)\}.
\end{equation}
The right side is the relative real closure inside \(\No\).
Indeed, a real algebraic root is singled out by its position in the
ordered finite root set of a polynomial.  Conversely, a formula reduces
to a Boolean combination of polynomial sign conditions over a field
generated by finitely many parameters.  At a point transcendental over
that field none of the finitely many nonzero polynomials vanishes.
All signs remain unchanged on some interval about the point, so that
formula cannot isolate a singleton.

In particular,
\[
 \dcl_{\mathfrak F}(\varnothing)=\R_{\mathrm{alg}},\qquad
 \dcl_{\mathfrak F}(\R)=\R.
\]
Thus allowing every real as an intrinsic ordered-field parameter still
does not define any infinite or nonzero infinitesimal surreal.
The field defines only ordinary integers among its parameter-free
definable omnific integers, because
\(\R_{\mathrm{alg}}\cap\Oz=\Z\).

The same elimination theorem makes every definable subset of the line
a finite union of points and intervals.  Consequently \(\Oz\), an
infinite discrete subclass, is not definable in this reduct with any
finite parameters.  The monomial class is also infinite and discrete:
the interval \((m/2,2m)\) contains no other monomial when \(m\) is a
monomial.  Hence the omega map is not definable there either.  These
arguments apply formula by formula to the class structure.

\subsection{The ordinary exponential and the omnific predicate}

The real exponential field is an elementary substructure of the
Gonshor surreal exponential field; this is a classical result of
van den Dries--Ehrlich, read with its erratum \cite{vdDE,vdDEerr}.
It follows that the parameter-free definable elements of that reduct
are real.  If a formula has a unique solution in the large structure,
elementarity supplies a solution in \(\R\), and uniqueness identifies
them.  This does not calculate which reals are definable in the real
exponential field.  It does show that naming the ordinary exponential
does not parameter-freely define \(\omega\).

The pair \(\Szero=(\mathfrak F,\Oz)\) behaves differently.  The next
section gives short formulas defining its canonical \(\R\), coefficient
at zero, and positive-support projection.  Yet its parameter-free
definable elements are still real.  One elementary witness is the
normal-form dilation
\[
 S_2\left(\sum_a r_a\omega^a\right)=\sum_a r_a\omega^{2a}.
\]
It is a field automorphism preserving \(\Oz\).  Its fixed field is
\(\R\): if a fixed series had a nonzero support exponent \(a\), invariance
under doubling and halving would give an infinite increasing sequence
in its reverse-well-ordered support.  For \(a>0\) use
\(a,2a,4a,\ldots\); for \(a<0\) use
\(a,a/2,a/4,\ldots\).  Thus any uniquely parameter-free definable point,
which must be fixed by \(S_2\), is real.  This observation and much
stronger parameterwise obstructions are in (Part~\ref{dsn:part:one}; \cite{RepoOpa}).

Because \(\R\) and \(\Z=\R\cap\Oz\) are definable in the pair, every
real definable in \((\R;+,-,\cdot,<,\Z)\) is definable there.  The exact
upper bound remains unresolved here.  Fixedness under automorphisms
alone cannot give an exact description of definable closure: even the
ordinary ordered real field is rigid while its parameter-free
definable elements are only the real algebraic numbers.

\subsection{Definability in the omnific ring itself}
\label{dsn:op:sec:native-omnific}

In this subsection work in the canonical universe, or externally in a
fixed transitive model.  There is also a natural intrinsic definition
that uses only the ring of omnific integers:
\[
 \mathfrak I=(\Oz;+,-,\cdot,0,1).
\]
Its definable elements and operations should be distinguished from
ambiently definable omnific integers.  The elementary fraction-field
construction makes the relationship with \(\Szero\) precise.

\begin{proposition}[The native ring and its fraction interpretation]
\label[proposition]{dsn:op:prop:native-omnific}
The numerical order on \(\Oz\) is definable in its pure unital ring
language.  The structure \(\mathfrak I\) and the pair \(\Szero\) are
bi-interpretable using the usual fraction quotient.  Moreover,
\[
 \dcl_{\mathfrak I}(\varnothing)=\Z,\qquad
 \dcl_{\Szero}(\R)=\R.
\]
\end{proposition}
\begin{proof}
The same-denominator construction in \cref{dsn:op:thm:ambient-hull} gives
\(\Frac(\Oz)=\No\).  For \(a,b\in\Oz\), use the formula
\[
 a<b\quad\Longleftrightarrow\quad
 \exists p\,\exists q\,
   \bigl(p\ne0\wedge q\ne0\wedge (b-a)q^2=p^2\bigr).
\]
If \(b-a>0\), express its positive square root in \(\No\) as a
nonzero omnific fraction \(p/q\).  Conversely, the displayed identity
makes \(b-a\) a nonzero square in the ordered fraction field, hence
positive.

Interpret \(\No\) by pairs \((p,q)\in\Oz^2\) with \(q\ne0\),
modulo \(pq'=p'q\).  The field operations and order have their usual
fraction definitions, and the interpreted omnific predicate is
\(\exists z\,(p=zq)\).  In the opposite direction the omnific domain
is already a definable subset of \(\Szero\).  The comparison map on
the ring is \(z\mapsto[(z,1)]\); on the field it is the relation
between \(x\) and the fraction class of \((p,q)\) specified by \(p=xq\).
These are definable bijections onto the relevant quotients and
preserve the operations.  This proves the asserted bi-interpretation.

The dilation \(S_2\) just constructed fixes exactly \(\R\) in \(\No\)
and exactly \(\R\cap\Oz=\Z\) in \(\Oz\).  Therefore every
parameter-free definable element of \(\mathfrak I\) is an ordinary
integer.  Conversely every ordinary integer is named by a finite ring
term using \(0,1\).  Since \(S_2\) fixes every real parameter, the same
argument gives \(\dcl_{\Szero}(\R)\subseteq\R\), and the reverse
inclusion follows from permission to name each real.
\end{proof}

Thus naming the order adds no expressive power to the native ring.
This fraction interpretation is elementary algebraic background,
using the repository's fraction-field and dilation results
(Part~\ref{dsn:part:one}; \cite{RepoOdg,RepoOpa}).  It is a different comparison from
the set-universe interpretation studied below.
\begin{writenote}[Antecedents, 4 October 2026]
The order formula, the fraction bi-interpretation and $\dcl_{\Szero}(\R)=\R$
are largely in the omnific Diophantine report
(\lbl{odg:def:thm:realrecovery}, \lbl{odg:def:cor:arithmetic}, with the same
dilation as \lbl{odg:def:ex:dilation}) and in \cref{dsn:prop:pair}~(i) of
Part~\ref{dsn:part:one}; the inclusion $\dcl_{\Szero}(\R)\subseteq\R$ also
follows from \lbl{opa:sc:cor:innerbound}. The manuscript credits these
reports in general terms.
\end{writenote}


A nonintegral real cannot literally be a parameter in a one-sorted
\(\Oz\)-structure, since it is not in its domain.  It can be used
as a parameter in the interpreted fraction sort, as an equivalence
class.  Naming a chosen numerator and denominator separately can add
information: the pair \((r\omega,\omega)\) represents \(r\) but
also names \(\omega\).  The second definable-closure equality above
shows that allowing real parameters in the fraction sort still does
not define a nonordinary omnific integer.

Finally the native ring has definable quotient and remainder operations.
For every \(a,b\in\Oz\) with \(b>0\), there are unique \(q,r\in\Oz\)
such that
\[
 a=bq+r,\qquad 0\le r<b.
\]
Existence follows by taking \(q=\lfloor a/b\rfloor_{\Oz}\) and
\(r=a-bq\).  For uniqueness, a nonzero difference of two quotients
has absolute value at least one, while the corresponding difference
of remainders has absolute value less than \(b\), a contradiction.
These equations are first-order graphs, using the definable order.
They do not assert termination of an iterated division procedure:
the positive omnific order is not well founded.  For example, starting
with \((a,b)=(\sqrt2\,\omega,\omega)\), the first quotient is \(1\)
and all later quotients are \(2\).  The successive positive remainders
are \((\sqrt2-1)^n\omega\), for \(n=1,2,\ldots\); they never become
zero at any finite step.


\begin{center}
\begin{tabular}{P{0.29\textwidth}P{0.6\textwidth}}
\toprule
Language on \(\No\) & Consequence relevant to definability \\
\midrule
Ordered field & Exact closure \eqref{dsn:op:eq:pure-dcl}; all unary definable sets
are semialgebraic.\\
Ordered field and \(\exp\) & No parameter-free infinite element; exact
real definable closure is not computed here.\\
Ordered field and \(\Oz\) & Defines canonical reals, zero coefficient,
positive truncation and floor; parameter-free definable elements are real.\\
Ordered field and \(\Omega\) & Defines \(\omega=\Omega(1)\), a natural
valuation, and bounded finite-tuple codes.\\
Ordered field, \(\Oz\), \(\Omega\) & Defines all coefficients, coded Hahn
sums and the quotient interpretation of set theory proved below.\\
Ordered field and birthday order & Bi-interpretation with set theory is
the announced Chen--Hamkins--Yang result; a different signature.\\
\bottomrule
\end{tabular}
\end{center}

Naming all ordinals or all reals in one of these intrinsic languages still
means allowing finitely many named elements in each formula.  The ambient
classification in \cref{dsn:op:sec:hierarchy} should not be transferred to a
reduct merely by retaining the same informal word ``definable.''



\section{Definable support calculus}
\label{dsn:op:sec:interpretation}

Write
\[
 \Szero=(\No;0,1,+,-,\cdot,<,\Oz),\qquad
 \Somega=(\Szero,\Omega),\qquad
 \Omega(a)=\omega^a.
\]
The distinction between the two structures is substantial.  The omnific
predicate suffices to recover the real coefficient field and the coefficient
at exponent zero.  Naming the omega map then permits a uniform shift to
every exponent.  This gives the entire support relation, and supports can
carry arbitrary set-sized relations.  We give the intermediate steps
explicitly before drawing the set-theoretic conclusion.

\subsection{Recovering the coefficient at zero}

Let $\Pi$ denote the additive subgroup of purely infinite surreals,
including zero: its elements have normal forms supported on strictly
positive exponents.  Let $\mathcal A$ be the subring of normal forms
supported on nonnegative exponents.  Thus
\[
 \Oz=\mathbb Z+\Pi,
 \qquad \mathcal A=\mathbb R+\Pi.
\]
These equalities refer to the canonical normal form.  The following
formulas recover these objects inside $\Szero$.

\begin{lemma}[A quadratic detector and its multiplier ring]
\label[lemma]{dsn:op:lem:real-recovery}
The classes $\Pi$, $\mathcal A$, and $\mathbb R$ are parameter-free
definable in $\Szero$.  Explicit definitions are
\begin{align}
 \operatorname{PI}(x)
 &\;\longleftrightarrow\;
 x\in\Oz\ \wedge\
 \exists y\in\Oz\ (y^2=2x^2),
 \label{dsn:op:eq:PI}\\
 \operatorname{A}(t)
 &\;\longleftrightarrow\;
 \forall p\,(\operatorname{PI}(p)\longrightarrow
                        \operatorname{PI}(tp)),
 \label{dsn:op:eq:A}\\
 \operatorname{Real}(t)
 &\;\longleftrightarrow\;
 t=0\ \vee\
 \bigl(t\ne0\wedge\operatorname{A}(t)
                     \wedge\operatorname{A}(t^{-1})\bigr).
 \label{dsn:op:eq:Real}
\end{align}
Here an inverse is an ordinary definable field operation, and all displayed
restricted quantifiers abbreviate first-order formulas.
\end{lemma}

\begin{proof}
The constant coefficient on $\Oz$ is a ring homomorphism into
$\mathbb Z$: in a product of two nonnegative-support series, the only
contribution to exponent zero comes from their constant terms.  If
$y^2=2x^2$ with $x,y\in\Oz$, their integer constant terms $n,m$
satisfy $m^2=2n^2$, whence $m=n=0$.  Conversely, if $x\in\Pi$, then
$\sqrt2\,x\in\Pi\subseteq\Oz$ witnesses the formula.  This proves
\eqref{dsn:op:eq:PI}.

Every element of $\mathcal A$ multiplies $\Pi$ into $\Pi$.  Conversely,
suppose the normal form of $t$ contains a nonzero term at an exponent
$g<0$.  Put $p=\Omega(-g/2)\in\Pi$.  Multiplication by this monomial shifts
every exponent by $-g/2$; in particular, $tp$ has a nonzero term at
$g/2<0$.  Thus $tp\notin\Pi$.  No cancellation at this exponent is
possible because translation of exponents is injective.  This proves
\eqref{dsn:op:eq:A}.

Every nonzero real and its inverse belong to $\mathcal A$.  If
$0\ne t\in\mathcal A$ is not real, its largest exponent is strictly
positive.  The largest exponent of $t^{-1}$ is its negative, so
$t^{-1}\notin\mathcal A$.  This proves \eqref{dsn:op:eq:Real}.
\end{proof}

The use of the omega map in the proof of the multiplier identity does not
put it into the defining formula.  It constructs a witness showing that
\eqref{dsn:op:eq:A}, a formula in the smaller language, has the stated
extension.  The quadratic detector and coefficient reconstruction are
already part of the repository's omnific Diophantine work \cite{RepoOdg};
the proof is
included to make the present interpretation independent of a lengthy
imported reconstruction argument.

\begin{proposition}[Global coefficient and positive-part operations]
\label[proposition]{dsn:op:prop:ct}
In $\Szero$, the infinitesimal class $\mathfrak m$, the finite-element
valuation ring $\mathcal O$, and the functions
\[
 x\longmapsto c_0(x),\qquad x\longmapsto x_{>0}
\]
are parameter-free definable.  Here $c_0(x)$ is the real coefficient of
$\omega^0$ in the full Conway normal form, and $x_{>0}$ is its positive
exponent truncation.  The graph of $c_0$ is given by
\[
 \operatorname{Real}(r)\ \wedge\
 \exists p\bigl(\operatorname{PI}(p)
             \wedge x-p-r\in\mathfrak m\bigr),
\]
where
\[
 \epsilon\in\mathfrak m
 \ \longleftrightarrow\
 \forall r\bigl(\operatorname{Real}(r)\wedge r>0
                         \longrightarrow |\epsilon|<r\bigr).
\]
\end{proposition}

\begin{proof}
The formula for $\mathfrak m$ expresses precisely that $\epsilon=0$ or
its largest exponent is negative.  Similarly,
\[
 x\in\mathcal O\ \longleftrightarrow\
 \exists r\bigl(\operatorname{Real}(r)\wedge r>0\wedge |x|<r\bigr).
\]
Split the support of a normal form at exponent zero.  This gives
$x=p+r+\epsilon$ with $p\in\Pi$, $r\in\mathbb R$, and
$\epsilon\in\mathfrak m$.  The splitting is unique because
$\Pi\cap\mathcal O=\{0\}$ and
$\mathbb R\cap\mathfrak m=\{0\}$.  Hence the displayed formula defines
$r=c_0(x)$ uniquely, and the same decomposition defines $p=x_{>0}$.
\end{proof}

This global coefficient projection is not a field homomorphism.  For
example, $c_0(\omega)=c_0(\omega^{-1})=0$ whereas
$c_0(\omega\omega^{-1})=1$.  On the finite valuation ring, its restriction
is the ordinary standard-part map and is a ring homomorphism.

\subsection{Every coefficient and every truncation}

\begin{theorem}[Uniform support calculus]
\label[theorem]{dsn:op:thm:support}
The following operations and relations are parameter-free definable in
$\Somega$, uniformly in the variables shown:
\begin{align*}
 c_a(x)&=c_0\bigl(x/\Omega(a)\bigr),\\
 T_{>a}(x)&=\Omega(a)\bigl(x/\Omega(a)\bigr)_{>0},\\
 T_{\ge a}(x)&=T_{>a}(x)+c_a(x)\Omega(a),\\
 \operatorname{Supp}(a,x)&\;\longleftrightarrow\;c_a(x)\ne0.
\end{align*}
Thus $T_{>a}(x)$ and $T_{\ge a}(x)$ are the canonical normal-form
truncations at any surreal exponent $a$, and
$\operatorname{Supp}(a,x)$ means that $a$ occurs in the Conway support of
$x$.  The complementary truncations are definable by subtraction.
\end{theorem}

\begin{proof}
The identity $\Omega(b)/\Omega(a)=\Omega(b-a)$ shifts all exponents by
$-a$.  The coefficient at zero after this shift is exactly the old
coefficient at $a$, and the positive part after the shift is exactly the
old part above $a$, scaled by $\Omega(-a)$.  Apply
Proposition~\ref{dsn:op:prop:ct} and compose definable functions.
\end{proof}

The theorem is a statement about the richer language.  The absence of a
definable omega map in $\Szero$ prevents transferring these uniform
exponent-indexed formulas to that reduct.  A monomial supplied as a
parameter still allows coefficient extraction at that particular scale.

\begin{remark}[Valuation already from the omega map]
\label[remark]{dsn:op:rem:valuation}
Even without the omnific predicate, $(\No;0,1,+,-,\cdot,<,\Omega)$
defines the leading exponent of a nonzero $x$.  It is the unique $a$ such
that
\[
 \forall b\left[
 (b<a\longrightarrow\Omega(b)<|x|)\ \wedge\
 (a<b\longrightarrow |x|<\Omega(b))\right].
\]
The actual largest exponent satisfies the formula because distinct
monomials have distinct Archimedean classes.  If a different $a$ did so,
choose $b$ strictly between $a$ and the largest exponent to obtain a
contradiction.  Negating the leading exponent gives the natural valuation.
In particular its valuation ring and infinitesimal ideal are definable
already in this language.  This observation does not by itself select the
canonical real residue representatives or recover all coefficients.
\end{remark}

\subsection{A bounded lexicographic pairing operation}

The parameter-free ordered-field map
\[
 \tau(x)=\frac12\left(1+\frac{x}{\sqrt{1+x^2}}\right)
\]
is an increasing bijection from $\No$ onto $(0,1)$, with inverse
\[
 \tau^{-1}(u)=\frac{2u-1}{\sqrt{1-(2u-1)^2}}.
\]
Define
\begin{align}
 B(a,b)&=\Omega\left(\frac{2+\tau(a)}4\right)
                +\Omega\left(\frac{\tau(b)}4\right),
 \label{dsn:op:eq:B}\\
 p(a,b)&=\tau\bigl(B(a,b)\bigr).
 \label{dsn:op:eq:p}
\end{align}

\begin{lemma}[Pairing and preservation of reverse well-orders]
\label[lemma]{dsn:op:lem:pair}
The function $B$ is a parameter-free definable injection
$\No^2\to\Oz\cap(1,\omega)$ in the omega-map expansion.
It is an order embedding for the lexicographic order with first coordinate
primary.  Consequently $p$ is a definable lexicographic order embedding
$\No^2\to(0,1)$.  If $D,E$ are reverse well-ordered sets of
surreals, then $p[D\times E]$, and each of its subsets, is reverse
well-ordered.
\end{lemma}

\begin{proof}
The first exponent in \eqref{dsn:op:eq:B} belongs to $(1/2,3/4)$ and the
second to $(0,1/4)$.  Thus both exponents are positive, strictly separated,
and below one.  Their sum is an omnific integer strictly between one and
$\omega$.  If $a<a'$, the leading monomial in $B(a',b')$ dominates both
terms of $B(a,b)$, whatever $b,b'$ are.  Hence
$B(a,b)<B(a',b')$.  When $a=a'$, comparison reduces to the second term and
therefore to comparison of $b,b'$.  This proves the lexicographic embedding
and in particular injectivity.  Composition with $\tau$ preserves it.

For any nonempty subset of $D\times E$, first choose the greatest first
coordinate appearing in it, and then the greatest second coordinate above
that first coordinate.  These exist by reverse well-ordering and give its
lexicographically greatest pair.  The same assertion therefore holds in
its image under $p$.
\end{proof}

The inverse of $B$ on its image is definable by its graph and uniqueness.
Iterating $B$ consequently compresses any fixed finite tuple to one
omnific integer below $\omega$.  In particular,
\[
 \operatorname{dcl}_{\Somega}(A,a,b)
 =\operatorname{dcl}_{\Somega}(A,B(a,b))
\]
for every allowed parameter set $A$.  This is an intrinsic version of the
repository's bounded coding principle.  No normal-form decoding symbol
needs to be added to the language for the finite-tuple assertion.

\subsection{Powersets of supports and arbitrary relations}

Call a surreal $u$ a \emph{unit-coefficient series} if
\[
 \operatorname{Unit}(u):\qquad
 \forall a\ (c_a(u)=0\ \vee\ c_a(u)=1).
\]
It may be zero.  Its support determines it uniquely.

\begin{lemma}[Exact powerset coding]
\label[lemma]{dsn:op:lem:powerset}
Let $d$ be a unit-coefficient series and put
$D=\{a:\operatorname{Supp}(a,d)\}$.  The first-order predicate
\[
 \operatorname{Sub}(u,d):\qquad
 \operatorname{Unit}(u)\ \wedge\
 \forall a\bigl(\operatorname{Supp}(a,u)
                  \longrightarrow\operatorname{Supp}(a,d)\bigr)
\]
has exactly one solution $u_X$ for each subset $X\subseteq D$, namely
\[
 u_X=\sum_{a\in X}\Omega(a).
\]
If $d,e$ have supports $D,E$, the analogous relation-code predicate
\[
 \begin{split}
 \operatorname{Rel}(q;d,e):\quad&\operatorname{Unit}(q)\ \wedge\\[-1mm]
 &\forall t\left(\operatorname{Supp}(t,q)\longrightarrow
  \exists a\exists b\,[\operatorname{Supp}(a,d)\wedge
  \operatorname{Supp}(b,e)\wedge t=p(a,b)]\right)
 \end{split}
\]
has exactly one solution for each relation $R\subseteq D\times E$:
\[
 q_R=\sum_{(a,b)\in R}\Omega(p(a,b)).
\]
For this solution,
$R(a,b)$ is expressed by $\operatorname{Supp}(p(a,b),q_R)$.
\end{lemma}

\begin{proof}
A subset of a set-sized reverse well-ordered support is again a set-sized
reverse well-ordered support.  Hence every displayed unary series exists
as a Conway normal form.  Its nonzero coefficients are exactly one, and
uniqueness of normal forms gives uniqueness of the code.  Conversely, any
solution of $\operatorname{Sub}$ has just this form for its support.
For the binary assertion use Lemma~\ref{dsn:op:lem:pair}: every subset of
$p[D\times E]$ is reverse well-ordered, so every relation has a legal
normal-form code.  Injectivity of $p$ and normal-form uniqueness give both
decoding and uniqueness.  All relation codes have support in $(0,1)$ and
therefore lie in $\Oz\cap[0,\omega)$.
\end{proof}

This lemma covers \emph{every} ambient subset of $D$ and every relation on
$D\times E$, not just subsets definable from the original parameters.
Definability is used for the membership test and the family of codes.
An individual code need not be parameter-free definable.  This distinction
is what allows the following first-order expression of genuine
well-foundedness.

\section{Interpreting set theory in bounded omnific monomials}
\label{dsn:op:sec:set-interpretation}

The support-subset mechanism has precedents in interpretations of
second-order arithmetic in Hahn fields \cite{Camacho}.  Here it applies
to supports of arbitrary set order type and therefore to arbitrary
set-sized membership graphs.

\subsection{Pointed well-founded extensional graphs}

A candidate set code is a triple $(d,q,r)$ with $d$ a unit-coefficient
series, $D=\operatorname{supp}(d)\subseteq(0,1)$ nonempty,
$\operatorname{Rel}(q;d,d)$, and $r\in D$.  Write
\[
 a\ E_q\ b\quad\longleftrightarrow\quad
                   \operatorname{Supp}(p(a,b),q)
\]
for its membership-oriented edge relation.  The following are all
first-order conditions in $\Somega$.

\begin{enumerate}
\item \emph{Well-foundedness:} for every $u$ with
$\operatorname{Sub}(u,d)$ and $u\ne0$, there is
$a\in\operatorname{supp}(u)$ with no $E_q$-predecessor in
$\operatorname{supp}(u)$.
\item \emph{Extensionality:} if $a,b\in D$ have the same
$E_q$-predecessors in $D$, then $a=b$.
\item \emph{Accessibility from the root:} every $u$ with
$\operatorname{Sub}(u,d)$ whose support contains $r$ and is closed under
$E_q$-predecessors has support equal to $D$.
\end{enumerate}

For instance, the first condition is literally
\[
 \begin{split}
 \forall u\bigl(&\operatorname{Sub}(u,d)\wedge u\ne0
 \longrightarrow\\[-1mm]
 &\exists a[\operatorname{Supp}(a,u)\wedge
  \forall b(\operatorname{Supp}(b,u)\longrightarrow\neg(bE_q a))]
 \bigr).
 \end{split}
\]
There are no quantifiers over sets of surreals in these formulas: $u$ is
one surreal number.  Lemma~\ref{dsn:op:lem:powerset} proves that ranging over
these numbers really ranges over all subsets needed in the displayed
well-foundedness test.

Call such a triple \emph{admissible}.  The set-sized Mostowski collapse
of $(D,E_q)$ sends its root to a unique set, denoted here by
$\operatorname{val}(d,q,r)$.  The collapse function is used in the proof
in the ambient universe, not named in the surreal language.

\begin{theorem}[Set theory in the omega--omnific structure]
\label[theorem]{dsn:op:thm:V-interpretation}
In ZFC, $\Somega$ has a parameter-free first-order quotient
interpretation of $(V,\in)$.  Its representatives may all be taken to be
single monomials in
$\Qrange=\{\Omega(a):0<a<1\}\subseteq\Oz\cap(1,\omega)$.
More precisely, there are fixed
first-order formulas for a class $C$, an equivalence relation $\equiv$ on
$C$, and a relation $\mathrel{E^*}$ invariant under $\equiv$, such that the
Mostowski-collapse valuation gives
\[
 (C/{\equiv},E^*)\ \cong\ (V,\in).
\]
The displayed quotient means a definable quotient presentation; it does
not require proper-class equivalence classes to be sets.
\end{theorem}

\begin{proof}
First use admissible triples as representatives.  The admissibility
conditions were given above and are first-order.  By exact powerset
coding, the asserted well-foundedness is actual well-foundedness of the
set relation $E_q$, so the set-sized Mostowski collapse applies.
Extensionality makes it an isomorphism onto a transitive set $T$.
Accessibility says that $T$ is exactly the transitive closure of the
singleton containing the collapsed root.  Indeed, the nodes reachable
from the root by a finite sequence of predecessor steps form a
predecessor-closed subset containing the root; by accessibility this is
all of $D$.

Two representatives are declared equivalent when there is a root-preserving
isomorphism between their graphs.  This is first-order: for supports
$D,D'$, quantify over a relation code $h$ satisfying
$\operatorname{Rel}(h;d,d')$ and require its decoded relation to be the
graph of a total bijection $D\to D'$ that takes root to root and preserves
edges in both directions.  These requirements are finite first-order
formulas.  The relation-coding lemma supplies every bijection required,
so this condition is equivalent to equality of the collapsed roots.

For membership, declare a representative $A$ to be $E^*$-related to a
representative $B$ if there are a node $b$ in the graph of $B$ and a coded
map $h$ with $\operatorname{Rel}(h;d_A,d_B)$ satisfying the following
conditions: $b$ is an immediate
$E_B$-predecessor of the root of $B$; $h$ is an injective map from the
entire domain of $A$ into that of $B$; it sends the root of $A$ to $b$;
it preserves and reflects edges; and its image is closed under
$E_B$-predecessors.  All these assertions are first-order statements about
a relation code.  The last requirement can be written
\[
 \forall a\in D_A\ \forall v\in D_B\,
 \bigl(vE_B h(a)\longrightarrow
          \exists a'\in D_A\ (h(a')=v)\bigr),
\]
using the graph of $h$ in place of function notation.

Such a map identifies $A$ with the full predecessor closure rooted at
$b$.  To see this, its predecessor-closed image contains $b$, so it
contains every node reachable downwards from $b$.  Conversely,
accessibility of $A$ puts every image node in that closure.  Applying
collapse, its existence is therefore equivalent to
$\operatorname{val}(A)\in\operatorname{val}(B)$.  If this membership holds,
the corresponding isomorphism of transitive closures supplies a map
$h$, and exact relation coding supplies its code.  Equality and
membership are consequently independent of the chosen representatives.

For surjectivity, fix any set $x$.  By Choice, the set
$\TC(\{x\})$ has an enumeration by an ordinal $\alpha$.
Label its elements bijectively by the reverse well-ordered set
\[
 D_\alpha=\left\{\frac1{\underline\beta+2}:\beta<\alpha\right\}
                 \subseteq(0,1),
\]
where $\underline\beta$ is the canonical ordinal surreal.  Use the
unit-coefficient series supported on $D_\alpha$ for $d$, use
Lemma~\ref{dsn:op:lem:powerset} to code the membership relation as $q$, and
take for $r$ the label of $x$.  This is an admissible triple whose
collapsed root is $x$.  The empty set causes no exception: its pointed
transitive closure is the singleton $\{\varnothing\}$.

Finally replace the triple $(d,q,r)$ by the single number
\[
Q\bigl(B(d,B(q,r))\bigr)
 =\Omega\bigl(\tau(B(d,B(q,r)))\bigr).
\]
Lemma~\ref{dsn:op:lem:pair} and \eqref{dsn:op:eq:bounded-U} place it in $\Qrange$
and provide a definable inverse on its image.
Pulling back the preceding domain, equality, and membership
formulas gives the asserted one-coordinate presentation.
\end{proof}

\begin{example}[A code for the ordinal two]
Put \(a_0=1/2\), \(a_1=1/3\), and \(a_2=1/4\), and take
\[
 d=\Omega(a_0)+\Omega(a_1)+\Omega(a_2),\qquad
 q=\Omega(p(a_0,a_1))+\Omega(p(a_0,a_2))+\Omega(p(a_1,a_2)).
\]
The root is \(r=a_2\).  The graph has edges \(a_0E_qa_1\),
\(a_0E_qa_2\), and \(a_1E_qa_2\).  Collapse sends the three nodes to
\(\varnothing\), \(\{\varnothing\}\), and
\(\{\varnothing,\{\varnothing\}\}\), respectively.  Thus its value
is the set-theoretic ordinal \(2\).  The single representative
\(Q(B(d,B(q,a_2)))\) is an omnific monomial between \(1\) and \(\omega\).
Different choices of labels give equivalent representatives; they
need not give the same surreal number.
\end{example}

\begin{corollary}[Uniform translation of set-theoretic formulas]
\label[corollary]{dsn:op:cor:translation}
For each fixed first-order set-theoretic formula
$\varphi(x_1,\ldots,x_n)$ there is an effectively obtained formula
$\varphi^*(z_1,\ldots,z_n)$ in the language of $\Somega$ such
that, for admissible codes,
\[
 \Somega\models\varphi^*(z_1,\ldots,z_n)
 \quad\Longleftrightarrow\quad
 V\models\varphi(\operatorname{val}(z_1),\ldots,
                            \operatorname{val}(z_n)).
\]
\end{corollary}

\begin{proof}
Replace atomic equality and membership by their formulas from
Theorem~\ref{dsn:op:thm:V-interpretation}; preserve Boolean connectives; and
relativize each quantifier to the admissible-code domain.  Induction on
the fixed formula, using surjectivity of the valuation map, proves the
claim.  This is a syntactic translation, not a truth predicate for
arbitrary formulas in $V$ or in the class structure $\Somega$.
\end{proof}

\paragraph{What this resolves.}
The repository asks how much ambient definability can be captured after
adjoining the omega map to the field with the omnific predicate.  The
theorem supplies a strong answer at the level of interpretation: this
expansion can express arbitrary first-order set theory through bounded
omnific codes.  It also gives explicit intrinsic definitions of all
normal-form coefficients and exponent truncations.  These conclusions
do not identify the exact definable closure of the empty set, nor do they
prove a bi-interpretation.  In particular, the quotient interpretation
does not supply a definable choice of one representative for each set,
or a definable identification between an original surreal and the
canonical surreal reconstructed inside the interpreted set universe.
Those are additional mathematical requirements.
\begin{writenote}[Effect on Part~\ref{dsn:part:one}, 4 October 2026]
This answers the omega-map half of \cref{dsn:q:languages} at the level of an
interpretation: with $\Oz$ and $\Omega$ named, first-order set theory is
expressible through bounded omnific codes. Source~02's
\cref{dsn:rp:prop:omega-code} is the special case of the bounded unary code.
The exact definable closures remain open (\cref{dsn:op:q:dcl}), as does a
bi-interpretation (\cref{dsn:op:q:biinterpretation}). The birthday-enriched
structures of the birthday-cutoff report are bi-interpretable with the
corresponding $H_\kappa$ (\lbl{hset:thm:biinterpretation},
\lbl{hset:thm:main-rounding}); the present signature names no birthday
relation.
\end{writenote}


\paragraph{Choice and the interpretation domain.}
The proof of exact powerset and relation coding does not use Choice.
Without Choice, every admissible graph still has a well-orderable domain,
since its domain is a reverse well-ordered support.  Its collapsed
transitive closure is therefore well-orderable.  Conversely, every set
whose transitive closure is well-orderable has a code by the same argument.
Thus over ZF the quotient presents exactly the hereditarily well-orderable
sets, with their actual membership relation; Choice makes this all of $V$.

\paragraph{Class semantics.}
All normal-form supports, graph domains, and relations used in a single
code are sets.  Consequently the proof uses ordinary set recursion and
the set-sized Mostowski collapse.  No collapse of a proper-class graph
is invoked.  The theorem can be read schematically over ZFC or externally
for the surreal structure of a transitive set model.  In the latter
reading, ``every subset'' means every subset in that model, and the
interpreted structure is precisely its set universe.

\section{Definable Hahn summation of coded families}
\label{dsn:op:subsec:summation}

The support calculus also makes an infinitary operation first-order
definable once its set-sized input family is represented by one number.
The coding of that family is part of the input: individual definability of
all of its values is not a substitute for definability of the family.

For a unit-coefficient series $d$ with support $D$, define
\[
 F_f(i,x)\quad\longleftrightarrow\quad
             \operatorname{Supp}(p(i,x),f).
\]
Let $\operatorname{Fam}(f,d)$ be the conjunction of
\begin{enumerate}
\item $\operatorname{Unit}(d)$ and $\operatorname{Unit}(f)$;
\item every element of $\operatorname{supp}(f)$ is $p(i,x)$ for some
$i\in D$ and some surreal $x$;
\item for every $i\in D$ there is exactly one $x$ with $F_f(i,x)$.
\end{enumerate}
These three conditions are a fixed first-order formula.  They say that
$f$ is the unit-coefficient support code of the graph of a function
$i\mapsto x_i$ on $D$.

\begin{lemma}[Every support-indexed family has a code]
\label[lemma]{dsn:op:lem:family}
For every set function $i\mapsto x_i$ on a reverse well-ordered set $D$,
the normal form
\[
 f=\sum_{i\in D}\Omega(p(i,x_i))
\]
exists and satisfies $\operatorname{Fam}(f,d)$, where
$d=\sum_{i\in D}\Omega(i)$.  Conversely every such code defines a
set-sized family.
\end{lemma}

\begin{proof}
Every nonempty subset of the graph has a largest first coordinate $i$,
and the corresponding second coordinate is unique.  Thus its image under
the lexicographic embedding $p$ is reverse well-ordered.  The normal-form
existence theorem and injectivity of $p$ give the code.  Conversely,
decoding the set-sized support of $f$ gives a set of pairs; the family
conditions make it the graph of a set function on $D$.
\end{proof}

\begin{lemma}[First-order finiteness of a support]
\label[lemma]{dsn:op:lem:finite}
For unit-coefficient $t$, the support of $t$ is finite if and only if
\[
 \begin{split}
 \operatorname{Fin}(t):\quad\forall u\bigl(
  \operatorname{Sub}(u,t)\wedge u\ne0\longrightarrow
  \exists i[\operatorname{Supp}(i,u)\wedge
   \forall j(\operatorname{Supp}(j,u)\longrightarrow i\le j)]
 \bigr).
 \end{split}
\]
\end{lemma}

\begin{proof}
Every nonempty subset of a finite linear order has a least element.
Conversely, an infinite reverse well-ordered set has a strictly decreasing
sequence consisting of its first $\omega$ elements in reverse order.
The set of those elements is an allowed sub-support and has no numerical
least element.  Its unit-coefficient code therefore falsifies the
displayed formula.  This uses the fact that all subsets are represented,
as proved in Lemma~\ref{dsn:op:lem:powerset}.
\end{proof}

\begin{theorem}[Definability of summability and the Hahn sum]
\label[theorem]{dsn:op:thm:summation}
There are parameter-free first-order formulas
$\operatorname{Summable}(f,d)$ and $\operatorname{Sum}(f,d,y)$ in
$\Somega$ which, for family codes, express respectively that
$(x_i)_{i\in D}$ is Hahn-summable and that
\[
 y=\sum_{i\in D}x_i.
\]
The sum formula defines a unique output exactly on the summable inputs.
\end{theorem}

\begin{proof}
Require $\operatorname{Fam}(f,d)$.  The union of the supports is expressed
by the formula
\[
 U_f(a):\quad
 \exists i\exists x\bigl(
 \operatorname{Supp}(i,d)\wedge F_f(i,x)
                       \wedge\operatorname{Supp}(a,x)\bigr).
\]
That this union is reverse well-ordered is equivalent to
\[
 \exists u\bigl(\operatorname{Unit}(u)\wedge
    \forall a(\operatorname{Supp}(a,u)\longleftrightarrow U_f(a))\bigr).
\]
Indeed an existing support is reverse well-ordered, and every set-sized
reverse well-ordered set has its unit-coefficient series.  The union is
a set because the decoded family is a set.

For each exponent $a$, require a code $t$ for precisely the indices at
which its coefficient is nonzero:
\[
 \begin{split}
 \operatorname{Fib}(a,t;f,d):\quad&\operatorname{Sub}(t,d)\ \wedge\\
 &\forall i\left(
   \operatorname{Supp}(i,t)\longleftrightarrow
   [\operatorname{Supp}(i,d)\wedge
     \exists x(F_f(i,x)\wedge c_a(x)\ne0)]\right).
 \end{split}
\]
Such a $t$ always exists uniquely by exact subset coding.  The condition
\[
 \forall a\ \exists t\bigl(
        \operatorname{Fib}(a,t;f,d)\wedge\operatorname{Fin}(t)\bigr)
\]
says exactly that every exponent occurs in only finitely many members of
the family.  This, together with the preceding union-support condition,
is the usual Hahn-summability criterion, so defines
$\operatorname{Summable}(f,d)$.

For completeness, the finite coefficient sum in the output specification
can itself be expressed by a fixed formula, without a variable-length
term.  For the finite fibre $I=\operatorname{supp}(t)$, write
\[
 \operatorname{Next}_t(j,i):\quad
 i,j\in I\ \wedge\ i<j\ \wedge\
 \neg\exists k\in I\ (i<k<j).
\]
This is the next step in descending order.  To say that $c$ is the sum of
the coefficients $c_a(x_i)$ for $i\in I$, use the following alternatives.
If $t=0$, require $c=0$.  Otherwise quantify a family code $g$ on $t$ and
require all of the following finite list of first-order conditions:
\begin{enumerate}
\item $\operatorname{Fam}(g,t)$;
\item for every $i\in I$, every $x,v$ with
$F_f(i,x)$ and $F_g(i,v)$, if $i$ is greatest in $I$, then $v=c_a(x)$;
\item for every $i,j\in I$ and every $x,v,w$ with
$\operatorname{Next}_t(j,i)$, $F_f(i,x)$, $F_g(i,v)$, and $F_g(j,w)$,
require $v=w+c_a(x)$;
\item for the least $i\in I$, require $F_g(i,c)$.
\end{enumerate}
Greatest and least are expressed by quantifying over $I$, a definable
support.  A finite linear order has unique descending partial sums, so
these conditions determine $c$ uniquely, with the intended value.  Their
witness family $g$ has a code by Lemma~\ref{dsn:op:lem:family}.  Denote this
fixed formula, including the finite-fibre requirement, by
$\operatorname{CoeffSum}(a,f,d,c)$.

Now define
\[
 \operatorname{Sum}(f,d,y):\quad
 \operatorname{Summable}(f,d)\ \wedge\
 \forall a\ \operatorname{CoeffSum}(a,f,d,c_a(y)).
\]
For summable inputs the standard Hahn construction forms these finite
real coefficient sums, whose nonzero support is contained in the reverse
well-ordered union already obtained.  Thus a surreal $y$ with exactly
these coefficients exists.  Normal-form uniqueness proves uniqueness.
\end{proof}

The formula does not assign a Hahn sum to a family that fails the
summability criterion even if some informal regrouping would cancel its
terms.  It defines precisely the usual unconditional Hahn operation.
Since the operation has a unique output and a parameter-free graph, it
preserves each intrinsic definable closure when both its domain code and
family code belong to that closure.  The corresponding ambient
definability-preservation statement follows by the same unique-output
argument.

\section{An intrinsically definable commuting family of derivations}
\label{dsn:op:subsec:derivations}

The support calculus provides explicit derivations without choosing a
Hamel basis or an unspecified derivation.  These are diagonal Euler
derivations, rather than the Berarducci--Mantova derivation.

\begin{theorem}[Coefficient-indexed derivations]
\label[theorem]{dsn:op:thm:derivations}
There is a parameter-free definable function
$(b,x)\mapsto D_b(x)$ in $\Somega$ characterized by
\[
 c_a(D_b(x))=c_b(a)c_a(x)\qquad(a\in\No).
\]
For every surreal $b$, the function $D_b$ is a strongly
$\mathbb R$-linear derivation.  The derivations commute pairwise, satisfy
$D_b(\Oz)\subseteq\Pi$, and have common constant field exactly
$\mathbb R$:
\[
 \{x:\forall b\ D_b(x)=0\}=\mathbb R.
\]
In particular $D_0$ is a nonzero parameter-free definable derivation and
$D_0(\omega)=\omega$.
\end{theorem}

\begin{proof}
For $x=\sum_a r_a\Omega(a)$, set
\[
 D_b(x)=\sum_a c_b(a)r_a\Omega(a).
\]
The new nonzero support is a subset of the old support, so the series
exists.  Its coefficients characterize it uniquely and yield the
displayed first-order definition.  The map $a\mapsto c_b(a)$ is
$\mathbb R$-linear on the additive surreal group: addition and real
scalar multiplication act coefficientwise on normal forms.  Thus $D_b$
is $\mathbb R$-linear.

For multiplication, the contribution to any exponent $g$ in $xy$ is a
finite sum over pairs $a+a'=g$ of exponents occurring in the two inputs.
Since
\[
 c_b(g)=c_b(a)+c_b(a'),
\]
its coefficient in $D_b(xy)$ equals the coefficient in
$D_b(x)y+xD_b(y)$.  Normal-form uniqueness proves the Leibniz rule.
The same coefficient argument proves strong linearity: for a summable
family $(x_i)$, each output support is contained in its original support,
so $(D_b(x_i))$ is summable, and the finite coefficient sums commute with
multiplication by $c_b(a)$.

The coefficient at $a$ in either $D_bD_{b'}(x)$ or $D_{b'}D_b(x)$ is
$c_b(a)c_{b'}(a)c_a(x)$, proving commutation.  An omnific input has no
negative exponents; the coefficient at zero in its derivative vanishes
because $c_b(0)=0$.  Hence its derivative belongs to $\Pi$.

Every real is annihilated by every $D_b$.  If $x$ is not real, choose a
nonzero exponent $a$ with $c_a(x)\ne0$.  Since $a\ne0$, its own normal
form has some exponent $b$ with $c_b(a)\ne0$.  Then $D_b(x)$ has nonzero
coefficient at $a$, so $x$ is not a common constant.  Finally,
$c_0(1)=1$, and the only term of $\omega$ has exponent one, yielding
$D_0(\omega)=\omega$.
\end{proof}

\begin{proposition}[Canonical partial integration]
\label[proposition]{dsn:op:prop:partial-integration}
For each $b$ there are uniformly definable $\mathbb R$-linear operators
$K_b,J_b$ such that
\[
 \begin{array}{ll}
 c_a(K_b(x))=c_a(x),&c_b(a)=0,\\
 c_a(K_b(x))=0,&c_b(a)\ne0,
 \end{array}
 \qquad
 c_a(J_b(x))=
 \begin{cases}
 c_a(x)/c_b(a),&c_b(a)\ne0,\\
 0,&c_b(a)=0.
 \end{cases}
\]
They satisfy
\[
 D_bJ_b=J_bD_b=\operatorname{id}-K_b,
 \qquad\operatorname{im}(K_b)=\ker(D_b),
 \qquad\operatorname{im}(D_b)=\ker(K_b).
\]
Consequently $D_b(y)=x$ has a solution exactly when $K_b(x)=0$; then
$J_b(x)$ is the unique solution with $K_b(y)=0$.
\end{proposition}

\begin{proof}
All displayed series have support contained in the support of their
input, so they exist and their coefficient specifications are first-order
definitions with unique outputs.  The identities follow coefficient by
coefficient.  If $D_b(y)=x$, the vanishing multiplier at every exponent
with $c_b(a)=0$ implies $K_b(x)=0$.  Conversely, for $K_b(x)=0$, the
identity $D_bJ_b(x)=x-K_b(x)$ gives a solution.  Two solutions differ by an
element of $\ker D_b=\operatorname{im}K_b$, so the stated normalization
makes it unique.
\end{proof}

For clarity, $D_0$ is not the standard surreal derivation normalized by
$\partial\omega=1$: here $D_0\omega=\omega$, and its kernel is much
larger than $\mathbb R$.  For example,
$D_0(\Omega(\omega))=0$ because $c_0(\omega)=0$.  The advantage of this
family is its explicit intrinsic definability in the stated language.
The general classification of strong real-linear derivations diagonal on
monomials is already discussed in the repository's real-vector-space
report \cite{RepoRvs}; the present result isolates a uniformly definable coefficient
family and supplies its definable partial inverses.

\section{Further consequences for parameters, fields, and operations}
\label{dsn:op:sec:consequences}

\subsection{Coefficientwise operations and full product rings}

Once coefficients are intrinsically available, many useful operations
have short definitions even when their normal forms are transfinite.
Coordinatewise linear algebra and diagonal operators are developed in
the repository's vector-space report \cite{RepoRvs}.  The point here
is the uniform intrinsic definability of their coefficient formulas.
The following is a general existence and preservation principle.

\begin{theorem}[Coefficientwise lifting]\label[theorem]{dsn:op:thm:coefficientwise}
Let \(f:\R^k\to\R\) be definable in \(\Somega\), with specified
parameters, and assume \(f(0,\ldots,0)=0\).  There is a unique operation
\(\widetilde f:\No^k\to\No\) given by
\[
 c_a\bigl(\widetilde f(x_1,\ldots,x_k)\bigr)
   =f\bigl(c_a(x_1),\ldots,c_a(x_k)\bigr)
       \qquad(a\in\No).
\]
It is definable with the same parameters.  If
\(f[\Z^k]\subseteq\Z\), then
\(\widetilde f[\Oz^k]\subseteq\Oz\).
\end{theorem}
\begin{proof}
The new support is contained in the finite union of the input supports,
because the all-zero coefficient tuple is sent to zero.  A finite union
of reverse-well-ordered sets is reverse well ordered: a nonempty subset
has a greatest member in each nonempty component, and the greatest of
those finitely many choices is its greatest member.  The normal-form
existence theorem therefore gives the output.  The coefficient equations
give uniqueness and a first-order graph, using
\cref{dsn:op:thm:support}.  For omnific inputs there are no negative-exponent
coefficients, and the coefficient at zero is an integer tuple; both
properties survive under the stated assumptions.
\end{proof}

For example, define the Hadamard product and support indicator by
\[
 c_a(x\odot y)=c_a(x)c_a(y),\qquad
 c_a(e_x)=\begin{cases}1,&c_a(x)\ne0,\\0,&c_a(x)=0.\end{cases}
\]
Both are definable operations on surreals and preserve omnific integers.
The Hadamard product differs from ordinary field multiplication:
\(\omega\odot\omega=\omega\), whereas
\(\omega\cdot\omega=\omega^2\).
Coefficientwise absolute value gives another example: it sends
\(\omega-1\) to \(\omega+1\), whereas numerical absolute value leaves
\(\omega-1\) unchanged.

For a unit-coefficient series \(d\) with support \(D\), set
\[
 H_D=\{x:\supp(x)\subseteq D\}.
\]
This is a definable domain with parameter \(d\).  Every coefficient
function \(D\to\R\) is allowed, because its nonzero support is a subset
of the reverse-well-ordered set \(D\).  Consequently
\begin{equation}\label{dsn:op:eq:product-ring}
 (H_D,+,\odot,d)\cong\R^D
\end{equation}
as unital rings.  If \(D\subseteq\No_{>0}\), the entire domain \(H_D\)
lies in \(\Oz\).  Thus omnific integers carry copies of arbitrary set
products of the real field under this additional definable multiplication.
There is no claim that \(H_D\) is a subring for ordinary multiplication.

The idempotents for \(\odot\) have only coefficients zero and one, so
they correspond exactly to \(\mathcal P(D)\).  Their Boolean operations
are
\[
 e\wedge f=e\odot f,\qquad e\vee f=e+f-e\odot f,
 \qquad \neg_D e=d-e.
\]
This is another way to see the exact powerset coding used in the
interpretation.  There is no global Hadamard identity on all of \(\No\):
it would need coefficient one at every surreal exponent, a proper-class
support.

\subsection{The inner-model transcendence consequence}

The next corollary applies the repository's support-amplification argument
(\cref{dsn:thm:amplify} of Part~\ref{dsn:part:one}) after the missing absoluteness hypothesis has
been established.  We include the short algebraic mechanism.

\begin{corollary}[A proper inner model misses a class of independent integers]
\label[corollary]{dsn:op:cor:inner-independent}
Let \(M\subsetneq V\) be a transitive inner model of ZFC with the same
ordinals, and put \(K=\No^M\).  There is an ordinal-indexed family of
omnific integers, all in \((\omega,\omega+\sqrt\omega)\), which is
algebraically independent over \(K\) in the sense that every finite
subfamily is algebraically independent.  No equality of the real fields
is required.
\end{corollary}
\begin{proof}
The field \(K\) is proper.  Otherwise every ordinal subset would have
its characteristic sign sequence in \(M\), and hence belong to \(M\);
ordinary pointed-relation coding would then give \(M=V\).
Choose \(t\in(0,1)\setminus K\) by applying \(\tau\) to an element
outside \(K\).  Its inverse is algebraic, so \(\tau\) cannot send that
element into \(K\).

For each ordinal \(\alpha\), put
\[
 e_\alpha=\omega^{-(\alpha+1)},\qquad
 u_\alpha=t e_\alpha,\qquad
 z_\alpha=\omega+\omega^{u_\alpha}.
\]
All \(e_\alpha\) belong to \(K\).  The cosets \(u_\alpha+K\) are
linearly independent over \(\Q\) in the additive quotient \(\No/K\).
Indeed, a nontrivial finite rational relation would give
\(t\sum q_\alpha e_\alpha\in K\).  The sum is a nonzero member of
\(K\), by uniqueness of normal form, so this would put \(t\) in \(K\).

Consider a nonzero polynomial over \(K\) in finitely many
\(\omega^{u_\alpha}\).  A term indexed by an exponent vector \(n\)
has its coefficient's normal-form support translated by
\(\sum n_\alpha u_\alpha\).  By
\cref{dsn:op:thm:all-reals-free-absoluteness}, every exponent of every
coefficient from \(K\) itself belongs to \(K\).  Distinct exponent
vectors therefore produce disjoint support cosets, by the preceding
linear independence.  Their nonzero series cannot cancel.  Thus the
monomials are algebraically independent, and translating each by
\(\omega\in K\) preserves this property.  Finally
\(0<u_\alpha<1/2\), giving the claimed omnific interval.
\end{proof}
The result states finite polynomial independence for a definable family
from the parameter \(t\).  It does not assign a cardinality to a
proper-class transcendence basis or require Global Choice.

\subsection{Exact preservation of definability levels by the bounded code}

\begin{corollary}[The bounded code does not raise a fixed Levy level]
\label[corollary]{dsn:op:cor:code-levels}
For every fixed standard \(n\ge1\), every fixed allowed parameter family,
and every surreal \(x\), the objects \(x\) and \(Q(x)\) have the same
\(\Sigma_n\)-, \(\Pi_n\)-, and \(\Delta_n\)-singleton definability
status in ambient set theory.
\end{corollary}
\begin{proof}
By \cref{dsn:op:thm:canonical-delta-one}, the graph of \(Q=\Omega\circ\tau\)
is \(\Delta_1\).  Its range is \(\Delta_1\): a normal form must have
exactly one term, coefficient one, and exponent in \((0,1)\).
The inverse uses that exponent and the algebraic formula for
\(\tau^{-1}\), so its graph is \(\Delta_1\) too.  Apply
\cref{dsn:op:prop:definition-complexity-transfer} in both directions.
\end{proof}

There is also a modest literal description-length conclusion once a
syntax convention is fixed.  For literal finite-alphabet length, fix a
finite parameter family \(A\) with named constant references; broader
parameter permissions are treated one finite tuple at a time.
Use a fixed finite alphabet, allow either
of two fixed symbols \(v_0,v_1\) as the sole free object variable, and
write permitted parameters as constant references with a fixed cost
convention.  Let \(L_A(x)\) be the minimum length among definitions with
either output variable.  If \(\varphi(v_0)\) defines \(x\), the context
\[
 \exists v_0\bigl(\varphi(v_0)\wedge
                  \operatorname{Graph}_Q(v_0,v_1)\bigr)
\]
defines \(Q(x)\), with \(\varphi\) inserted unchanged.  All variables
bound inside \(\varphi\) retain their scopes.  Reverse the roles of the
two free symbols when necessary, using two precompiled versions of
the fixed graph formula.  The inverse graph gives the reverse bound.
Thus for every \(A\)-definable \(x\),
\begin{equation}\label{dsn:op:eq:length-bound}
 |L_A(x)-L_A(Q(x))|\le c
\end{equation}
for a constant depending only on the coding formulas and the syntax
convention.  This is a statement about minimum lengths of successful
definitions, not an algorithm for recognizing success or minimizing
length.  It supplies an additive overhead bound for this specific
reversible code; it is not an optimal proof-length or birthday bound.

\subsection{Elementary theories and elementary inclusions}

\begin{theorem}[The expansion retains the full elementary theory of sets]
\label[theorem]{dsn:op:thm:theory-faithful}
For canonical surreal structures of models \(M,N\models\mathrm{ZFC}\),
\[
 M\equiv N\quad\Longleftrightarrow\quad
 \Somega^M\equiv\Somega^N.
\]
For nested transitive models \(M\subseteq W\models\mathrm{ZFC}\),
with the canonical inclusions,
\begin{equation}\label{dsn:op:eq:elementary-exact}
 M\preccurlyeq W\quad\Longleftrightarrow\quad
 \Somega^M\preccurlyeq\Somega^W.
\end{equation}
For proper classes each elementarity assertion is a scheme over standard
formulas.
\end{theorem}
\begin{proof}
The class \(\No\), all the field operations, \(\Oz\), and \(\Omega\)
have fixed set-theoretic definitions.  Relativizing quantifiers therefore
translates every sentence of the enriched surreal language into a
set-theoretic sentence.  This proves the forward implication for
elementary equivalence.  The interpretation of
\cref{dsn:op:thm:V-interpretation} gives the reverse translation, and hence
the reverse implication.  In a nonstandard model the graph codes and
collapses are interpreted internally, so this argument remains the
ordinary syntactic argument for interpretations.

The same translation proves the forward implication of
\eqref{dsn:op:eq:elementary-exact}, now with parameters.  For the converse,
given any finite tuple of parameters \(\bar x\in M\), choose in \(M\)
admissible surreal codes \(\bar c\) whose collapsed values are
\(\bar x\).  Normal-form absoluteness preserves these codes and their
relations in \(W\).  Their old pointed graphs remain well founded and
extensional.  Accessibility is preserved because it is equivalent to
every node having a finite predecessor path from the root, and such
paths are absolute.  Mostowski collapse is absolute, so the codes have
the same values \(\bar x\) there.  Elementarity of the surreal inclusion applied
to the translation of any formula \(\varphi(\bar x)\) now gives
\(M\models\varphi(\bar x)\) iff \(W\models\varphi(\bar x)\).
\end{proof}

This sharply separates preservation of operations from preservation of
all formulas.  Every canonical inclusion in the last theorem is
elementary in the ordered-field reduct by
\cref{dsn:op:cor:inner-rcf-elementarity}.  The enriched inclusion fails to be
elementary whenever the set universes disagree on a first-order formula
with parameters from the smaller model.  For
example, the translated continuum-hypothesis sentence distinguishes
canonical structures built in universes disagreeing about CH.  More
generally the translation detects any first-order set-theoretic
difference, including differences above the reals.

The theorem still does not prove a bi-interpretation.  A choice of codes
for parameters in its proof is an existential choice for a finite
tuple, not a definable assignment of canonical codes to all original
surreals.  Nor does elementary equivalence classify isomorphism types.
\begin{writenote}[On the nested clause, 4 October 2026]
For nested models of the same height, \eqref{dsn:op:eq:elementary-exact} holds
only in the trivial case: both of its sides are equivalent to $M=W$. This
observation is the write's, and is proved as follows; it does not contradict
the theorem.

\emph{Observation.} Let $M\subseteq W$ be transitive models of ZFC (set models,
or inner classes with elementarity read as a scheme) with
$\Ord\cap M=\Ord\cap W$. If $M\preccurlyeq W$, or equivalently
$\Somega^M\preccurlyeq\Somega^W$, then $M=W$.

\emph{Proof.} Suppose $M\preccurlyeq W$ and take $x\in W\setminus M$. By Choice
in $W$, $x$ is the root of the Mostowski collapse of a well-founded extensional
relation on an ordinal, which an ordinal pairing codes by a set $A\subseteq\alpha$
of ordinals. If $A$ were in $M$, the collapse, which is absolute between
transitive models, would put $x$ in $M$; so $A\in W\setminus M$, while
$\alpha\in M$ because the ordinals agree. Let $P=\mathcal P(\alpha)^M\in M$.
Then $M\models\forall y\,(y\subseteq\alpha\rightarrow y\in P)$, whereas
$W\models\exists y\,(y\subseteq\alpha\wedge y\notin P)$, witnessed by $A$. One
formula with the parameters $\alpha,P\in M$ thus refutes $M\preccurlyeq W$.
\qed

This is the argument of \lbl{hset:thm:outer} of the birthday-cutoff report
(there for birthday structures, with codes for $\alpha$ and
$\mathcal P(\alpha)^M$ as parameters). The nested clause of the theorem
therefore has content only for models of different heights, for example
$M=V_\kappa$ and $W=V_\lambda$ when $\kappa<\lambda$ are inaccessible cardinals
with $V_\kappa\prec V_\lambda$. The clause on elementary equivalence is
unaffected; it is the analogue for $\Somega$ of
$M\equiv N\iff\mathfrak B^M\equiv\mathfrak B^N$ in \lbl{hset:rem:internal}
and \lbl{hset:cor:equiv}.
\end{writenote}


\subsection{An operation calculus has no universal self-evaluator}

There are two different meanings of ``all definable operations.''  One
can prove a theorem scheme that each specified formula gives a
parameter-preserving operation, or supply a uniform translation of
finite syntax.  One cannot conclude that there is an internal evaluator
for truth of arbitrary formulas in the same universe.

Here is an elementary diagonal version in the enriched surreal language.
Its ordinary naturals are definable by
\(\N=\{n\in\Oz\cap\R:n\ge0\}\).  Suppose a parameter-free definable
total function \(U:\N\times\No\to\No\) enumerated all parameter-free
definable total unary surreal functions, in the sense that each such
function was \(x\mapsto U(n,x)\) for some natural \(n\).  Then the
parameter-free function
\[
 g(x)=\begin{cases}U(x,x)+1,&x\in\N,\\0,&x\notin\N\end{cases}
\]
would have an index \(k\).  Evaluation at \(k\) would give
\(U(k,k)=U(k,k)+1\), a contradiction.  The support and summation
formulas evade this obstruction because their inputs are explicitly
coded mathematical families, not an alleged enumeration of every
operation definable in their own language.

\section{Size, first omissions, and the limits of collection}
\label{dsn:op:sec:first-omissions}

The parameter distinctions also control what a counting argument can
show.  For an externally fixed set model \(M\), a set of parameters
\(A\subseteq M\) gives
\[
 |D_A^M|_{\mathrm{external}}\le
          \max(\aleph_0,|A|_{\mathrm{external}}).
\]
There are only countably many finite formulas and only the indicated
number of finite parameter tuples.  In particular a fixed finite
parameter family \(A\) produces an externally countable definability hull.  This does
not prove that the hull is a countable set \emph{inside} \(M\), or that
\(M\) contains a sequence enumerating its successful definitions.
Pointwise definable models illustrate why the distinction matters
\cite{HLR}: every element may have an external parameter-free definition
while the model internally satisfies the usual uncountability theorems.

For the full ambient universe \(V\), allowing all ordinals gives a
proper-class parameter supply and the preceding set bound does not apply.
For a fixed set model \(M\), its ordinals are still an external set and
the bound continues to apply.  With
individual real parameters, every real is trivially definable from
itself, and \cref{dsn:op:cor:countable-birthday-real} gives much more.
Nevertheless a particular formula uses only one finite tuple.  Neither
permission is an infinitary conjunction of names.

\subsection{Two different first-omission ordinals}

The repository distinguishes a first ordinal without an allowed
definition from the first birthday at which a new sign sequence appears
(\cref{dsn:sec:core,dsn:sec:firstbirthday} of Part~\ref{dsn:part:one}).  They measure different phenomena.

First fix a transitive model \(M\) and an external parameter family \(A\).
Suppose there is a least ordinal \(\rho\in\Ord^M\) not in \(D_A^M\).
Suppressing the superscript, the
largest initial subclass of \(K_A\), for sign simplicity, is
\begin{equation}\label{dsn:op:eq:initial-core}
 K_A\cap\No_{<\rho}.
\end{equation}
Indeed, every prefix of a definable \(x\) of birthday below \(\rho\)
is definable from \(x\) and its length, and that length is \(A\)-definable.
Conversely, if \(x\in K_A\) has birthday at least \(\rho\), its prefix
of length \(\rho\) cannot belong to \(K_A\), since its definable birthday
operation would then define \(\rho\).  No initial subclass of \(K_A\)
can contain such an \(x\).

The ordinal \(\rho\), when it exists, is an epsilon number.  It is not
finite.  If it were not a fixed point of ordinal
\(\alpha\mapsto\omega^\alpha\), its finite Cantor normal form would use
only exponents strictly below \(\rho\).  Each such exponent and each
finite coefficient is \(A\)-definable, so the normal form would define
\(\rho\), a contradiction.  The stronger field and exponential
properties of \eqref{dsn:op:eq:initial-core} are proved in the existing report
using the classical birthday-cutoff theorems.  This discussion does not
assert that such a \(\rho\) exists for every parameter budget; with all
ordinals permitted there is none.

Second let \(M\subsetneq V\) be a transitive inner model of ZFC with all
ordinals.  Define
\[
 \lambda_M=\min\{\alpha:\mathcal P(\alpha)^M\ne
                              \mathcal P(\alpha)^V\}.
\]
This exists: otherwise all sets of ordinals, and then all sets by
well-founded relation coding, would belong to \(M\).
Characteristic sign sequences show immediately that \(\lambda_M\)
is the first birthday of a surreal not in \(\No^M\).

For clarity, the standard cardinal argument is included.  If
\(\lambda_M\) were not a \(V\)-cardinal, let
\(\kappa=|\lambda_M|^V<\lambda_M\), and choose a bijection
\(f:\kappa\to\lambda_M\).  Pull ordinal order back along \(f\) to a
well-order \(E\) on \(\kappa\).  The usual ordinal pairing codes \(E\)
by a subset of \(\kappa\); by minimality this code belongs to \(M\).
The unique ordinal collapse of \(E\), computed in \(M\), is absolute.
It therefore recovers the original \(f\) and its range
\(\lambda_M\).  Now every \(A\subseteq\lambda_M\) has pullback
\(f^{-1}[A]\subseteq\kappa\) in \(M\), and hence belongs to \(M\),
a contradiction.  Thus \(\lambda_M\) is a \(V\)-cardinal, and
therefore also an \(M\)-cardinal.  The case \(M=\HOD\) is the
repository's first non-OD birthday theorem.
\begin{writenote}[Antecedents, 4 October 2026]
The epsilon-number argument is the lemma behind \cref{dsn:thm:core}
(\cref{dsn:sec:core} of Part~\ref{dsn:part:one}), and the cardinal threshold is
\lbl{univ:prop:beta} of the universes report, proved there for arbitrary
transitive $M\subseteq N$; \cref{dsn:thm:firstbirthday} is its case $M=\HOD$.
\end{writenote}


These observations organize different questions without identifying
them.  The ordinal \(\rho\) concerns definitions using a fixed external
budget; \(\lambda_{\HOD}\) concerns hereditary ordinal definability;
and \(\omega_1^{\HOD}=\omega_1\) concerns the ability to choose real
codes uniformly.  Cardinal correctness can hold in a proper inner model,
as the Cohen example in \cref{dsn:op:sec:real-parameter-extension} shows.

\section{Further research questions and proposed program}
\label{dsn:op:sec:research}

The following questions are deliberately stated beyond what has been
proved.  The first three concern expressive power; the next group concerns
parameters and complexity; the final group concerns operations and
formalization.  Their order reflects dependencies rather than a claim
about difficulty.

\begin{question}[The missing comparison map]
\label[question]{dsn:op:q:biinterpretation}
Is \(\Somega\) bi-interpretable with \((V,\in)\), or can one rule this
out?  Specifically, is there a first-order definable relation associating
an original surreal \(x\) with exactly the equivalence class coding its
canonical set-theoretic sign sequence in the interpretation?
\end{question}

Theorem~\ref{dsn:op:thm:V-interpretation} gives a quotient universe, and the
ordinary set construction gives surreals within that universe.  A
bi-interpretation requires the second composite to be definably
isomorphic to the original surreal structure.  An abstract external
isomorphism does not meet this requirement.  An explicit sign-prefix
relation, or a sufficiently strong intrinsic birthday relation, would
be a promising route.  The birthday-enriched Chen--Hamkins--Yang
announcement \cite{CHYKobe} identifies a useful comparison target but does
not supply the missing definition in the present signature.

\begin{writenote}[Added 4 October 2026, batch~95]
Part~IX of
\repopath{foundations-and-computation/birthday-cutoffs-and-hereditary-sets}
\cite{RepoHset} (its sources~13 and~14, batch-95 manuscripts~01 and~02, merged;
AI-assisted, unrefereed, not formalized) is the birthday-enriched comparison
target named above. In $(\No;0,1,+,-,\cdot,<,\prec_b)$, and natively in
$(\Oz;0,1,+,-,\cdot,<,\prec_b)$, with $x\prec_b y\iff\bd(x)<\bd(y)$, it
interprets $(V,\in)$ parameter-free by codes of pointed well-founded
extensional graphs, each stored in a single number, and proves the missing
comparison: a parameter-free formula sending each original number to the class
of codes of its own sign sequence (its Theorems~71.1 and~72.2,
\lbl{hset:sf:thm:graph}, \lbl{hset:sf:thm:roundtrip}). Every parameter-free
set-theoretically definable relation on $\No$ is therefore parameter-free
definable in that language (its Corollary~74.1, \lbl{hset:sf:cor:definability}),
in particular $\Oz$ and the graph of $\Omega$, so $\Somega$ is a reduct of a
definitional expansion of it. Its Proposition~74.2 (\lbl{hset:sf:prop:reducts},
added in that merge) shows that the specific form of this question, a
definable relation sending each original surreal to the class of codes of its
own sign sequence, holds for $\Somega$ exactly when birthday precedence is
definable in $\Somega$, with the same parameters; by its last clause this
applies to the interpretation of \cref{dsn:op:thm:V-interpretation}, whose
isomorphism onto $(V,\in)$ is the set-theoretically definable Mostowski
collapse. Whether $\prec_b$ is so definable is not decided there, and this
question stays open. Its Remark~67.5 (\lbl{hset:sf:rem:dsn}) compares the two
codings: sign blocks there, Hahn supports here.
\end{writenote}

\begin{question}[Exact intrinsic definable closures]
\label[question]{dsn:op:q:dcl}
Determine
\[
 \dcl_{\Szero}(\varnothing),\qquad
 \dcl_{\Somega}(\varnothing),\qquad
 \dcl_{\Somega}(\R),
\]
and their ordinal-parameter analogues, in specified ambient models.
\end{question}

For the smaller pair, the parameter-free closure lies in \(\R\) and
contains the reals definable in \((\R;+,\cdot,<,\Z)\).
For the richer structure, \(\omega\), all the operations proved here,
and bounded tuple codes are available.  Neither statement identifies the
closure.  An answer to \cref{dsn:op:q:biinterpretation} may connect this problem
to ambient definability; without that comparison map, the interpretation
theorem cannot simply be substituted for an exact closure theorem.
This is the remaining part of repository \cref{dsn:q:languages} (Part~\ref{dsn:part:one}).

\begin{question}[How much of the signature is necessary?]
Can the named omega map be replaced by its range predicate, or by a
weaker normal-form operation?  Can the omnific predicate be removed
while preserving an interpretation of set theory?
\end{question}

The proof uses \(\Oz\) to recover canonical real coefficients and
\(\Omega\) to index them by arbitrary exponents and to build an
order-compatible pairing.  A monomial predicate names scales but does
not automatically identify their surreal exponents.  Conversely,
\(\Omega\) already defines a natural valuation
(\cref{dsn:op:rem:valuation}), but a valuation alone does not provide a
canonical residue-field section.  A useful first target is an exact
replacement for one of these two missing ingredients, keeping the
other unchanged.

\begin{writenote}[Added 4 October 2026, batch~95]
For interpretations that come with a comparison sending each number to codes
of itself, Proposition~74.2 of
\repopath{foundations-and-computation/birthday-cutoffs-and-hereditary-sets}
(\lbl{hset:sf:prop:reducts}) reduces the analogue of this question to
definability: a structure on $\No$ or $\Oz$ whose finitely many primitives are
given by parameter-free set-theoretic formulas on sign sequences, a reduct of
$\Somega$ among them, has such a comparison exactly when it defines
$0,1,+,-,\cdot,<,\prec_b$. Its Question~79.3 (\lbl{hset:sf:q:reducts}) asks which
reducts of $\{0,1,+,-,\cdot,<,\prec_b\}$ have one, and records the overlap with
this question. This question asks for interpretations without a comparison,
and is not affected.
\end{writenote}

\begin{question}[Automorphisms compatible with the coding structure]
Classify automorphisms of \(\Somega\), and determine how they can act
on representatives of the interpreted quotient universe.
\end{question}

Kaplan--Krapp--Serra ask about nontrivial omega-preserving automorphisms
in Question~5.6 and a related extension problem in Question~5.7
\cite{KKS}.  Their pointwise-valuation
rigidity result must be distinguished from preserving a valuation
equivariantly while also acting on its value group.  The present
interpretation places an additional constraint on automorphisms when
\(\Oz\) is also named.  Even a rigid interpreted quotient would leave
the possibility of movement within individual code classes.  Studying
that action is a concrete way to test whether the missing comparison
map in \cref{dsn:op:q:biinterpretation} can exist.

\begin{question}[The parameter degree of a uniform section]
For a fixed universe, classify the sets of ordinals \(p\), or the reals
\(r\), for which
\[
 \omega_1^{\HOD(p)}=\omega_1,\qquad
 (\kappa^+)^{\HOD(p)}=\kappa^+.
\]
What ordering of parameter strengths is induced by the existence of
the corresponding definable coding sections?
\end{question}

Theorems~\ref{dsn:op:thm:real-code-section} and
\ref{dsn:op:thm:kappa-code-section} turn this into an exact inner-model
question.  Natural variants restrict the input birthdays, require
sections to commute with embeddings, or require one parameter to work
at every cardinal.  The all-cardinal criterion is agreement on cardinals,
not equality with \(V\).  Further work should distinguish changes in
that criterion from changes in which individual surreals are OD.

\begin{question}[Optimal logical complexity]
\label[question]{dsn:op:q:complexity}
When a real-code section exists, is the \(\Delta_2(p)\) upper bound
optimal?  Which additional hypotheses permit a \(\Delta_1(p)\) section?
What are the sharp complexities for variants with a prescribed
enumeration convention?
\end{question}

The conditional totality in the existing \(\Delta_2\) bound matters.
If an eligible ordinal has no successful OD enumeration, adding a
default for that additional failure is an unbounded negative search;
the current argument does not establish the same bound for such an
unconditional extension.  This is a concrete place where a complexity
lower bound could be informative, rather than merely shortening a
formula.

\begin{question}[Birthday and description costs]
\label[question]{dsn:op:q:cost}
Find sharp birthday bounds for \(Q(x)\), the lexicographic pair \(B(a,b)\),
and the pointed-graph code of a set.  How do the bounds change when the
leading term, coefficient alphabet, or target interval is prescribed?
\end{question}

The fixed-level preservation in \cref{dsn:op:cor:code-levels} and the additive
literal length bound \eqref{dsn:op:eq:length-bound} concern logical syntax.
They do not bound the ordinal length of the output sign sequence.
For graph codes, the chosen enumeration of the transitive closure may
also affect the cost.  An initial project is to combine corrected
classical birthday estimates \cite{vdDE,vdDEerr} with the explicit
formulas for \(\tau,B,Q\), before asking for optimal constructions.
This develops repository \cref{dsn:q:cost-birthday,dsn:q:cost} (Part~\ref{dsn:part:one}).

\begin{question}[Interpretations in bounded birthday structures]
For which \(\lambda\) does an appropriate expansion of
\(\No_{<\lambda}\) interpret a substantial set-theoretic structure such
as \(H_\kappa\) or \(V_\alpha\)?  What is the exact interpreted domain?
\end{question}

The unrestricted proof quantifies over codes for \emph{all} subsets
of a support.  At a birthday cutoff, those subset codes or the paired
relation codes may lie outside the structure.  In that case the
well-foundedness formula need not express ambient well-foundedness.
Determining precisely which closure and size hypotheses restore the
proof is the first necessary step; merely restricting the quantifiers
of the present formulas is not a solution.

\begin{question}[Strong definable operations beyond diagonal calculus]
Classify the intrinsically definable strong real-linear maps and
derivations in \(\Somega\).  Which are determined by a definable
action on monomials, and which require more general support changes?
\end{question}

The family \(D_b\) and its partial integration operators are
support-preserving, which makes existence transparent.
Support-changing operators require a summability criterion for their
monomial images.  The coded-family theorem supplies a way to express
that criterion for a given family.  One can then ask which such
operators are definable uniformly, and how their kernels, images and
commutants are described.  The Berarducci--Mantova derivation is a
separate target: its compatibility with exponential structure and its
normalization cannot be inferred from the diagonal construction.

\begin{question}[Coefficientwise algebra and definable product structures]
Which definable real operations have coefficientwise lifts compatible
with ordinary surreal multiplication?  What substructures of the
definable product rings \(H_D\cong\R^D\) can be characterized uniformly?
\end{question}

Hadamard multiplication gives all Boolean support algebra directly.
Ordinary multiplication instead convolves exponents.  Requiring
compatibility with both operations should impose strong constraints,
but the present coefficientwise theorem assumes no such compatibility.
Finite-support ideals, support filters, and the interactions between
coordinate projections and convolution provide explicit starting
examples.

\begin{question}[First omissions and inner support fields]
\label[question]{dsn:op:q:support}
Which first-omission ordinals \(\rho\) and first new birthdays
\(\lambda_M\) can occur under specified hypotheses on \(M\) and the
parameter family?  Which intermediate support fields arise as
\(\No^M\), and which do not?
\end{question}

The absoluteness theorem proves that every transitive inner model of ZFC
gives a support field, including models with fewer reals.  It does not
classify all support fields, or all possible first-omission spectra.
One can now compare combinatorial closure of a field with set-theoretic
closure of a candidate inner model, without carrying the obsolete
same-reals restriction.  This extends the still-open classification
parts of repository \cref{dsn:q:support,dsn:q:rho} (Part~\ref{dsn:part:one}).

\begin{question}[Axiom strength and formal verification]
Which fragments of set theory suffice for the canonical-operation
certificates, the powerset coding, and the interpretation?  Can their
formula translations and correctness proofs be formalized with an
explicit dependency ledger?
\end{question}

There are several distinct obligations: existence of the relevant
normal forms, Collection for evaluation certificates, the set-sized
Mostowski collapse, and Choice for coding every set by a well-orderable
domain.  The ZF version identifies the hereditarily well-orderable
quotient, but does not by itself calibrate weaker theories.  Existing
formal work on Conway normal forms \cite{PK} is useful background.
Formalizing a finite syntactic translator alone would verify the
translation algorithm, not the semantic adequacy of its graph codes.

\begin{writenote}[Status of the questions, 4 October 2026, batch~93]
All twelve questions stay open as printed. \Cref{dsn:op:q:dcl} is the remainder
of \cref{dsn:q:languages} of Part~\ref{dsn:part:one} and contains source~02's
question on definable closure versus fixed elements in
\cref{dsn:rp:sec:agenda}. The question on automorphisms meets
\lbl{sse:sr:q:omegafield} of the self-embeddings report, which also starts from
Kaplan--Krapp--Serra's Questions~5.6--5.7. The question on birthday and
description costs (\cref{dsn:op:q:cost}) develops
\cref{dsn:q:cost-birthday,dsn:q:cost}, as do source~02's question on the
birthday cost of countable codes and source~05's Question~13.8 in
\cref{dsn:rp:sec:agenda}. \Cref{dsn:op:q:support} extends
\cref{dsn:q:support,dsn:q:rho}; \cref{dsn:rp:rem:support} supplies one more
support field.
\end{writenote}

\subsection{Further questions and research added at the write \writetag}
\label{dsn:op:sub:further}

The following items are recorded under Vladimir's standing rule of
4~October 2026: claims made by the manuscript with an argument that it
itself presents as a sketch or an upper bound are kept as open questions,
with the argument and what is missing.

\begin{question}[The complexity of the decoder and of its sections \writetag]
\label[question]{dsn:op:q:remark45}
Prove in full the bounds of the remark after
\cref{dsn:op:thm:real-code-section}: that the real decoder $D$ has a
$\Delta_1$ graph over ZFC, and that under the hypotheses of
\cref{dsn:op:thm:real-code-section,dsn:op:thm:kappa-code-section} the
canonical section has a $\Sigma_2(p)$ graph and a total $\Delta_2(p)$ (or
$\Delta_2(p,\kappa)$) extension.
\end{question}
The manuscript argues this in a remark, not a proof: existential set witnesses
for the graph of $D$ and its complement (an ordinal collapse; a bounded failure
of the order axioms; a nonempty subset without least element); and, for the
section, the standard $\Delta_2(p)$ surjection from $\Ord$ onto $\OD(p)$ of
\cite[Section~4]{DoraisHamkins} with a least-code search. What is missing is a
written verification that the least-code search, the choice of a real
well-order of the required type and the default off the intended domain stay
within the stated classes; the manuscript itself calls these upper bounds and
asks for optimality in \cref{dsn:op:q:complexity}.

\begin{question}[The description-length bound \writetag]
\label[question]{dsn:op:q:length}
With the syntax convention stated before \eqref{dsn:op:eq:length-bound}, write
out the two precompiled graph formulas of $Q$ and $Q^{-1}$ and the resulting
constant $c$, and prove \eqref{dsn:op:eq:length-bound}.
\end{question}
The manuscript's argument inserts a definition $\varphi(v_0)$ unchanged into a
fixed context $\exists v_0(\varphi(v_0)\wedge\operatorname{Graph}_Q(v_0,v_1))$,
renaming the free variable when necessary. What is missing is the explicit
cost accounting for variable capture and the constant-reference convention;
the bound depends on that convention, as the manuscript states.

\section{Proof dependencies, contribution ledger, and verification status}
\label{dsn:op:sec:status}

\subsection{What is proved here and what is imported}

The following ledger records the scope of the claims.  An extension
means an addition to the audited repository treatment with a proof in
this article.  It is not a certification that no earlier source
contains an equivalent result.

\begingroup
\small
\setlength{\tabcolsep}{5pt}
\begin{longtable}{P{0.27\textwidth}P{0.18\textwidth}P{0.47\textwidth}}
\toprule
Claim or construction & Status & Dependency and precise scope\\
\midrule
\endfirsthead
\toprule
Claim or construction & Status & Dependency and precise scope\\
\midrule
\endhead
\bottomrule
\endfoot
Surreal field and normal forms & Classical &
Conway--Gonshor theory; legal supports are set-sized and reverse well
ordered.\\
Real closed ambient hulls, floor, fraction representation & Repository
background &
Unique-output closure and classical field facts; rederived here to fix
the parameter budget.\\
Bounded unary monomial code \(Q\) & Repository background &
Parameter-free algebraic compression followed by \(\Omega\).\\
Canonical real field and coefficient at zero from \(\Oz\) &
Repository background &
Quadratic detector and multiplier ring; no omega symbol is needed in
the defining formulas.\\
Global real consolidation & Extension &
Choice of a well-order of \(\R\), then substitution of definitions;
equivalent to \(V=\HOD(r)\) for one real.\\
Uniform real and \(\kappa\)-code sections & Extension &
Exact HOD successor-cardinal criteria; \(\Delta_2\) encoder bound is
conditional on the agreement hypothesis.\\
Normal-form absoluteness without common reals & Extension &
Absolute cut tables and persistent sign prefixes, plus the classical
normal-form sign-sum identity.\\
Set theory in bounded omnific monomials & Extension &
Coefficient recovery, lexicographic pairing, exact support subsets and
relations, and set-sized Mostowski collapse.\\
Hahn summability and sum for coded families & Extension &
One fixed formula uses support-union existence, finite fibres and
finite partial-sum codes.\\
Coefficient-indexed derivations and partial integration &
Definable instance and extension &
General diagonal Euler derivations have precedents; this family and its
normalized partial inverses are given by intrinsic coefficient formulas.\\
Fixed Levy levels and literal description length &
Levels preserved; length bound proposed &
\(\Delta_1\) graph, range and inverse for \(Q\).
The finite-alphabet convention alone does not finish the additive-bound
accounting; see \cref{dsn:op:q:length}. The earlier row called both an
``Extension'' with ``a stated finite-alphabet syntax convention for the
additive bound''; the length claim remains source~06's credited proposal.\\
Elementary-theory and inclusion equivalences & Consequence &
Two syntactic translations, and absoluteness of old graph codes for
the inclusion statement.\\
Exact intrinsic definable closure; bi-interpretation &
Not proved &
The missing comparison map and exact closure problems remain research
questions.\\
\end{longtable}
\endgroup

The main proof chain can be read in four layers.  Classical surreal
theory supplies legal series and uniqueness of coefficients.  The
omnific predicate and omega map expose these coefficients to a fixed
first-order language.  Ordered support coding then exposes all subsets
and relations on each set support.  Finally, ordinary well-founded
collapse identifies the quotient with set theory.  The real-code
section criterion is independent of this intrinsic interpretation: it
is an ambient definability theorem about HOD and ordinal enumerations.

\subsection{A concrete formalization sequence}

A useful formalization effort would proceed in the following order.
\begin{enumerate}
\item Fix canonical sign representations and prove the bounded
simplest-separator test.  Construct the finite-input, set-sized
dependency tables for arithmetic and \(\Omega\).
\item Connect the sign and normal-form representations.  Prove that
the persistent-prefix limit evaluates legal normal forms, with explicit
set witnesses at limit stages.  Establish the \(\Delta_1\) graph results
in the chosen set-theoretic coding.
\item Formalize the three reconstruction formulas for \(\Pi\),
\(\mathcal A\) and \(\R\), then coefficient extraction and truncation.
These are small interfaces on which the later interpretation depends.
\item Verify the lexicographic pairing, the exact subset and relation
coding lemmas, and the three admissibility conditions.  Define the
isomorphism and membership formulas without using collapse as a
primitive of the surreal language.
\item Prove the semantic correspondence by the set-sized Mostowski
collapse theorem.  Only after that, verify the recursive translation of
formulas and its preservation of truth, one standard formula at a time.
\item Add family codes, finite coefficient folds, Hahn sums and
derivations.  Develop the HOD and cardinal-correctness section in a
separate set-theory layer, because those proofs have different
foundational dependencies.
\end{enumerate}

This article contains written proofs and was subjected to independent
mathematical checks during preparation, particularly at the support
coding, inner-model absoluteness, cardinal selection and syntax-cost
steps.  It is not a Lean, Mizar or other proof-assistant development.
No finite numerical experiment is offered as evidence for a transfinite
existence theorem.  The supplied source compiles without an external
bibliography database; the accompanying archive contains the PDF and
the precise source used to generate it.

\begin{writenote}[4 October 2026]
In this report the manuscript's source, its PDF and its delivery README are
not shipped; they survive in the arrival commit's archive. Only
\fn{10-definable-operations-PROOF_STATUS.txt} is shipped
(\cref{dsn:op:sub:provenance}).
\end{writenote}



\clearpage
\begin{thebibliography}{99}
\addcontentsline{toc}{section}{References}

\bibitem{Conway}
J. H. Conway, \emph{On Numbers and Games}, second edition,
A K Peters, 2001; first edition, Academic Press, 1976.

\bibitem{Gonshor}
H. Gonshor, \emph{An Introduction to the Theory of Surreal Numbers},
London Mathematical Society Lecture Note Series 110,
Cambridge University Press, 1986.

\bibitem{vdDE}
L. van den Dries and P. Ehrlich,
Fields of surreal numbers and exponentiation,
\emph{Fundamenta Mathematicae} \textbf{167} (2001), 173--188.
\href{https://doi.org/10.4064/fm167-2-3}{doi:10.4064/fm167-2-3}.
In particular Theorem 2.1, Lemmas 4.1--4.2, Proposition 4.7, and
Corollary 5.5. To be read with \cite{vdDEerr}.

\bibitem{vdDEerr}
L. van den Dries and P. Ehrlich, Erratum to ``Fields of surreal numbers and
exponentiation'', \emph{Fundamenta Mathematicae} \textbf{168} (2001),
295--297.
\href{https://doi.org/10.4064/fm168-3-5}{doi:10.4064/fm168-3-5}.

\bibitem{BG}
O. Bournez and Q. Guilmant, Surreal fields stable under exponential and
logarithmic functions, arXiv:2201.08199, 2022.
\url{https://arxiv.org/abs/2201.08199}.

\bibitem{LM}
S. L'Innocente and V. Mantova,
\emph{A factorisation theory for generalised power series and omnific
integers}, arXiv:1710.07304, version 5, January 22, 2024.
\url{https://arxiv.org/abs/1710.07304}.
Journal version: \href{https://doi.org/10.1016/j.aim.2024.109513}
{doi:10.1016/j.aim.2024.109513}.

\bibitem{Marker}
D. Marker, \emph{Model Theory: An Introduction}, Graduate Texts in
Mathematics 217, Springer, 2002.
For the quantifier-elimination argument, see also the author's
\emph{Introductory Lectures on Model Theory II: Quantifier Elimination},
October 14, 2020,
\url{https://homepages.math.uic.edu/~marker/MSRI-2.pdf}.

\bibitem{vddTame}
L. van den Dries, \emph{Tame Topology and o-Minimal Structures}, London
Mathematical Society Lecture Note Series 248, Cambridge University Press,
1998.

\bibitem{EhrlichKaplan}
P. Ehrlich and E. Kaplan, \emph{Number systems with simplicity hierarchies:
a generalization of Conway's theory of surreal numbers II}, preprint,
arXiv:1512.04001, version 1, 2015. \url{https://arxiv.org/abs/1512.04001v1}.

\bibitem{Jech}
T. Jech, \emph{Set Theory: The Third Millennium Edition, Revised and
Expanded}, Springer Monographs in Mathematics, Springer, 2003.
\href{https://doi.org/10.1007/3-540-44761-X}{doi:10.1007/3-540-44761-X}.

\bibitem{Shepherdson}
J. C. Shepherdson,
A non-standard model for a free variable fragment of number theory,
\emph{Bulletin de l'Acad\'emie Polonaise des Sciences, S\'erie des
sciences math\'ematiques, astronomiques et physiques}
\textbf{12} (1964), 79--86.
The integer-part application to $\Oz$ is also stated explicitly
in \cite[Section 1.1]{LM}.

\bibitem{HLR}
J. D. Hamkins, D. Linetsky, and J. Reitz,
Pointwise definable models of set theory,
\emph{Journal of Symbolic Logic} \textbf{78} (2013), no.~1, 139--156.
\url{https://arxiv.org/abs/1105.4597}.
Author's discussion:
\url{https://jdh.hamkins.org/pointwisedefinablemodelsofsettheory/}.

\bibitem{HamkinsDensity}
J. D. Hamkins, answer to ``Definable map from all the ordinals to the
surreal numbers with a dense image'', \emph{MathOverflow}, 8 April 2012.
\url{https://mathoverflow.net/questions/93468/definable-map-from-all-the-ordinals-to-the-surreal-numbers-with-a-dense-image}.

\bibitem{HamkinsLeahy}
J. D. Hamkins and C. Leahy,
Algebraicity and implicit definability in set theory,
\emph{Notre Dame Journal of Formal Logic} \textbf{57} (2016),
no.~3, 431--439.
\url{https://arxiv.org/abs/1305.5953};
\href{https://doi.org/10.1215/00294527-3542326}
{doi:10.1215/00294527-3542326}.

\bibitem{GoldbergPoveda}
G. Goldberg and A. Poveda,
Compactness phenomena in HOD,
\emph{Forum of Mathematics, Sigma} \textbf{13} (2025), e118.
\href{https://doi.org/10.1017/fms.2025.10070}
{doi:10.1017/fms.2025.10070}.
Author's version:
\url{https://math.berkeley.edu/~goldberg/Papers/CompactnessPhenomenaInHOD.pdf}.

\bibitem{BCHGP2026}
T. Benhamou, J. Cummings, G. Goldberg, Y. Hayut, and A. Poveda,
\emph{Compactness phenomena in HOD and the Optimality of Magidor's
Covering theorem}, arXiv:2608.24190, version 1, August 25, 2026.
\url{https://arxiv.org/abs/2608.24190}.
Theorems 1.1--1.2 supply the cited first-disagreement results;
these are external inputs, not newly proved claims of this report.

\bibitem{CHY}
J. D. Hamkins, \emph{Surreal arithmetic is bi-interpretable with set
theory}, CUNY Logic Workshop announcement, posted 4 February 2026 for a
talk on 13 March 2026; joint work with Junhong Chen and Ruizhi Yang,
described as work in progress.
\url{https://jdh.hamkins.org/surreal-arithmetic-cuny-logic-workshop-march-2026/}.
Cited as an announcement, not a refereed proof source.

\bibitem{WikiDef}
Wikipedia contributors, \emph{Definable real number},
user-specified starting reference, accessed September 23, 2026.
\url{https://en.wikipedia.org/wiki/Definable_real_number}.
No theorem depends on this secondary reference.

\bibitem{Repo}
V. Reshetnikov and contributors, \emph{Surreal}, GitHub repository;
the sources inspected it at commit \repoSHA, 23 September 2026.
\url{https://github.com/VladimirReshetnikov/Surreal}.
The repository describes its manuscripts as AI-assisted research drafts and
tracks formalization separately.

\bibitem{RepoHset}
\emph{Birthday Cutoffs and Hereditary Sets}, merged research report in
\cite{Repo},
\repopath{docs/foundations-and-computation/birthday-cutoffs-and-hereditary-sets/}
(label prefix \texttt{hset:}).

\bibitem{RepoUniv}
\emph{Surreal Fields Across Set-Theoretic Universes}, merged research report
in \cite{Repo},
\repopath{docs/foundations-and-computation/surreal-fields-across-universes/}
(label prefix \texttt{univ:}).

\bibitem{RepoOdg}
\emph{Omnific Integers and Omnific--Diophantine Geometry}, merged research
report in \cite{Repo},
\repopath{docs/surreal/omnific-diophantine-geometry/}
(label prefix \texttt{odg:}).

\bibitem{RepoOsq}
\emph{Set-Sized Quotients of Omnific Integers}, merged research report in
\cite{Repo},
\repopath{docs/surreal/set-sized-quotients-of-omnific-integers/}
(label prefix \texttt{osq:}).

\bibitem{RepoOpa}
\emph{Omnific-Preserving Automorphisms}, merged research report in
\cite{Repo},
\repopath{docs/surreal/omnific-preserving-automorphisms/}
(label prefix \texttt{opa:}).


% ----- added for Parts II and III (batch 93) -----
\bibitem{RepoSse}
\emph{Surreal Self-Embeddings}, merged research report in \cite{Repo},
\repopath{docs/surreal/surreal-self-embeddings/} (label prefix
\texttt{sse:}). Its question \texttt{sse:sr:q:stabilizers} asks for the
common fixed field with or without omnific or strong-sum constraints.

\bibitem{RepoLce}
\emph{Large-Cardinal Embeddings and Normal Forms}, research report in
\cite{Repo},
\repopath{docs/foundations-and-computation/large-cardinal-embeddings-and-normal-forms/}
(label prefix \texttt{lce:}).

\bibitem{RepoRvs}
\emph{The Surreal Numbers as a Real Vector Space: Conway--Hahn coordinates,
Hamel bases, strong duality, canonical operators, and the Euler quotient},
research report in \cite{Repo},
\repopath{docs/surreal/real-vector-space-structure/} (label prefix
\texttt{rvs:}). Cited in Part~III only for the general diagonal-derivation
framework.

\bibitem{DoraisHamkins}
F. G. Dorais and J. D. Hamkins,
When does every definable nonempty set have a definable element?,
\emph{Mathematical Logic Quarterly} \textbf{65} (2019), no.~4, 407--411;
\href{https://arxiv.org/abs/1706.07285}{arXiv:1706.07285};
\href{https://doi.org/10.1002/malq.201700035}{doi:10.1002/malq.201700035}.
Sections~2 and~4: the reflected definition of $\OD(X)$ and the $\Delta_2$
ordinal enumeration. Used for the rank-local framework, not for an assertion
that $\HOD(\R)$ satisfies Choice.

\bibitem{Kunen}
K. Kunen, \emph{Set Theory}, Studies in Logic 34, College Publications,
2011. Used for standard forcing, homogeneity, and hereditary definability.

\bibitem{KaragilaForcing}
A. Karagila, \emph{Lecture Notes: Forcing \& Symmetric Extensions}, version
dated 24 August 2026, especially Section~9.2 and Theorem~9.10.
\url{https://karagila.org/files/Forcing-2023.pdf}. The filename predates the
version consulted by source~02.

\bibitem{Prikry}
K. L. Prikry, Changing measurable into accessible cardinals,
\emph{Dissertationes Mathematicae} \textbf{68} (1970), 5--52.

\bibitem{Benhamou}
T. Benhamou, Prikry forcing and tree Prikry forcing of various filters,
\emph{Archive for Mathematical Logic} \textbf{58} (2019), 787--817.
\url{https://doi.org/10.1007/s00153-019-00660-3};
\url{https://arxiv.org/abs/1801.04424}. The classical normal-measure Prikry
property and its no-new-bounded-subsets consequence are the inputs source~05
uses.

\bibitem{LietzNamba2024}
A. Lietz, \emph{The Basics of Namba Forcing}, author-written exposition,
22 March 2024. \url{https://andreas-lietz.github.io/blog/NambaBlog.html}.
The opening distinguishes the two tree presentations; the CH no-new-reals
proof is stated to apply to both.

\bibitem{MillerForcing2009}
A. W. Miller, \emph{Forcing Tidbits}, Math 873, Fall 2009, version dated
21 December 2009, Theorem~23 and the comparison table on p.~4.
\url{https://people.math.wisc.edu/~awmille1/old/m873-09/forcing.pdf}.

\bibitem{ShelahPIF}
S. Shelah, \emph{Proper and Improper Forcing}, second edition, Perspectives in
Mathematical Logic, Springer, 1998. Chapter~XI treats Namba forcing. Added at
the write as a standard reference; not consulted (see
\cref{dsn:rp:q:nambainput}).

\bibitem{FKW}
V. Fischer, M. Koelbing and W. Wohofsky, Fresh function spectra,
\emph{Annals of Pure and Applied Logic} \textbf{174} (2023), no.~9, 103300.
\url{https://doi.org/10.1016/j.apal.2023.103300}. Section~7.6 (Namba forcing
under CH), as imported by the universes report.

\bibitem{EhrlichKaplanJSL}
P. Ehrlich and E. Kaplan, Surreal ordered exponential fields,
\emph{Journal of Symbolic Logic} \textbf{86} (2021), no.~3, 1066--1115.
\url{https://arxiv.org/abs/2002.07739}. Background on canonical surreal
structures, not a proof source for the assembly theorems.

\bibitem{Solovay}
R. M. Solovay, A model of set-theory in which every set of reals is Lebesgue
measurable, \emph{Annals of Mathematics} (2) \textbf{92} (1970), no.~1,
1--56. \url{https://www.jstor.org/stable/1970696}. The countable-ordinal
definability construction is in III.2.2--III.2.6, pp.~51--52.

\bibitem{KL}
V. Kanovei and V. Lyubetsky, \emph{A countable equivalence relation of second
projective class not generated by projective functions}, arXiv:2605.03126v1,
2026. \url{https://arxiv.org/html/2605.03126v1}. Proposition~2.2(i) records
countable $\ROD$ assembly in the Levy--Solovay extension, citing Solovay. The
unversioned record has since used the title \emph{Locally countable graphs of
second projective class not generated by countably many projective
functions}; the cited proposition is from version~1.

\bibitem{LippariniMezo}
P. Lipparini and I. Mez\H{o}, \emph{Surreal limits}, arXiv:1603.09289v3,
8 June 2026. \url{https://arxiv.org/abs/1603.09289}. Definition~2.1 and
Corollary~3.2: the persistent-prefix sign limit and legal normal-form sums.

\bibitem{PK}
K. P\k{a}k and C. Kaliszyk, \emph{Conway Normal Form: Bridging Approaches for
Comprehensive Formalization of Surreal Numbers}, arXiv:2410.00065, 2024.
\url{https://arxiv.org/abs/2410.00065}.

\bibitem{Camacho}
S. Camacho, Truncation in Hahn fields is undecidable and wild,
\emph{South American Journal of Logic} \textbf{5} (2019), no.~1, 39--47.
\url{https://arxiv.org/abs/1706.03722}.

\bibitem{BKMMe}
A. Berarducci, S. Kuhlmann, V. Mantova, and M. Matusinski,
Exponential fields and Conway's omega-map,
\emph{Proceedings of the American Mathematical Society} \textbf{151} (2023),
no.~6, 2655--2669.
\href{https://doi.org/10.1090/proc/14577}{doi:10.1090/proc/14577};
\url{https://arxiv.org/abs/1810.03029}.

\bibitem{KKS}
E. Kaplan, L. S. Krapp, and M. Serra, \emph{Decomposing the automorphism group
of the surreal numbers}, author-hosted manuscript, accessed 4 October 2026 by
source~06.
\url{https://elliotakaplan.github.io/Decomposing_the_automorphism_group_of_the_surreal_numbers.pdf}.
See Proposition~5.5 and Questions~5.6--5.7.

\bibitem{CHYKobe}
J. D. Hamkins, \emph{The elementary theory of surreal arithmetic is
bi-interpretable with set theory}, Kobe conference announcement, posted
20 August 2025 for a September 2025 talk; joint work with Junhong Chen and
Ruizhi Yang, described as work in progress.
\url{https://jdh.hamkins.org/theory-of-surreal-arithmetic-bi-interpretable-with-set-theory-japan-september-2025/}.
Source~06 also consulted the November 2025 Notre Dame slides. The signature
includes birthday order; no theorem of Part~III depends on the announcement.
Compare \cite{CHY}.

\end{thebibliography}
\end{document}
